\documentclass[11pt,reqno]{amsart}
\usepackage{iftex}
\ifPDFTeX
  \usepackage[T1]{fontenc}
  \usepackage[utf8]{inputenc}
  \usepackage{lmodern}
\else
  \usepackage[no-math]{fontspec}
\fi
\usepackage{amsmath}
\usepackage{amssymb}
\usepackage{amsfonts}
\usepackage{mathrsfs}
\usepackage{amsthm}
\usepackage{mathtools}
\usepackage{geometry}
\usepackage{enumitem}
\usepackage{longtable}
\usepackage{booktabs}
\usepackage{array}
\usepackage{microtype}
\usepackage{graphicx}
\graphicspath{{../figures/}{figures/}}
\usepackage{listings}
\lstset{basicstyle=\ttfamily\scriptsize,breaklines=true,columns=fullflexible,keepspaces=true,
  frame=single,framerule=0.3pt,showstringspaces=false,keywordstyle=\color{NavyBlue},
  commentstyle=\color{OliveGreen},stringstyle=\color{BrickRed},upquote=true}
\lstdefinestyle{out}{basicstyle=\ttfamily\scriptsize,language={},keywordstyle={},commentstyle={},stringstyle={}}
\usepackage{float}
\usepackage[dvipsnames]{xcolor}
\usepackage[pagebackref=true]{hyperref}
\usepackage{bookmark}
\usepackage{pdfcomment}

% Hover previews for internal references (no layout change).
\makeatletter
\let\CG@origref\ref
\let\CG@origeqref\eqref
\newcommand{\CG@gettip}[1]{\@ifundefined{CGTip@#1}{Reference: #1}{\csname CGTip@#1\endcsname}}
\DeclareRobustCommand{\CGref}[1]{\pdftooltip{\CG@origref{#1}}{\CG@gettip{#1}}}
\DeclareRobustCommand{\CGeqref}[1]{\pdftooltip{\CG@origeqref{#1}}{\CG@gettip{#1}}}
\makeatother
% Tooltip texts shown when hovering over internal references.
\makeatletter
\expandafter\gdef\csname CGTip@sec:notation\endcsname{Section : Notation (p. 3)}
\expandafter\gdef\csname CGTip@sec:introduction\endcsname{Section 1: Introduction (p. 5)}
\expandafter\gdef\csname CGTip@subsec:background\endcsname{Section 1.1: Background (p. 5)}
\expandafter\gdef\csname CGTip@eq:mertens-def\endcsname{Eq. (1.1), p. 6: M(x):=sum\_n<= xmu(n).}
\expandafter\gdef\csname CGTip@subsec:midpoint-system\endcsname{Section 1.2: The midpoint RVM system (p. 6)}
\expandafter\gdef\csname CGTip@eq:intro-clock\endcsname{Eq. (1.2), p. 6: N\_RVM(T)=(T)/(2pi)log(T)/(2pi e)+7/8,}
\expandafter\gdef\csname CGTip@eq:intro-lambert\endcsname{Eq. (1.3), p. 6: tau\_k=(2pi(k-11/8))/(W\_0((k-11/8)/e)) .}
\expandafter\gdef\csname CGTip@eq:intro-zeta-rvm\endcsname{Eq. (1.4), p. 6: P\_RVM\textasciicircum{}crit(s)=prod\_k>=1(1+(s(s-1))/(1/4+tau\_k\textasciicircum{}2)), zeta\_RVM(s)=P\_RVM\textasciicircum{}crit(s)G(s), G(s)=s(s-1)Gamma(s/2)pi\textasciicircum{}-s/2.}
\expandafter\gdef\csname CGTip@eq:intro-pi\endcsname{Eq. (1.5), p. 7: pi\_RVM(x)=sum\_1<= m<=log\_2x(mu(m))/m J\_RVM\textasciicircum{}*(x\textasciicircum{}1/m), J\_RVM\textasciicircum{}*(x)=sum\_1<= m<=log\_2x1/m pi\_RVM(x\textasciicircum{}1/m).}
\expandafter\gdef\csname CGTip@eq:intro-real-primes\endcsname{Eq. (1.6), p. 7: p\_j\textasciicircum{}RVM:=pi\_RVM\textasciicircum{}-1(j), pi\_RVM(p\_j\textasciicircum{}RVM)=j,}
\expandafter\gdef\csname CGTip@eq:intro-transition\endcsname{Eq. (1.7), p. 7: zeta\_RVM(s)=varpi\_RVM(1+w(s))(1+O(1/|t|)), w(s)=exp-2pi i N\_RVM(-i(s-1/2)),}
\expandafter\gdef\csname CGTip@eq:intro-Mbound\endcsname{Eq. (1.8), p. 7: M\_RVM\textasciicircum{}crit(x)=L+O\_eps(x\textasciicircum{}1/2+eps)}
\expandafter\gdef\csname CGTip@eq:intro-muh\endcsname{Eq. (1.9), p. 7: mu\_h(n)=mu(n)\textasciicircum{}2varsigma\_n, varsigma\_n in -1,+1, M\_h(n):=sum\_m<= nmu\_h(m),}
\expandafter\gdef\csname CGTip@eq:intro-quantizer\endcsname{Eq. (1.10), p. 7: varsigma\_n= +1, M\_h(n-1)<M(n),; -1, M\_h(n-1)>=M(n),}
\expandafter\gdef\csname CGTip@subsec:core\endcsname{Section 1.3: The core of the method (p. 8)}
\expandafter\gdef\csname CGTip@eq:intro-family\endcsname{Eq. (1.11), p. 8: gamma\_k(lambda)=(1-lambda)gamma\_k+lambdatau\_k, zeta\_lambda(s)=1/(G(s))prod\_k>=1(1+(s(s-1))/(1/4+gamma\_k(lambda)\textasciicircum{}2)),}
\expandafter\gdef\csname CGTip@eq:intro-hamburger\endcsname{Eq. (1.12), p. 8: zeta\_Z1=zeta <=> (D): zeta\_Z1 is an absolutely convergent Dirichlet series in Re s>1 .}
\expandafter\gdef\csname CGTip@eq:intro-varpi-Z1\endcsname{Eq. (1.13), p. 8: varpi\_Z1=lim\_sigma->+infzeta\_Z1(sigma)=varpi\_RVMprod\_k>=1(1/4+tau\_k\textasciicircum{}2)/(1/4+gamma\_k\textasciicircum{}2) >=1,}
\expandafter\gdef\csname CGTip@eq:intro-step3\endcsname{Eq. (1.14), p. 9: zeta\_Z1=zeta => M(x)=O\_eps(x\textasciicircum{}1/2+eps) (eps>0),}
\expandafter\gdef\csname CGTip@eq:intro-prime-chain\endcsname{Eq. (1.15), p. 9: tau\_k->zeta\_RVM-> J\_RVM\_RVM->(dpi\_RVM, p\_j\textasciicircum{}RVM, c\_n\textasciicircum{}RVM),}
\expandafter\gdef\csname CGTip@eq:intro-M-chain\endcsname{Eq. (1.16), p. 9: F\_RVM=zeta\_RVM\textasciicircum{}-1-> M\_RVM\textasciicircum{}crit->(M,mu\_h,M\_h),}
\expandafter\gdef\csname CGTip@eq:intro-family-chain\endcsname{Eq. (1.17), p. 9: zeta\_RVM=zeta\_1->zeta\_0=zeta\_Z1=?zeta => M(x)=O\_eps(x\textasciicircum{}1/2+eps).}
\expandafter\gdef\csname CGTip@subsec:status\endcsname{Section 1.4: Originality, scope and status (p. 9)}
\expandafter\gdef\csname CGTip@subsec:organization\endcsname{Section 1.5: Organization (p. 10)}
\expandafter\gdef\csname CGTip@sec:clock\endcsname{Section 2: The midpoint Riemann--von Mangoldt clock (p. 11)}
\expandafter\gdef\csname CGTip@eq:zero-count\endcsname{Eq. (2.1), p. 11: N\_zeta(T):=no.rho: 0<Imrho<= T=no.k:gamma\_k<= T.}
\expandafter\gdef\csname CGTip@eq:RvM-full\endcsname{Eq. (2.2), p. 11: N\_zeta(T)=(T)/(2pi)log(T)/(2pi e)+7/8+S(T)+O(1/T), S(T):=1pi(1/2+iT)=O(log T),}
\expandafter\gdef\csname CGTip@eq:clock\endcsname{Eq. (2.3), p. 11: N\_RVM(T):=(T)/(2pi)log(T)/(2pi e)+7/8 .}
\expandafter\gdef\csname CGTip@eq:midpoint-condition\endcsname{Eq. (2.4), p. 11: N\_RVM(tau\_k)=k-1/2, k>=1,}
\expandafter\gdef\csname CGTip@prop:lambert\endcsname{Proposition 2.1 (Lambert inversion) (p. 11): For every k>=1, tau\_k=(2pi(k-11/8))/(W\_0((k-11/8)/e)), The quotient is positive also for k=1.}
\expandafter\gdef\csname CGTip@eq:lambert\endcsname{Eq. (2.5), p. 11: tau\_k=(2pi(k-11/8))/(W\_0((k-11/8)/e)),}
\expandafter\gdef\csname CGTip@eq:lambert-proof-y\endcsname{Eq. (2.6), p. 11: y(log y-1)=k-(11)/(8).}
\expandafter\gdef\csname CGTip@eq:clock-derivative\endcsname{Eq. (2.7), p. 11: N\_RVM'(T)=1/(2pi)log(T)/(2pi),}
\expandafter\gdef\csname CGTip@eq:clock-growth\endcsname{Eq. (2.8), p. 11: tau\_k (k)/(log k) (k>=2).}
\expandafter\gdef\csname CGTip@eq:RvM-formula\endcsname{Eq. (2.9), p. 11: N\_zeta(T)=N\_RVM(T)+S(T)+O(T\textasciicircum{}-1), S(T)=O(log T) (T>=2).}
\expandafter\gdef\csname CGTip@eq:symmetric-count\endcsname{Eq. (2.10), p. 12: N\_zeta\textasciicircum{}*(T):=1/2(N\_zeta(T+0)+N\_zeta(T-0))}
\expandafter\gdef\csname CGTip@eq:eta-def\endcsname{Eq. (2.11), p. 12: eta\_k:=tau\_k-gamma\_k .}
\expandafter\gdef\csname CGTip@lem:bounded-displacement\endcsname{Lemma 2.2 (Bounded displacement) (p. 12): There is a constant C\_eta such that |eta\_k|<= C\_eta (k>=1).}
\expandafter\gdef\csname CGTip@eq:bounded-displacement\endcsname{Eq. (2.12), p. 12: |eta\_k|<= C\_eta (k>=1).}
\expandafter\gdef\csname CGTip@eq:block-count\endcsname{Eq. (2.13), p. 12: |N\_zeta\textasciicircum{}*(gamma\_k)-(k-1/2)|<=(m(gamma\_k))/2 .}
\expandafter\gdef\csname CGTip@eq:clock-discrepancy\endcsname{Eq. (2.14), p. 12: |N\_RVM(gamma\_k)-N\_RVM(tau\_k)|=O(loggamma\_k).}
\expandafter\gdef\csname CGTip@sec:zeta-rvm\endcsname{Section 3: The pure critical RVM zeta function (p. 12)}
\expandafter\gdef\csname CGTip@eq:q-def\endcsname{Eq. (3.1), p. 12: q(s):=s(s-1),}
\expandafter\gdef\csname CGTip@eq:Pcrit\endcsname{Eq. (3.2), p. 12: P\_RVM\textasciicircum{}crit(s):=prod\_k>=1(1+(q(s))/(1/4+tau\_k\textasciicircum{}2)),}
\expandafter\gdef\csname CGTip@eq:Gdef\endcsname{Eq. (3.3), p. 12: G(s):=s(s-1)Gamma(s/2)pi\textasciicircum{}-s/2,}
\expandafter\gdef\csname CGTip@eq:zetaRVM\endcsname{Eq. (3.4), p. 12: xi\_RVM(s):=1/2P\_RVM\textasciicircum{}crit(s), zeta\_RVM(s):=P\_RVM\textasciicircum{}crit(s)G(s).}
\expandafter\gdef\csname CGTip@thm:canonical\endcsname{Theorem 3.1 (Canonical characterization of zeta\_RVM) (p. 12): The product CGeqrefeq:Pcrit converges absolutely and locally uniformly in s in C. It defines an entire function with P\_RVM\textasciicircum{}crit(1-s)=P\_RVM\textasciicircum{}crit(s), P\_RVM\textasciicircum{}crit(conj s)=conj P\_RVM\textasciicircum{}crit(s), P\_RVM\textasciicircum{}crit(0)=P\_RVM\textasciicircum{}crit(1)=1 . Its zeros are exactly the simple zeros rho\_k\textasciicircum{}RVM=1/2+itau\_k, conj rho\_k\textasciicircum{}RVM=1/...}
\expandafter\gdef\csname CGTip@eq:P-symmetry\endcsname{Eq. (3.5), p. 12: P\_RVM\textasciicircum{}crit(1-s)=P\_RVM\textasciicircum{}crit(s), P\_RVM\textasciicircum{}crit(conj s)=conj P\_RVM\textasciicircum{}crit(s), P\_RVM\textasciicircum{}crit(0)=P\_RVM\textasciicircum{}crit(1)=1 .}
\expandafter\gdef\csname CGTip@eq:synthetic-divisor\endcsname{Eq. (3.6), p. 12: rho\_k\textasciicircum{}RVM=1/2+itau\_k, conj rho\_k\textasciicircum{}RVM=1/2-itau\_k, k>=1 .}
\expandafter\gdef\csname CGTip@eq:canonical-normal\endcsname{Eq. (3.7), p. 12: sup\_s in K|(q(s))/(1/4+tau\_k\textasciicircum{}2)|<<\_Ktau\_k\textasciicircum{}-2,}
\expandafter\gdef\csname CGTip@eq:factor-zero\endcsname{Eq. (3.8), p. 12: (s-1/2)\textasciicircum{}2+tau\_k\textasciicircum{}2=0 .}
\expandafter\gdef\csname CGTip@sec:exp-sums\endcsname{Section 4: RVM exponential sums, centroids, and the height transform (p. 13)}
\expandafter\gdef\csname CGTip@eq:S-RVM\endcsname{Eq. (4.1), p. 13: S\_RVM(N;u):=sum\_k<= Ne\textasciicircum{}iutau\_k, C\_RVM\textasciicircum{}(N;u):=S\_RVM(N;u)N,}
\expandafter\gdef\csname CGTip@eq:path-centroid\endcsname{Eq. (4.2), p. 13: C\_RVM\textasciicircum{}(N;u):=1/Nsum\_n<= NS\_RVM(n;u) =sum\_k<= N(N-k+1)/(N)e\textasciicircum{}iutau\_k.}
\expandafter\gdef\csname CGTip@thm:exp-sum\endcsname{Theorem 4.1 (Uniform logarithmic exponential-sum bound) (p. 13): Let 0<u\_1<= u\_2<inf. Then sup\_u in [u\_1,u\_2]|S\_RVM(N;u)|<<\_u\_1,u\_21+log N .}
\expandafter\gdef\csname CGTip@eq:exp-sum-bound\endcsname{Eq. (4.3), p. 13: sup\_u in [u\_1,u\_2]|S\_RVM(N;u)|<<\_u\_1,u\_21+log N .}
\expandafter\gdef\csname CGTip@eq:inverse-clock\endcsname{Eq. (4.4), p. 13: T'(y)=(2pi)/(log(T(y)/2pi)), T''(y)=-(4pi\textasciicircum{}2)/(T(y)log\textasciicircum{}3(T(y)/2pi)) .}
\expandafter\gdef\csname CGTip@eq:rvm-gap-bounds\endcsname{Eq. (4.5), p. 13: 2pi\_k+1<=tau\_k+1-tau\_k<=(2pi)/(\_k),}
\expandafter\gdef\csname CGTip@eq:rvm-second-difference\endcsname{Eq. (4.6), p. 13: |tau\_k+2-2tau\_k+1+tau\_k|<=(4pi\textasciicircum{}2)/(tau\_k\_k\textasciicircum{}3).}
\expandafter\gdef\csname CGTip@eq:tau-ell-k\endcsname{Eq. (4.7), p. 13: tau\_k\_k k (k>=2).}
\expandafter\gdef\csname CGTip@eq:Ek-identity\endcsname{Eq. (4.8), p. 13: E\_k=b\_k(E\_k+1-E\_k).}
\expandafter\gdef\csname CGTip@eq:summation-by-parts\endcsname{Eq. (4.9), p. 13: sum\_k=k\_0\textasciicircum{}NE\_k=b\_NE\_N+1-b\_k\_0E\_k\_0+sum\_k=k\_0+1\textasciicircum{}N(b\_k-1-b\_k)E\_k .}
\expandafter\gdef\csname CGTip@eq:kappa-bounds\endcsname{Eq. (4.10), p. 13: |b\_k|<=(pi)/(2Delta\_k)<=\_k+14u\_1, |b\_k-1-b\_k|<=|Delta\_k-Delta\_k-1| |b\_k-1||b\_k| <=(pi\textasciicircum{}2u\_2)/(4u\_1\textasciicircum{}2) \_k+1\textasciicircum{}2tau\_k-1\_k-1\textasciicircum{}3<<\_u\_1,u\_21/(tau\_k\_k)<<1/k,}
\expandafter\gdef\csname CGTip@cor:point-centroid\endcsname{Corollary 4.2 (Point centroid) (p. 14): Uniformly for u in [u\_1,u\_2](0,inf), C\_RVM\textasciicircum{}(N;u)=O\_u\_1,u\_2 ((1+log N)/(N))->0 . Theorem CGrefthm:exp-sum alone does not imply convergence of the path centroid CGeqrefeq:path-centroid.}
\expandafter\gdef\csname CGTip@eq:point-centroid\endcsname{Eq. (4.11), p. 14: C\_RVM\textasciicircum{}(N;u)=O\_u\_1,u\_2 ((1+log N)/(N))->0 .}
\expandafter\gdef\csname CGTip@eq:n-RVM\endcsname{Eq. (4.12), p. 14: n\_RVM(H):=no.k:tau\_k<= H.}
\expandafter\gdef\csname CGTip@eq:MH-RVM\endcsname{Eq. (4.13), p. 14: M\_H\textasciicircum{}RVM(x):=1n\_RVM(H)sum\_tau\_k<= Hx\textasciicircum{}itau\_k,}
\expandafter\gdef\csname CGTip@eq:MH-genuine\endcsname{Eq. (4.14), p. 14: M\_H(x):=1/(N\_zeta(H))sum\_gamma\_k<= Hx\textasciicircum{}igamma\_k.}
\expandafter\gdef\csname CGTip@cor:MH-decay\endcsname{Corollary 4.3 (Synthetic height-transform decay) (p. 14): For every fixed x>1, with N=n\_RVM(H), M\_H\textasciicircum{}RVM(x)=O\_x ((1+log N)/(N))->0 (H->inf).}
\expandafter\gdef\csname CGTip@eq:MH-decay\endcsname{Eq. (4.15), p. 14: M\_H\textasciicircum{}RVM(x)=O\_x ((1+log N)/(N))->0 (H->inf).}
\expandafter\gdef\csname CGTip@rem:MH-not-mertens\endcsname{Remark 4.4 (p. 14): Neither CGeqrefeq:MH-RVM nor CGeqrefeq:MH-genuine is the Mertens function; both are normalized phase averages.}
\expandafter\gdef\csname CGTip@sec:explicit-formula\endcsname{Section 5: The RVM explicit formula, J\_RVMJ\_RVM, and pi\_RVMpi\_RVM (p. 14)}
\expandafter\gdef\csname CGTip@eq:Ei-def\endcsname{Eq. (5.1), p. 14: Ei(w):=int\_-inf\textasciicircum{}ae\textasciicircum{}v+ibv+ib dv, Ei(conj w):=conj Ei(w) .}
\expandafter\gdef\csname CGTip@eq:J-RVM\endcsname{Eq. (5.2), p. 14: J\_RVM(x):=li(x)-2Resum\_k>=1Ei ((1/2+itau\_k)log x)-log2+int\_x\textasciicircum{}inf(dy)/(y(y\textasciicircum{}2-1)log y),}
\expandafter\gdef\csname CGTip@lem:Ei\endcsname{Lemma 5.1 (Asymptotics of Ei off the real axis) (p. 14): For w=a+ib, b>0, Ei(w)=(e\textasciicircum{}w)/(w)+varrho(w), |varrho(w)|<=(e\textasciicircum{}a)/(b\textasciicircum{}2).}
\expandafter\gdef\csname CGTip@eq:Ei-remainder\endcsname{Eq. (5.3), p. 14: Ei(w)=(e\textasciicircum{}w)/(w)+varrho(w), |varrho(w)|<=(e\textasciicircum{}a)/(b\textasciicircum{}2).}
\expandafter\gdef\csname CGTip@thm:J-exists\endcsname{Theorem 5.2 (Compact-uniform existence of the RVM zero sum) (p. 14): The series sum\_kReEi((1/2+itau\_k)u) converges uniformly for u in every compact interval [u\_1,u\_2](0,inf). Hence J\_RVM is a continuous function on (1,inf).}
\expandafter\gdef\csname CGTip@eq:Ei-asymptotic\endcsname{Eq. (5.4), p. 15: ReEi((1/2+itau\_k)u)=e\textasciicircum{}u/2u (sin(utau\_k))/(tau\_k)+O\_u\_1,u\_2(tau\_k\textasciicircum{}-2),}
\expandafter\gdef\csname CGTip@eq:abel-weights\endcsname{Eq. (5.5), p. 15: 1/(tau\_K)-1tau\_K+1<=tau\_K+1-tau\_Ktau\_K\textasciicircum{}2<=(2pi)/(\_Ktau\_K\textasciicircum{}2)<<(log K)/(K\textasciicircum{}2)}
\expandafter\gdef\csname CGTip@eq:J-star\endcsname{Eq. (5.6), p. 15: J\_RVM\textasciicircum{}*(x):= J\_RVM(x), x>=2,; 0, 0<x<2,}
\expandafter\gdef\csname CGTip@eq:pi-RVM\endcsname{Eq. (5.7), p. 15: pi\_RVM(x):=sum\_1<= m<=log\_2x(mu(m))/m J\_RVM\textasciicircum{}*(x\textasciicircum{}1/m).}
\expandafter\gdef\csname CGTip@prop:finite-inversion\endcsname{Proposition 5.3 (Exact finite prime-power inversion) (p. 15): For every x>=2, J\_RVM\textasciicircum{}*(x)=sum\_1<= m<=log\_2x1/m pi\_RVM(x\textasciicircum{}1/m).}
\expandafter\gdef\csname CGTip@eq:J-pi-inverse\endcsname{Eq. (5.8), p. 15: J\_RVM\textasciicircum{}*(x)=sum\_1<= m<=log\_2x1/m pi\_RVM(x\textasciicircum{}1/m).}
\expandafter\gdef\csname CGTip@eq:mobius-coefficient\endcsname{Eq. (5.9), p. 15: 1/rsum\_d| rmu(d),}
\expandafter\gdef\csname CGTip@subsec:square-root\endcsname{Section 5.1: The square-root estimate (p. 15)}
\expandafter\gdef\csname CGTip@eq:A-RVM\endcsname{Eq. (5.10), p. 15: A\_RVM(u):=sum\_k>=1e\textasciicircum{}iutau\_ktau\_k,}
\expandafter\gdef\csname CGTip@lem:A-RVM\endcsname{Lemma 5.4 (The weighted clock sum) (p. 15): For every u>=log2 the series CGeqrefeq:A-RVM converges, and |A\_RVM(u)|<=(log\textasciicircum{}2(2+2u))/(4pi)+(log(2+2u))/(2pi)+C\_A, with an absolute constant C\_A.}
\expandafter\gdef\csname CGTip@eq:A-RVM-bound\endcsname{Eq. (5.11), p. 15: |A\_RVM(u)|<=(log\textasciicircum{}2(2+2u))/(4pi)+(log(2+2u))/(2pi)+C\_A,}
\expandafter\gdef\csname CGTip@eq:A-head\endcsname{Eq. (5.12), p. 15: sum\_tau\_k<= Y\_01/(tau\_k)=(n(Y\_0))/(Y\_0)+int\_tau\_1\textasciicircum{}Y\_0(n(t))/(t\textasciicircum{}2) dt<=((Y\_0))/(2pi)+int\_tau\_1\textasciicircum{}Y\_0((t))/(2pi t) dt<=((Y\_0))/(2pi)+((Y\_0)\textasciicircum{}2)/(4pi),}
\expandafter\gdef\csname CGTip@eq:phi-second\endcsname{Eq. (5.13), p. 16: phi''(y)=-(2pi u)/(T (T)\textasciicircum{}3), T=T(y-1/2).}
\expandafter\gdef\csname CGTip@eq:block-vdc\endcsname{Eq. (5.14), p. 16: |sum\_a<k<= be\textasciicircum{}iutau\_k|<= A(16(Y\_j(Y\_j+1)2pi+1)kappa\_j\textasciicircum{}1/2+kappa\_j\textasciicircum{}-1/2)=:P\_j,}
\expandafter\gdef\csname CGTip@eq:block-weighted\endcsname{Eq. (5.15), p. 16: (P\_j)/(Y\_j)<= A((16)/(2pi)(pi uY\_j(Y\_j+1))\textasciicircum{}1/2+(16)/(Y\_j)+((Y\_j+1)\textasciicircum{}3pi uY\_j)\textasciicircum{}1/2).}
\expandafter\gdef\csname CGTip@thm:Lambert-estimate\endcsname{Theorem 5.5 (Square-root estimate for J\_RVM and pi\_RVM) (p. 16): For x>=8, J\_RVM(x)-li(x)=O (( x(loglog x)\textasciicircum{}2)/(log x)), pi\_RVM(x)-li(x)=O (( x(loglog x)\textasciicircum{}2)/(log x)), with absolute implied constants. No hypothesis on zeta is used.}
\expandafter\gdef\csname CGTip@eq:J-baseline\endcsname{Eq. (5.16), p. 16: J\_RVM(x)-li(x)=O (( x(loglog x)\textasciicircum{}2)/(log x)),}
\expandafter\gdef\csname CGTip@eq:pi-baseline\endcsname{Eq. (5.17), p. 16: pi\_RVM(x)-li(x)=O (( x(loglog x)\textasciicircum{}2)/(log x)),}
\expandafter\gdef\csname CGTip@eq:Ei-split\endcsname{Eq. (5.18), p. 16: Ei(rho\_ku)=e\textasciicircum{}rho\_kurho\_ku+varrho\_k, |varrho\_k|<=e\textasciicircum{}u/2tau\_k\textasciicircum{}2u\textasciicircum{}2.}
\expandafter\gdef\csname CGTip@eq:Ei-sum-A\endcsname{Eq. (5.19), p. 16: sum\_k>=1Ei(rho\_ku)=e\textasciicircum{}u/2u(-i A\_RVM(u)+sum\_kr\_ke\textasciicircum{}iutau\_k)+sum\_kvarrho\_k =e\textasciicircum{}u/2u(-i A\_RVM(u)+O(1)).}
\expandafter\gdef\csname CGTip@eq:pi-minus-li\endcsname{Eq. (5.20), p. 16: pi\_RVM(x)-li(x)=(J\_RVM(x)-li(x))+sum\_2<= m<=log\_2x(mu(m))/m J\_RVM\textasciicircum{}*(x\textasciicircum{}1/m).}
\expandafter\gdef\csname CGTip@cor:three-limits\endcsname{Corollary 5.6 (The three limits) (p. 17): As x->inf, li(x)pi\_RVM(x)->1, (pi(x))/(li(x))->1, pi(x)pi\_RVM(x)->1 . More precisely, for some absolute c>0, pi(x)-pi\_RVM(x)=O (x e\textasciicircum{}-csqrt(log x)).}
\expandafter\gdef\csname CGTip@eq:three-limits\endcsname{Eq. (5.21), p. 17: li(x)pi\_RVM(x)->1, (pi(x))/(li(x))->1, pi(x)pi\_RVM(x)->1 .}
\expandafter\gdef\csname CGTip@eq:pi-minus-piRVM\endcsname{Eq. (5.22), p. 17: pi(x)-pi\_RVM(x)=O (x e\textasciicircum{}-csqrt(log x)).}
\expandafter\gdef\csname CGTip@sec:integer-weights\endcsname{Section 6: Integer weights and the continuous-to-atomic problem (p. 17)}
\expandafter\gdef\csname CGTip@eq:cn\endcsname{Eq. (6.1), p. 17: c\_n\textasciicircum{}RVM:=pi\_RVM(n)-pi\_RVM(n-1), n>=2,}
\expandafter\gdef\csname CGTip@eq:weighted-zeta\endcsname{Eq. (6.2), p. 17: zeta\_Z,RVM(s):=prod\_n>=2(1-n\textasciicircum{}-s)\textasciicircum{}-c\_n\textasciicircum{}RVM.}
\expandafter\gdef\csname CGTip@eq:JZ\endcsname{Eq. (6.3), p. 17: J\_Z,RVM(x):=sum\_1<= m<=log\_2x1/m pi\_RVM(floor( x\textasciicircum{}1/m)),}
\expandafter\gdef\csname CGTip@eq:J-star-Z\endcsname{Eq. (6.4), p. 17: J\_RVM\textasciicircum{}*,Z(x):=J\_RVM\textasciicircum{}*(floor( x)).}
\expandafter\gdef\csname CGTip@thm:unit-cell\endcsname{Theorem 6.1 (Lipschitz control of J\_RVM) (p. 17): There is an absolute constant C\_J such that |J\_RVM(x+theta)-J\_RVM(x)|<= C\_J theta (x>=2, 0<theta<= x).}
\expandafter\gdef\csname CGTip@eq:unit-cell\endcsname{Eq. (6.5), p. 17: |J\_RVM(x+theta)-J\_RVM(x)|<= C\_J theta (x>=2, 0<theta<= x).}
\expandafter\gdef\csname CGTip@eq:Dk-def\endcsname{Eq. (6.6), p. 17: D\_k:=Ei(rho\_k\textasciicircum{}RVMu\_2)-Ei(rho\_k\textasciicircum{}RVMu\_1)=int\_u\_1\textasciicircum{}u\_2e\textasciicircum{}v/2e\textasciicircum{}ivtau\_kv dv .}
\expandafter\gdef\csname CGTip@eq:damped-sum\endcsname{Eq. (6.7), p. 17: sum\_ke\textasciicircum{}-epstau\_kD\_k=int\_u\_1\textasciicircum{}u\_2e\textasciicircum{}v/2v S\_eps(v) dv, S\_eps(v):=sum\_k>=1g(tau\_k), g(tau):=e\textasciicircum{}(iv-eps)tau.}
\expandafter\gdef\csname CGTip@eq:S-eps-split\endcsname{Eq. (6.8), p. 17: S\_eps(v)=int\_tau\_0\textasciicircum{}inf g dN\_RVM+g(tau\_0)N\_RVM(tau\_0)+int\_tau\_0\textasciicircum{}infpsi(N\_RVM(tau)+1/2)g'(tau) dtau .}
\expandafter\gdef\csname CGTip@eq:K-harmonics\endcsname{Eq. (6.9), p. 18: int\_tau\_0\textasciicircum{}infpsi(N\_RVM+1/2)g' dtau=-sum\_m>=1((-1)\textasciicircum{}m)/(2pi i m)(K\_m\textasciicircum{}+-K\_m\textasciicircum{}-), K\_m\textasciicircum{}+-:=int\_tau\_0\textasciicircum{}inf(iv-eps)e\textasciicircum{}-epstaue\textasciicircum{}iPhi\_m\textasciicircum{}+-(tau) dtau,}
\expandafter\gdef\csname CGTip@eq:K-minus-bound\endcsname{Eq. (6.10), p. 18: |K\_m\textasciicircum{}-|<=2(v+1)((8)/(mlog2)+8sqrt((2tau\_m)/m)) =2(v+1)((8)/(mlog2)+16 e\textasciicircum{}v/(2m) m).}
\expandafter\gdef\csname CGTip@eq:S-eps-bound\endcsname{Eq. (6.11), p. 18: |S\_eps(v)|<= C (v+1) e\textasciicircum{}v/2 (v>=log2, 0<eps<=1),}
\expandafter\gdef\csname CGTip@eq:Dk-sum-bound\endcsname{Eq. (6.12), p. 18: |sum\_ke\textasciicircum{}-epstau\_kD\_k|<= C'(u\_2-u\_1) x<= C'theta,}
\expandafter\gdef\csname CGTip@cor:cn-bounded\endcsname{Corollary 6.2 (Bounded cell masses) (p. 18): pi\_RVM(y)-pi\_RVM(floor( y))=O(1), c\_n\textasciicircum{}RVM=O(1).}
\expandafter\gdef\csname CGTip@eq:unit-cell-consequence\endcsname{Eq. (6.13), p. 18: pi\_RVM(y)-pi\_RVM(floor( y))=O(1), c\_n\textasciicircum{}RVM=O(1).}
\expandafter\gdef\csname CGTip@cor:ZRVM-Dirichlet\endcsname{Corollary 6.3 (The integer-weighted zeta is a Dirichlet series) (p. 18): The product CGeqrefeq:weighted-zeta converges absolutely for Re s>1 and zeta\_Z,RVM(s)=sum\_N>=1a\_N\textasciicircum{}ZN\textasciicircum{}s, a\_1\textasciicircum{}Z=1, sum\_N>=1|a\_N\textasciicircum{}Z|N\textasciicircum{}sigma<inf (sigma>1). The inverse is also an absolutely convergent Dirichlet series on the whole half-plane: 1zeta\_Z,RVM(s)=sum\_N>=1conj a\_N\textasciicircum{}ZN\textasciicircum{}s, conj a\_1\textasciicircum{}Z=1, sum\_N>=...}
\expandafter\gdef\csname CGTip@eq:ZRVM-Dirichlet\endcsname{Eq. (6.14), p. 18: zeta\_Z,RVM(s)=sum\_N>=1a\_N\textasciicircum{}ZN\textasciicircum{}s, a\_1\textasciicircum{}Z=1, sum\_N>=1|a\_N\textasciicircum{}Z|N\textasciicircum{}sigma<inf (sigma>1).}
\expandafter\gdef\csname CGTip@eq:ZRVM-inverse-Dirichlet\endcsname{Eq. (6.15), p. 18: 1zeta\_Z,RVM(s)=sum\_N>=1conj a\_N\textasciicircum{}ZN\textasciicircum{}s, conj a\_1\textasciicircum{}Z=1, sum\_N>=1|conj a\_N\textasciicircum{}Z|N\textasciicircum{}sigma<inf (sigma>1).}
\expandafter\gdef\csname CGTip@eq:log-ZRVM\endcsname{Eq. (6.16), p. 19: logzeta\_Z,RVM(s)=sum\_n>=2c\_n\textasciicircum{}RVMsum\_m>=1n\textasciicircum{}-msm=sum\_N>=2(b\_N)/(N\textasciicircum{}s), b\_N:=sum\_n\textasciicircum{}m=Nc\_n\textasciicircum{}RVMm,}
\expandafter\gdef\csname CGTip@thm:MS\endcsname{Theorem 6.4 (Mellin--Stieltjes identification) (p. 19): There is a function J\_RVM\textasciicircum{} M on (1,inf) with J\_RVM\textasciicircum{} M(x)x\textasciicircum{}-c in L\textasciicircum{}2((1,inf),dx/x) for every c>1 such that logzeta\_RVM(s)=sint\_1\textasciicircum{}inf J\_RVM\textasciicircum{} M(x) x\textasciicircum{}-s-1 dx (Re s>1), and J\_RVM\textasciicircum{} M-J\_RVM is constant on (1,inf). In particular J\_RVM\textasciicircum{} M=J\_RVM\textasciicircum{}*+C\_RVM on [2,inf) for a constant C\_RVM.}
\expandafter\gdef\csname CGTip@eq:MS-RVM\endcsname{Eq. (6.17), p. 19: logzeta\_RVM(s)=sint\_1\textasciicircum{}inf J\_RVM\textasciicircum{} M(x) x\textasciicircum{}-s-1 dx (Re s>1),}
\expandafter\gdef\csname CGTip@cor:weighted-factorization\endcsname{Corollary 6.5 (Weighted factorization) (p. 19): J\_Z,RVM(x)-J\_RVM\textasciicircum{}*(x)=O(loglog x), and zeta\_Z,RVM(s)=zeta\_RVM(s) e\textasciicircum{}h\_RVM(s), where h\_RVM(s):=sint\_1\textasciicircum{}inf(J\_Z,RVM(x)-J\_RVM\textasciicircum{} M(x))x\textasciicircum{}-s-1 dx is holomorphic for Re s>0. The two functions have the same zeros in Re s>0; compare the modified Euler products of Grosswald and Schnitzer [ref].}
\expandafter\gdef\csname CGTip@eq:JZ-diff\endcsname{Eq. (6.18), p. 19: J\_Z,RVM(x)-J\_RVM\textasciicircum{}*(x)=O(loglog x),}
\expandafter\gdef\csname CGTip@eq:weighted-factorization\endcsname{Eq. (6.19), p. 19: zeta\_Z,RVM(s)=zeta\_RVM(s) e\textasciicircum{}h\_RVM(s),}
\expandafter\gdef\csname CGTip@eq:hRVM\endcsname{Eq. (6.20), p. 19: h\_RVM(s):=sint\_1\textasciicircum{}inf(J\_Z,RVM(x)-J\_RVM\textasciicircum{} M(x))x\textasciicircum{}-s-1 dx}
\expandafter\gdef\csname CGTip@lem:cell-cancel\endcsname{Lemma 6.6 (Exact first-harmonic cell cancellation) (p. 19): For every integer n, int\_n-1\textasciicircum{}n2cos(2pi x) dx=0 .}
\expandafter\gdef\csname CGTip@eq:cell-cancel\endcsname{Eq. (6.21), p. 19: int\_n-1\textasciicircum{}n2cos(2pi x) dx=0 .}
\expandafter\gdef\csname CGTip@hyp:cell-remainder\endcsname{Hypothesis 6.7 (Uniform cell remainder) (p. 19): c\_n\textasciicircum{}RVM-int\_n-1\textasciicircum{}n(1/(log x)+2cos(2pi x))dx=o (1/(log n)).}
\expandafter\gdef\csname CGTip@eq:cell-remainder\endcsname{Eq. (6.22), p. 19: c\_n\textasciicircum{}RVM-int\_n-1\textasciicircum{}n(1/(log x)+2cos(2pi x))dx=o (1/(log n)).}
\expandafter\gdef\csname CGTip@prop:eventual-positive\endcsname{Proposition 6.8 (Conditional eventual positivity) (p. 19): Hypothesis CGrefhyp:cell-remainder implies c\_n\textasciicircum{}RVM 1/(log n), hence 0<c\_n\textasciicircum{}RVM<1 for all large n.}
\expandafter\gdef\csname CGTip@eq:eventual-positive\endcsname{Eq. (6.23), p. 19: c\_n\textasciicircum{}RVM 1/(log n),}
\expandafter\gdef\csname CGTip@sec:mertens-transform\endcsname{Section 7: The critical RVM Mertens transform (p. 20)}
\expandafter\gdef\csname CGTip@eq:F-RVM\endcsname{Eq. (7.1), p. 20: F\_RVM(s)=hat M\_RVM\textasciicircum{}crit(s):=1zeta\_RVM(s)=G(s)P\_RVM\textasciicircum{}crit(s) .}
\expandafter\gdef\csname CGTip@eq:Mcrit-formal\endcsname{Eq. (7.2), p. 20: M\_RVM\textasciicircum{}crit(x)=1/(2pi i)int\_(c)F\_RVM(s)(x\textasciicircum{}s)/(s) ds ;}
\expandafter\gdef\csname CGTip@prop:Mcrit-poles\endcsname{Proposition 7.1 (Synthetic pole abscissa) (p. 20): F\_RVM is holomorphic in Re s>1/2. Its poles are the simple poles 1/2+- itau\_k (k>=1) and s=-2,-4,-6,... and s=1 is a simple zero of F\_RVM.}
\expandafter\gdef\csname CGTip@eq:Mcrit-poles\endcsname{Eq. (7.3), p. 20: 1/2+- itau\_k (k>=1) and s=-2,-4,-6,...}
\expandafter\gdef\csname CGTip@prop:Gaussian\endcsname{Proposition 7.2 (Gaussian square-root bound under vertical growth) (p. 20): Fix c>1 and 0<eps<c-1/2. Assume, uniformly for 1/2+eps<=sigma<= c, F\_RVM(sigma+it)=exp(o(t\textasciicircum{}2)) (|t|->inf). Define M\_RVM,G\textasciicircum{}crit(x):=1/(2pi i)int\_(c)e\textasciicircum{}s\textasciicircum{}2F\_RVM(s)(x\textasciicircum{}s)/(s) ds . Then M\_RVM,G\textasciicircum{}crit(x)=O\_eps(x\textasciicircum{}1/2+eps).}
\expandafter\gdef\csname CGTip@eq:subgaussian\endcsname{Eq. (7.4), p. 20: F\_RVM(sigma+it)=exp(o(t\textasciicircum{}2)) (|t|->inf).}
\expandafter\gdef\csname CGTip@eq:MG\endcsname{Eq. (7.5), p. 20: M\_RVM,G\textasciicircum{}crit(x):=1/(2pi i)int\_(c)e\textasciicircum{}s\textasciicircum{}2F\_RVM(s)(x\textasciicircum{}s)/(s) ds .}
\expandafter\gdef\csname CGTip@eq:MG-bound\endcsname{Eq. (7.6), p. 20: M\_RVM,G\textasciicircum{}crit(x)=O\_eps(x\textasciicircum{}1/2+eps).}
\expandafter\gdef\csname CGTip@sec:transition\endcsname{Section 8: The transition theorem and the unsmoothed synthetic Mertens bound (p. 20)}
\expandafter\gdef\csname CGTip@eq:chi\endcsname{Eq. (8.1), p. 20: chi(s):=2\textasciicircum{}spi\textasciicircum{}s-1sin ((pi s)/2)Gamma(1-s)=pi\textasciicircum{}s-1/2(Gamma((1-s)/2))/(Gamma(s/2)),}
\expandafter\gdef\csname CGTip@eq:transition-zw\endcsname{Eq. (8.2), p. 20: s=1/2+z, z=sigma'+it, tau\_*:=-iz=t-isigma', w(s):=exp-2pi i N\_RVM(tau\_*).}
\expandafter\gdef\csname CGTip@eq:chi-size\endcsname{Eq. (8.3), p. 20: |chi(sigma+it)|=((|t|)/(2pi))\textasciicircum{}1/2-sigma(1+O(|t|\textasciicircum{}-1)).}
\expandafter\gdef\csname CGTip@thm:transition\endcsname{Theorem 8.1 (Transition evaluation of the midpoint product) (p. 20): There is a constant varpi\_RVM>0 with the following properties. For every B>0 there are t\_0(B) and C\_B such that, for 0<sigma'<= B and t>= t\_0(B), zeta\_RVM(s)=varpi\_RVM(1+w(s))(1+E\_0(s)), |E\_0(s)|<=(C\_B)/t . For every 0<sigma'\_1<sigma'\_2<1/2 there are t\_0 and C such that, for sigma'\_1<=sigma'<=sig...}
\expandafter\gdef\csname CGTip@eq:transition-w\endcsname{Eq. (8.4), p. 20: zeta\_RVM(s)=varpi\_RVM(1+w(s))(1+E\_0(s)), |E\_0(s)|<=(C\_B)/t .}
\expandafter\gdef\csname CGTip@eq:transition-chi\endcsname{Eq. (8.5), p. 21: zeta\_RVM(s)=varpi\_RVM(1+chi(s))(1+E(s)), |E(s)|<=C/t .}
\expandafter\gdef\csname CGTip@subsec:transition-proof\endcsname{Section 8.1: Proof of Theorem thm:transition (p. 21)}
\expandafter\gdef\csname CGTip@eq:transition-product\endcsname{Eq. (8.6), p. 21: P\_RVM\textasciicircum{}crit(s)P\_RVM\textasciicircum{}crit(1/2)=prod\_k>=1(1+(z\textasciicircum{}2)/(tau\_k\textasciicircum{}2))=expLambda(z), Lambda(z):=sum\_k>=1f(tau\_k), f(tau):=Log(1+(z\textasciicircum{}2)/(tau\textasciicircum{}2)).}
\expandafter\gdef\csname CGTip@eq:fprime\endcsname{Eq. (8.7), p. 21: f'(tau)=-(2z\textasciicircum{}2)/(tau(tau\textasciicircum{}2+z\textasciicircum{}2))=-2tau+1/(tau-iz)+1/(tau+iz).}
\expandafter\gdef\csname CGTip@eq:sawtooth-count\endcsname{Eq. (8.8), p. 21: n\_RVM(tau)=floor(N\_RVM(tau)+1/2)=N\_RVM(tau)-psi(N\_RVM(tau)+1/2), n\_RVM(tau\_0)=0 .}
\expandafter\gdef\csname CGTip@eq:transition-split\endcsname{Eq. (8.9), p. 21: Lambda(z)=Lambda\_ sm(z)+Lambda\_ saw(z), Lambda\_ sm(z):=int\_tau\_0\textasciicircum{}inf f dN\_RVM+f(tau\_0)N\_RVM(tau\_0),; Lambda\_ saw(z):=int\_tau\_0\textasciicircum{}infpsi(N\_RVM(tau)+1/2)f'(tau) dtau .}
\expandafter\gdef\csname CGTip@lem:TR-smooth\endcsname{Lemma 8.2 (Smooth part) (p. 21): With log G(s):=Log(z\textasciicircum{}2-1/4)+LogGamma(s/2)-s/2logpi, Lambda\_ sm(z)-log G(s)=nu\_3+O(t\textasciicircum{}-1), nu\_3 in R independent of z .}
\expandafter\gdef\csname CGTip@eq:smooth-difference\endcsname{Eq. (8.10), p. 21: Lambda\_ sm(z)-log G(s)=nu\_3+O(t\textasciicircum{}-1), nu\_3 in R independent of z .}
\expandafter\gdef\csname CGTip@eq:smooth-J-ibp\endcsname{Eq. (8.11), p. 21: Lambda\_ sm(z)=-int\_tau\_0\textasciicircum{}infN\_RVM f' dtau=int\_tau\_0\textasciicircum{}infN\_RVM(tau)(2tau-(2tau)/(tau\textasciicircum{}2+z\textasciicircum{}2))dtau .}
\expandafter\gdef\csname CGTip@eq:J2\endcsname{Eq. (8.12), p. 21: 7/8Log(1+(z\textasciicircum{}2)/(tau\_0\textasciicircum{}2))=7/4Log z-7/4logtau\_0+O(t\textasciicircum{}-2).}
\expandafter\gdef\csname CGTip@eq:I-def\endcsname{Eq. (8.13), p. 21: I(z):=1\_0\textasciicircum{}inflog(tau)/(2pi e) (z\textasciicircum{}2)/(tau\textasciicircum{}2+z\textasciicircum{}2) dtau .}
\expandafter\gdef\csname CGTip@eq:J-asymp\endcsname{Eq. (8.14), p. 22: Lambda\_ sm(z)=z/2Logz/(2pi e)+7/4Log z+nu\_1+O(t\textasciicircum{}-2), nu\_1=-7/4logtau\_0-1\_0\textasciicircum{}tau\_0log2pi e dtau .}
\expandafter\gdef\csname CGTip@eq:stirling-loggamma\endcsname{Eq. (8.15), p. 22: LogGamma(s/2)=z/2Logz/2-1/4Log z+1/4log2-z/2+1/2log2pi+O(t\textasciicircum{}-1),}
\expandafter\gdef\csname CGTip@eq:G-asymp\endcsname{Eq. (8.16), p. 22: log G(s)=z/2Logz/(2pi e)+7/4Log z+nu\_2+O(t\textasciicircum{}-1), nu\_2=1/4log2+1/2log2pi-1/4logpi .}
\expandafter\gdef\csname CGTip@lem:TR-clock\endcsname{Lemma 8.3 (The clock as a holomorphic phase) (p. 22): N\_RVM(tau)=2pi(Log2pi-1)+7/8 is holomorphic on Retau>0, N\_RVM(conj tau)=conj N\_RVM(tau), and for Retau>=tau\_0, 2pi ImN\_RVM(tau)>=(log2) Imtau (Imtau>=0), 2pi ImN\_RVM(tau)<=(log2) Imtau (Imtau<=0).}
\expandafter\gdef\csname CGTip@eq:clock-halfplane-bound\endcsname{Eq. (8.17), p. 22: 2pi ImN\_RVM(tau)>=(log2) Imtau (Imtau>=0), 2pi ImN\_RVM(tau)<=(log2) Imtau (Imtau<=0).}
\expandafter\gdef\csname CGTip@eq:clock-im-identity\endcsname{Eq. (8.18), p. 22: 2\_RVM(tau)=Ylog(|tau|)/(2pi e)+X =Y[log(X)/(2pi e)+1/2log(1+v\textasciicircum{}2)+( v)/v]>= Y[log2+varphi(v)],}
\expandafter\gdef\csname CGTip@lem:TR-harmonic\endcsname{Lemma 8.4 (Abel-regularized harmonic sum) (p. 22): Lambda\_ saw(z)=Log(1+w(s))+nu\_4+O(t\textasciicircum{}-2), nu\_4 in R independent of z .}
\expandafter\gdef\csname CGTip@eq:Psi-asymp\endcsname{Eq. (8.19), p. 22: Lambda\_ saw(z)=Log(1+w(s))+nu\_4+O(t\textasciicircum{}-2), nu\_4 in R independent of z .}
\expandafter\gdef\csname CGTip@eq:psi-r\endcsname{Eq. (8.20), p. 22: psi\_r(y):=-sum\_m>=1r\textasciicircum{}m(sin2pi my)/(pi m)=1(1-re\textasciicircum{}2pi iy).}
\expandafter\gdef\csname CGTip@eq:Psi-r-limit\endcsname{Eq. (8.21), p. 22: Lambda\_ saw(z)=lim\_r1Lambda\_r(z), Lambda\_r(z):=int\_tau\_0\textasciicircum{}infpsi\_r(N\_RVM+1/2)f' dtau .}
\expandafter\gdef\csname CGTip@eq:Psi-r-series\endcsname{Eq. (8.22), p. 22: Lambda\_r(z)=-sum\_m>=1((-r)\textasciicircum{}m)/(2pi i m)(I\_m\textasciicircum{}+-I\_m\textasciicircum{}-), I\_m\textasciicircum{}+-:=int\_tau\_0\textasciicircum{}inf e\textasciicircum{}+-2pi imN\_RVM(tau)f'(tau) dtau .}
\expandafter\gdef\csname CGTip@eq:ray-integrals\endcsname{Eq. (8.23), p. 23: R\_m\textasciicircum{}+:=iint\_0\textasciicircum{}inf e\textasciicircum{}2pi imN\_RVM(tau\_0+iy)f'(tau\_0+iy) dy, R\_m\textasciicircum{}-:=-iint\_0\textasciicircum{}inf e\textasciicircum{}-2pi imN\_RVM(tau\_0-iy)f'(tau\_0-iy) dy .}
\expandafter\gdef\csname CGTip@eq:lower-contour\endcsname{Eq. (8.24), p. 23: R\_m\textasciicircum{}--I\_m\textasciicircum{}-=2pi i e\textasciicircum{}-2pi imN\_RVM(tau\_*)=2pi i w\textasciicircum{}m .}
\expandafter\gdef\csname CGTip@eq:Psi-r-closed\endcsname{Eq. (8.25), p. 23: Lambda\_r(z)=-sum\_m>=1((-rw)\textasciicircum{}m)/m+Phi\_r(z)=Log(1+rw)+Phi\_r(z), Phi\_r(z):=-sum\_m>=1((-r)\textasciicircum{}m)/(2pi i m)(R\_m\textasciicircum{}+-R\_m\textasciicircum{}-).}
\expandafter\gdef\csname CGTip@eq:ray-bound\endcsname{Eq. (8.26), p. 23: |R\textasciicircum{}+-,(z)\_m|<=4/(t\textasciicircum{}2)int\_0\textasciicircum{}inf2\textasciicircum{}-my(tau\_0+y) dy=4/(t\textasciicircum{}2)((tau\_0)/(mlog2)+1/((mlog2)\textasciicircum{}2)).}
\expandafter\gdef\csname CGTip@eq:assembly\endcsname{Eq. (8.27), p. 23: zeta\_RVM(s) =P\_RVM\textasciicircum{}crit(s)G(s)=P\_RVM\textasciicircum{}crit(1/2)exp(Lambda\_ sm(z)-log G(s)+Lambda\_ saw(z)); =P\_RVM\textasciicircum{}crit(1/2) e\textasciicircum{}nu\_3+nu\_4(1+w)(1+O(t\textasciicircum{}-1)),}
\expandafter\gdef\csname CGTip@eq:varpi-def\endcsname{Eq. (8.28), p. 23: varpi\_RVM:=P\_RVM\textasciicircum{}crit(1/2) e\textasciicircum{}nu\_3+nu\_4>0,}
\expandafter\gdef\csname CGTip@eq:theta-def\endcsname{Eq. (8.29), p. 23: vartheta(v):=1/(2i)(LogGamma(1/4+(iv)/2)-LogGamma(1/4-(iv)/2))-v/2logpi,}
\expandafter\gdef\csname CGTip@eq:clock-theta\endcsname{Eq. (8.30), p. 23: N\_RVM(v)=vartheta(v)pi+1+O(|v|\textasciicircum{}-1).}
\expandafter\gdef\csname CGTip@eq:w-chi\endcsname{Eq. (8.31), p. 23: w=e\textasciicircum{}-2ivartheta(tau\_*)-2pi i+O(1/t)=chi(s)(1+O(t\textasciicircum{}-1)).}
\expandafter\gdef\csname CGTip@lem:TR-boundary\endcsname{Lemma 8.5 (Extension to the critical line) (p. 24): The function E\_0=zeta\_RVM/(varpi\_RVM(1+w))-1 extends holomorphically to |sigma'|<1/2, t>t\_0, and |E\_0|<= C/t for 0<=sigma'<=1/4.}
\expandafter\gdef\csname CGTip@cor:derivative-zeros\endcsname{Corollary 8.6 (Derivative at the zeros) (p. 24): For k with tau\_k>t\_0, zeta\_RVM'(rho\_k\textasciicircum{}RVM)=varpi\_RVMlog(tau\_k)/(2pi)(1+O(1/(tau\_k))).}
\expandafter\gdef\csname CGTip@eq:derivative-zeros\endcsname{Eq. (8.32), p. 24: zeta\_RVM'(rho\_k\textasciicircum{}RVM)=varpi\_RVMlog(tau\_k)/(2pi)(1+O(1/(tau\_k))).}
\expandafter\gdef\csname CGTip@rem:TR-audit\endcsname{Remark 8.7 (Structure of the proof) (p. 24): The delicate points are handled as follows: the smooth part is computed with principal branches by the identity theorem (Lemma CGreflem:TR-smooth); the clock estimate is quantified (Lemma CGreflem:TR-clock); the sawtooth series is replaced by its Abel means, termwise integration is justified for ...}
\expandafter\gdef\csname CGTip@cor:vertical-limit\endcsname{Corollary 8.8 (Vertical constant limit) (p. 24): For every fixed sigma>1/2, zeta\_RVM(sigma+it)->varpi\_RVM, F\_RVM(sigma+it)-> L:=1varpi\_RVM (|t|->inf). More precisely, with sigma'=sigma-1/2, |w(sigma+it)|=O\_sigma(|t|\textasciicircum{}-sigma'), and, with alpha\_c:=min(c-1/2,1) for c>1/2, F\_RVM(c+it)-L=O\_c(|t|\textasciicircum{}-alpha\_c).}
\expandafter\gdef\csname CGTip@eq:vertical-limit\endcsname{Eq. (8.33), p. 24: zeta\_RVM(sigma+it)->varpi\_RVM, F\_RVM(sigma+it)-> L:=1varpi\_RVM (|t|->inf).}
\expandafter\gdef\csname CGTip@eq:w-decay\endcsname{Eq. (8.34), p. 24: |w(sigma+it)|=O\_sigma(|t|\textasciicircum{}-sigma'),}
\expandafter\gdef\csname CGTip@eq:FminusL-all\endcsname{Eq. (8.35), p. 24: F\_RVM(c+it)-L=O\_c(|t|\textasciicircum{}-alpha\_c).}
\expandafter\gdef\csname CGTip@subsec:halfplane\endcsname{Section 8.2: Bounds on the whole half-plane (p. 24)}
\expandafter\gdef\csname CGTip@eq:saw-primitive\endcsname{Eq. (8.36), p. 24: s\_1(tau):=int\_tau\_0\textasciicircum{}taupsi(N\_RVM(v)+1/2) dv (tau>=tau\_0).}
\expandafter\gdef\csname CGTip@lem:saw-primitive\endcsname{Lemma 8.9 (Bounded primitive of the midpoint sawtooth) (p. 24): |s\_1(tau)|<=(pi)/(4log2) for all tau>=tau\_0.}
\expandafter\gdef\csname CGTip@thm:halfplane-RVM\endcsname{Theorem 8.10 (Two-sided bound and limit at +inf) (p. 24): There is an absolute constant R\_0 such that, for s=1/2+z, z=sigma'+it, sigma'>1/2 and |z|>= R\_0, zeta\_RVM(s)=varpi\_RVM e\textasciicircum{}theta\_RVM(s), |theta\_RVM(s)|<=(pi\textasciicircum{}2)/(2log2)*1/(sigma')+(C)/(|z|). Consequently, for every sigma\_0>1 there is C(sigma\_0) with C(sigma\_0)\textasciicircum{}-1<=|zeta\_RVM(s)|<= C(sigma\_0) (Re s>=s...}
\expandafter\gdef\csname CGTip@eq:halfplane-RVM\endcsname{Eq. (8.37), p. 24: zeta\_RVM(s)=varpi\_RVM e\textasciicircum{}theta\_RVM(s), |theta\_RVM(s)|<=(pi\textasciicircum{}2)/(2log2)*1/(sigma')+(C)/(|z|).}
\expandafter\gdef\csname CGTip@eq:halfplane-two-sided\endcsname{Eq. (8.38), p. 25: C(sigma\_0)\textasciicircum{}-1<=|zeta\_RVM(s)|<= C(sigma\_0) (Re s>=sigma\_0),}
\expandafter\gdef\csname CGTip@eq:nu4-closed\endcsname{Eq. (8.39), p. 25: nu\_4=-2int\_tau\_0\textasciicircum{}inf(s\_1(tau))/(tau\textasciicircum{}2) dtau .}
\expandafter\gdef\csname CGTip@eq:smooth-exact\endcsname{Eq. (8.40), p. 25: Lambda\_ sm(z)=z/2Logz/(2pi e)+7/8Log(1+(z\textasciicircum{}2)/(tau\_0\textasciicircum{}2))-1\_0\textasciicircum{}tau\_0log2pi e (z\textasciicircum{}2)/(tau\textasciicircum{}2+z\textasciicircum{}2) dtau .}
\expandafter\gdef\csname CGTip@eq:saw-ibp\endcsname{Eq. (8.41), p. 25: Lambda\_ saw(z)=-int\_tau\_0\textasciicircum{}infs\_1f'' dtau, f''(tau)=2/(tau\textasciicircum{}2)-1/((tau-iz)\textasciicircum{}2)-1/((tau+iz)\textasciicircum{}2).}
\expandafter\gdef\csname CGTip@eq:saw-bound\endcsname{Eq. (8.42), p. 25: |Lambda\_ saw(z)+2int\_tau\_0\textasciicircum{}inf(s\_1(tau))/(tau\textasciicircum{}2) dtau|<=(pi)/(4log2)*(2pi)/(sigma') .}
\expandafter\gdef\csname CGTip@thm:no-Dirichlet-RVM\endcsname{Theorem 8.11 (No absolutely convergent Dirichlet expansion) (p. 25): Neither zeta\_RVM nor F\_RVM=1/zeta\_RVM is represented by a Dirichlet series sum\_na\_nn\textasciicircum{}-s converging absolutely in some half-plane Re s>sigma\_0.}
\expandafter\gdef\csname CGTip@eq:mean-square\endcsname{Eq. (8.43), p. 25: 1/(2T)int\_-T\textasciicircum{}T|F\_RVM(sigma+it)-L|\textasciicircum{}2dt =sum\_m,n a\_mconj a\_n(mn)\textasciicircum{}-sigma 1/(2T)int\_-T\textasciicircum{}T(n/m)\textasciicircum{}itdt ,}
\expandafter\gdef\csname CGTip@eq:mean-square-limit\endcsname{Eq. (8.44), p. 25: lim\_T->inf1/(2T)int\_-T\textasciicircum{}T|F\_RVM(sigma+it)-L|\textasciicircum{}2dt=sum\_n>=1| a\_n|\textasciicircum{}2n\textasciicircum{}-2sigma.}
\expandafter\gdef\csname CGTip@thm:unsmoothed-RVM\endcsname{Theorem 8.12 (Unsmoothed synthetic Mertens bound) (p. 26): For every x>1 and every c>1/2 the symmetric Perron integral M\_RVM\textasciicircum{}crit(x):=lim\_U->inf1/(2pi i)int\_c-iU\textasciicircum{}c+iUF\_RVM(s)(x\textasciicircum{}s)/(s) ds exists, is independent of c>1/2, and equals the absolutely convergent expression M\_RVM\textasciicircum{}crit(x)=L+1/(2pi i)int\_(c)(F\_RVM(s)-L)(x\textasciicircum{}s)/(s) ds . For every eps in (0,1/2), M\_R...}
\expandafter\gdef\csname CGTip@eq:Mcrit-symmetric\endcsname{Eq. (8.45), p. 26: M\_RVM\textasciicircum{}crit(x):=lim\_U->inf1/(2pi i)int\_c-iU\textasciicircum{}c+iUF\_RVM(s)(x\textasciicircum{}s)/(s) ds}
\expandafter\gdef\csname CGTip@eq:Mcrit-renorm\endcsname{Eq. (8.46), p. 26: M\_RVM\textasciicircum{}crit(x)=L+1/(2pi i)int\_(c)(F\_RVM(s)-L)(x\textasciicircum{}s)/(s) ds .}
\expandafter\gdef\csname CGTip@eq:Mcrit-unsmoothed-bound\endcsname{Eq. (8.47), p. 26: M\_RVM\textasciicircum{}crit(x)=L+O\_eps(x\textasciicircum{}1/2+eps).}
\expandafter\gdef\csname CGTip@eq:bromwich\endcsname{Eq. (8.48), p. 26: lim\_U->inf1/(2pi i)int\_c-iU\textasciicircum{}c+iU(x\textasciicircum{}s)/(s) ds=1 .}
\expandafter\gdef\csname CGTip@eq:Mcrit-final-estimate\endcsname{Eq. (8.49), p. 26: |M\_RVM\textasciicircum{}crit(x)-L|<=x\textasciicircum{}1/2+eps2piint\_-inf\textasciicircum{}inf|F\_RVM(1/2+eps+it)-L||1/2+eps+it| dt,}
\expandafter\gdef\csname CGTip@rem:unsmoothed-scope\endcsname{Remark 8.13 (Scope) (p. 26): Theorem CGrefthm:unsmoothed-RVM concerns 1/zeta\_RVM. The transfer to M goes through the endpoint of the lambda-family and is Theorem CGrefthm:step3.}
\expandafter\gdef\csname CGTip@sec:comparison\endcsname{Section 9: The exact Hadamard product and the comparison with zetazeta (p. 26)}
\expandafter\gdef\csname CGTip@eq:Pzeta\endcsname{Eq. (9.1), p. 26: P\_zeta(s):=G(s)zeta(s)=2xi(s), xi(s)=1/2s(s-1)pi\textasciicircum{}-s/2Gamma(s/2)zeta(s).}
\expandafter\gdef\csname CGTip@thm:Hadamard\endcsname{Theorem 9.1 (Exact Hadamard product) (p. 26): Let rho run over the nontrivial zeros of zeta with Imrho>0, with multiplicity, ordered by Imrho. The product P\_zeta(s)=prod\_Imrho>0(1+(q(s))/(rho(1-rho))) converges absolutely and locally uniformly in s in C, and the identity holds with no exponential or constant factor.}
\expandafter\gdef\csname CGTip@eq:Hadamard-exact\endcsname{Eq. (9.2), p. 26: P\_zeta(s)=prod\_Imrho>0(1+(q(s))/(rho(1-rho)))}
\expandafter\gdef\csname CGTip@eq:hadamard-general\endcsname{Eq. (9.3), p. 26: 2xi(s)=e\textasciicircum{}Bsprod\_rho(1- srho)e\textasciicircum{}s/rho,}
\expandafter\gdef\csname CGTip@eq:pair-factor\endcsname{Eq. (9.4), p. 26: (1- srho)(1-s/(1-rho))e\textasciicircum{}s/rho+s/(1-rho) =(1+(q(s))/(rho(1-rho)))e\textasciicircum{}s/(rho(1-rho)).}
\expandafter\gdef\csname CGTip@eq:hadamard-exp\endcsname{Eq. (9.5), p. 27: 2xi(s)=e\textasciicircum{}B's Pi\_zeta(s), Pi\_zeta(s):=prod\_Imrho>0(1+(q(s))/(rho(1-rho))), B'=B+sum\_Imrho>01/(rho(1-rho)) .}
\expandafter\gdef\csname CGTip@eq:Q-RVM\endcsname{Eq. (9.6), p. 27: Q\_RVM(s):=zeta(s)zeta\_RVM(s)=P\_zeta(s)P\_RVM\textasciicircum{}crit(s).}
\expandafter\gdef\csname CGTip@prop:M-transform-relation\endcsname{Proposition 9.2 (Exact transform relation) (p. 27): For Re s>1, and by meromorphic continuation, F\_RVM(s)=Q\_RVM(s) 1/(zeta(s)) . If the inverse Mellin transform Q\_RVM\textasciicircum{}v of Q\_RVM and the inverse defining M\_RVM\textasciicircum{}crit exist in a common distribution space with Mellin convolution, then M\_RVM\textasciicircum{}crit=Q\_RVM\textasciicircum{}v*Mel M .}
\expandafter\gdef\csname CGTip@eq:M-transform-relation\endcsname{Eq. (9.7), p. 27: F\_RVM(s)=Q\_RVM(s) 1/(zeta(s)) .}
\expandafter\gdef\csname CGTip@eq:M-convolution\endcsname{Eq. (9.8), p. 27: M\_RVM\textasciicircum{}crit=Q\_RVM\textasciicircum{}v*Mel M .}
\expandafter\gdef\csname CGTip@cor:Q-control\endcsname{Corollary 9.3 (All growth of zeta to the right of 1/2 lies in Q\_RVM) (p. 27): Let 0<eps<c'-1/2. There are t\_0 and 0<c\_-<= c\_+ such that c\_-<=|zeta\_RVM(sigma+it)|<= c\_+ (1/2+eps<=sigma<= c', |t|>= t\_0), and therefore, in the same region, c\_-|Q\_RVM(s)|<=|zeta(s)|<= c\_+|Q\_RVM(s)|, 1c\_+|Q\_RVM(s)|<=|1/(zeta(s))|<=1c\_-|Q\_RVM(s)|. Consequently: the Lindel"of hypothesis holds if a...}
\expandafter\gdef\csname CGTip@eq:zetaRVM-two-sided\endcsname{Eq. (9.9), p. 27: c\_-<=|zeta\_RVM(sigma+it)|<= c\_+ (1/2+eps<=sigma<= c', |t|>= t\_0),}
\expandafter\gdef\csname CGTip@eq:zeta-asymp-Q\endcsname{Eq. (9.10), p. 27: c\_-|Q\_RVM(s)|<=|zeta(s)|<= c\_+|Q\_RVM(s)|, 1c\_+|Q\_RVM(s)|<=|1/(zeta(s))|<=1c\_-|Q\_RVM(s)|.}
\expandafter\gdef\csname CGTip@eq:prime-defect\endcsname{Eq. (9.11), p. 27: d\_n\textasciicircum{}RVM:=1\_n prime-c\_n\textasciicircum{}RVM.}
\expandafter\gdef\csname CGTip@eq:prime-defect-cumulative\endcsname{Eq. (9.12), p. 27: sum\_2<= n<= Nd\_n\textasciicircum{}RVM=pi(N)-pi\_RVM(N).}
\expandafter\gdef\csname CGTip@sec:lambda-family\endcsname{Section 10: The lambdalambda-family from zeta\_RVMzeta\_RVM to zeta\_Z1zeta\_Z1 (p. 28)}
\expandafter\gdef\csname CGTip@eq:gamma-lambda\endcsname{Eq. (10.1), p. 28: gamma\_k(lambda):=(1-lambda)gamma\_k+lambdatau\_k=gamma\_k+lambdaeta\_k,}
\expandafter\gdef\csname CGTip@eq:P-lambda\endcsname{Eq. (10.2), p. 28: P\_lambda(s):=prod\_k>=1(1+(q(s))/(1/4+gamma\_k(lambda)\textasciicircum{}2)),}
\expandafter\gdef\csname CGTip@eq:zeta-lambda\endcsname{Eq. (10.3), p. 28: zeta\_lambda(s):=(P\_lambda(s))/(G(s)), F\_lambda(s):=1/(zeta\_lambda(s)),}
\expandafter\gdef\csname CGTip@eq:family-endpoints\endcsname{Eq. (10.4), p. 28: P\_Z1:=P\_0, zeta\_Z1:=zeta\_0, zeta\_1=zeta\_RVM.}
\expandafter\gdef\csname CGTip@thm:path\endcsname{Theorem 10.1 (The lambda-family) (p. 28): The product CGeqrefeq:P-lambda converges absolutely and locally uniformly jointly on [0,1]xC. For every lambda, P\_lambda is entire of order exactly one, P\_lambda(1-s)=P\_lambda(s), P\_lambda(conj s)=conj P\_lambda(s), P\_lambda(0)=P\_lambda(1)=1, and its zeros, with multiplicity, are exactly 1/2+- iga...}
\expandafter\gdef\csname CGTip@eq:family-divisor\endcsname{Eq. (10.5), p. 28: 1/2+- igamma\_k(lambda), k>=1 .}
\expandafter\gdef\csname CGTip@eq:lambda-normal\endcsname{Eq. (10.6), p. 28: sum\_k>=1sup\_s in K, 0<=lambda<=1|(q(s))/(1/4+gamma\_k(lambda)\textasciicircum{}2)|<inf,}
\expandafter\gdef\csname CGTip@eq:order-one\endcsname{Eq. (10.7), p. 28: logmax\_|s|<= R\_0|P\_lambda(s)|<=sum\_k>=1log(1+(R\_0\textasciicircum{}2+R\_0)/(1/4+gamma\_k(lambda)\textasciicircum{}2))=O(R\_0log R\_0),}
\expandafter\gdef\csname CGTip@eq:ordinate-gaps\endcsname{Eq. (10.8), p. 28: gamma\_k+1(lambda)-gamma\_k(lambda)=(1-lambda)(gamma\_k+1-gamma\_k)+lambda(tau\_k+1-tau\_k)>=lambda(tau\_k+1-tau\_k)>0,}
\expandafter\gdef\csname CGTip@rem:only-gamma-tau\endcsname{Remark 10.2 (All differences lie in gamma\_k-tau\_k) (p. 28): For lambda,lambda' in [0,1] and s off the divisors, zeta\_lambda(s)zeta\_lambda'(s)=prod\_k>=1((1/4+(gamma\_k+lambda'eta\_k)\textasciicircum{}2)(1/4+(gamma\_k+lambdaeta\_k)\textasciicircum{}2+q(s)))/((1/4+(gamma\_k+lambdaeta\_k)\textasciicircum{}2)(1/4+(gamma\_k+lambda'eta\_k)\textasciicircum{}2+q(s))) . The gamma factor cancels; the members differ only through eta\_k=tau\_k-...}
\expandafter\gdef\csname CGTip@eq:ratio-members\endcsname{Eq. (10.9), p. 28: zeta\_lambda(s)zeta\_lambda'(s)=prod\_k>=1((1/4+(gamma\_k+lambda'eta\_k)\textasciicircum{}2)(1/4+(gamma\_k+lambdaeta\_k)\textasciicircum{}2+q(s)))/((1/4+(gamma\_k+lambdaeta\_k)\textasciicircum{}2)(1/4+(gamma\_k+lambda'eta\_k)\textasciicircum{}2+q(s))) .}
\expandafter\gdef\csname CGTip@thm:generator\endcsname{Theorem 10.3 (Infinitesimal generator) (p. 28): Put A\_k,lambda:=1/4+gamma\_k(lambda)\textasciicircum{}2. Off the divisor of P\_lambda, V\_lambda(s):=d\_ P\_lambda(s)=-2q(s)sum\_k>=1gamma\_k(lambda) eta\_kA\_k,lambda(A\_k,lambda+q(s)), the series converging locally uniformly on compact sets separated from the divisor. Since G does not depend on lambda, d\_\_lambda=V\_lambda...}
\expandafter\gdef\csname CGTip@eq:generator\endcsname{Eq. (10.10), p. 28: V\_lambda(s):=d\_ P\_lambda(s)=-2q(s)sum\_k>=1gamma\_k(lambda) eta\_kA\_k,lambda(A\_k,lambda+q(s)),}
\expandafter\gdef\csname CGTip@eq:generator-zeta\endcsname{Eq. (10.11), p. 29: d\_\_lambda=V\_lambda, d\_ F\_lambda=-V\_lambda .}
\expandafter\gdef\csname CGTip@eq:mode-derivative\endcsname{Eq. (10.12), p. 29: d\_(1+ qA\_k,lambda)=-2q gamma\_k(lambda)eta\_kA\_k,lambda(A\_k,lambda+q) .}
\expandafter\gdef\csname CGTip@prop:integrated-multiplier\endcsname{Proposition 10.4 (Integrated multiplier) (p. 29): On every component avoiding the moving divisor, zeta\_lambda(s)=zeta\_Z1(s)exp(int\_0\textasciicircum{}\_lambda'(s) dlambda'), F\_lambda(s)=F\_Z1(s)exp(-int\_0\textasciicircum{}\_lambda'(s) dlambda'), and the total multiplier is Pi\_ tr(s):=exp(-int\_0\textasciicircum{}1V\_lambda'(s) dlambda')=P\_Z1(s)P\_RVM\textasciicircum{}crit(s), zeta\_Z1=zeta\_RVM Pi\_ tr . The identities e...}
\expandafter\gdef\csname CGTip@eq:integrated-zeta\endcsname{Eq. (10.13), p. 29: zeta\_lambda(s)=zeta\_Z1(s)exp(int\_0\textasciicircum{}\_lambda'(s) dlambda'),}
\expandafter\gdef\csname CGTip@eq:integrated-inverse\endcsname{Eq. (10.14), p. 29: F\_lambda(s)=F\_Z1(s)exp(-int\_0\textasciicircum{}\_lambda'(s) dlambda'),}
\expandafter\gdef\csname CGTip@eq:transport-multiplier\endcsname{Eq. (10.15), p. 29: Pi\_ tr(s):=exp(-int\_0\textasciicircum{}1V\_lambda'(s) dlambda')=P\_Z1(s)P\_RVM\textasciicircum{}crit(s), zeta\_Z1=zeta\_RVM Pi\_ tr .}
\expandafter\gdef\csname CGTip@eq:mode-integral\endcsname{Eq. (10.16), p. 29: int\_0\textasciicircum{}1d\_lambda'log(1+ qA\_k,lambda')dlambda'==log(1+q/(1/4+tau\_k\textasciicircum{}2))/(1+q/(1/4+gamma\_k\textasciicircum{}2))+-od2pi i.}
\expandafter\gdef\csname CGTip@cor:endpoint-expansion\endcsname{Corollary 10.5 (Expansion at the endpoint) (p. 29): Let K be compact and disjoint from the zeros of P\_Z1. Uniformly for s in K, as lambda0, zeta\_lambda(s)=zeta\_Z1(s)(1+\_0(s)+O\_K(lambda\textasciicircum{}2)), F\_lambda(s)=F\_Z1(s)(1-\_0(s)+O\_K(lambda\textasciicircum{}2)).}
\expandafter\gdef\csname CGTip@eq:endpoint-zeta-expansion\endcsname{Eq. (10.17), p. 29: zeta\_lambda(s)=zeta\_Z1(s)(1+\_0(s)+O\_K(lambda\textasciicircum{}2)),}
\expandafter\gdef\csname CGTip@eq:endpoint-inverse-expansion\endcsname{Eq. (10.18), p. 29: F\_lambda(s)=F\_Z1(s)(1-\_0(s)+O\_K(lambda\textasciicircum{}2)).}
\expandafter\gdef\csname CGTip@prop:residues\endcsname{Proposition 10.6 (Residues along the family) (p. 29): Suppose gamma\_k is a simple ordinate (it occurs once in the list). Put rho\_k,lambda=1/2+igamma\_k(lambda) and, for 0<lambda<=1, r\_k,lambda:=Res\_s=rho\_k,lambda1/(s zeta\_lambda(s))=G(rho\_k,lambda)rho\_k,lambdaP\_lambda'(rho\_k,lambda). Then r\_k,lambda->G(1/2+igamma\_k)(1/2+igamma\_k)P\_Z1'(1/2+igamma\_k) (...}
\expandafter\gdef\csname CGTip@eq:residue-transport\endcsname{Eq. (10.19), p. 29: r\_k,lambda:=Res\_s=rho\_k,lambda1/(s zeta\_lambda(s))=G(rho\_k,lambda)rho\_k,lambdaP\_lambda'(rho\_k,lambda).}
\expandafter\gdef\csname CGTip@eq:residue-limit\endcsname{Eq. (10.20), p. 29: r\_k,lambda->G(1/2+igamma\_k)(1/2+igamma\_k)P\_Z1'(1/2+igamma\_k) (lambda0).}
\expandafter\gdef\csname CGTip@rem:multiplicity\endcsname{Remark 10.7 (Multiplicity) (p. 29): For 0<lambda<=1 every zero of zeta\_lambda is simple. At lambda=0 a repeated ordinate of zeta appears in zeta\_Z1 with the same multiplicity, and the residues of Proposition CGrefprop:residues are replaced by the full principal part. Multiplicity plays no role in Step 3: Theorem CGrefthm:step3 uses...}
\expandafter\gdef\csname CGTip@sec:Z1-step3\endcsname{Section 11: The endpoint zeta\_Z1zeta\_Z1, Hamburger's theorem, and Step 3 (p. 30)}
\expandafter\gdef\csname CGTip@subsec:Z1-class\endcsname{Section 11.1: The analytic class of the endpoint (p. 30)}
\expandafter\gdef\csname CGTip@prop:Z1-class\endcsname{Proposition 11.1 (Analytic class of zeta\_Z1) (p. 30): The endpoint P\_Z1(s)=prod\_k>=1(1+(q(s))/(1/4+gamma\_k\textasciicircum{}2)), zeta\_Z1(s)=P\_Z1(s)G(s), has the following properties. P\_Z1 is entire of order exactly one, P\_Z1(1-s)=P\_Z1(s), P\_Z1(0)=P\_Z1(1)=1; the nontrivial zeros of zeta\_Z1 are exactly 1/2+- igamma\_k, k>=1, with the multiplicities of the ordinate list...}
\expandafter\gdef\csname CGTip@eq:Z1-product\endcsname{Eq. (11.1), p. 30: P\_Z1(s)=prod\_k>=1(1+(q(s))/(1/4+gamma\_k\textasciicircum{}2)), zeta\_Z1(s)=P\_Z1(s)G(s),}
\expandafter\gdef\csname CGTip@eq:Z1-divisor\endcsname{Eq. (11.2), p. 30: 1/2+- igamma\_k, k>=1,}
\expandafter\gdef\csname CGTip@eq:Z1-FE\endcsname{Eq. (11.3), p. 30: pi\textasciicircum{}-s/2Gamma(s/2)zeta\_Z1(s)=pi\textasciicircum{}-(1-s)/2Gamma((1-s)/2)zeta\_Z1(1-s);}
\expandafter\gdef\csname CGTip@eq:G-near-one\endcsname{Eq. (11.4), p. 30: G(s)=(s-1)+O((s-1)\textasciicircum{}2) (s->1).}
\expandafter\gdef\csname CGTip@eq:Z1-FE-proof\endcsname{Eq. (11.5), p. 30: pi\textasciicircum{}-s/2Gamma(s/2)zeta\_Z1(s)=P\_Z1(s)q(s),}
\expandafter\gdef\csname CGTip@eq:Z1-entire-removal\endcsname{Eq. (11.6), p. 30: (s-1)zeta\_Z1(s)=P\_Z1(s) pi\textasciicircum{}s/2s Gamma(s/2),}
\expandafter\gdef\csname CGTip@eq:Z1-defect-ratio\endcsname{Eq. (11.7), p. 30: R\_Z1(s):=zeta\_Z1(s)zeta(s)=P\_Z1(s)P\_zeta(s),}
\expandafter\gdef\csname CGTip@subsec:Hamburger\endcsname{Section 11.2: Hamburger's theorem at the endpoint (p. 30)}
\expandafter\gdef\csname CGTip@thm:Hamburger-Z1\endcsname{Theorem 11.2 (Hamburger closure) (p. 30): Consider the condition (D) zeta\_Z1(s)=sum\_n>=1(a\_n)/(n\textasciicircum{}s) absolutely for Re s>1 . Then (D) <=> R\_Z1==1 <=> zeta\_Z1=zeta <=> RH.}
\expandafter\gdef\csname CGTip@eq:Z1-D\endcsname{Eq. (11.8), p. 30: (D) zeta\_Z1(s)=sum\_n>=1(a\_n)/(n\textasciicircum{}s) absolutely for Re s>1 .}
\expandafter\gdef\csname CGTip@eq:Hamburger-equivalence\endcsname{Eq. (11.9), p. 30: (D) <=> R\_Z1==1 <=> zeta\_Z1=zeta <=> RH.}
\expandafter\gdef\csname CGTip@eq:Hamburger-conclusion\endcsname{Eq. (11.10), p. 31: zeta\_Z1=C\_H zeta}
\expandafter\gdef\csname CGTip@eq:RH-product\endcsname{Eq. (11.11), p. 31: P\_zeta(s)=prod\_k>=1(1+(q(s))/(1/4+gamma\_k\textasciicircum{}2))=P\_Z1(s)}
\expandafter\gdef\csname CGTip@subsec:D-multiplier\endcsname{Section 11.3: Condition (D) and the integer-weighted zeta (p. 31)}
\expandafter\gdef\csname CGTip@eq:K-multiplier\endcsname{Eq. (11.12), p. 31: K(s):=zeta\_Z1(s)zeta\_Z,RVM(s) .}
\expandafter\gdef\csname CGTip@thm:D-multiplier\endcsname{Theorem 11.3 (Reduction of (D) to the multiplier) (p. 31): If K has an expansion K(s)=sum\_n>=1k\_nn\textasciicircum{}-s that converges absolutely in Re s>1, then CGeqrefeq:Z1-D holds. If CGeqrefeq:Z1-D holds, then K(s)=sum\_n>=1k\_nn\textasciicircum{}-s absolutely on the whole half-plane Re s>1. Thus (D) holds if and only if K is an absolutely convergent Dirichlet series in Re s>1. Hence, b...}
\expandafter\gdef\csname CGTip@eq:K-explicit\endcsname{Eq. (11.13), p. 31: K(s)=Pi\_ tr(s) e\textasciicircum{}-h\_RVM(s)=P\_Z1(s)P\_RVM\textasciicircum{}crit(s) e\textasciicircum{}-h\_RVM(s) ,}
\expandafter\gdef\csname CGTip@rem:D-atomic\endcsname{Remark 11.4 (Continuous-to-atomic meaning of (D)) (p. 31): The prime-power count J\_Z1\textasciicircum{} M of zeta\_Z1 (Theorem CGrefthm:mellin-counts) is a continuous object built from the zeros, and (D) makes its measure the atomic measure sum\_p,mm\textasciicircum{}-1delta\_p\textasciicircum{}m. The final step of the method is therefore the passage of zeta\_Z1 into the discrete Dirichlet category. Theorem ...}
\expandafter\gdef\csname CGTip@subsec:step3\endcsname{Section 11.4: Step 3: the Mertens bound from zeta\_Z1=zetazeta\_Z1 = zeta (p. 31)}
\expandafter\gdef\csname CGTip@lem:zeta-growth\endcsname{Lemma 11.5 (Vertical growth of zeta) (p. 31): For 1/4<=sigma<=4 and |t|>=2, |zeta(sigma+it)|<=15|t| .}
\expandafter\gdef\csname CGTip@eq:zeta-growth\endcsname{Eq. (11.14), p. 31: |zeta(sigma+it)|<=15|t| .}
\expandafter\gdef\csname CGTip@eq:zeta-integral-rep\endcsname{Eq. (11.15), p. 31: zeta(s)=s/(s-1)-sint\_1\textasciicircum{}infyy\textasciicircum{}-s-1 dy .}
\expandafter\gdef\csname CGTip@lem:log-zeta\endcsname{Lemma 11.6 (Logarithm of zeta off the critical line) (p. 32): Assume zeta(s)!=0 for Re s>1/2. Let 0<delta<=1/4 and omega(delta):=(log(3-2delta))/(log(3-delta))<1 . Then, with the branch of logzeta equal to sum\_n>=2Lambda(n)n\textasciicircum{}-s/log n on Re s>1, |logzeta(sigma+it)|<=(48|t|)\textasciicircum{}omega(delta) (sigma>=1/2+delta, |t|>=4). Consequently, for every eps>0 there is t\_2=t...}
\expandafter\gdef\csname CGTip@eq:three-circles-exponent\endcsname{Eq. (11.16), p. 32: omega(delta):=(log(3-2delta))/(log(3-delta))<1 .}
\expandafter\gdef\csname CGTip@eq:log-zeta-bound\endcsname{Eq. (11.17), p. 32: |logzeta(sigma+it)|<=(48|t|)\textasciicircum{}omega(delta) (sigma>=1/2+delta, |t|>=4).}
\expandafter\gdef\csname CGTip@eq:inv-zeta-bound\endcsname{Eq. (11.18), p. 32: |1/(zeta(sigma+it))|<=|t|\textasciicircum{}eps (sigma>=1/2+delta, |t|>= t\_2).}
\expandafter\gdef\csname CGTip@eq:borel-caratheodory\endcsname{Eq. (11.19), p. 32: max\_|s-s\_0|<= r\_3|logzeta(s)|<=(2r\_3)/(R-r\_3)max\_|s-s\_0|=RRelogzeta(s)+(R+r\_3)/(R-r\_3)|logzeta(s\_0)|,}
\expandafter\gdef\csname CGTip@eq:BC-bound\endcsname{Eq. (11.20), p. 32: max\_|s-s\_0|<= r\_3|logzeta(s)|<=12delta(log(30|t|)+logzeta(2))<=48|t| .}
\expandafter\gdef\csname CGTip@eq:three-circles\endcsname{Eq. (11.21), p. 32: M\_2<= M\_1\textasciicircum{}1-omegaM\_3\textasciicircum{}omega, omega=(log(r\_2/r\_1))/(log(r\_3/r\_1))<=(log(3-2delta))/(log(3-delta))=omega(delta).}
\expandafter\gdef\csname CGTip@lem:perron-kernel\endcsname{Lemma 11.7 (Perron kernel) (p. 32): For y>0, y!=1, c>0, T>0 [ref], |1/(2pi i)int\_c-iT\textasciicircum{}c+iT(y\textasciicircum{}s)/(s) ds-1\_y>1|<= y\textasciicircum{}cmin(1,1/(T|log y|)).}
\expandafter\gdef\csname CGTip@eq:perron-kernel\endcsname{Eq. (11.22), p. 32: |1/(2pi i)int\_c-iT\textasciicircum{}c+iT(y\textasciicircum{}s)/(s) ds-1\_y>1|<= y\textasciicircum{}cmin(1,1/(T|log y|)).}
\expandafter\gdef\csname CGTip@thm:step3\endcsname{Theorem 11.8 (Step 3) (p. 32): If zeta\_Z1=zeta, then for every eps>0 M(x)=O\_eps(x\textasciicircum{}1/2+eps) (x>=1).}
\expandafter\gdef\csname CGTip@eq:step3-bound\endcsname{Eq. (11.23), p. 32: M(x)=O\_eps(x\textasciicircum{}1/2+eps) (x>=1).}
\expandafter\gdef\csname CGTip@eq:perron-truncated\endcsname{Eq. (11.24), p. 32: M(x)=1/(2pi i)int\_c-iT\textasciicircum{}c+iT(x\textasciicircum{}s)/(s zeta(s)) ds+O(sum\_n>=1(x/n)\textasciicircum{}cmin(1,1/(T|log(x/n)|))).}
\expandafter\gdef\csname CGTip@eq:perron-error\endcsname{Eq. (11.25), p. 33: M(x)=1/(2pi i)int\_c-iT\textasciicircum{}c+iT(x\textasciicircum{}s)/(s zeta(s)) ds+O(log x).}
\expandafter\gdef\csname CGTip@eq:contour-shift\endcsname{Eq. (11.26), p. 33: int\_c-iT\textasciicircum{}c+iT=int\_1/2+delta-iT\textasciicircum{}1/2+delta+iT+int\_1/2+delta+iT\textasciicircum{}c+iT-int\_1/2+delta-iT\textasciicircum{}c-iT.}
\expandafter\gdef\csname CGTip@eq:horizontal-bound\endcsname{Eq. (11.27), p. 33: 1/(2pi)int\_1/2+delta\textasciicircum{}cx\textasciicircum{}sigma T\textasciicircum{}epsT dsigma<= x\textasciicircum{}cT\textasciicircum{}eps-1=e x\textasciicircum{}eps.}
\expandafter\gdef\csname CGTip@eq:left-bound\endcsname{Eq. (11.28), p. 33: 1/(2pi)int\_-T\textasciicircum{}Tx\textasciicircum{}1/2+delta|1/2+delta+it| |zeta(1/2+delta+it)| dt<= x\textasciicircum{}1/2+delta(C'(delta,eps)+int\_t\_2\textasciicircum{}Tt\textasciicircum{}eps-1dt)<<\_delta,epsx\textasciicircum{}1/2+delta+eps.}
\expandafter\gdef\csname CGTip@thm:transfer\endcsname{Theorem 11.9 (Transfer of the synthetic Mertens bound to M) (p. 33): Assume zeta\_Z1=zeta and let 0<delta<=1/4, eps>0. Then: on Re s>1/2, 1/(zeta(s))=F\_RVM(s) Pi\_ tr(s)\textasciicircum{}-1, Pi\_ tr=P\_Z1P\_RVM\textasciicircum{}crit=zetazeta\_RVM, where Pi\_ tr is holomorphic and zero-free in Re s>1/2 and Pi\_ tr(1)=1; there is t\_3=t\_3(delta,eps) such that, for 1/2+delta<=sigma<=2 and |t|>= t\_3, |F\_RVM(si...}
\expandafter\gdef\csname CGTip@eq:transfer-factorization\endcsname{Eq. (11.29), p. 33: 1/(zeta(s))=F\_RVM(s) Pi\_ tr(s)\textasciicircum{}-1, Pi\_ tr=P\_Z1P\_RVM\textasciicircum{}crit=zetazeta\_RVM,}
\expandafter\gdef\csname CGTip@eq:transfer-multiplier-bound\endcsname{Eq. (11.30), p. 33: |F\_RVM(sigma+it)|<=2L, |Pi\_ tr(sigma+it)\textasciicircum{}+-1|<=|t|\textasciicircum{}eps;}
\expandafter\gdef\csname CGTip@eq:transfer-conclusion\endcsname{Eq. (11.31), p. 33: M(x)=O\_eps(x\textasciicircum{}1/2+eps).}
\expandafter\gdef\csname CGTip@eq:log-zetaRVM-bounded\endcsname{Eq. (11.32), p. 33: logzeta\_RVM=logvarpi\_RVM+Log(1+w)+Log(1+E\_0)=O(1).}
\expandafter\gdef\csname CGTip@eq:log-Pi-tr\endcsname{Eq. (11.33), p. 33: |logPi\_ tr(sigma+it)|=|logzeta-logzeta\_RVM|<=(48|t|)\textasciicircum{}omega(delta)+O(1)|t|}
\expandafter\gdef\csname CGTip@rem:transfer-mellin\endcsname{Remark 11.10 (The transfer in Mellin form) (p. 34): Formally, CGeqrefeq:transfer-factorization reads M=(Pi\_ tr\textasciicircum{}-1)\textasciicircum{}v*Mel M\_RVM\textasciicircum{}crit: the Mertens function is the synthetic Mertens function convolved with the inverse Mellin transform of the inverse transport multiplier. The synthetic factor F\_RVM contributes the exponent 1/2 and is bounded (Theorem ...}
\expandafter\gdef\csname CGTip@cor:step3-D\endcsname{Corollary 11.11 (Step 3 through (D)) (p. 34): Condition CGeqrefeq:Z1-D implies M(x)=O\_eps(x\textasciicircum{}1/2+eps) for every eps>0.}
\expandafter\gdef\csname CGTip@rem:Step3-status\endcsname{Remark 11.12 (Status of Step 3) (p. 34): The implication zeta\_Z1=zeta M(x)=O\_eps(x\textasciicircum{}1/2+eps) is proved completely above, and it requires nothing beyond the identity: no separate statement about the zeros of zeta enters, because zeta=zeta\_Z1 already carries the divisor CGeqrefeq:Z1-divisor. The identity zeta\_Z1=zeta itself, equivalently (...}
\expandafter\gdef\csname CGTip@subsec:reciprocal\endcsname{Section 11.5: Reciprocal endpoint and finite-height Perron integrals (p. 34)}
\expandafter\gdef\csname CGTip@eq:FZ1\endcsname{Eq. (11.34), p. 34: F\_Z1(s):=1zeta\_Z1(s)=F\_RVM(s) Pi\_ tr(s)\textasciicircum{}-1,}
\expandafter\gdef\csname CGTip@eq:truncated-Perron\endcsname{Eq. (11.35), p. 34: M\_lambda,U(x):=1/(2pi i)int\_c-iU\textasciicircum{}c+iUF\_lambda(s)(x\textasciicircum{}s)/(s) ds, M\_Z1,U:=M\_0,U.}
\expandafter\gdef\csname CGTip@prop:finite-Perron\endcsname{Proposition 11.13 (Finite-height Perron integrals along the family) (p. 34): For fixed c>1, U<inf and x>0, M\_lambda,U(x)-> M\_Z1,U(x) (lambda0).}
\expandafter\gdef\csname CGTip@eq:finite-Perron-limit\endcsname{Eq. (11.36), p. 34: M\_lambda,U(x)-> M\_Z1,U(x) (lambda0).}
\expandafter\gdef\csname CGTip@thm:coefficient-criterion\endcsname{Theorem 11.14 (Reciprocal-coefficient criterion) (p. 34): Assume that in Re s>1 F\_Z1(s)=sum\_n>=1mu\_Z1(n)n\textasciicircum{}s absolutely, and let M\_Z1(x)=sum\_n<= xmu\_Z1(n), so that mu\_Z1(n)=M\_Z1(n)-M\_Z1(n-1). Then the following are equivalent: (i) M\_Z1(n)=M(n) for every n>=1; (ii) mu\_Z1=mu; (iii) F\_Z1=1/zeta in Re s>1; (iv) zeta\_Z1=zeta; (v) RH.}
\expandafter\gdef\csname CGTip@eq:Z1-reciprocal-Dirichlet\endcsname{Eq. (11.37), p. 34: F\_Z1(s)=sum\_n>=1mu\_Z1(n)n\textasciicircum{}s}
\expandafter\gdef\csname CGTip@eq:Z1-increments\endcsname{Eq. (11.38), p. 34: mu\_Z1(n)=M\_Z1(n)-M\_Z1(n-1).}
\expandafter\gdef\csname CGTip@rem:negative-convolution\endcsname{Remark 11.15 (A route that does not close Step 3) (p. 35): The following approach does not work. One may try to expand both factors of F\_Z1=F\_RVM Pi\_ tr\textasciicircum{}-1 in CGeqrefeq:FZ1 as absolutely convergent Dirichlet series and identify their Dirichlet convolution with mu. By Theorem CGrefthm:no-Dirichlet-RVM the expansion of F\_RVM does not exist in any half-plan...}
\expandafter\gdef\csname CGTip@sec:euler-products\endcsname{Section 12: The endpoint at infinity and the Euler products (p. 35)}
\expandafter\gdef\csname CGTip@subsec:value-infinity\endcsname{Section 12.1: The value of the endpoint at infinity (p. 35)}
\expandafter\gdef\csname CGTip@eq:frak-e\endcsname{Eq. (12.1), p. 35: e(tau):=N\_zeta(tau)-N\_RVM(tau), e\_1(tau):=int\_tau\_0\textasciicircum{}taue(v) dv (tau>=tau\_0=4pi).}
\expandafter\gdef\csname CGTip@eq:frak-e1-bound\endcsname{Eq. (12.2), p. 35: |e\_1(tau)|<= C\_elogtau (tau>=tau\_0).}
\expandafter\gdef\csname CGTip@thm:Z1-clock\endcsname{Theorem 12.1 (Clock decomposition of the endpoint) (p. 35): Let s=1/2+z, z=sigma'+it, sigma'>0. Then logzeta\_Z1(s)=logvarpi\_Z1-int\_tau\_0\textasciicircum{}infe\_1(tau)(1/((tau-iz)\textasciicircum{}2)+1/((tau+iz)\textasciicircum{}2))dtau+O(1/(|z|)) (|z|>= R\_0), where varpi\_Z1:=P\_Z1(1/2)exp(nu\_3+2int\_tau\_0\textasciicircum{}inf(e\_1(tau))/(tau\textasciicircum{}2) dtau)>0 . Consequently: for sigma'>0 and |z|>= R\_0, |logzeta\_Z1(s)-logvarpi\_Z1|<=2...}
\expandafter\gdef\csname CGTip@eq:Z1-clock\endcsname{Eq. (12.3), p. 35: logzeta\_Z1(s)=logvarpi\_Z1-int\_tau\_0\textasciicircum{}infe\_1(tau)(1/((tau-iz)\textasciicircum{}2)+1/((tau+iz)\textasciicircum{}2))dtau+O(1/(|z|)) (|z|>= R\_0),}
\expandafter\gdef\csname CGTip@eq:varpi-Z1\endcsname{Eq. (12.4), p. 35: varpi\_Z1:=P\_Z1(1/2)exp(nu\_3+2int\_tau\_0\textasciicircum{}inf(e\_1(tau))/(tau\textasciicircum{}2) dtau)>0 .}
\expandafter\gdef\csname CGTip@eq:Z1-growth\endcsname{Eq. (12.5), p. 35: |logzeta\_Z1(s)-logvarpi\_Z1|<=2(pi+2)C\_esigma'(1+log(2|t|+2sigma'+tau\_0))+C/(|z|);}
\expandafter\gdef\csname CGTip@eq:varpi-Z1-product\endcsname{Eq. (12.6), p. 35: varpi\_Z1=varpi\_RVMprod\_k>=1(1/4+tau\_k\textasciicircum{}2)/(1/4+gamma\_k\textasciicircum{}2):=varpi\_RVMlim\_K->infprod\_k<= K(1/4+tau\_k\textasciicircum{}2)/(1/4+gamma\_k\textasciicircum{}2).}
\expandafter\gdef\csname CGTip@eq:Z1-split\endcsname{Eq. (12.7), p. 35: sum\_kf(gamma\_k)=Lambda\_ sm(z)+int\_tau\_0\textasciicircum{}infe\_1f'' dtau,}
\expandafter\gdef\csname CGTip@eq:log-kernel\endcsname{Eq. (12.8), p. 35: int\_tau\_0\textasciicircum{}inf(logtau)/((tau-a)\textasciicircum{}2+sigma'\textasciicircum{}2) dtau<=sigma'log X+4 (1+log X)/X<=(pi+2)/(sigma')(1+log X).}
\expandafter\gdef\csname CGTip@eq:varpi-ratio\endcsname{Eq. (12.9), p. 36: varpi\_Z1varpi\_RVM=P\_Z1(1/2)P\_RVM\textasciicircum{}crit(1/2)exp(2int\_tau\_0\textasciicircum{}inf(e\_1(tau)+s\_1(tau))/(tau\textasciicircum{}2) dtau).}
\expandafter\gdef\csname CGTip@eq:count-difference-integral\endcsname{Eq. (12.10), p. 36: int\_tau\_0\textasciicircum{}inf(e\_1+s\_1)/(tau\textasciicircum{}2) dtau=lim\_X->infint\_tau\_0\textasciicircum{}XN\_zeta(tau)-n\_RVM(tau)tau dtau=lim\_K->infsum\_k<= Klog(tau\_k)/(gamma\_k).}
\expandafter\gdef\csname CGTip@thm:varpi-Z1\endcsname{Theorem 12.2 (The value at infinity decides the endpoint identity) (p. 36): varpi\_Z1>=1, and varpi\_Z1=1 if and only if zeta\_Z1=zeta.}
\expandafter\gdef\csname CGTip@eq:pair-ratio\endcsname{Eq. (12.11), p. 36: (W\textasciicircum{}2)/((W-delta\textasciicircum{}2)\textasciicircum{}2+4delta\textasciicircum{}2gamma\textasciicircum{}2)* c\_gamma, c\_gamma=(|rho(1-rho)|\textasciicircum{}2)/((1/4+gamma\textasciicircum{}2)\textasciicircum{}2)=1+(2delta\textasciicircum{}2(gamma\textasciicircum{}2-1/4)+delta\textasciicircum{}4)/((1/4+gamma\textasciicircum{}2)\textasciicircum{}2)>1 .}
\expandafter\gdef\csname CGTip@cor:product-identity\endcsname{Corollary 12.3 (The endpoint identity as an identity of the displacements) (p. 36): The following are equivalent: (D); zeta\_Z1=zeta; varpi\_Z1=1; prod\_k>=1(1/4+gamma\_k\textasciicircum{}2)/(1/4+tau\_k\textasciicircum{}2)=varpi\_RVM (ordered product); and lim\_sigma->+infPi\_ tr(sigma)=L=1/varpi\_RVM.}
\expandafter\gdef\csname CGTip@eq:product-identity\endcsname{Eq. (12.12), p. 36: prod\_k>=1(1/4+gamma\_k\textasciicircum{}2)/(1/4+tau\_k\textasciicircum{}2)=varpi\_RVM}
\expandafter\gdef\csname CGTip@prop:Z1-bounded\endcsname{Proposition 12.4 (Boundedness of the endpoint) (p. 36): Let sigma\_0>1. Then |zeta\_Z1| is bounded above and below on Re s>=sigma\_0 if and only if sup\_sigma'>=sigma\_0-1/2, t>=0|Reint\_tau\_0\textasciicircum{}inf(e\_1(tau)+s\_1(tau))/((tau-t+isigma')\textasciicircum{}2) dtau|<inf .}
\expandafter\gdef\csname CGTip@eq:conjugate-poisson\endcsname{Eq. (12.13), p. 36: sup\_sigma'>=sigma\_0-1/2, t>=0|Reint\_tau\_0\textasciicircum{}inf(e\_1(tau)+s\_1(tau))/((tau-t+isigma')\textasciicircum{}2) dtau|<inf .}
\expandafter\gdef\csname CGTip@subsec:euler-products\endcsname{Section 12.2: The Euler products of zeta\_RVMzeta\_RVM and zeta\_Z1zeta\_Z1 (p. 37)}
\expandafter\gdef\csname CGTip@thm:mellin-counts\endcsname{Theorem 12.5 (Prime-power counts as Mellin transforms) (p. 37): Let X in RVM,Z1. There is a unique function J\_X\textasciicircum{} M on (1,inf) with J\_X\textasciicircum{} M(x)x\textasciicircum{}-c in L\textasciicircum{}2((1,inf),dx/x) for every c>1 and logzeta\_X(s)=sint\_1\textasciicircum{}inf J\_X\textasciicircum{} M(x) x\textasciicircum{}-s-1 dx (Re s>1), the branch of logzeta\_X being real on (1,inf). Equivalently zeta\_X(s)=exp(int\_[1,inf)x\textasciicircum{}-s dJ\_X\textasciicircum{} M(x)).}
\expandafter\gdef\csname CGTip@eq:mellin-count\endcsname{Eq. (12.14), p. 37: logzeta\_X(s)=sint\_1\textasciicircum{}inf J\_X\textasciicircum{} M(x) x\textasciicircum{}-s-1 dx (Re s>1),}
\expandafter\gdef\csname CGTip@thm:distributional-EF\endcsname{Theorem 12.6 (Distributional explicit formula) (p. 37): Let X in RVM,Z1, with ordinates kappa\_k=tau\_k for X=RVM and kappa\_k=gamma\_k for X=Z1. In D'(1,inf), (log x) dJ\_X\textasciicircum{} Mdx=1-2/( x) Resum\_k>=1x\textasciicircum{}ikappa\_k-1/(x(x\textasciicircum{}2-1)), the series converging in D'(1,inf). The explicit-formula count J\_X(x):=li(x)-2Resum\_k>=1Ei((1/2+ikappa\_k)log x)-log2+int\_x\textasciicircum{}inf(dy)/(y(y...}
\expandafter\gdef\csname CGTip@eq:distributional-EF\endcsname{Eq. (12.15), p. 37: (log x) dJ\_X\textasciicircum{} Mdx=1-2/( x) Resum\_k>=1x\textasciicircum{}ikappa\_k-1/(x(x\textasciicircum{}2-1)),}
\expandafter\gdef\csname CGTip@eq:J-X\endcsname{Eq. (12.16), p. 37: J\_X(x):=li(x)-2Resum\_k>=1Ei((1/2+ikappa\_k)log x)-log2+int\_x\textasciicircum{}inf(dy)/(y(y\textasciicircum{}2-1)log y)}
\expandafter\gdef\csname CGTip@eq:log-derivative-X\endcsname{Eq. (12.17), p. 37: -(zeta\_X')/(zeta\_X)(s)=1/(s-1)-sum\_k>=1(1/(s-rho\_k)+1/(s-conj rho\_k))+sum\_n>=1(1/(2n)-1/(s+2n))-(gamma\_E)/2-1/2logpi .}
\expandafter\gdef\csname CGTip@thm:cell-form\endcsname{Theorem 12.7 (Cell form of the endpoint identity) (p. 37): The following are equivalent. (D), that is, zeta\_Z1=zeta. J\_Z1\textasciicircum{} M is constant on every interval (n,n+1), n>=1, and its jumps b\_n:=v\_n-v\_n-1, where v\_n is its value on (n,n+1) and v\_0:=0, satisfy sum\_n<= X| b\_n|=O\_eps(X\textasciicircum{}1+eps) for every eps>0. For every n>=1, in D'(n,n+1), 2 Resum\_k>=1(x\textasciicircum{}igamma\_k-...}
\expandafter\gdef\csname CGTip@eq:cell-identity\endcsname{Eq. (12.18), p. 38: 2 Resum\_k>=1(x\textasciicircum{}igamma\_k-x\textasciicircum{}itau\_k)= x log x J\_RVM'(x),}
\expandafter\gdef\csname CGTip@eq:cell-difference\endcsname{Eq. (12.19), p. 38: (log x)(dJ\_Z1\textasciicircum{} Mdx-J\_RVM'(x))=-2/( x) Resum\_k>=1(x\textasciicircum{}igamma\_k-x\textasciicircum{}itau\_k),}
\expandafter\gdef\csname CGTip@rem:cell-form\endcsname{Remark 12.8 (The final step in cell form) (p. 38): Theorem CGrefthm:cell-form is the exact form of the continuous-to-atomic passage of Remark CGrefrem:D-atomic. Both sides of CGeqrefeq:cell-identity are built from the midpoint construction: the right side from the clock alone (J\_RVM' is bounded on [2,inf) by Theorem CGrefthm:unit-cell, equal to 1...}
\expandafter\gdef\csname CGTip@sec:atomized\endcsname{Section 13: The atomized family and the programme for the final step (p. 38)}
\expandafter\gdef\csname CGTip@subsec:family-infinity\endcsname{Section 13.1: The family at infinity (p. 38)}
\expandafter\gdef\csname CGTip@prop:family-clock\endcsname{Proposition 13.1 (Clock decomposition along the family) (p. 38): There is an absolute constant C with |e\_1\textasciicircum{}(lambda)(tau)|<= Clogtau for tau>=tau\_0 and all lambda in [0,1]. Consequently Theorem CGrefthm:Z1-clock holds for every zeta\_lambda with e\_1 replaced by e\_1\textasciicircum{}(lambda), uniformly in lambda; in particular varpi\_lambda:=lim\_sigma->+infzeta\_lambda(sigma) exist...}
\expandafter\gdef\csname CGTip@eq:varpi-lambda\endcsname{Eq. (13.1), p. 38: varpi\_lambda=varpi\_RVMprod\_k>=1(1/4+tau\_k\textasciicircum{}2)/(1/4+gamma\_k(lambda)\textasciicircum{}2), varpi\_1=varpi\_RVM, varpi\_0=varpi\_Z1.}
\expandafter\gdef\csname CGTip@eq:family-count-primitive\endcsname{Eq. (13.2), p. 39: int\_tau\_0\textasciicircum{}X(N\_lambda-n\_RVM)=sum\_k((X-gamma\_k(lambda))\textasciicircum{}+-(X-tau\_k)\textasciicircum{}+)=(1-lambda)sum\_gamma\_k(lambda),tau\_k<= Xeta\_k+O(log X),}
\expandafter\gdef\csname CGTip@eq:Upsilon\endcsname{Eq. (13.3), p. 39: (lambda):=sum\_k>=1(2gamma\_k(lambda) eta\_k)/(1/4+gamma\_k(lambda)\textasciicircum{}2) (summed in the order of k).}
\expandafter\gdef\csname CGTip@prop:family-flow\endcsname{Proposition 13.2 (The value at infinity along the family) (p. 39): The series CGeqrefeq:Upsilon converges uniformly on [0,1], lambda\_lambda is C\textasciicircum{}1 on [0,1], and (d)/(dlambda)logvarpi\_lambda=-(lambda)=lim\_sigma->+infV\_lambda(sigma), '(lambda)=sum\_k>=1(2eta\_k\textasciicircum{}2(1/4-gamma\_k(lambda)\textasciicircum{}2))/((1/4+gamma\_k(lambda)\textasciicircum{}2)\textasciicircum{}2)<0 , with V\_lambda the generator of Theorem CGrefthm:...}
\expandafter\gdef\csname CGTip@eq:varpi-derivative\endcsname{Eq. (13.4), p. 39: (d)/(dlambda)logvarpi\_lambda=-(lambda)=lim\_sigma->+infV\_lambda(sigma), '(lambda)=sum\_k>=1(2eta\_k\textasciicircum{}2(1/4-gamma\_k(lambda)\textasciicircum{}2))/((1/4+gamma\_k(lambda)\textasciicircum{}2)\textasciicircum{}2)<0 ,}
\expandafter\gdef\csname CGTip@eq:flow-integral\endcsname{Eq. (13.5), p. 39: int\_0\textasciicircum{}+\textasciicircum{}1(lambda) dlambda=sum\_k>=1log(1/4+tau\_k\textasciicircum{}2)/(1/4+gamma\_k\textasciicircum{}2)=log L+logvarpi\_Z1,}
\expandafter\gdef\csname CGTip@prop:varpi-verified\endcsname{Proposition 13.3 (The value at infinity and the verified zeros) (p. 39): If every zero beta+igamma with 0<gamma<= T\_0 has beta=1/2, then 0<=logvarpi\_Z1<=0.26sum\_gamma>T\_01/(gamma\textasciicircum{}2). With T\_0=3*10\textasciicircum{}12 [ref], the right side is (0.26)/(2pi T\_0)(log(T\_0)/(2pi)+1)<4*10\textasciicircum{}-13.}
\expandafter\gdef\csname CGTip@eq:varpi-verified\endcsname{Eq. (13.6), p. 39: 0<=logvarpi\_Z1<=0.26sum\_gamma>T\_01/(gamma\textasciicircum{}2).}
\expandafter\gdef\csname CGTip@subsec:smooth-cells\endcsname{Section 13.2: Atomization by smooth cells (p. 39)}
\expandafter\gdef\csname CGTip@def:atomized\endcsname{Definition 13.4 (Atomized family) (p. 40): The cell masses, the atomized zeta function and its coefficients are c\_n(lambda):= dpi\_lambda,varphi\_n, zeta\_Z,lambda(s):=prod\_n>=2(1-n\textasciicircum{}-s)\textasciicircum{}- c\_n(lambda)=sum\_N>=1( a\_N(lambda))/(N\textasciicircum{}s), a\_N(lambda)=sum\_prod\_nn\textasciicircum{}r\_n=N prod\_n: r\_n>=1 c\_n(lambda)+r\_n-1r\_n, a\_1(lambda)=1 .}
\expandafter\gdef\csname CGTip@eq:atomized-family\endcsname{Eq. (13.7), p. 40: c\_n(lambda):= dpi\_lambda,varphi\_n, zeta\_Z,lambda(s):=prod\_n>=2(1-n\textasciicircum{}-s)\textasciicircum{}- c\_n(lambda)=sum\_N>=1( a\_N(lambda))/(N\textasciicircum{}s),}
\expandafter\gdef\csname CGTip@eq:atomized-coefficients\endcsname{Eq. (13.8), p. 40: a\_N(lambda)=sum\_prod\_nn\textasciicircum{}r\_n=N prod\_n: r\_n>=1 c\_n(lambda)+r\_n-1r\_n, a\_1(lambda)=1 .}
\expandafter\gdef\csname CGTip@eq:J-lambda-derivative\endcsname{Eq. (13.9), p. 40: (d J\_lambda)/(dlambda)(x)=2sum\_k>=1eta\_k Imx\textasciicircum{}rho\_k(lambda)rho\_k(lambda), (d)/(dlambda)((log x)(dJ\_lambda)/(dx))=(2log x)/( x)sum\_k>=1eta\_ksin(gamma\_k(lambda)log x).}
\expandafter\gdef\csname CGTip@subsec:joining-factor\endcsname{Section 13.3: The joining factor cell by cell (p. 40)}
\expandafter\gdef\csname CGTip@eq:atomization-kernel\endcsname{Eq. (13.10), p. 40: a\_s(x):=sum\_n>=2varphi\_n(x) n\textasciicircum{}-s-x\textasciicircum{}-s, h\_n(lambda;s):= dpi\_lambda, varphi\_n(x)(n\textasciicircum{}-s-x\textasciicircum{}-s) .}
\expandafter\gdef\csname CGTip@eq:first-cell\endcsname{Eq. (13.11), p. 40: u\_lambda(s):=C\_lambda 2\textasciicircum{}-s+sint\_1\textasciicircum{}2J\textasciicircum{} M\_lambda(x) x\textasciicircum{}-s-1 dx, C\_lambda:=J\textasciicircum{} M\_lambda-J\_lambda ,}
\expandafter\gdef\csname CGTip@prop:first-cell\endcsname{Proposition 13.5 (The first cell) (p. 40): Let lambda in [0,1]. u\_lambda is entire, and for Re s>1 logzeta\_lambda(s)=u\_lambda(s)+ dJ\textasciicircum{}*\_lambda,x\textasciicircum{}-s, dJ\textasciicircum{}*\_lambda,x\textasciicircum{}-s:=sint\_2\textasciicircum{}inf J\_lambda(x) x\textasciicircum{}-s-1 dx . u\_lambda(sigma+it)=o(|t|) as |t|->inf, uniformly for sigma in compact sets. u\_lambda is constant if and only if J\_lambda=0 on (1,2); in tha...}
\expandafter\gdef\csname CGTip@eq:first-cell-split\endcsname{Eq. (13.12), p. 40: logzeta\_lambda(s)=u\_lambda(s)+ dJ\textasciicircum{}*\_lambda,x\textasciicircum{}-s, dJ\textasciicircum{}*\_lambda,x\textasciicircum{}-s:=sint\_2\textasciicircum{}inf J\_lambda(x) x\textasciicircum{}-s-1 dx .}
\expandafter\gdef\csname CGTip@thm:joining-factor\endcsname{Theorem 13.6 (The joining factor) (p. 41): Let h\_lambda:=log zeta\_Z,lambda-logzeta\_lambda on the half-plane where zeta\_Z,lambda converges absolutely (Theorem CGrefthm:atomized-continuation). Then zeta\_Z,lambda=zeta\_lambda e\textasciicircum{} h\_lambda, h\_lambda(s)=-u\_lambda(s)+sum\_m>=11/msum\_n>=2h\_n(lambda;ms), and, explicitly, h\_lambda(s)=-u\_lambda(s)+sum...}
\expandafter\gdef\csname CGTip@eq:joining-factor\endcsname{Eq. (13.13), p. 41: zeta\_Z,lambda=zeta\_lambda e\textasciicircum{} h\_lambda, h\_lambda(s)=-u\_lambda(s)+sum\_m>=11/msum\_n>=2h\_n(lambda;ms),}
\expandafter\gdef\csname CGTip@eq:joining-explicit\endcsname{Eq. (13.14), p. 41: h\_lambda(s)=-u\_lambda(s)+sum\_m>=11/msum\_j>=1(mu(j))/jint\_2\textasciicircum{}infa\_ms(y\textasciicircum{}j) 1/(log y)(1-2/( y)Resum\_k>=1y\textasciicircum{}igamma\_k(lambda)-1/(y(y\textasciicircum{}2-1)))dy,}
\expandafter\gdef\csname CGTip@eq:joining-derivative\endcsname{Eq. (13.15), p. 41: d\_lambda h\_lambda(s)=-d\_\_lambda(s)+sum\_m>=11/msum\_j>=1(mu(j))/jint\_2\textasciicircum{}infa\_ms(y\textasciicircum{}j) 2/( y)sum\_k>=1eta\_ksin(gamma\_k(lambda)log y) dy .}
\expandafter\gdef\csname CGTip@prop:cell-explicit\endcsname{Proposition 13.7 (Explicit cell defects) (p. 41): For every n>=2, lambda in [0,1] and Re s>1, h\_n(lambda;s)=sum\_j>=1(mu(j))/jint\_2\textasciicircum{}infvarphi\_n(y\textasciicircum{}j)(n\textasciicircum{}-s-y\textasciicircum{}-js) 1/(log y)(1-2/( y)Resum\_k>=1y\textasciicircum{}igamma\_k(lambda)-1/(y(y\textasciicircum{}2-1)))dy, d\_\_n(lambda;s)=sum\_j>=1(mu(j))/jint\_2\textasciicircum{}infvarphi\_n(y\textasciicircum{}j)(n\textasciicircum{}-s-y\textasciicircum{}-js) 2/( y)sum\_k>=1eta\_ksin(gamma\_k(lambda)log y) dy , the su...}
\expandafter\gdef\csname CGTip@eq:cell-explicit\endcsname{Eq. (13.16), p. 41: h\_n(lambda;s)=sum\_j>=1(mu(j))/jint\_2\textasciicircum{}infvarphi\_n(y\textasciicircum{}j)(n\textasciicircum{}-s-y\textasciicircum{}-js) 1/(log y)(1-2/( y)Resum\_k>=1y\textasciicircum{}igamma\_k(lambda)-1/(y(y\textasciicircum{}2-1)))dy,}
\expandafter\gdef\csname CGTip@eq:cell-explicit-derivative\endcsname{Eq. (13.17), p. 41: d\_\_n(lambda;s)=sum\_j>=1(mu(j))/jint\_2\textasciicircum{}infvarphi\_n(y\textasciicircum{}j)(n\textasciicircum{}-s-y\textasciicircum{}-js) 2/( y)sum\_k>=1eta\_ksin(gamma\_k(lambda)log y) dy ,}
\expandafter\gdef\csname CGTip@prop:cell-limit\endcsname{Proposition 13.8 (The limit lambda->0 and the meaning of a vanishing cell) (p. 42): For every n>=2 and Re s>1, h\_n(lambda;s)->h\_n(0;s) and c\_n(lambda)-> c\_n(0) as lambda->0\textasciicircum{}+, and h\_n(0;s)=h\_n(1;s)-int\_0\textasciicircum{}+\textasciicircum{}1d\_\_n(lambda;s) dlambda . Moreover h\_n(0;s)=0 for all Re s>1 if and only if varphi\_n dpi\_Z1= c\_n(0) delta\_n.}
\expandafter\gdef\csname CGTip@eq:cell-flow\endcsname{Eq. (13.18), p. 42: h\_n(0;s)=h\_n(1;s)-int\_0\textasciicircum{}+\textasciicircum{}1d\_\_n(lambda;s) dlambda .}
\expandafter\gdef\csname CGTip@subsec:mean-square\endcsname{Section 13.4: Mean square of the cell masses (p. 42)}
\expandafter\gdef\csname CGTip@eq:VW\endcsname{Eq. (13.19), p. 42: V\_n(gamma):=sum\_j>=1(mu(j))/jint\_2\textasciicircum{}infvarphi\_n(y\textasciicircum{}j) y\textasciicircum{}-1/2e\textasciicircum{}i y dy, W\_n(gamma):=sum\_j>=1(mu(j))/jint\_2\textasciicircum{}infvarphi\_n(y\textasciicircum{}j) y\textasciicircum{}-1/2e\textasciicircum{}i ylog y dy ,}
\expandafter\gdef\csname CGTip@eq:cell-mass-W\endcsname{Eq. (13.20), p. 42: c\_n(lambda)=m\_n-2 Resum\_k>=1W\_n(gamma\_k(lambda)), d\_lambda c\_n(lambda)=2sum\_k>=1eta\_k Im V\_n(gamma\_k(lambda)),}
\expandafter\gdef\csname CGTip@lem:mean-value-ordinates\endcsname{Lemma 13.9 (Mean value over the moving ordinates) (p. 42): There is C such that for all lambda in [0,1], T>=2, u\_0 in R and complex a\_k with |a\_k|<=1, int\_u\_0-1\textasciicircum{}u\_0+1|sum\_gamma\_k(lambda)<= Ta\_k e\textasciicircum{}igamma\_k(lambda)u|\textasciicircum{}2du<= C Tlog\textasciicircum{}2T .}
\expandafter\gdef\csname CGTip@eq:mean-value-ordinates\endcsname{Eq. (13.21), p. 42: int\_u\_0-1\textasciicircum{}u\_0+1|sum\_gamma\_k(lambda)<= Ta\_k e\textasciicircum{}igamma\_k(lambda)u|\textasciicircum{}2du<= C Tlog\textasciicircum{}2T .}
\expandafter\gdef\csname CGTip@eq:fejer-pairs\endcsname{Eq. (13.22), p. 42: <<no.(k,k'):gamma\_k(lambda)<= T, |gamma\_k(lambda)-gamma\_k'(lambda)|<=1.}
\expandafter\gdef\csname CGTip@thm:mean-square-cells\endcsname{Theorem 13.10 (Mean square of the cell masses along the family) (p. 42): For every eps>0 there is C\_eps such that for all X>=2 and all lambda in [0,1] sum\_X<n<=2X| c\_n(lambda)|\textasciicircum{}2<= C\_eps X\textasciicircum{}1+eps, sum\_X<n<=2X|d\_lambda c\_n(lambda)|\textasciicircum{}2<= C\_eps X\textasciicircum{}1+eps. In particular, letting lambda->0\textasciicircum{}+ (each c\_n is continuous in lambda and the constant does not depend on lambda), sum\_X<n...}
\expandafter\gdef\csname CGTip@eq:mean-square-cells\endcsname{Eq. (13.23), p. 42: sum\_X<n<=2X| c\_n(lambda)|\textasciicircum{}2<= C\_eps X\textasciicircum{}1+eps, sum\_X<n<=2X|d\_lambda c\_n(lambda)|\textasciicircum{}2<= C\_eps X\textasciicircum{}1+eps.}
\expandafter\gdef\csname CGTip@eq:A-nj\endcsname{Eq. (13.24), p. 43: A\_n,j=sum\_kint\_2\textasciicircum{}infvarphi\_n(y\textasciicircum{}j) y\textasciicircum{}-1/2+igamma\_k(lambda)log y dy , 1<= j<= J:=floor((log(2X+1))/(log2)) .}
\expandafter\gdef\csname CGTip@eq:cell-CS\endcsname{Eq. (13.25), p. 43: |intvarphi\_n(y\textasciicircum{}j)S(y) dy|\textasciicircum{}2<=|varphi|\_inf\textasciicircum{}2 |I\_n,j|int\_I\_n,j|S|\textasciicircum{}2dy .}
\expandafter\gdef\csname CGTip@eq:cell-block\endcsname{Eq. (13.26), p. 43: sum\_X<n<=2X|intvarphi\_n(y\textasciicircum{}j)S dy|\textasciicircum{}2<< X\textasciicircum{}1/j-1((j)/(log X))\textasciicircum{}2int\_u\_0-1\textasciicircum{}u\_0+1|sum\_gamma\_k(lambda)<= Te\textasciicircum{}igamma\_k(lambda)u|\textasciicircum{}2du<< j\textasciicircum{}2X\textasciicircum{}1/j+2eps.}
\expandafter\gdef\csname CGTip@cor:abscissa-one\endcsname{Corollary 13.11 (Abscissa 1 for every member) (p. 43): For every lambda in [0,1] and sigma>1, sum\_n>=2| c\_n(lambda)| n\textasciicircum{}-sigma<inf ; hence zeta\_Z,lambda and 1/ zeta\_Z,lambda are absolutely convergent Dirichlet series in Re s>1, for every lambda, including the endpoint lambda=0. Moreover, uniformly in lambda, no.n<= X:| c\_n(lambda)|>n\textasciicircum{}delta<<\_eps X\textasciicircum{}1-2de...}
\expandafter\gdef\csname CGTip@eq:abscissa-one\endcsname{Eq. (13.27), p. 43: sum\_n>=2| c\_n(lambda)| n\textasciicircum{}-sigma<inf ;}
\expandafter\gdef\csname CGTip@subsec:atomized-continuation\endcsname{Section 13.5: Analytic continuation of the atomized family (p. 44)}
\expandafter\gdef\csname CGTip@thm:atomized-continuation\endcsname{Theorem 13.12 (Continuation of the atomized family) (p. 44): Let lambda in [0,1]. For every c>1/2, (pi\_lambda-L) x\textasciicircum{}-c in L\textasciicircum{}2((2,inf),dx/x). | c\_n(lambda)|<= C(1+I\_n\textasciicircum{}1/2) with I\_n as below, and sum\_n| c\_n(lambda)|n\textasciicircum{}-sigma<inf for sigma>3/2; by Corollary CGrefcor:abscissa-one this holds for sigma>1, and zeta\_Z,lambda is an absolutely convergent Dirichlet ser...}
\expandafter\gdef\csname CGTip@eq:cell-defect-bound\endcsname{Eq. (13.28), p. 44: |h\_n(lambda;s)|<= C|s| n\textasciicircum{}-sigma-1(1+I\_n\textasciicircum{}1/2), sum\_nI\_n\textasciicircum{}1/2n\textasciicircum{}-sigma-1<inf (sigma>c).}
\expandafter\gdef\csname CGTip@thm:closing\endcsname{Theorem 13.13 (Closing criterion; via Hamburger's theorem [ref]) (p. 44): The following are equivalent. (D), that is, zeta\_Z1=zeta. h\_0==0. J\_Z1=0 on (1,2), and for every n>=2 and every Re s>1, int\_0\textasciicircum{}+\textasciicircum{}1d\_\_n(lambda;s) dlambda=h\_n(1;s). In this case c\_n(0)=1\_ prime(n), a\_N(0)=1 for all N, and varpi\_Z1=1.}
\expandafter\gdef\csname CGTip@eq:closing-cells\endcsname{Eq. (13.29), p. 44: int\_0\textasciicircum{}+\textasciicircum{}1d\_\_n(lambda;s) dlambda=h\_n(1;s).}
\expandafter\gdef\csname CGTip@rem:A6-scale\endcsname{Remark 13.14 (The pointwise size of the cell masses) (p. 45): Theorem CGrefthm:closing needs no pointwise bound on c\_n(0): its proof uses the cell masses only through CGeqrefeq:abscissa-one. The bound c\_n(0)=O(n\textasciicircum{}eps) for every n is not proved directly; Corollary CGrefcor:abscissa-one restricts the exceptions to a set of n<= X of size O(X\textasciicircum{}1-2delta+eps) for t...}
\expandafter\gdef\csname CGTip@subsec:atomized-endpoints\endcsname{Section 13.6: The two endpoints of the atomized family (p. 45)}
\expandafter\gdef\csname CGTip@prop:atomized-clock\endcsname{Proposition 13.15 (Atomization at the clock endpoint) (p. 45): At lambda=1: c\_n(1)=O(1), and zeta\_Z,1 is an absolutely convergent Dirichlet series in Re s>1; h\_1(sigma)->-logvarpi\_RVM=log L!=0 as sigma->+inf; h\_1 is not constant: for every sigma>1, zeta\_RVM(sigma+it)->varpi\_RVM as |t|->inf, whereas zeta\_Z,1(sigma+it) has no limit as |t|->inf.}
\expandafter\gdef\csname CGTip@eq:duality\endcsname{Eq. (13.30), p. 45: zeta\_lambda\_functional equation zeros 1/2+- igamma\_k(lambda)* e\textasciicircum{} h\_lambda = zeta\_Z,lambda\_Dirichlet series Euler product, h\_1 not constant, (D) h\_0==0 .}
\expandafter\gdef\csname CGTip@eq:atomized-mobius\endcsname{Eq. (13.31), p. 45: 1 zeta\_Z,lambda(s)=prod\_n>=2(1-n\textasciicircum{}-s)\textasciicircum{} c\_n(lambda)=sum\_N>=1( mu\_lambda(N))/(N\textasciicircum{}s), mu\_lambda(N)=sum\_prod\_nn\textasciicircum{}r\_n=N prod\_n: r\_n>=1(-1)\textasciicircum{}r\_n c\_n(lambda)r\_n,}
\expandafter\gdef\csname CGTip@eq:atomized-mertens\endcsname{Eq. (13.32), p. 45: M\_lambda(x):=sum\_N<= x mu\_lambda(N), 1 zeta\_Z,lambda=e\textasciicircum{}- h\_lambdaF\_lambda ,}
\expandafter\gdef\csname CGTip@prop:Mellin-Mertens\endcsname{Proposition 13.16 (Mellin form of the endpoint Mertens function) (p. 45): There are c\_1>=2 and a unique measurable B\_Z1 on (0,inf), vanishing on (0,1), with x\textasciicircum{}-cB\_Z1 in L\textasciicircum{}2((1,inf),dx/x) and F\_Z1(s)=sint\_1\textasciicircum{}inf B\_Z1(x) x\textasciicircum{}-s-1 dx=exp(-int\_[1,inf)x\textasciicircum{}-s dJ\textasciicircum{} M\_Z1(x)) (Re s>c\_1), for every c>c\_1. Moreover (D) holds if and only if B\_Z1=M almost everywhere; and if F\_Z1 has an a...}
\expandafter\gdef\csname CGTip@eq:Mellin-Mertens\endcsname{Eq. (13.33), p. 45: F\_Z1(s)=sint\_1\textasciicircum{}inf B\_Z1(x) x\textasciicircum{}-s-1 dx=exp(-int\_[1,inf)x\textasciicircum{}-s dJ\textasciicircum{} M\_Z1(x)) (Re s>c\_1),}
\expandafter\gdef\csname CGTip@subsec:hamburger-table\endcsname{Section 13.7: Hamburger's hypotheses at the two endpoints (p. 46)}
\expandafter\gdef\csname CGTip@eq:complement\endcsname{Eq. (13.34), p. 47: zeta\_lambda\_keeps the functional equation and the zeros on Re s=1/2 zeta\_Z,lambda\_keeps the Dirichlet series and the Euler product, zeta\_Z,lambda=zeta\_lambda e\textasciicircum{} h\_lambda.}
\expandafter\gdef\csname CGTip@eq:closure-common\endcsname{Eq. (13.35), p. 47: h\_0==0 J\_Z1=0 on (1,2), h\_n(0;s)=0 for all n>=2, Re s>1, and varpi\_Z1=1 .}
\expandafter\gdef\csname CGTip@subsec:programme\endcsname{Section 13.8: Programme (p. 47)}
\expandafter\gdef\csname CGTip@rem:input-A7\endcsname{Remark 13.17 (Input for A7) (p. 48): Steps A1--A6 use only two properties of an ordinate list: bounded displacement from the clock and a primitive of the count difference of size O(logtau). The list (tau\_k) has both properties, and for it the corresponding function zeta\_RVM is not a Dirichlet series (Theorem CGrefthm:no-Dirichlet-RV...}
\expandafter\gdef\csname CGTip@sec:criteria\endcsname{Section 14: Criteria for the final step and the clock end of the family (p. 48)}
\expandafter\gdef\csname CGTip@eq:cell-moments\endcsname{Eq. (14.1), p. 48: (xi):=int\_Rvarphi(t)e\textasciicircum{}-2pi ixi t dt, M\_r:=int\_Rt\textasciicircum{}rvarphi(t) dt, C\_r:=int\_Rt\textasciicircum{}rvarphi(t)cos(2pi t) dt .}
\expandafter\gdef\csname CGTip@eq:partition-fourier\endcsname{Eq. (14.2), p. 48: (q)=int\_0\textasciicircum{}1sum\_m in Zvarphi(t-m) e\textasciicircum{}-2pi iqt dt=delta\_q0 (q in Z),}
\expandafter\gdef\csname CGTip@subsec:open-path\endcsname{Section 14.1: The open path (p. 49)}
\expandafter\gdef\csname CGTip@prop:nonconstant\endcsname{Proposition 14.1 (No member of the open path is a Dirichlet series) (p. 49): Let 0<lambda<=1. Then zeta\_lambda is not a constant multiple of a Dirichlet series sum\_na\_nn\textasciicircum{}-s converging absolutely in Re s>1. In particular the joining factor h\_lambda is not constant.}
\expandafter\gdef\csname CGTip@subsec:convexity\endcsname{Section 14.2: Convexity along the family (p. 49)}
\expandafter\gdef\csname CGTip@prop:convexity\endcsname{Proposition 14.2 (Convexity on the real axis) (p. 49): For every real sigma>1 the function lambda\_lambda(sigma) is strictly convex on [0,1]: d\_lambda\textasciicircum{}2logzeta\_lambda(sigma)=sum\_k>=1eta\_k\textasciicircum{}2(C(sigma-1/2,gamma\_k(lambda))-C(1/2,gamma\_k(lambda)))>0, C(a,gamma):=(2(a\textasciicircum{}2-gamma\textasciicircum{}2))/((a\textasciicircum{}2+gamma\textasciicircum{}2)\textasciicircum{}2). Consequently, with the generator V\_lambda of Theorem CGreft...}
\expandafter\gdef\csname CGTip@eq:convexity-second\endcsname{Eq. (14.3), p. 49: d\_lambda\textasciicircum{}2logzeta\_lambda(sigma)=sum\_k>=1eta\_k\textasciicircum{}2(C(sigma-1/2,gamma\_k(lambda))-C(1/2,gamma\_k(lambda)))>0, C(a,gamma):=(2(a\textasciicircum{}2-gamma\textasciicircum{}2))/((a\textasciicircum{}2+gamma\textasciicircum{}2)\textasciicircum{}2).}
\expandafter\gdef\csname CGTip@eq:convexity-sandwich\endcsname{Eq. (14.4), p. 49: logzeta\_RVM(sigma)-V\_1(sigma)<logzeta\_Z1(sigma)<logzeta\_RVM(sigma)-V\_0(sigma) (sigma>1).}
\expandafter\gdef\csname CGTip@subsec:pair\endcsname{Section 14.3: Hamburger's theorem for the pair (p. 49)}
\expandafter\gdef\csname CGTip@thm:pair-criterion\endcsname{Theorem 14.3 (Hamburger's theorem for the pair (zeta\_lambda, zeta\_Z,lambda)) (p. 49): Let lambda in [0,1]. Suppose that the joining factor e\textasciicircum{} h\_lambda extends to a meromorphic function on C of finite order with finitely many poles, and that e\textasciicircum{} h\_lambda(s)=e\textasciicircum{} h\_lambda(1-s) . Then zeta\_Z,lambda=zeta. For lambda=0, moreover, zeta\_Z1=zeta and h\_0==0.}
\expandafter\gdef\csname CGTip@eq:pair-symmetry\endcsname{Eq. (14.5), p. 49: e\textasciicircum{} h\_lambda(s)=e\textasciicircum{} h\_lambda(1-s) .}
\expandafter\gdef\csname CGTip@subsec:positivity\endcsname{Section 14.4: Positivity and the value at infinity (p. 50)}
\expandafter\gdef\csname CGTip@lem:summation-positive\endcsname{Lemma 14.4 (Summation formula for a representing measure) (p. 50): Let lambda in [0,1] and let A be a positive measure on [1,inf) with zeta\_lambda(s)=int\_[1,inf)u\textasciicircum{}-s d A(u) (Re s>1). Then for every even Schwartz function g g(0)+2int\_[1,inf) g d A= g(0)+2int\_[1,inf) g d A .}
\expandafter\gdef\csname CGTip@eq:representing-measure\endcsname{Eq. (14.6), p. 50: zeta\_lambda(s)=int\_[1,inf)u\textasciicircum{}-s d A(u) (Re s>1).}
\expandafter\gdef\csname CGTip@eq:summation-positive\endcsname{Eq. (14.7), p. 50: g(0)+2int\_[1,inf) g d A= g(0)+2int\_[1,inf) g d A .}
\expandafter\gdef\csname CGTip@eq:tate-local\endcsname{Eq. (14.8), p. 50: pi\textasciicircum{}-s/2Gamma(s/2) g\textasciicircum{}(1-s)=pi\textasciicircum{}-(1-s)/2Gamma((1-s)/2) g\textasciicircum{}(s)}
\expandafter\gdef\csname CGTip@eq:summation-shift\endcsname{Eq. (14.9), p. 50: int g d A=1/2 g(0)-1/2 g(0)+1/(2pi i)int\_(1-c)zeta\_lambda(s) g\textasciicircum{}(s) ds .}
\expandafter\gdef\csname CGTip@thm:positivity\endcsname{Theorem 14.5 (Positivity criterion) (p. 50): Let lambda in [0,1]. If zeta\_lambda has a representation CGeqrefeq:representing-measure with a positive measure A, then varpi\_lambda= A(1)<=1 . In particular, if zeta\_Z1 has such a representation, then varpi\_Z1=1 and zeta\_Z1=zeta.}
\expandafter\gdef\csname CGTip@eq:positivity-bound\endcsname{Eq. (14.10), p. 50: varpi\_lambda= A(1)<=1 .}
\expandafter\gdef\csname CGTip@eq:extremal-function\endcsname{Eq. (14.11), p. 50: g(u):=((sinpi u)/(pi u))\textasciicircum{}4(1-u\textasciicircum{}2), g(xi)= B(xi)+1/(4pi\textasciicircum{}2) B''(xi) .}
\expandafter\gdef\csname CGTip@eq:lp-inequality\endcsname{Eq. (14.12), p. 51: 1 >= g\_epsilon(0)+2int g\_epsilon d A= g\_epsilon(0)+2int g\_epsilon d A >= g\_epsilon(0)+2varpi\_lambda g\_epsilon(1) .}
\expandafter\gdef\csname CGTip@eq:lp-value\endcsname{Eq. (14.13), p. 51: varpi\_lambda<=1- g(0)2 g(1)=(1/3+1/(2pi\textasciicircum{}2))/(1/3+1/(2pi\textasciicircum{}2))=1 .}
\expandafter\gdef\csname CGTip@cor:positivity\endcsname{Corollary 14.6 (Positive form of the final step) (p. 51): The following are equivalent: up(i) (D); up(ii) zeta\_Z1 is the Mellin--Stieltjes transform CGeqrefeq:representing-measure of a positive measure on [1,inf); up(iii) J\textasciicircum{} M\_Z1, extended by 0 to (0,1], is non-decreasing on (0,inf).}
\expandafter\gdef\csname CGTip@subsec:clock-density\endcsname{Section 14.5: The density of the clock (p. 51)}
\expandafter\gdef\csname CGTip@eq:U-def\endcsname{Eq. (14.14), p. 51: U(M,t\_0):=int\_t\_0\textasciicircum{}infe\textasciicircum{}-iMPsi(t) dt:=lim\_T->infint\_t\_0\textasciicircum{}Te\textasciicircum{}-iMPsi(t) dt .}
\expandafter\gdef\csname CGTip@lem:stationary-U\endcsname{Lemma 14.7 (Stationary phase for Psi) (p. 51): The limit CGeqrefeq:U-def exists, and there is an absolute constant C\_U with |U(M,t\_0)-sqrt((2pi)/M) e\textasciicircum{}i(M-pi/4)|<=C\_UM (M>=1, 0<t\_0<=1/2).}
\expandafter\gdef\csname CGTip@eq:U-asymptotic\endcsname{Eq. (14.15), p. 51: |U(M,t\_0)-sqrt((2pi)/M) e\textasciicircum{}i(M-pi/4)|<=C\_UM (M>=1, 0<t\_0<=1/2).}
\expandafter\gdef\csname CGTip@eq:U-substitution\endcsname{Eq. (14.16), p. 51: p( w):=(dt)/(d w)=( w)/(log t), p(0)=1, q( w):=( p( w)-1)/( w)=1/(log t)-1/( w).}
\expandafter\gdef\csname CGTip@eq:U-split\endcsname{Eq. (14.17), p. 52: e\textasciicircum{}-iM U(M,t\_0)=int\_ w\_0\textasciicircum{}inf e\textasciicircum{}-iM w\textasciicircum{}2/2 d w+int\_ w\_0\textasciicircum{}inf e\textasciicircum{}-iM w\textasciicircum{}2/2 w q( w) d w .}
\expandafter\gdef\csname CGTip@eq:clock-harmonic\endcsname{Eq. (14.18), p. 52: R\_m(y):=-2(-1)\textasciicircum{}m m\textasciicircum{}-3/2 y\textasciicircum{}1/(2m)-1/2cos(2pi m y\textasciicircum{}1/m-((7m+1)pi)/4),}
\expandafter\gdef\csname CGTip@thm:clock-density\endcsname{Theorem 14.8 (The density of the clock) (p. 52): J\_RVM is continuously differentiable on (2,inf), and for y>=16 J\_RVM'(y)=1/(log y)+2cos(2pi y)-y\textasciicircum{}-1/42cos(4pi y+4)+O(y\textasciicircum{}-1/3). On (2,inf) the measure dpi\_RVM is the sum of the atoms (mu(j))/jJ\_RVM(2) delta\_2\textasciicircum{}j, j>=1, and of a density that is continuous off the points 2\textasciicircum{}j and equals J\_RVM'(y)+O(y\textasciicircum{}-...}
\expandafter\gdef\csname CGTip@eq:clock-density\endcsname{Eq. (14.19), p. 52: J\_RVM'(y)=1/(log y)+2cos(2pi y)-y\textasciicircum{}-1/42cos(4pi y+4)+O(y\textasciicircum{}-1/3).}
\expandafter\gdef\csname CGTip@eq:H-harmonic\endcsname{Eq. (14.20), p. 52: H\_m\textasciicircum{}+-(v):=int\_tau\_0\textasciicircum{}inf e\textasciicircum{}iPhi\_m\textasciicircum{}+-(tau) dtau ,}
\expandafter\gdef\csname CGTip@eq:S-zero\endcsname{Eq. (14.21), p. 52: sum\_kD\_k=int\_u\_1\textasciicircum{}u\_2e\textasciicircum{}v/2v S\_0(v) dv, S\_0(v):=S\_ sm(v)+sum\_m>=1((-1)\textasciicircum{}m v)/(2pi m)(H\_m\textasciicircum{}-(v)-H\_m\textasciicircum{}+(v)),}
\expandafter\gdef\csname CGTip@eq:J-prime-exact\endcsname{Eq. (14.22), p. 52: J\_RVM'(y)=1/(log y)(1-1/(y(y\textasciicircum{}2-1))-2/( y) Re S\_0(log y)).}
\expandafter\gdef\csname CGTip@eq:H-by-U\endcsname{Eq. (14.23), p. 52: H\_m\textasciicircum{}-(v)= t\_m e\textasciicircum{}-7pi im/4 U(m t\_m,tau\_0/ t\_m), tau\_0/ t\_m=2e\textasciicircum{}-v/m<=1/2 (y>=4\textasciicircum{}m).}
\expandafter\gdef\csname CGTip@eq:H-asymptotic\endcsname{Eq. (14.24), p. 53: ((-1)\textasciicircum{}mv)/(2pi m) H\_m\textasciicircum{}-(v)=(-1)\textasciicircum{}m vm\textasciicircum{}3/2 y\textasciicircum{}1/(2m)exp(i(2pi m y\textasciicircum{}1/m-((7m+1)pi)/4))+O(v/(m\textasciicircum{}2)).}
\expandafter\gdef\csname CGTip@cor:clock-cells\endcsname{Corollary 14.9 (Cell masses and cell defects at the clock end) (p. 53): For n>=17, with 1\_2\textasciicircum{}j(n)=1 if n=2\textasciicircum{}j and 0 otherwise, c\_n\textasciicircum{}RVM=int\_n-1\textasciicircum{}n(1/(log x)+2cos(2pi x))dx+1\_2\textasciicircum{}j(n) (mu(j))/j J\_RVM(2)+O(n\textasciicircum{}-1/4), c\_n(1)=1/(log n)+1\_2\textasciicircum{}j(n) (mu(j))/j J\_RVM(2)+O(n\textasciicircum{}-1/4), and for 0<Re s<=sigma\_1 |h\_n(1;s)|<= C\_sigma\_1(|s| n\textasciicircum{}-sigma-4/3+|s|(1+|s|) n\textasciicircum{}-sigma-2). In particular Hypo...}
\expandafter\gdef\csname CGTip@eq:clock-cell-sharp\endcsname{Eq. (14.25), p. 53: c\_n\textasciicircum{}RVM=int\_n-1\textasciicircum{}n(1/(log x)+2cos(2pi x))dx+1\_2\textasciicircum{}j(n) (mu(j))/j J\_RVM(2)+O(n\textasciicircum{}-1/4),}
\expandafter\gdef\csname CGTip@eq:clock-cell-smooth\endcsname{Eq. (14.26), p. 53: c\_n(1)=1/(log n)+1\_2\textasciicircum{}j(n) (mu(j))/j J\_RVM(2)+O(n\textasciicircum{}-1/4),}
\expandafter\gdef\csname CGTip@eq:clock-defect\endcsname{Eq. (14.27), p. 53: |h\_n(1;s)|<= C\_sigma\_1(|s| n\textasciicircum{}-sigma-4/3+|s|(1+|s|) n\textasciicircum{}-sigma-2).}
\expandafter\gdef\csname CGTip@subsec:height-kernels\endcsname{Section 14.6: The cell masses in the heights of the ordinates (p. 53)}
\expandafter\gdef\csname CGTip@prop:height-kernels\endcsname{Proposition 14.10 (Height kernel of the cell masses) (p. 53): For every eps>0, lambda in [0,1] and n>=3 that is not a power of 2, c\_n(lambda)=m\_n-2/( n log n)sum\_k>=1((gamma\_k(lambda))/(2pi n))cos(gamma\_k(lambda)log n)+E\_n(lambda), where sum\_X<n<=2X|E\_n(lambda)|\textasciicircum{}2<= C\_eps X\textasciicircum{}1/2+eps uniformly in lambda (sum over the n that are not powers of 2). The kernel (g...}
\expandafter\gdef\csname CGTip@eq:height-kernel\endcsname{Eq. (14.28), p. 53: c\_n(lambda)=m\_n-2/( n log n)sum\_k>=1((gamma\_k(lambda))/(2pi n))cos(gamma\_k(lambda)log n)+E\_n(lambda),}
\expandafter\gdef\csname CGTip@eq:Xi-rotation\endcsname{Eq. (14.29), p. 54: Xi\_n(Omega):=sum\_k>=1((tau\_k)/(2pi n))n\textasciicircum{}itau\_ke\textasciicircum{}-iOmegaeta\_k (0<=Omega<=log n).}
\expandafter\gdef\csname CGTip@prop:flow-count\endcsname{Proposition 14.11 (The flow of a cell mass and the count difference) (p. 54): For every n>=3 that is not a power of 2, c\_n(0)- c\_n(1)=-2int\_0\textasciicircum{}+\textasciicircum{}1sum\_k>=1eta\_k Im V\_n(gamma\_k(lambda)) dlambda=-2 Imint\_tau\_0\textasciicircum{}inf V\_n(t)(N\_zeta(t)-n\_RVM(t)) dt .}
\expandafter\gdef\csname CGTip@eq:flow-count\endcsname{Eq. (14.30), p. 54: c\_n(0)- c\_n(1)=-2int\_0\textasciicircum{}+\textasciicircum{}1sum\_k>=1eta\_k Im V\_n(gamma\_k(lambda)) dlambda=-2 Imint\_tau\_0\textasciicircum{}inf V\_n(t)(N\_zeta(t)-n\_RVM(t)) dt .}
\expandafter\gdef\csname CGTip@subsec:gaussian-defect\endcsname{Section 14.7: The Gaussian defect of the atomization (p. 54)}
\expandafter\gdef\csname CGTip@lem:vertical-growth\endcsname{Lemma 14.12 (Vertical growth of the joining factor) (p. 54): For every lambda in [0,1] and 1/2<sigma\_0<sigma\_1, h\_lambda(sigma+it)=o(|t|) as |t|->inf, uniformly for sigma\_0<=sigma<=sigma\_1.}
\expandafter\gdef\csname CGTip@eq:cell-RL\endcsname{Eq. (14.31), p. 54: h\_n(lambda;ms)= intvarphi\_n(x)(n\textasciicircum{}-ms-x\textasciicircum{}-ms)dL(x)-int E(x)varphi\_n'(x)(n\textasciicircum{}-ms-x\textasciicircum{}-ms)dx; -msint E(x)varphi\_n(x) x\textasciicircum{}-ms-1dx .}
\expandafter\gdef\csname CGTip@eq:gaussian-defect\endcsname{Eq. (14.32), p. 54: Theta\_lambda(y):=1/(4pi i)int\_(c)zeta\_lambda(s) Gamma(s/2) pi\textasciicircum{}-s/2 y\textasciicircum{}-s/2 ds, G\_lambda(y):=sum\_N>=1 a\_N(lambda) e\textasciicircum{}-pi N\textasciicircum{}2y-Theta\_lambda(y).}
\expandafter\gdef\csname CGTip@prop:gaussian-defect\endcsname{Proposition 14.13 (Gaussian defect) (p. 54): Let lambda in [0,1]. Theta\_lambda(y)=y\textasciicircum{}-1/2 Theta\_lambda(1/y)+1/2(y\textasciicircum{}-1/2-1) for y>0, and Theta\_lambda(y)=O\_A(y\textasciicircum{}-A) as y->inf for every A. The coefficients a\_N(lambda) satisfy the theta relation of zeta, sum\_N>=1 a\_N(lambda) e\textasciicircum{}-pi N\textasciicircum{}2y=y\textasciicircum{}-1/2sum\_N>=1 a\_N(lambda) e\textasciicircum{}-pi N\textasciicircum{}2/y+1/2(y\textasciicircum{}-1/2-1) (y>0), if...}
\expandafter\gdef\csname CGTip@eq:theta-atomized\endcsname{Eq. (14.33), p. 55: sum\_N>=1 a\_N(lambda) e\textasciicircum{}-pi N\textasciicircum{}2y=y\textasciicircum{}-1/2sum\_N>=1 a\_N(lambda) e\textasciicircum{}-pi N\textasciicircum{}2/y+1/2(y\textasciicircum{}-1/2-1) (y>0),}
\expandafter\gdef\csname CGTip@eq:defect-selfdual\endcsname{Eq. (14.34), p. 55: G\_lambda(y)=y\textasciicircum{}-1/2 G\_lambda(1/y) (y>0).}
\expandafter\gdef\csname CGTip@eq:defect-expansion\endcsname{Eq. (14.35), p. 55: G\_lambda(y)=1/2(e\textasciicircum{} h\_lambda(1)-1) y\textasciicircum{}-1/2+O\_delta(y\textasciicircum{}-1/4-delta),}
\expandafter\gdef\csname CGTip@cor:residue-condition\endcsname{Corollary 14.14 (The density of the atomized integers) (p. 55): sum\_N<= x a\_N(lambda) has mean density e\textasciicircum{} h\_lambda(1) in the sense of CGeqrefeq:defect-expansion, and (D) implies h\_0(1)=0, that is, u\_0(1)=sum\_m>=11/msum\_n>=2h\_n(0;m).}
\expandafter\gdef\csname CGTip@eq:residue-condition\endcsname{Eq. (14.36), p. 55: h\_0(1)=0, that is, u\_0(1)=sum\_m>=11/msum\_n>=2h\_n(0;m).}
\expandafter\gdef\csname CGTip@sec:numerics\endcsname{Section 15: Numerical illustration with 2 001 0522001052 zeros (p. 55)}
\expandafter\gdef\csname CGTip@subsec:num-clock\endcsname{Section 15.1: The clock and the zeros (p. 56)}
\expandafter\gdef\csname CGTip@exp:clock\endcsname{Experiment 15.1 (Clock, midpoints and displacements) (p. 56): Object: N\_RVM, N\_zeta, tau\_k, eta\_k=tau\_k-gamma\_k. Range: k<= N. Expected: CGeqrefeq:midpoint-condition to rounding, and |eta\_k|<= C\_eta (Lemma CGreflem:bounded-displacement). Observed: max\_k|N\_RVM(tau\_k)-(k-1/2)|=6.985*10\textasciicircum{}-10; max\_k|eta\_k|=0.9937 (at k=34), mean |eta\_k|=0.1512, root mean square ...}
\expandafter\gdef\csname CGTip@fig:clock-zoom\endcsname{Figure 1 (p. 56): The midpoint clock N\_RVM (red), the zero staircase N\_zeta (black), the midpoint ordinates tau\_k (red points) and the points (gamma\_k,k-1/2) (squares). Left: the first nine zeros. Right: the last twelve zeros of the file.}
\expandafter\gdef\csname CGTip@fig:remainder\endcsname{Figure 2 (p. 56): The remainder N\_zeta(T)-N\_RVM(T)=S(T)+O(1/T) of CGeqrefeq:RvM-full for T<=200.}
\expandafter\gdef\csname CGTip@subsec:num-family\endcsname{Section 15.2: The lambdalambda-family and the endpoint (p. 56)}
\expandafter\gdef\csname CGTip@fig:family\endcsname{Figure 3 (p. 57): Top: eta\_k for every 50th k<= N and its distribution. Bottom: moving ordinates in clock units, k<=30, and the functions Z\_lambda(t) on the critical line.}
\expandafter\gdef\csname CGTip@exp:Z1\endcsname{Experiment 15.2 (The endpoint zeta\_Z1 against zeta) (p. 56): Object: zeta\_Z1=P\_Z1/G computed from the N ordinates gamma\_k, and zeta computed by mpmath. Range: sigma in 2,1.25,0.75, 1<= t<=2000 (499 points per line), and the critical line 0.5<= t<=60. Expected: if zeta\_Z1=zeta, the ratio is 1 up to the truncation of the product. Observed: max|zeta\_Z1/zeta-1...}
\expandafter\gdef\csname CGTip@fig:Z1\endcsname{Figure 4 (p. 57): Left: |zeta\_Z1/zeta-1| on three vertical lines. Right: Z\_Z1 and Hardy's Z on the critical line.}
\expandafter\gdef\csname CGTip@exp:product-identity\endcsname{Experiment 15.3 (The value at infinity: the product identity) (p. 57): Object: the ordered partial products P\_K=prod\_k<= K(1/4+gamma\_k\textasciicircum{}2)/(1/4+tau\_k\textasciicircum{}2) of CGeqrefeq:product-identity. Range: K<= N. Expected: by Corollary CGrefcor:product-identity, zeta\_Z1=zeta is the statement P\_K->varpi\_RVM. Observed: P\_K=0.96256, 0.96248, 0.96229, 0.961951, 0.961930, 0.9619418189, ...}
\expandafter\gdef\csname CGTip@fig:product-identity\endcsname{Figure 5 (p. 58): The partial products P\_K of CGeqrefeq:product-identity against varpi\_RVM (dashed). Left: K<=10\textasciicircum{}4. Right: (P\_K-varpi\_RVM)*10\textasciicircum{}6 for 10\textasciicircum{}5<= K<= N. The value of the endpoint at +inf is varpi\_Z1=varpi\_RVM/lim\_KP\_K (Theorem CGrefthm:Z1-clock).}
\expandafter\gdef\csname CGTip@subsec:num-transition\endcsname{Section 15.3: The transition theorem and zeta\_RVMzeta\_RVM against zetazeta (p. 58)}
\expandafter\gdef\csname CGTip@exp:transition\endcsname{Experiment 15.4 (Transition theorem) (p. 58): Object: zeta\_RVM, w(s), varpi\_RVM, L. Range: sigma' in 0.25,0.5,1.5, 10<= t<=10\textasciicircum{}4. Expected: Theorem CGrefthm:transition and Corollary CGrefcor:vertical-limit: t |zeta\_RVM/(varpi\_RVM(1+w))-1| bounded, |w|=(t/2pi)\textasciicircum{}-sigma'(1+O(1/t)). Observed: Richardson extrapolation on Re s=2 gives varpi\_RVM=0.96...}
\expandafter\gdef\csname CGTip@fig:transition\endcsname{Figure 6 (p. 59): Top: the transition error multiplied by t, and the rescaled vertical limit t\textasciicircum{}min(sigma',1)|1/zeta\_RVM-L|. Bottom: |zeta| and |zeta\_RVM| on sigma=2 and sigma=0.75.}
\expandafter\gdef\csname CGTip@subsec:num-expsums\endcsname{Section 15.4: Exponential sums (p. 58)}
\expandafter\gdef\csname CGTip@exp:expsums\endcsname{Experiment 15.5 (Exponential sums over the two spectra) (p. 58): Object: S(N;log x)=sum\_k<= Nx\textasciicircum{}io\_k for o\_k=tau\_k and o\_k=gamma\_k. Range: N<=2 001 052, x in 2,3,4,6,10. Expected: |S\_RVM|<<1+log N (Theorem CGrefthm:exp-sum); for gamma\_k, Landau's term -T/(2pi)(Lambda(x))/( x). Observed: max\_N|S\_RVM(N;log x)|=3.57,3.98,3.95,4.53,9.70; at N=2 001 052, |sum x\textasciicircum{}igam...}
\expandafter\gdef\csname CGTip@fig:expsums\endcsname{Figure 7 (p. 59): The midpoint spectrum carries no prime frequencies (left); the spectrum of zeta resonates at prime powers (right).}
\expandafter\gdef\csname CGTip@subsec:num-primes\endcsname{Section 15.5: Prime counts: pipi and pi\_RVMpi\_RVM (p. 58)}
\expandafter\gdef\csname CGTip@exp:primes\endcsname{Experiment 15.6 (Prime counts and RVM primes) (p. 58): Object: J\_RVM, pi\_RVM, c\_n\textasciicircum{}RVM, p\_j\textasciicircum{}RVM against J, pi, the primes. Range: 2<= x<=100 with all N ordinates of each list (3000 points); 2<= x<=1000 with the first K=17 061 ordinates and a smooth cut (40 000 points). Expected: the explicit formula with the gamma\_k reproduces pi(x); pi\_RVM is continu...}
\expandafter\gdef\csname CGTip@fig:prime-counts\endcsname{Figure 8 (p. 60): pi(x) (black staircase), the explicit formula with the zeros of zeta (blue), pi\_RVM(x) (red) and li(x)-li(2) (green). The distance between the red and green curves is the quantity bounded in Theorem CGrefthm:Lambert-estimate; the ratios are in Figure CGreffig:pi-ratios.}
\expandafter\gdef\csname CGTip@fig:rvm-primes\endcsname{Figure 9 (p. 60): Left: cell masses c\_n\textasciicircum{}RVM, the prime indicator and 1/log n. Right: real RVM prime levels against the primes up to 1000.}
\expandafter\gdef\csname CGTip@exp:pi-ratios\endcsname{Experiment 15.7 (The three ratios and the square-root scale) (p. 59): Object: li/pi\_RVM, pi/li, pi/pi\_RVM, and (pi\_RVM-li)/B, (pi-li)/B with B(x)= x(loglog x)\textasciicircum{}2/log x. Range: the grid of Experiment CGrefexp:primes on [8,1000] (K=17 061 ordinates). Expected: the three ratios tend to 1 (Corollary CGrefcor:three-limits) and (pi\_RVM-li)/B stays bounded (Theorem CGrefth...}
\expandafter\gdef\csname CGTip@fig:pi-ratios\endcsname{Figure 10 (p. 60): Left: the three ratios of Corollary CGrefcor:three-limits. Right: pi\_RVM-li and pi-li in units of the scale B(x) of Theorem CGrefthm:Lambert-estimate. The ratio pi/pi\_RVM (black) is the closest of the three to 1; the ratios with li approach 1 at the rate 1/2li( x)/li(x) of the bias J-pi.}
\expandafter\gdef\csname CGTip@subsec:num-mertens\endcsname{Section 15.6: The Mertens functions and the ternary division (p. 61)}
\expandafter\gdef\csname CGTip@eq:ternary-update\endcsname{Eq. (15.1), p. 61: M\_h(n)=M\_h(n-1)+mu(n)\textasciicircum{}2 varsigma\_n, |M\_h(n)-M(n)| small,}
\expandafter\gdef\csname CGTip@tab:ternary\endcsname{Table 3 (p. 61): The ternary division for n<=15 (the first 40 values are in Appendix CGrefapp:out-s05-mertens). The zero columns of mu and mu\_h coincide; the signs of mu\_h are chosen by comparing M\_h(n-1) with M(n). For n<=12 the two walks agree; at n=13 the target -1.85 lies above M\_h(12)=-2, so mu\_h(13)=+1 whil...}
\expandafter\gdef\csname CGTip@exp:mertens\endcsname{Experiment 15.8 (Mertens functions and the ternary walk) (p. 61): Object: M, M\_RVM\textasciicircum{}crit, M\_h; the explicit formulas. Range: n<=10\textasciicircum{}6; explicit formulas on [2,2000] with all N ordinates of each list. Expected: the explicit formula with the zeros of zeta reproduces M; M\_RVM\textasciicircum{}crit(x)=O(x\textasciicircum{}1/2+eps) (Theorem CGrefthm:unsmoothed-RVM); M\_h tracks M\_RVM\textasciicircum{}crit with the dens...}
\expandafter\gdef\csname CGTip@fig:mertens\endcsname{Figure 11 (p. 62): Left: M (black), M\_RVM\textasciicircum{}crit (red) and M\_h (green) at every integer n<=2000. Right: range of M and mean of M\_RVM\textasciicircum{}crit in blocks of 500 integers up to 10\textasciicircum{}6, with +- n.}
\expandafter\gdef\csname CGTip@fig:ternary\endcsname{Figure 12 (p. 62): Left: mu(n) and mu\_h(n) for n<=60: same zeros, different signs. Right: M, the target M\_RVM\textasciicircum{}crit and M\_h for n<=300.}
\expandafter\gdef\csname CGTip@fig:explicit\endcsname{Figure 13 (p. 63): Left: M(x), its explicit formula from the zeros of zeta and M\_RVM\textasciicircum{}crit(x) from the tau\_k on [2,300]. Right: normalized derivatives at the zeros: erratic for zeta, constant varpi\_RVM for zeta\_RVM.}
\expandafter\gdef\csname CGTip@sec:conclusions\endcsname{Section 16: Conclusions and open problems (p. 62)}
\expandafter\gdef\csname CGTip@rem:open-problems\endcsname{Remark 16.1 (Open problems) (p. 64): Prove the identity zeta\_Z1=zeta, equivalently condition (D), equivalently an absolutely convergent Dirichlet expansion of K=zeta\_Z1/zeta\_Z,RVM in Re s>1, equivalently the product identity CGeqrefeq:product-identity, equivalently the cell identity CGeqrefeq:cell-identity with the jump bound of The...}
\expandafter\gdef\csname CGTip@app:code\endcsname{Section A: Python programs (p. 64)}
\expandafter\gdef\csname CGTip@app:code-rvm-common\endcsname{Section A.1: ttrvm\_common.py (p. 64)}
\expandafter\gdef\csname CGTip@app:code-s01-clock-and-family\endcsname{Section A.2: tts01\_clock\_and\_family.py (p. 67)}
\expandafter\gdef\csname CGTip@app:code-s02-transition\endcsname{Section A.3: tts02\_transition.py (p. 70)}
\expandafter\gdef\csname CGTip@app:code-s03-exponential-sums\endcsname{Section A.4: tts03\_exponential\_sums.py (p. 71)}
\expandafter\gdef\csname CGTip@app:code-s04-primes\endcsname{Section A.5: tts04\_primes.py (p. 72)}
\expandafter\gdef\csname CGTip@app:code-s05-mertens\endcsname{Section A.6: tts05\_mertens.py (p. 75)}
\expandafter\gdef\csname CGTip@app:code-s06-prime-ratios\endcsname{Section A.7: tts06\_prime\_ratios.py (p. 80)}
\expandafter\gdef\csname CGTip@app:code-s07-product-identity\endcsname{Section A.8: tts07\_product\_identity.py (p. 81)}
\expandafter\gdef\csname CGTip@app:output\endcsname{Section B: Console output of the programs (p. 82)}
\expandafter\gdef\csname CGTip@app:out-s01-clock-and-family\endcsname{Section B.1: tts01\_clock\_and\_family (p. 82)}
\expandafter\gdef\csname CGTip@app:out-s01b-clock-zoom\endcsname{Section B.2: tts01b\_clock\_zoom (p. 82)}
\expandafter\gdef\csname CGTip@app:out-s02-transition\endcsname{Section B.3: tts02\_transition (p. 83)}
\expandafter\gdef\csname CGTip@app:out-s03-exponential-sums\endcsname{Section B.4: tts03\_exponential\_sums (p. 83)}
\expandafter\gdef\csname CGTip@app:out-s04-primes\endcsname{Section B.5: tts04\_primes (p. 83)}
\expandafter\gdef\csname CGTip@app:out-s05-mertens\endcsname{Section B.6: tts05\_mertens (p. 84)}
\expandafter\gdef\csname CGTip@app:out-s06-prime-ratios\endcsname{Section B.7: tts06\_prime\_ratios (p. 85)}
\expandafter\gdef\csname CGTip@app:out-s07-product-identity\endcsname{Section B.8: tts07\_product\_identity (p. 86)}
\makeatother


\hypersetup{
  colorlinks=false,
  pdfborder={0 0 1},
  linkbordercolor={1 0 0},
  citebordercolor={0 0.3 1},
  urlbordercolor={0 0.6 0.6},
  filebordercolor={0 0.6 0.6},
  linktoc=all,
  bookmarksnumbered=true,
  bookmarksopen=true,
  bookmarksopenlevel=1,
  pdfstartview=FitH,
  pdftitle={The midpoint Riemann-von Mangoldt synthetic zeta system},
  pdfauthor={Rotkin Aynasac},
  pdfsubject={Midpoint RVM zeta function, the lambda-family to the endpoint Z1, Hamburger closure, and the Mertens bound},
  pdfkeywords={Riemann zeta function, Riemann-von Mangoldt clock, Hadamard product, Hamburger theorem, Mertens function}
}
% Back-references in the bibliography (placed after the DOI).
\renewcommand*{\backref}[1]{}
\renewcommand*{\backrefalt}[4]{%
  \ifcase #1 %
    \relax
  \or
    \ (Cited on p.~#2.)%
  \else
    \ (Cited on pp.~#2.)%
  \fi}

% In \cite[note]{key} the note (e.g. a section number) is linked too.
\makeatletter
\let\CG@origcite\cite
\renewcommand*{\cite}[2][]{\def\CG@citekey{#2}%
  \ifx\relax#1\relax\CG@origcite{#2}\else\CG@origcite[#1]{#2}\fi}
\def\@cite#1#2{{\@citestyle[\citeform{#1}\if@tempswa, \hyper@link{cite}{cite.\CG@citekey}{#2}\fi]}}
\makeatother

\geometry{margin=1in}
\numberwithin{equation}{section}
\setcounter{tocdepth}{3}
\setcounter{secnumdepth}{3}
\setlength{\emergencystretch}{1.5em}

\theoremstyle{plain}
\newtheorem{theorem}{Theorem}[section]
\newtheorem{proposition}[theorem]{Proposition}
\newtheorem{lemma}[theorem]{Lemma}
\newtheorem{corollary}[theorem]{Corollary}
\newtheorem{conjecture}[theorem]{Conjecture}
\theoremstyle{definition}
\newtheorem{definition}[theorem]{Definition}
\newtheorem{hypothesis}[theorem]{Hypothesis}
\newtheorem{heuristic}[theorem]{Heuristic}
\newtheorem{experiment}[theorem]{Experiment}
\theoremstyle{remark}
\newtheorem{remark}[theorem]{Remark}
\newtheorem{warning}[theorem]{Warning}

\newcommand{\R}{\mathbb R}
\newcommand{\C}{\mathbb C}
\newcommand{\N}{\mathbb N}
\newcommand{\Z}{\mathbb Z}
\newcommand{\eps}{\varepsilon}
\newcommand{\li}{\operatorname{li}}
\newcommand{\Log}{\operatorname{Log}}
\newcommand{\Ei}{\operatorname{Ei}}
\newcommand{\Res}{\operatorname*{Res}}
\newcommand{\Mell}{\ast_{\rm Mel}}
\newcommand{\RVM}{\mathrm{RVM}}
\newcommand{\crit}{\mathrm{crit}}
\newcommand{\pts}{\mathrm{pts}}
\newcommand{\pathc}{\mathrm{path}}
\newcommand{\NR}{\mathcal N_{\RVM}}
\newcommand{\vpi}{\varpi_{\RVM}}
\newcommand{\Zo}{Z1}
\providecommand{\NoCaseChange}[1]{\lowercase{#1}}

\begin{document}

\title[The midpoint RVM synthetic zeta system]{The midpoint Riemann--von Mangoldt synthetic zeta system:\\
canonical products, exponential sums, prime counts, Mertens transforms,\\
and the transformation to $\zeta_{Z1}$}
\author{Rotkin Aynasac\\[2pt]\NoCaseChange{\normalfont\texttt{rotkivanyasac@gmail.com}}}
\date{4 October 2026}

\begin{abstract}
We study the zeta-like system generated by the midpoint inverse of the smooth Riemann--von Mangoldt counting law.  Its nontrivial zeros are $\tfrac12\pm i\tau_k$, where $\tau_k$ is the midpoint ordinate of the clock.  We prove normal convergence of the canonical product, the functional symmetry, the exact zero divisor, a uniform logarithmic bound for the exponential sums of the $\tau_k$, compact-uniform existence of the explicit-formula prime-power count with an explicit branch of $\Ei$, exact finite M\"obius inversion to a synthetic prime count $\pi_{\RVM}$, and the unconditional square-root estimate $\pi_{\RVM}(x)-\li(x)=O(\sqrt x(\log\log x)^2/\log x)$; with the prime number theorem this gives $\li/\pi_{\RVM}\to1$, $\pi/\li\to1$ and $\pi/\pi_{\RVM}\to1$.  The transition theorem $\zeta_{\RVM}=\varpi_{\RVM}(1+w)(1+O(1/|t|))$ is proved in full; it gives a constant vertical limit of $1/\zeta_{\RVM}$ on every line to the right of $\tfrac12$, excludes absolutely convergent Dirichlet expansions of $\zeta_{\RVM}^{\pm1}$, and proves the unsmoothed synthetic Mertens bound $M_{\RVM}^{\crit}(x)=L+O_\eps(x^{1/2+\eps})$, $L=1/\varpi_{\RVM}$.  A contour-free form of the same decomposition shows that $\zeta_{\RVM}$ is bounded above and below on every half-plane $\Re s\ge\sigma_0>1$ and tends to $\varpi_{\RVM}$ as $\Re s\to+\infty$.

The core of the method is a transformation of $\zeta_{\RVM}$ depending on a parameter $\lambda\in[0,1]$: the ordinates move along $\gamma_k(\lambda)=(1-\lambda)\gamma_k+\lambda\tau_k$, where $\gamma_k$ are the ordinates of the zeros of the Riemann zeta function, every member $\zeta_\lambda$ has all its nontrivial zeros on $\Re s=\tfrac12$ by construction, $\zeta_1=\zeta_{\RVM}$, and $\zeta_0=\zeta_{Z1}$.  It is impossible for $\zeta_{Z1}$, or for any $\zeta_\lambda$, to have a zero off the critical line other than the trivial zeros.  All differences between the moving zetas lie in the numbers $\gamma_k-\tau_k$.  The final step of the method is the identity $\zeta_{Z1}=\zeta$.  This identity is the single and complete test: while $\zeta_{Z1}=\zeta$ there is nothing outside $\Re s=\tfrac12$ to speak of, and $\zeta_{Z1}\neq\zeta$ is the only possible obstruction.

We prove Step~3 in the form: if $\zeta_{Z1}=\zeta$, then $M(x)=O_\eps(x^{1/2+\eps})$ for the Mertens function $M$, with a self-contained proof (Borel--Carath\'eodory, three circles, truncated Perron).  By Hamburger's theorem the identity $\zeta_{Z1}=\zeta$ is equivalent to the condition (D) that $\zeta_{Z1}$ is an absolutely convergent Dirichlet series in $\Re s>1$.  We prove three further forms of (D): the value of the endpoint at infinity is $\varpi_{Z1}=\lim_{\sigma\to+\infty}\zeta_{Z1}(\sigma)=\varpi_{\RVM}\prod_k(\frac14+\tau_k^2)/(\frac14+\gamma_k^2)\ge1$, and (D) holds exactly when $\varpi_{Z1}=1$, that is, when $\prod_k(\frac14+\gamma_k^2)/(\frac14+\tau_k^2)=\varpi_{\RVM}$; the prime-power counts of $\zeta_{\RVM}$ and $\zeta_{Z1}$ are the Mellin transforms of their logarithms and satisfy a distributional explicit formula, which proves the Mellin--Stieltjes identification for the integer weights; and (D) holds exactly when the prime-power count of $\zeta_{Z1}$ is constant between consecutive integers, an identity on each cell between the displacements $\gamma_k-\tau_k$ and the clock.  Along the family, $\log\varpi_\lambda$ has derivative $-\Upsilon(\lambda)$, the generator at infinity, and (D) holds exactly when $\int_{0^+}^1\Upsilon\,d\lambda=\log(1/\varpi_{\RVM})$.  Finally we atomize every member of the family by smooth cells: the cell masses have mean square $X^{1+\eps}$ on every dyadic block, uniformly in $\lambda$, so every atomized function $\tilde\zeta_{\Z,\lambda}=\zeta_\lambda e^{\tilde h_\lambda}$, the endpoint $\lambda=0$ included, is an absolutely convergent Dirichlet series in $\Re s>1$; the joining factor $\tilde h_\lambda$ is written cell by cell and is holomorphic in $\Re s>\frac12$, and (D) is reduced to a single statement at $\lambda=0$, the vanishing of every cell defect, which is the remaining step of the programme.  We prove further that no member $\zeta_\lambda$ with $0<\lambda\le1$ is a Dirichlet series, that $\lambda\mapsto\log\zeta_\lambda(\sigma)$ is strictly convex for every real $\sigma>1$, that Hamburger's theorem applies to the pair $(\zeta_\lambda,\tilde\zeta_{\Z,\lambda})$ as soon as the joining factor is symmetric under $s\mapsto1-s$, and that (D) holds exactly when $\zeta_{Z1}$ is the Mellin--Stieltjes transform of a positive measure, in which case the functional equation gives $\varpi_{Z1}\le1$ by an extremal summation argument; at the clock end the density of $d\pi_{\RVM}$ is $\frac1{\log y}+2\cos(2\pi y)-2^{-1/2}y^{-1/4}\cos(4\pi\sqrt y+\frac\pi4)+O(y^{-1/3})$.  The identity $\zeta_{Z1}=\zeta$, equivalently (D), is not proved in this paper; it is the open step of the method.
\end{abstract}

\maketitle
\markboth{ROTKIN AYNASAC}{THE MIDPOINT RVM SYNTHETIC ZETA SYSTEM}

\tableofcontents

\section*{Notation}
\label{sec:notation}

Every symbol below has a single meaning throughout the paper.  The last column links to the defining equation; hovering over the link shows the formula.  Symbols introduced inside a single proof (indices $j,k_0,m$, auxiliary variables $v,y$, thresholds $t_0,t_1,t_2$, radii $r_1,r_2,r_3$) are local to that proof and do not appear elsewhere.

\medskip
{\small
\begin{longtable}{>{\raggedright\arraybackslash}p{0.27\textwidth}>{\raggedright\arraybackslash}p{0.53\textwidth}>{\raggedright\arraybackslash}p{0.12\textwidth}}
\toprule
Symbol & Meaning & Defined\\
\midrule
\endfirsthead
\toprule
Symbol & Meaning & Defined\\
\midrule
\endhead
\bottomrule
\endfoot
\multicolumn{3}{l}{\emph{Variables and the Riemann zeta function}}\\
$s=\sigma+it$ & complex variable, $\sigma=\Re s$, $t=\Im s$ & --\\
$z=s-\tfrac12=\sigma'+it$ & centred variable, $\sigma'=\sigma-\tfrac12$ & \CGeqref{eq:transition-zw}\\
$q(s)=s(s-1)$ & symmetric quadratic, $q(1-s)=q(s)$ & \CGeqref{eq:q-def}\\
$\zeta$, $\rho$ & Riemann zeta function; $\rho$ runs over its nontrivial zeros with multiplicity & \S\CGref{sec:clock}\\
$\gamma_k$ & $k$th positive ordinate of the nontrivial zeros of $\zeta$, with multiplicity, nondecreasing & \CGeqref{eq:zero-count}\\
$N_\zeta(T)$, $N_\zeta^{*}(T)$ & zero count up to height $T$; its symmetric version & \CGeqref{eq:zero-count}, \CGeqref{eq:symmetric-count}\\
$S(T)$ & remainder in the Riemann--von Mangoldt formula & \CGeqref{eq:RvM-formula}\\
$G(s)$ & $s(s-1)\Gamma(s/2)\pi^{-s/2}$ & \CGeqref{eq:Gdef}\\
$P_\zeta=2\xi$ & completed product $G\zeta$ & \CGeqref{eq:Pzeta}\\
$\mu$, $M(x)$ & M\"obius function; Mertens function $\sum_{n\le x}\mu(n)$ & \CGeqref{eq:mertens-def}\\
$\pi(x)$, $\li$ & prime-counting function; logarithmic integral & --\\
$\chi(s)$, $\vartheta$ & functional-equation factor; Riemann--Siegel theta function & \CGeqref{eq:chi}, \CGeqref{eq:theta-def}\\
\midrule
\multicolumn{3}{l}{\emph{The midpoint clock}}\\
$\NR(T)$ & midpoint Riemann--von Mangoldt clock & \CGeqref{eq:clock}\\
$\tau_k$ & midpoint ordinates, $\NR(\tau_k)=k-\tfrac12$ & \CGeqref{eq:midpoint-condition}\\
$W_0$ & principal branch of Lambert's $W$-function & \CGeqref{eq:lambert}\\
$\mathcal T$ & inverse function of the clock & \CGeqref{eq:inverse-clock}\\
$\ell_k$ & $\log(\tau_k/2\pi)$ & \CGeqref{eq:rvm-gap-bounds}\\
$n_{\RVM}(H)$ & number of $\tau_k\le H$ & \CGeqref{eq:n-RVM}\\
$\eta_k$, $C_\eta$ & $\eta_k=\tau_k-\gamma_k$; a bound $|\eta_k|\le C_\eta$ & \CGeqref{eq:eta-def}, \CGeqref{eq:bounded-displacement}\\
\midrule
\multicolumn{3}{l}{\emph{The synthetic zeta function and its sums}}\\
$P_{\RVM}^{\crit}$, $\xi_{\RVM}$, $\zeta_{\RVM}$ & canonical product, completed and synthetic zeta & \CGeqref{eq:Pcrit}, \CGeqref{eq:zetaRVM}\\
$\rho_k^{\RVM}$ & synthetic zero $\tfrac12+i\tau_k$ & \CGeqref{eq:synthetic-divisor}\\
$S_{\RVM}(N;u)$, $C_{\RVM}^{\pts}$, $C_{\RVM}^{\pathc}$ & exponential sum; point and path centroids & \CGeqref{eq:S-RVM}, \CGeqref{eq:path-centroid}\\
$u_1,u_2$ & a frequency interval $[u_1,u_2]\subset(0,\infty)$ & \CGeqref{eq:exp-sum-bound}\\
$M_H(x)$, $M_H^{\RVM}(x)$ & finite-height phase averages (not Mertens functions) & \CGeqref{eq:MH-genuine}, \CGeqref{eq:MH-RVM}\\
$\Ei$ & exponential integral on $\Im w\ne0$ & \CGeqref{eq:Ei-def}\\
$J_{\RVM}$, $J_{\RVM}^{*}$, $\pi_{\RVM}$ & explicit-formula prime-power count; truncation; synthetic prime count & \CGeqref{eq:J-RVM}, \CGeqref{eq:J-star}, \CGeqref{eq:pi-RVM}\\
$\mathcal A_{\RVM}(u)$, $\ell(t)$ & weighted clock sum $\sum_ke^{iu\tau_k}/\tau_k$; $\log(t/2\pi)$ & \CGeqref{eq:A-RVM}\\
$\mathfrak s_1$; $\mathfrak e$, $\mathfrak e_1$ & primitive of the midpoint sawtooth; $N_\zeta-\NR$ and its primitive & \CGeqref{eq:saw-primitive}, \CGeqref{eq:frak-e}\\
$\varpi_{Z1}$ & value of $\zeta_{Z1}$ at $+\infty$ & \CGeqref{eq:varpi-Z1}\\
$J_X^{\rm M}$, $J_X$, $C_X$ & Mellin and explicit-formula prime-power counts of $\zeta_X$, $X\in\{\RVM,Z1\}$; their difference & \CGeqref{eq:mellin-count}, \CGeqref{eq:J-X}\\
$\mathsf b_n$ & jumps of $J_{Z1}^{\rm M}$ at the integers & Thm.~\CGref{thm:cell-form}\\
$\eta_k$; $\varpi_\lambda$; $\Upsilon(\lambda)$ & $\tau_k-\gamma_k$; value of $\zeta_\lambda$ at $+\infty$; flow at infinity & \CGeqref{eq:varpi-lambda}, \CGeqref{eq:Upsilon}\\
$\varphi_n$, $\tilde c_n(\lambda)$, $\tilde a_N(\lambda)$ & smooth cells; atomized masses and coefficients & \CGeqref{eq:atomized-family}, \CGeqref{eq:atomized-coefficients}\\
$\tilde\zeta_{\Z,\lambda}$, $\tilde h_\lambda$ & atomized family; joining factor & \CGeqref{eq:atomized-family}, \CGeqref{eq:joining-factor}\\
$\mathfrak a_s$, $\mathfrak h_n(\lambda;s)$, $\mathfrak L$ & atomization kernel; cell defects; $\int_1^x(1-y^{-1})dy/\log y$ & \CGeqref{eq:atomization-kernel}\\
$V_n$, $W_n$ & ordinate kernels of the cell $n$ & \CGeqref{eq:VW}\\
$V_n$, $W_n$ & ordinate kernels of the cell $n$ & \CGeqref{eq:VW}\\
$\tilde\mu_\lambda$, $\tilde M_\lambda$ & atomized M\"obius coefficients; summatory function & \CGeqref{eq:atomized-mobius}\\
$B_{Z1}$ & Mellin form of the endpoint Mertens function & \CGeqref{eq:Mellin-Mertens}\\
$c_n^{\RVM}$, $\zeta_{\Z,\RVM}$ & integer weights; weighted Euler product & \CGeqref{eq:cn}, \CGeqref{eq:weighted-zeta}\\
$J_{\Z,\RVM}$, $J_{\RVM}^{*,\Z}$, $h_{\RVM}$ & integer-step counts; Mellin correction & \CGeqref{eq:JZ}, \CGeqref{eq:J-star-Z}, \CGeqref{eq:hRVM}\\
$\theta$ & unit-cell step, $0<\theta\le1$ & \CGeqref{eq:unit-cell}\\
$p_j^{\RVM}$ & real RVM prime levels $\pi_{\RVM}^{-1}(j)$ on monotone branches & \CGeqref{eq:intro-real-primes}\\
$\mathcal M(n)$, $\varsigma_n$, $\mu_h$, $M_h$ & target samples of $M_{\RVM}^{\crit}$; signs; ternary walk and its sum (companion object; $h$ is a label) & \CGeqref{eq:intro-muh}, \CGeqref{eq:intro-quantizer}\\
$a_N^{\Z}$ & Dirichlet coefficients of $\zeta_{\Z,\RVM}$ & \CGeqref{eq:ZRVM-Dirichlet}\\
$C_J$ & Lipschitz constant of $J_{\RVM}$ & \CGeqref{eq:unit-cell}\\
$\mathfrak d_n^{\RVM}$ & prime defect $\mathbf 1_{\rm prime}(n)-c_n^{\RVM}$ & \CGeqref{eq:prime-defect}\\
\midrule
\multicolumn{3}{l}{\emph{Mertens transforms and the transition theorem}}\\
$F_{\RVM}=\widehat M_{\RVM}^{\crit}$ & $1/\zeta_{\RVM}$ & \CGeqref{eq:F-RVM}\\
$c$, $U$ & Perron abscissa; truncation height & \CGeqref{eq:Mcrit-symmetric}\\
$M_{\RVM}^{\crit}$, $M_{\RVM,G}^{\crit}$ & symmetric Perron inverse; Gaussian-smoothed inverse & \CGeqref{eq:Mcrit-symmetric}, \CGeqref{eq:MG}\\
$\tau_*$, $w(s)$ & transition point $-iz=t-i\sigma'$; transition variable $e^{-2\pi i\NR(\tau_*)}$ & \CGeqref{eq:transition-zw}\\
$B$; $\sigma'_1,\sigma'_2$ & upper bound for $\sigma'$; a range $[\sigma'_1,\sigma'_2]\subset(0,\tfrac12)$ & \CGeqref{eq:transition-w}, \CGeqref{eq:transition-chi}\\
$\mathcal E_0,\mathcal E$ & error factors of the transition theorem & \CGeqref{eq:transition-w}, \CGeqref{eq:transition-chi}\\
$\varpi_{\RVM}$, $L$ & transition constant $0.9619415776\ldots$; $L=1/\varpi_{\RVM}$ & \CGeqref{eq:varpi-def}, \CGeqref{eq:vertical-limit}\\
$f,\Lambda,\Lambda_{\rm sm},\Lambda_{\rm saw}$ & logarithm of the product and its smooth and sawtooth parts (\S\CGref{sec:transition} only) & \CGeqref{eq:transition-product}, \CGeqref{eq:transition-split}\\
$\psi,\psi_r,\Phi_r,\nu_1,\dots,\nu_4$ & sawtooth, Abel means, ray remainder, real constants (\S\CGref{sec:transition} only) & \CGeqref{eq:sawtooth-count}, \CGeqref{eq:psi-r}\\
$\alpha_c$ & $\min(c-\tfrac12,1)$ & \CGeqref{eq:FminusL-all}\\
\midrule
\multicolumn{3}{l}{\emph{Comparison with $\zeta$}}\\
$\mathcal Q_{\RVM}$, $\mathcal Q_{\RVM}^{\vee}$ & quotient $\zeta/\zeta_{\RVM}$; its inverse Mellin transform & \CGeqref{eq:Q-RVM}, \CGeqref{eq:M-convolution}\\
\midrule
\multicolumn{3}{l}{\emph{The $\lambda$-family and the endpoint $\zeta_{Z1}$}}\\
$\lambda$ & family parameter; $\lambda=1$: $\zeta_{\RVM}$, $\lambda=0$: $\zeta_{Z1}$ & \CGeqref{eq:gamma-lambda}\\
$\gamma_k(\lambda)$ & moving ordinates $(1-\lambda)\gamma_k+\lambda\tau_k$ & \CGeqref{eq:gamma-lambda}\\
$P_\lambda$, $\zeta_\lambda$, $F_\lambda$ & family product, family zeta, $1/\zeta_\lambda$ & \CGeqref{eq:P-lambda}, \CGeqref{eq:zeta-lambda}\\
$A_{k,\lambda}$ & $\tfrac14+\gamma_k(\lambda)^2$ & \CGeqref{eq:generator}\\
$\mathcal V_\lambda$ & infinitesimal generator $\partial_\lambda\log P_\lambda$ & \CGeqref{eq:generator}\\
$\Pi_{\rm tr}$ & total multiplier $P_{Z1}/P_{\RVM}^{\crit}$ & \CGeqref{eq:transport-multiplier}\\
$\rho_{k,\lambda}$, $\mathfrak r_{k,\lambda}$ & moving zero $\tfrac12+i\gamma_k(\lambda)$; residue of $1/(s\zeta_\lambda)$ & \CGeqref{eq:residue-transport}\\
$P_{Z1}$, $\zeta_{Z1}$, $F_{Z1}$ & endpoint product, endpoint zeta, $1/\zeta_{Z1}$ & \CGeqref{eq:Z1-product}, \CGeqref{eq:FZ1}\\
$\mathscr R_{Z1}$ & ratio $\zeta_{Z1}/\zeta$ & \CGeqref{eq:Z1-defect-ratio}\\
$\mathcal K$ & multiplier $\zeta_{Z1}/\zeta_{\Z,\RVM}$ & \CGeqref{eq:K-multiplier}\\
(D), $a_n$ & Dirichlet condition for $\zeta_{Z1}$ and its coefficients & \CGeqref{eq:Z1-D}\\
$C_H$ & constant in Hamburger's theorem & \CGeqref{eq:Hamburger-conclusion}\\
$\mu_{Z1}$, $M_{Z1}$ & coefficients of $F_{Z1}$; summatory function & \CGeqref{eq:Z1-reciprocal-Dirichlet}\\
$M_{\lambda,U}$, $M_{Z1,U}$ & finite-height Perron integrals & \CGeqref{eq:truncated-Perron}\\
$\delta$, $\eps$, $\omega(\delta)$ & margins $\sigma\ge\tfrac12+\delta$, exponent $\eps>0$, three-circle exponent & \CGeqref{eq:three-circles-exponent}\\
\end{longtable}
}

\medskip\noindent\textbf{Collision check.}
The table was checked symbol by symbol against the whole text: no symbol carries two meanings.  In particular $c$ denotes only a Perron abscissa and $c_n^{\RVM}$ only the integer weights; $\NR$ is the clock and $n_{\RVM}$ the count of the $\tau_k$; $\mathcal V_\lambda$ is the generator of the family and $\Phi_r$ a remainder inside one proof; $\eps$ is only an exponent; $\theta$ is the unit-cell step and $\vartheta$ the Riemann--Siegel function; $E_k=e^{iu\tau_k}$ while $\mathcal E_0,\mathcal E$ are error factors; (D) is only the Dirichlet condition; $T$ is a height, and the target of the ternary walk is $\mathcal M(n)$ with signs $\varsigma_n$; the subscript $h$ of $\mu_h,M_h$ is a label, distinct from $h_{\RVM}$.


\section{Introduction}
\label{sec:introduction}

\subsection{Background}
\label{subsec:background}

The theory of the Riemann zeta function joins four structures that are usually introduced in different languages.  On the spectral side there are the nontrivial zeros and their Riemann--von Mangoldt counting law; on the analytic side the completed function $\xi$, its functional equation and its Hadamard product; on the prime side the prime-power measure and the explicit formula; and on the inverse side the reciprocal Dirichlet series $1/\zeta$ and the Mertens function
\begin{equation}\label{eq:mertens-def}
 M(x):=\sum_{n\le x}\mu(n).
\end{equation}
Riemann's memoir made the first three structures inseparable, and the work of von Mangoldt, Hadamard and Hamburger supplied, respectively, the rigorous zero count and explicit formula, the entire-product factorization, and a converse characterization of $\zeta$ from its Dirichlet-series and functional-equation data \cite{Riemann1859,vonMangoldt1895,Hadamard1896,Hamburger1921,Titchmarsh1986,Edwards1974}.  This paper asks what remains of that architecture when the zero staircase is replaced, at the outset, by the exact midpoint inverse of its smooth Riemann--von Mangoldt clock, and how the resulting object is carried back to $\zeta$.

\subsection{The midpoint RVM system}
\label{subsec:midpoint-system}

The point of departure is the smooth counting function \CGeqref{eq:clock},
\begin{equation}\label{eq:intro-clock}
 \NR(T)=\frac{T}{2\pi}\log\frac{T}{2\pi e}+\frac78,
\end{equation}
and the real ordinates $\tau_k$ determined by the midpoint condition $\NR(\tau_k)=k-\tfrac12$.  The half-integer convention is the midpoint convention for a staircase and leads to the exact Lambert inversion \CGeqref{eq:lambert},
\begin{equation}\label{eq:intro-lambert}
 \tau_k=\frac{2\pi(k-11/8)}{W_0\bigl((k-11/8)/e\bigr)} .
\end{equation}
Lambert-$W$ approximations to individual zero ordinates and smooth inversions of the Riemann--von Mangoldt law occur in the literature, including explicit equations for the $n$th zero \cite{FrancaLeClair2014,Corless1996,DLMF2024}.  Here the midpoint inverse is treated as the exact spectrum of a deterministic system and is propagated, without the fluctuating term $S(T)$, through canonical products, explicit prime counting and Mertens inversion.

The first object is the canonical product and the synthetic zeta function
\begin{equation}\label{eq:intro-zeta-rvm}
 P_{\RVM}^{\crit}(s)=\prod_{k\ge1}\left(1+\frac{s(s-1)}{\tfrac14+\tau_k^2}\right),
 \qquad
 \zeta_{\RVM}(s)=\frac{P_{\RVM}^{\crit}(s)}{G(s)},
 \qquad G(s)=s(s-1)\Gamma(s/2)\pi^{-s/2}.
\end{equation}
Because $\tau_k\asymp k/\log k$, the factors form a normally convergent product in the symmetric variable $s(s-1)$; $P_{\RVM}^{\crit}$ is entire and invariant under $s\mapsto1-s$, and $\zeta_{\RVM}$ has a residue-one pole at $1$, the trivial zeros, and the simple nontrivial zeros $\rho_k^{\RVM}=\tfrac12\pm i\tau_k$ (Theorem~\CGref{thm:canonical}).  The real part $\tfrac12$ is part of the definition.

\emph{Generalized prime systems.}  In Beurling's theory the generalized primes are arbitrary reals $1<p_1\le p_2\le\cdots\to\infty$; their multiplicative semigroup defines generalized integers and a zeta function by a Dirichlet series, an Euler product or a Mellin--Stieltjes transform, separating prime distribution from the additive structure of the integers \cite{Beurling1937,BatemanDiamond1969,Diamond1970}.  Beurling systems with prescribed analytic behaviour, critical zeros and large oscillation show that zeta-like zero data alone do not determine the arithmetic realization \cite{DiamondMontgomeryVorhauer2006,BrouckeDebruyneVindas2020}.  The RVM system runs in the opposite direction: it starts from a prescribed zero clock, constructs an entire completed product first, obtains a continuous prime-power count from its explicit formula, and only afterwards asks for an atomic realization on the integers.

\emph{Continuous and discrete prime systems.}  The passage from a continuous prime system to a discrete one is classical in the theory of Beurling generalized primes \cite{Beurling1937,DiamondMontgomeryVorhauer2006,BrouckeVindas2024,DiamondZhang2016}: the atoms are placed at suitably chosen real points, and the discretization changes the zeta function by a factor holomorphic and zero-free in $\Re s>\frac12$ (compare Theorem~\CGref{thm:atomized-continuation}(iii)).  Measures with a continuous part and non-integer multiplicities also appear in the theory of complex dimensions \cite[Ch.~3]{LapidusvanFrankenhuysen2000}.  Grosswald and Schnitzer \cite{GrosswaldSchnitzer1978} replaced each prime $p_n$ by a real $q_n$ with $p_n\le q_n\le p_{n+1}$ and proved that the modified Euler product is $\zeta$ times a factor holomorphic and zero-free in $\Re s>0$, so that it has the same zeros as $\zeta$ there; Corollary~\CGref{cor:weighted-factorization} and Theorem~\CGref{thm:atomized-continuation}(iv) are of the same type for $\zeta_{\RVM}$.  Here the masses of $d\pi_\lambda$ are placed at the integers, with real exponents $\tilde c_n(\lambda)$; this yields an ordinary Dirichlet series for every $\lambda$, absolutely convergent in $\Re s>1$ (Corollary~\CGref{cor:abscissa-one}), and makes Hamburger's theorem \cite{Hamburger1921,Siegel1922} available at the endpoint.  Hamburger's theorem and its extensions \cite{BochnerChandrasekharan1956,KaczorowskiPerelli1999} need absolute convergence in $\Re s>1$, and Knopp \cite{Knopp1994} showed that a half-plane of absolute convergence further to the right does not suffice.  We are not aware of an earlier use of this combination.

\emph{Prime coordinate.}  From the RVM zeros we define the prime-power count $J_{\RVM}$ by \CGeqref{eq:J-RVM}, with the branch \CGeqref{eq:Ei-def} of $\Ei$.  Finite M\"obius inversion in the exponent variable, which relates Riemann's prime-power count to $\pi(x)$ \cite{Riemann1859,Titchmarsh1986,Edwards1974}, is used here as the \emph{definition} of the synthetic prime count \CGeqref{eq:pi-RVM},
\begin{equation}\label{eq:intro-pi}
 \pi_{\RVM}(x)=\sum_{1\le m\le\log_2x}\frac{\mu(m)}m\,J_{\RVM}^{*}(x^{1/m}),
 \qquad
 J_{\RVM}^{*}(x)=\sum_{1\le m\le\log_2x}\frac1m\,\pi_{\RVM}(x^{1/m}).
\end{equation}
The function $\pi_{\RVM}$ is continuous, not a staircase.  It supports two notions of RVM prime: the continuous measure $d\pi_{\RVM}$, and, on every interval where $\pi_{\RVM}$ is continuous and strictly increasing, the real level crossings
\begin{equation}\label{eq:intro-real-primes}
 p_j^{\RVM}:=\pi_{\RVM}^{-1}(j),\qquad \pi_{\RVM}(p_j^{\RVM})=j,
\end{equation}
which need not be integers.  The synthetic count has square-root accuracy without any hypothesis on $\zeta$: $J_{\RVM}(x)-\li(x)$ and $\pi_{\RVM}(x)-\li(x)$ are $O(\sqrt x(\log\log x)^2/\log x)$ (Theorem~\CGref{thm:Lambert-estimate}).  The proof reduces the zero sum to $\sum_ke^{iu\tau_k}/\tau_k$, which is $O(\log^2u)$ by the second-derivative test on dyadic blocks of the clock (Lemma~\CGref{lem:A-RVM}).  Hence $\li(x)/\pi_{\RVM}(x)\to1$, and with the prime number theorem $\pi(x)/\li(x)\to1$ and $\pi(x)/\pi_{\RVM}(x)\to1$ (Corollary~\CGref{cor:three-limits}, Figure~\CGref{fig:pi-ratios}).  Sampling on unit cells gives the masses $c_n^{\RVM}=\pi_{\RVM}(n)-\pi_{\RVM}(n-1)$ and the weighted Euler object $\zeta_{\Z,\RVM}$ of \CGeqref{eq:weighted-zeta}.  The masses are not $0$ or $1$ and need not be positive, but they are bounded: $J_{\RVM}$ is Lipschitz on $[2,\infty)$ (Theorem~\CGref{thm:unit-cell}), so $c_n^{\RVM}=O(1)$ and $\zeta_{\Z,\RVM}$ is an absolutely convergent Dirichlet series in $\Re s>1$ (Corollaries~\CGref{cor:cn-bounded} and~\CGref{cor:ZRVM-Dirichlet}).  The unit-periodic first harmonic of the clock integrates to zero on every integer cell (Lemma~\CGref{lem:cell-cancel}); under the cell hypothesis one obtains $c_n^{\RVM}\sim1/\log n$ (Proposition~\CGref{prop:eventual-positive}).

\emph{Mertens coordinate.}  The Mertens coordinate is the symmetric inverse Mellin object of $F_{\RVM}=1/\zeta_{\RVM}$.  A smooth/sawtooth Stieltjes decomposition of the clock product with Abel regularization of the sawtooth gives on every band to the right of $\tfrac12$
\begin{equation}\label{eq:intro-transition}
 \zeta_{\RVM}(s)=\varpi_{\RVM}\bigl(1+w(s)\bigr)\bigl(1+O(1/|t|)\bigr),
 \qquad w(s)=\exp\bigl\{-2\pi i\,\NR\bigl(-i(s-\tfrac12)\bigr)\bigr\},
\end{equation}
and $\zeta_{\RVM}=\varpi_{\RVM}(1+\chi)(1+O(1/|t|))$ on closed substrips of $\tfrac12<\Re s<1$ (Theorem~\CGref{thm:transition}).  Hence $F_{\RVM}(\sigma+it)\to L=1/\varpi_{\RVM}$ for every fixed $\sigma>\tfrac12$; neither $\zeta_{\RVM}$ nor $F_{\RVM}$ is an absolutely convergent Dirichlet series in any half-plane (Theorem~\CGref{thm:no-Dirichlet-RVM}); and after subtracting the vertical limit the Perron inverse converges absolutely, with
\begin{equation}\label{eq:intro-Mbound}
 M_{\RVM}^{\crit}(x)=L+O_\eps(x^{1/2+\eps})
\end{equation}
(Theorem~\CGref{thm:unsmoothed-RVM}).  The RVM Mertens function therefore exists as a continuous spectral inverse although no coefficient sequence for $F_{\RVM}$ exists.  Since the primitive of the midpoint sawtooth is bounded, the same decomposition without contours gives $C^{-1}\le|\zeta_{\RVM}|\le C$ on every half-plane $\Re s\ge\sigma_0>1$ and $\zeta_{\RVM}(s)\to\varpi_{\RVM}$ as $\Re s\to+\infty$ (Theorem~\CGref{thm:halfplane-RVM}).

\emph{Ternary discretization.}  An integer-valued representative of $M_{\RVM}^{\crit}$ is obtained by sampling a smooth target $\mathcal M(n)\approx M_{\RVM}^{\crit}(n)$ and choosing signs on the squarefree support:
\begin{equation}\label{eq:intro-muh}
 \mu_h(n)=\mu(n)^2\varsigma_n,\qquad\varsigma_n\in\{-1,+1\},\qquad M_h(n):=\sum_{m\le n}\mu_h(m),
\end{equation}
\begin{equation}\label{eq:intro-quantizer}
 \varsigma_n=\begin{cases}+1,&M_h(n-1)<\mathcal M(n),\\-1,&M_h(n-1)\ge\mathcal M(n),\end{cases}\qquad(\mu(n)^2=1).
\end{equation}
This is a greedy sigma--delta feedback quantization; feedback quantization and noise shaping are standard themes of sigma--delta theory \cite{DaubechiesDeVore2003,GunturkThao2003,BenedettoPowellYilmaz2006}.  Here the target is a zero-generated Mertens inverse and the admissible times are the squarefree integers.  The walk $\mu_h$ is a comparison object: it is not claimed to be multiplicative or equal to $\mu$.  The two discretizations are distinct: the cell masses $c_n^{\RVM}$ come from $\pi_{\RVM}$, the signs $\varsigma_n$ track $M_{\RVM}^{\crit}$.

\emph{Comparison with $\zeta$.}  At the level of transforms the defect is the exact quotient $\mathcal Q_{\RVM}=\zeta/\zeta_{\RVM}$ of \CGeqref{eq:Q-RVM}, with $F_{\RVM}=\mathcal Q_{\RVM}\cdot(1/\zeta)$.  Since $|\zeta_{\RVM}|$ is bounded above and below on $\frac12+\eps\le\sigma\le c'$, $|t|\ge t_0$ (Corollary~\CGref{cor:Q-control}), every bound for $\zeta$ or $1/\zeta$ to the right of $\frac12$ is equivalent to the same bound for $\mathcal Q_{\RVM}$.  For $M$ itself the bound $M(x)=O_\eps(x^{1/2+\eps})$ is equivalent to the Riemann hypothesis \cite[\S14.25]{Titchmarsh1986}, and the disproof of the Mertens conjecture shows that the $\eps$ cannot simply be dropped \cite{OdlyzkoTeRiele1985}.

\subsection{The core of the method}
\label{subsec:core}

The transfer from $\zeta_{\RVM}$ to $\zeta$ is a transformation of $\zeta_{\RVM}$ depending on a parameter $\lambda\in[0,1]$.  The ordinates move along
\begin{equation}\label{eq:intro-family}
 \gamma_k(\lambda)=(1-\lambda)\gamma_k+\lambda\tau_k,
 \qquad
 \zeta_\lambda(s)=\frac1{G(s)}\prod_{k\ge1}\left(1+\frac{s(s-1)}{\tfrac14+\gamma_k(\lambda)^2}\right),
\end{equation}
where $\gamma_k$ are the ordinates of the nontrivial zeros of $\zeta$; thus $\zeta_1=\zeta_{\RVM}$ and $\zeta_0=\zeta_{Z1}$ (Section~\CGref{sec:lambda-family}).
\begin{enumerate}[label=\textup{(\roman*)}]
\item Every member $\zeta_\lambda$ has its nontrivial zeros exactly at $\tfrac12\pm i\gamma_k(\lambda)$.  The real part $\tfrac12$ is part of the definition of the family, like the gamma factor, and never changes along it: it is impossible for $\zeta_{Z1}$, or for any $\zeta_\lambda$, to have a zero off the critical line other than the trivial zeros.  For $0<\lambda\le1$ all zeros are simple.
\item All differences between the moving zetas lie in the numbers $\gamma_k-\tau_k$, which are bounded (Lemma~\CGref{lem:bounded-displacement}); the quotient of two members is an explicit product in $\gamma_k$ and $\lambda(\tau_k-\gamma_k)$ (Remark~\CGref{rem:only-gamma-tau}).
\item The final step is the identity $\zeta_{Z1}=\zeta$.  It is the single and complete test: while $\zeta_{Z1}=\zeta$ there is nothing outside $\Re s=\tfrac12$ to speak of, since the identity itself places every nontrivial zero of $\zeta$ at $\tfrac12\pm i\gamma_k$, and $\zeta_{Z1}\ne\zeta$ is the only possible obstruction.
\end{enumerate}
The endpoint has the functional equation, the pole, the residue and the order of $\zeta$ (Proposition~\CGref{prop:Z1-class}).  Hamburger's theorem belongs to the family of converse theorems that recover Dirichlet series from analytic continuation, functional equation and growth \cite{Hamburger1921,Perelli2016}; applied to $\zeta_{Z1}$ it gives (Theorem~\CGref{thm:Hamburger-Z1})
\begin{equation}\label{eq:intro-hamburger}
 \zeta_{Z1}=\zeta
 \quad\Longleftrightarrow\quad
 \text{(D): }\zeta_{Z1}\text{ is an absolutely convergent Dirichlet series in }\Re s>1 .
\end{equation}
With the integer-weighted zeta $\zeta_{\Z,\RVM}$, which is a Dirichlet series, (D) is equivalent to an absolutely convergent Dirichlet expansion of the explicit multiplier $\mathcal K=\zeta_{Z1}/\zeta_{\Z,\RVM}=\Pi_{\rm tr}e^{-h_{\RVM}}$ (Theorem~\CGref{thm:D-multiplier}); since $1/\zeta_{\Z,\RVM}$ is also an absolutely convergent Dirichlet series in $\Re s>1$, the equivalence holds on the whole half-plane.  Section~\CGref{sec:euler-products} gives two sharper forms of (D).  First, the clock decomposition applied to the ordinates $\gamma_k$ gives the value of the endpoint at infinity,
\begin{equation}\label{eq:intro-varpi-Z1}
 \varpi_{Z1}=\lim_{\sigma\to+\infty}\zeta_{Z1}(\sigma)=\varpi_{\RVM}\prod_{k\ge1}\frac{\frac14+\tau_k^2}{\frac14+\gamma_k^2}\ \ge1,
\end{equation}
and $\zeta_{Z1}=\zeta$ holds if and only if $\varpi_{Z1}=1$ (Theorems~\CGref{thm:Z1-clock} and~\CGref{thm:varpi-Z1}, Corollary~\CGref{cor:product-identity}): the final step is the single identity $\prod_k(\frac14+\gamma_k^2)/(\frac14+\tau_k^2)=\varpi_{\RVM}$ between the displacements and the transition constant.  Second, the prime-power counts $J_{\RVM}^{\rm M}$ and $J_{Z1}^{\rm M}$ are the Mellin transforms of $\log\zeta_{\RVM}$ and $\log\zeta_{Z1}$ and satisfy a distributional explicit formula (Theorems~\CGref{thm:mellin-counts} and~\CGref{thm:distributional-EF}); this proves the Mellin--Stieltjes identification of the integer weights (Theorem~\CGref{thm:MS}), and (D) holds if and only if $J_{Z1}^{\rm M}$ is constant between consecutive integers (Theorem~\CGref{thm:cell-form}).  Third, along the family $\frac d{d\lambda}\log\varpi_\lambda=-\Upsilon(\lambda)$, the generator at infinity, and (D) is the identity $\int_{0^+}^1\Upsilon(\lambda)\,d\lambda=\log L$ (Proposition~\CGref{prop:family-flow}).  Section~\CGref{sec:atomized} atomizes every member: $\tilde\zeta_{\Z,\lambda}=\zeta_\lambda e^{\tilde h_\lambda}$ keeps the Euler product and the Dirichlet series while $\zeta_\lambda$ keeps the functional equation; the joining factor $\tilde h_\lambda$ is written cell by cell and is holomorphic in $\Re s>\frac12$ (Theorems~\CGref{thm:joining-factor} and~\CGref{thm:atomized-continuation}); the cell masses have mean square $X^{1+\eps}$ on $(X,2X]$ uniformly in $\lambda$, so every $\tilde\zeta_{\Z,\lambda}$, the endpoint included, is an absolutely convergent Dirichlet series in $\Re s>1$ (Theorem~\CGref{thm:mean-square-cells}, Corollary~\CGref{cor:abscissa-one}); and (D) is equivalent to the vanishing of every cell defect at $\lambda=0$ (Theorem~\CGref{thm:closing}).  At the clock endpoint the joining factor $\tilde h_1$ is not constant, and the endpoint Mertens function has a Mellin form $B_{Z1}$ with $B_{Z1}=M$ exactly when (D) holds (Propositions~\CGref{prop:atomized-clock} and~\CGref{prop:Mellin-Mertens}).  The remaining steps form the programme of Subsection~\CGref{subsec:programme}.  In RVM language, the final step asks whether the zero-generated analytic object $\zeta_{Z1}$ belongs to the discrete Dirichlet category; if it does, Hamburger uniqueness makes it $\zeta$, and the prime and Mertens coordinates become those of the integers.  Step~3 is Theorem~\CGref{thm:step3}:
\begin{equation}\label{eq:intro-step3}
 \zeta_{Z1}=\zeta
 \quad\Longrightarrow\quad
 M(x)=O_\eps\bigl(x^{1/2+\eps}\bigr)\qquad(\eps>0),
\end{equation}
with a self-contained proof by the Borel--Carath\'eodory lemma, Hadamard's three-circle theorem and the truncated Perron formula.  Theorem~\CGref{thm:transfer} gives the same conclusion as a transfer from the synthetic side: under $\zeta_{Z1}=\zeta$, $1/\zeta=F_{\RVM}\,\Pi_{\rm tr}^{-1}$ with $F_{\RVM}$ bounded and $\Pi_{\rm tr}^{\pm1}\ll|t|^{\eps}$ to the right of $\frac12$, so the bound $M_{\RVM}^{\crit}(x)=L+O_\eps(x^{1/2+\eps})$ carries over to $M(x)=O_\eps(x^{1/2+\eps})$.  The identity $\zeta_{Z1}=\zeta$ itself, equivalently (D), is not proved in this paper; it is the open step of the method (Remark~\CGref{rem:Step3-status}).

The relations among the principal objects are
\begin{equation}\label{eq:intro-prime-chain}
 \{\tau_k\}\longrightarrow\zeta_{\RVM}\longrightarrow J_{\RVM}\longleftrightarrow\pi_{\RVM}\longrightarrow\bigl(d\pi_{\RVM},\,p_j^{\RVM},\,c_n^{\RVM}\bigr),
\end{equation}
\begin{equation}\label{eq:intro-M-chain}
 F_{\RVM}=\zeta_{\RVM}^{-1}\longrightarrow M_{\RVM}^{\crit}\longrightarrow(\mathcal M,\mu_h,M_h),
\end{equation}
\begin{equation}\label{eq:intro-family-chain}
 \zeta_{\RVM}=\zeta_1\xrightarrow{\ \lambda\downarrow0\ }\zeta_0=\zeta_{Z1}\overset{?}{=}\zeta\ \Longrightarrow\ M(x)=O_\eps(x^{1/2+\eps}).
\end{equation}
The first chain is the prime coordinate, the second the Mertens coordinate, the third the $\lambda$-family; they share the midpoint spectrum and solve different realization problems.

\subsection{Originality, scope and status}
\label{subsec:status}

The ingredients are standard or belong to neighbouring theories: the Riemann--von Mangoldt law, Lambert inversion, Hadamard products, explicit formulas, M\"obius and Perron inversion, Beurling generalized primes, sigma--delta feedback, and Hamburger's converse theorem.  The contribution is their assembly into one midpoint RVM space with four coordinated levels: (i) an exact midpoint zero clock and its canonical zeta-like product; (ii) a continuous prime-power and prime coordinate, with real prime levels and non-integral cell masses; (iii) a continuous synthetic Mertens inverse with its ternary feedback discretization; (iv) the $\lambda$-family from the midpoint spectrum to the ordinates of $\zeta$, ending in the single identity $\zeta_{Z1}=\zeta$ and the Mertens bound that follows from it.  To the author's knowledge the cited literature contains the individual mechanisms but not this chain.  Standard material is cited at the point of use; statements without citation are original.

\begin{center}
\small
\begin{tabular}{p{0.50\textwidth}p{0.40\textwidth}}
\toprule
Statement & Status\\
\midrule
Midpoint clock, Lambert inversion, bounded displacement & theorem\\
Canonical product $P_{\RVM}^{\crit}$ and $\zeta_{\RVM}$ & theorem\\
$S_{\RVM}(N;u)\ll1+\log N$; centroid and $M_H^{\RVM}$ decay & theorem\\
Compact-uniform existence of $J_{\RVM}$; finite inversion & theorem\\
Real prime levels $p_j^{\RVM}$ & defined on monotone branches of $\pi_{\RVM}$\\
Square-root estimate for $J_{\RVM},\pi_{\RVM}$ & theorem\\
$\li/\pi_{\RVM}\to1$; $\pi/\pi_{\RVM}\to1$ & theorem (with the prime number theorem)\\
$J_{\RVM}$ Lipschitz, $c_n^{\RVM}=O(1)$ & theorem\\
$\zeta_{\Z,\RVM}$ absolutely convergent Dirichlet series in $\Re s>1$ & theorem\\
Mellin--Stieltjes identification; $\zeta_{\Z,\RVM}=\zeta_{\RVM}e^{h_{\RVM}}$ & theorem\\
$|\zeta|\asymp|\mathcal Q_{\RVM}|$ for $\sigma\ge\frac12+\eps$, $|t|\ge t_0$ & theorem\\
Transition theorem $\varpi_{\RVM}(1+w)$, $\varpi_{\RVM}(1+\chi)$ & theorem\\
$C^{-1}\le|\zeta_{\RVM}|\le C$ on $\Re s\ge\sigma_0>1$; $\zeta_{\RVM}\to\varpi_{\RVM}$ at $+\infty$ & theorem\\
Unsmoothed $M_{\RVM}^{\crit}=L+O_\eps(x^{1/2+\eps})$ & theorem\\
Dirichlet expansion of $\zeta_{\RVM}^{\pm1}$ & impossible (theorem)\\
Ternary walk $\mu_h$, $M_h$ & definition; comparison object\\
Exact Hadamard product of $2\xi$ & theorem\\
$\lambda$-family, generator, multiplier, simplicity for $\lambda>0$ & theorem\\
Analytic class of $\zeta_{Z1}$ & theorem\\
$\zeta_{Z1}=\zeta\Leftrightarrow$ (D) $\Leftrightarrow$ RH & theorem\\
(D) $\Leftrightarrow$ Dirichlet expansion of $\mathcal K=\zeta_{Z1}/\zeta_{\Z,\RVM}$ in $\Re s>1$ & theorem\\
$\varpi_{Z1}=\varpi_{\RVM}\prod_k\frac{1/4+\tau_k^2}{1/4+\gamma_k^2}$; (D) $\Leftrightarrow\varpi_{Z1}=1$ & theorem\\
Mellin prime-power counts; distributional explicit formula & theorem\\
(D) $\Leftrightarrow$ $J_{Z1}^{\rm M}$ constant between integers & theorem\\
$\frac d{d\lambda}\log\varpi_\lambda=-\Upsilon$; (D) $\Leftrightarrow\int_{0^+}^1\Upsilon=\log L$ & theorem\\
Atomized family; joining factor holomorphic in $\Re s>\frac12$ & theorem\\
Atomized family: Dirichlet series in $\Re s>1$ for every $\lambda$ & theorem\\
(D) $\Leftrightarrow$ vanishing cell defects at $\lambda=0$ & theorem\\
$\zeta_\lambda$ not a Dirichlet series in $\Re s>1$ for $0<\lambda\le1$ & theorem\\
$\lambda\mapsto\log\zeta_\lambda(\sigma)$ strictly convex for real $\sigma>1$ & theorem\\
Hamburger's theorem for the pair $(\zeta_\lambda,\tilde\zeta_{\Z,\lambda})$ with symmetric joining factor & theorem\\
(D) $\Leftrightarrow$ $\zeta_{Z1}$ has a positive representing measure; then $\varpi_{Z1}\le1$ & theorem\\
Density of $d\pi_{\RVM}$; cell masses and defects at $\lambda=1$ & theorem\\
Height kernel of the cell masses; Gaussian defect of the atomization & theorem\\
$\zeta_{Z1}=\zeta\Rightarrow M(x)=O_\eps(x^{1/2+\eps})$ & theorem (Step~3)\\
$\zeta_{Z1}=\zeta$ & open\\
\bottomrule
\end{tabular}
\end{center}

\subsection{Organization}
\label{subsec:organization}

Section~\CGref{sec:clock} treats the clock, its Lambert inversion and the bounded displacement; Section~\CGref{sec:zeta-rvm} the canonical product; Section~\CGref{sec:exp-sums} the exponential sums; Sections~\CGref{sec:explicit-formula} and~\CGref{sec:integer-weights} the prime coordinate and the integer weights; Sections~\CGref{sec:mertens-transform} and~\CGref{sec:transition} the Mertens transform and the transition theorem; Section~\CGref{sec:comparison} the exact Hadamard product and the comparison with $\zeta$; Section~\CGref{sec:lambda-family} the $\lambda$-family; Section~\CGref{sec:Z1-step3} the endpoint, Hamburger's theorem and Step~3; Section~\CGref{sec:euler-products} the value of the endpoint at infinity and the Euler products of $\zeta_{\RVM}$ and $\zeta_{Z1}$; Section~\CGref{sec:atomized} the atomized family, Hamburger's hypotheses at the two endpoints and the programme for the final step; Section~\CGref{sec:criteria} the open path, convexity, Hamburger's theorem for the pair, the positivity criterion, the density of the clock and the Gaussian defect; Section~\CGref{sec:numerics} the numerical illustration; Section~\CGref{sec:conclusions} the conclusions and open problems.

\section{The midpoint Riemann--von Mangoldt clock}
\label{sec:clock}

Let $\rho$ run over the nontrivial zeros of $\zeta$, counted with multiplicity.  No nontrivial zero is real, so the positive ordinates can be listed with multiplicity as $0<\gamma_1\le\gamma_2\le\cdots$.  Put
\begin{equation}\label{eq:zero-count}
 N_\zeta(T):=\#\{\rho:\ 0<\Im\rho\le T\}=\#\{k:\gamma_k\le T\}.
\end{equation}
The Riemann--von Mangoldt formula \cite[Thms.~9.3--9.4]{Titchmarsh1986}, \cite{vonMangoldt1895} states that for $T\ge2$, $T\ne\gamma_k$,
\begin{equation}\label{eq:RvM-full}
 N_\zeta(T)=\frac{T}{2\pi}\log\frac{T}{2\pi e}+\frac78+S(T)+O\Bigl(\frac1T\Bigr),
 \qquad
 S(T):=\frac1\pi\arg\zeta\bigl(\tfrac12+iT\bigr)=O(\log T),
\end{equation}
where $\arg\zeta(\frac12+iT)$ is obtained by continuous variation of the argument along the segments from $2$ to $2+iT$ and from $2+iT$ to $\frac12+iT$, starting from $\arg\zeta(2)=0$.  The $O(1/T)$ term comes from Stirling's formula for $\Gamma$ and is explicit; $S(T)$ is the fluctuating part, with a jump $+m$ at an ordinate of multiplicity $m$.  The \emph{midpoint RVM clock} is the smooth part of \CGeqref{eq:RvM-full}:
\begin{equation}\label{eq:clock}
 \NR(T):=\frac{T}{2\pi}\log\frac{T}{2\pi e}+\frac78 .
\end{equation}
It is strictly increasing on $[2\pi,\infty)$.  The midpoint ordinates $\tau_k$ are the positive solutions of
\begin{equation}\label{eq:midpoint-condition}
 \NR(\tau_k)=k-\frac12,
 \qquad k\ge1,
\end{equation}
that is, the abscissae where the smooth count crosses the midpoints $k-\frac12$ of the steps of a unit staircase.  Thus $\NR$ replaces the staircase $N_\zeta$ by its smooth part, and $\tau_k$ is the point where the $k$th step would be crossed half-way (Figure~\CGref{fig:clock-zoom}).

The Lambert $W$-function is the multivalued inverse of $v\mapsto ve^v$: $W(z)e^{W(z)}=z$.  On $[-1/e,\infty)$ the equation $ve^v=z$ has exactly one solution $v\ge-1$; this defines the principal branch $W_0:[-1/e,\infty)\to[-1,\infty)$, which is real, increasing, with $W_0(0)=0$ and $W_0(z)\sim\log z-\log\log z$ as $z\to\infty$ \cite[\S4.13]{DLMF2024}, \cite{Corless1996}.

\begin{proposition}[Lambert inversion]\label{prop:lambert}
For every $k\ge1$,
\begin{equation}\label{eq:lambert}
 \tau_k=\frac{2\pi(k-11/8)}{W_0\bigl((k-11/8)/e\bigr)},
\end{equation}
The quotient is positive also for $k=1$.
\end{proposition}

\begin{proof}
Put $y=\tau/(2\pi)$.  Equation \CGeqref{eq:midpoint-condition} reads
\begin{equation}\label{eq:lambert-proof-y}
 y(\log y-1)=k-\frac{11}{8}.
\end{equation}
With $v=\log y-1$ one has $y=e^{v+1}$ and $ve^v=(k-11/8)/e$, so $v=W_0((k-11/8)/e)$ and $y=(k-11/8)/W_0((k-11/8)/e)$.  For $k=1$, $(k-11/8)/e=-3/(8e)\in(-1/e,0)$, where $W_0$ is real and negative; numerator and denominator are both negative.
\end{proof}

The clock derivative and the resulting growth are
\begin{equation}\label{eq:clock-derivative}
 \NR'(T)=\frac1{2\pi}\log\frac{T}{2\pi},
\end{equation}
\begin{equation}\label{eq:clock-growth}
 \tau_k\asymp\frac{k}{\log k}\qquad(k\ge2).
\end{equation}
In terms of the clock, \CGeqref{eq:RvM-full} reads
\begin{equation}\label{eq:RvM-formula}
 N_\zeta(T)=\NR(T)+S(T)+O(T^{-1}),\qquad S(T)=O(\log T)\qquad(T\ge2).
\end{equation}
Let
\begin{equation}\label{eq:symmetric-count}
 N_\zeta^{*}(T):=\tfrac12\bigl(N_\zeta(T+0)+N_\zeta(T-0)\bigr)
\end{equation}
be the symmetric count, and define the displacements
\begin{equation}\label{eq:eta-def}
 \eta_k:=\tau_k-\gamma_k .
\end{equation}

\begin{lemma}[Bounded displacement]\label{lem:bounded-displacement}
There is a constant $C_\eta$ such that
\begin{equation}\label{eq:bounded-displacement}
 |\eta_k|\le C_\eta\qquad(k\ge1).
\end{equation}
\end{lemma}

\begin{proof}
Let $\gamma=\gamma_k$ and let $m=m(\gamma)$ be the number of indices $j$ with $\gamma_j=\gamma$; these indices form a block $j_0+1,\dots,j_0+m$ containing $k$.  Then $N_\zeta(\gamma-0)=j_0$ and $N_\zeta(\gamma+0)=j_0+m$, so $N_\zeta^{*}(\gamma)=j_0+m/2$ and
\begin{equation}\label{eq:block-count}
 \bigl|N_\zeta^{*}(\gamma_k)-(k-\tfrac12)\bigr|\le\frac{m(\gamma_k)}2 .
\end{equation}
By \CGeqref{eq:RvM-formula}, $m(\gamma)\le N_\zeta(\gamma+1)-N_\zeta(\gamma-1)=O(\log\gamma)$, and $N_\zeta^{*}(\gamma)=\NR(\gamma)+O(\log\gamma)$.  With \CGeqref{eq:midpoint-condition},
\begin{equation}\label{eq:clock-discrepancy}
 \bigl|\NR(\gamma_k)-\NR(\tau_k)\bigr|=O(\log\gamma_k).
\end{equation}
Both $\gamma_k$ and $\tau_k$ tend to infinity and \CGeqref{eq:clock-discrepancy} gives $\log\gamma_k\asymp\log\tau_k\asymp\log k$.  By the mean-value theorem and \CGeqref{eq:clock-derivative}, $\NR'$ is $\gg\log k$ between $\gamma_k$ and $\tau_k$ for large $k$, so $|\gamma_k-\tau_k|=O(1)$.  Finitely many initial indices are absorbed into $C_\eta$.
\end{proof}

\section{The pure critical RVM zeta function}
\label{sec:zeta-rvm}

Put
\begin{equation}\label{eq:q-def}
 q(s):=s(s-1),
\end{equation}
\begin{equation}\label{eq:Pcrit}
 P_{\RVM}^{\crit}(s):=\prod_{k\ge1}\left(1+\frac{q(s)}{\tfrac14+\tau_k^2}\right),
\end{equation}
\begin{equation}\label{eq:Gdef}
 G(s):=s(s-1)\Gamma(s/2)\pi^{-s/2},
\end{equation}
and
\begin{equation}\label{eq:zetaRVM}
 \xi_{\RVM}(s):=\frac12P_{\RVM}^{\crit}(s),
 \qquad
 \zeta_{\RVM}(s):=\frac{P_{\RVM}^{\crit}(s)}{G(s)}.
\end{equation}

\begin{theorem}[Canonical characterization of $\zeta_{\RVM}$]\label{thm:canonical}
The product \CGeqref{eq:Pcrit} converges absolutely and locally uniformly in $s\in\C$.  It defines an entire function with
\begin{equation}\label{eq:P-symmetry}
 P_{\RVM}^{\crit}(1-s)=P_{\RVM}^{\crit}(s),
 \qquad
 P_{\RVM}^{\crit}(\bar s)=\overline{P_{\RVM}^{\crit}(s)},
 \qquad
 P_{\RVM}^{\crit}(0)=P_{\RVM}^{\crit}(1)=1 .
\end{equation}
Its zeros are exactly the simple zeros
\begin{equation}\label{eq:synthetic-divisor}
 \rho_k^{\RVM}=\frac12+i\tau_k,\qquad \overline{\rho_k^{\RVM}}=\frac12-i\tau_k,
 \qquad k\ge1 .
\end{equation}
Consequently $\zeta_{\RVM}$ has a simple pole at $s=1$ with residue $1$, the trivial zeros $s=-2,-4,-6,\dots$, and nontrivial zeros exactly at \CGeqref{eq:synthetic-divisor}; moreover $\xi_{\RVM}(s)=\xi_{\RVM}(1-s)$.
\end{theorem}

\begin{proof}
By \CGeqref{eq:clock-growth}, $\sum_k\tau_k^{-2}<\infty$.  On a compact set $K$,
\begin{equation}\label{eq:canonical-normal}
 \sup_{s\in K}\left|\frac{q(s)}{\tfrac14+\tau_k^2}\right|\ll_K\tau_k^{-2},
\end{equation}
so the Weierstrass criterion gives normal convergence.  Each factor depends on $s$ only through $q(s)$ and equals $1$ at $s=0,1$, which gives \CGeqref{eq:P-symmetry}; real coefficients give the conjugation symmetry.  The $k$th factor vanishes exactly when
\begin{equation}\label{eq:factor-zero}
 \Bigl(s-\frac12\Bigr)^2+\tau_k^2=0 .
\end{equation}
The clock is strictly increasing on $[2\pi,\infty)$ and $\tau_1>2\pi$, so the $\tau_k$ are distinct; the derivative $2s-1=\pm2i\tau_k\ne0$ of the vanishing factor and the non-vanishing of the other factors make the zeros simple.  At $s=1$, $P_{\RVM}^{\crit}(1)=1$ and $G(s)=(s-1)(1+O(s-1))$, so the residue is $1$.  At $s=-2n$ ($n\ge1$) the pole of $\Gamma(s/2)$ gives a zero of $1/G$, while $P_{\RVM}^{\crit}(-2n)\ne0$ because $q(-2n)>0$.  At $s=0$, $s\Gamma(s/2)\to2$, so $G(0)\ne0$ and $s=0$ is not a zero.
\end{proof}

\section{RVM exponential sums, centroids, and the height transform}
\label{sec:exp-sums}

For $u\in\R$ and $N\ge1$ define
\begin{equation}\label{eq:S-RVM}
 S_{\RVM}(N;u):=\sum_{k\le N}e^{iu\tau_k},
 \qquad
 C_{\RVM}^{\pts}(N;u):=\frac{S_{\RVM}(N;u)}{N},
\end{equation}
and
\begin{equation}\label{eq:path-centroid}
 C_{\RVM}^{\pathc}(N;u):=\frac1N\sum_{n\le N}S_{\RVM}(n;u)
 =\sum_{k\le N}\frac{N-k+1}{N}e^{iu\tau_k}.
\end{equation}

\begin{theorem}[Uniform logarithmic exponential-sum bound]\label{thm:exp-sum}
Let $0<u_1\le u_2<\infty$.  Then
\begin{equation}\label{eq:exp-sum-bound}
 \sup_{u\in[u_1,u_2]}|S_{\RVM}(N;u)|\ll_{u_1,u_2}1+\log N .
\end{equation}
\end{theorem}

\begin{proof}
Let $\mathcal T:[-\tfrac18,\infty)\to[2\pi,\infty)$ be the inverse of the clock, so $\tau_k=\mathcal T(k-\tfrac12)$, and put $\ell_k=\log(\tau_k/2\pi)>0$.  Differentiating $\NR(\mathcal T(y))=y$,
\begin{equation}\label{eq:inverse-clock}
 \mathcal T'(y)=\frac{2\pi}{\log(\mathcal T(y)/2\pi)},
 \qquad
 \mathcal T''(y)=-\frac{4\pi^2}{\mathcal T(y)\log^3(\mathcal T(y)/2\pi)} .
\end{equation}
Since $\mathcal T'$ decreases and $|\mathcal T''|$ decreases,
\begin{equation}\label{eq:rvm-gap-bounds}
 \frac{2\pi}{\ell_{k+1}}\le\tau_{k+1}-\tau_k\le\frac{2\pi}{\ell_k},
\end{equation}
\begin{equation}\label{eq:rvm-second-difference}
 |\tau_{k+2}-2\tau_{k+1}+\tau_k|\le\frac{4\pi^2}{\tau_k\ell_k^3}.
\end{equation}
From \CGeqref{eq:midpoint-condition}, $\tau_k\ell_k=2\pi(k-\tfrac{11}8)+\tau_k$, hence
\begin{equation}\label{eq:tau-ell-k}
 \tau_k\ell_k\asymp k\qquad(k\ge2).
\end{equation}
Fix $u\in[u_1,u_2]$ and put $\Delta_k=u(\tau_{k+1}-\tau_k)$, $E_k=e^{iu\tau_k}$.  Choose $k_0=k_0(u_2)$ with $\ell_{k_0}\ge2u_2$; then $0<\Delta_k\le\pi$ for $k\ge k_0$.  With $b_k:=(e^{i\Delta_k}-1)^{-1}$ one has $E_{k+1}=e^{i\Delta_k}E_k$, hence
\begin{equation}\label{eq:Ek-identity}
 E_k=b_k(E_{k+1}-E_k).
\end{equation}
Summing \CGeqref{eq:Ek-identity} over $k_0\le k\le N$ and shifting the index in the first part gives the exact identity
\begin{equation}\label{eq:summation-by-parts}
 \sum_{k=k_0}^{N}E_k=b_NE_{N+1}-b_{k_0}E_{k_0}+\sum_{k=k_0+1}^{N}(b_{k-1}-b_k)E_k .
\end{equation}
For $0<x\le\pi$, $|e^{ix}-1|=2\sin(x/2)\ge2x/\pi$.  Using \CGeqref{eq:rvm-gap-bounds} and \CGeqref{eq:rvm-second-difference},
\begin{equation}\label{eq:kappa-bounds}
 |b_k|\le\frac{\pi}{2\Delta_k}\le\frac{\ell_{k+1}}{4u_1},
 \qquad
 |b_{k-1}-b_k|\le|\Delta_k-\Delta_{k-1}|\,|b_{k-1}||b_k|
 \le\frac{\pi^2u_2}{4u_1^2}\,\frac{\ell_{k+1}^2}{\tau_{k-1}\ell_{k-1}^3}\ll_{u_1,u_2}\frac1{\tau_k\ell_k}\ll\frac1k,
\end{equation}
the last step by \CGeqref{eq:tau-ell-k}.  Inserting \CGeqref{eq:kappa-bounds} into \CGeqref{eq:summation-by-parts} gives $|\sum_{k_0}^{N}E_k|\ll\ell_{N+1}+\ell_{k_0+1}+\sum_{k\le N}k^{-1}\ll1+\log N$.  The indices $k<k_0$ contribute at most $k_0$.  All constants depend only on $u_1,u_2$.
\end{proof}

\begin{corollary}[Point centroid]\label{cor:point-centroid}
Uniformly for $u\in[u_1,u_2]\subset(0,\infty)$,
\begin{equation}\label{eq:point-centroid}
 C_{\RVM}^{\pts}(N;u)=O_{u_1,u_2}\!\left(\frac{1+\log N}{N}\right)\longrightarrow0 .
\end{equation}
Theorem~\CGref{thm:exp-sum} alone does not imply convergence of the path centroid \CGeqref{eq:path-centroid}.
\end{corollary}

\begin{proof}
The first assertion is \CGeqref{eq:exp-sum-bound} divided by $N$.  For the second, the identity $C_{\RVM}^{\pathc}=N^{-1}\sum_{n\le N}S_{\RVM}(n;u)$ with $S_{\RVM}(n;u)=O(\log n)$ gives only $O(\log N)$.
\end{proof}

Define the midpoint count
\begin{equation}\label{eq:n-RVM}
 n_{\RVM}(H):=\#\{k:\tau_k\le H\}.
\end{equation}
For $x>0$ and $n_{\RVM}(H)\ge1$ define
\begin{equation}\label{eq:MH-RVM}
 M_H^{\RVM}(x):=\frac1{n_{\RVM}(H)}\sum_{\tau_k\le H}x^{i\tau_k},
\end{equation}
and for the ordinates of $\zeta$
\begin{equation}\label{eq:MH-genuine}
 M_H(x):=\frac1{N_\zeta(H)}\sum_{\gamma_k\le H}x^{i\gamma_k}.
\end{equation}

\begin{corollary}[Synthetic height-transform decay]\label{cor:MH-decay}
For every fixed $x>1$, with $N=n_{\RVM}(H)$,
\begin{equation}\label{eq:MH-decay}
 M_H^{\RVM}(x)=O_x\!\left(\frac{1+\log N}{N}\right)\longrightarrow0\qquad(H\to\infty).
\end{equation}
\end{corollary}

\begin{proof}
Take $u=\log x>0$ in Theorem~\CGref{thm:exp-sum}; \CGeqref{eq:MH-RVM} is the point centroid at $N=n_{\RVM}(H)$.
\end{proof}

\begin{remark}\label{rem:MH-not-mertens}
Neither \CGeqref{eq:MH-RVM} nor \CGeqref{eq:MH-genuine} is the Mertens function; both are normalized phase averages.
\end{remark}

\section{The RVM explicit formula, \texorpdfstring{$J_{\RVM}$}{J\_RVM}, and \texorpdfstring{$\pi_{\RVM}$}{pi\_RVM}}
\label{sec:explicit-formula}

For $w=a+ib$ with $b>0$ define
\begin{equation}\label{eq:Ei-def}
 \Ei(w):=\int_{-\infty}^{a}\frac{e^{v+ib}}{v+ib}\,dv,
 \qquad
 \Ei(\bar w):=\overline{\Ei(w)} .
\end{equation}
The path is the horizontal line $\Im=b$, which avoids $0$; this is the branch of $\li(x^\rho)=\Ei(\rho\log x)$ used in Riemann's explicit formula \cite{Edwards1974}.  For $x>1$ put
\begin{equation}\label{eq:J-RVM}
 J_{\RVM}(x):=\li(x)-2\Re\sum_{k\ge1}\Ei\!\left(\Bigl(\frac12+i\tau_k\Bigr)\log x\right)-\log2+\int_x^\infty\frac{dy}{y(y^2-1)\log y},
\end{equation}
the zero sum being taken in the order of $k$.

\begin{lemma}[Asymptotics of $\Ei$ off the real axis]\label{lem:Ei}
For $w=a+ib$, $b>0$,
\begin{equation}\label{eq:Ei-remainder}
 \Ei(w)=\frac{e^w}{w}+\varrho(w),\qquad|\varrho(w)|\le\frac{e^a}{b^2}.
\end{equation}
\end{lemma}

\begin{proof}
Integrate \CGeqref{eq:Ei-def} by parts: $\int_{-\infty}^ae^{v+ib}(v+ib)^{-1}dv=e^w/w+\int_{-\infty}^ae^{v+ib}(v+ib)^{-2}dv$, the boundary term at $-\infty$ being $0$.  The last integral is bounded by $\int_{-\infty}^ae^v b^{-2}dv=e^a/b^2$.
\end{proof}

\begin{theorem}[Compact-uniform existence of the RVM zero sum]\label{thm:J-exists}
The series $\sum_k\Re\Ei\bigl((\tfrac12+i\tau_k)u\bigr)$ converges uniformly for $u$ in every compact interval $[u_1,u_2]\subset(0,\infty)$.  Hence $J_{\RVM}$ is a continuous function on $(1,\infty)$.
\end{theorem}

\begin{proof}
Apply Lemma~\CGref{lem:Ei} with $w=(\tfrac12+i\tau_k)u$, $a=u/2$, $b=\tau_ku$.  Since $\Re\bigl(e^{iu\tau_k}/(\tfrac12+i\tau_k)\bigr)=(\tfrac12\cos u\tau_k+\tau_k\sin u\tau_k)/(\tfrac14+\tau_k^2)$,
\begin{equation}\label{eq:Ei-asymptotic}
 \Re\Ei\bigl((\tfrac12+i\tau_k)u\bigr)=\frac{e^{u/2}}u\,\frac{\sin(u\tau_k)}{\tau_k}+O_{u_1,u_2}(\tau_k^{-2}),
\end{equation}
uniformly for $u\in[u_1,u_2]$.  The error series converges absolutely and uniformly.  For the main series, Abel summation with $S_{\RVM}(K;u)=O_{u_1,u_2}(1+\log K)$ (Theorem~\CGref{thm:exp-sum}) and
\begin{equation}\label{eq:abel-weights}
 \frac1{\tau_K}-\frac1{\tau_{K+1}}\le\frac{\tau_{K+1}-\tau_K}{\tau_K^2}\le\frac{2\pi}{\ell_K\tau_K^2}\ll\frac{\log K}{K^2}
\end{equation}
gives a uniformly convergent majorant $\sum_K(\log K)^2K^{-2}$, while the boundary term $S_{\RVM}(K;u)/\tau_K\to0$ uniformly.  Hence $\sum_ke^{iu\tau_k}/\tau_k$ converges uniformly, and so does its imaginary part.
\end{proof}

Define the support-truncated version
\begin{equation}\label{eq:J-star}
 J_{\RVM}^{*}(x):=\begin{cases}J_{\RVM}(x),&x\ge2,\\0,&0<x<2,\end{cases}
\end{equation}
set $\pi_{\RVM}(x)=0$ for $0<x<2$, and for $x\ge2$
\begin{equation}\label{eq:pi-RVM}
 \pi_{\RVM}(x):=\sum_{1\le m\le\log_2x}\frac{\mu(m)}m\,J_{\RVM}^{*}(x^{1/m}).
\end{equation}

\begin{proposition}[Exact finite prime-power inversion]\label{prop:finite-inversion}
For every $x\ge2$,
\begin{equation}\label{eq:J-pi-inverse}
 J_{\RVM}^{*}(x)=\sum_{1\le m\le\log_2x}\frac1m\,\pi_{\RVM}(x^{1/m}).
\end{equation}
\end{proposition}

\begin{proof}
All sums are finite because $x^{1/m}<2$ for $m>\log_2x$.  Substituting \CGeqref{eq:pi-RVM} into the right side of \CGeqref{eq:J-pi-inverse}, the coefficient of $J_{\RVM}^{*}(x^{1/r})$ is
\begin{equation}\label{eq:mobius-coefficient}
 \frac1r\sum_{d\mid r}\mu(d),
\end{equation}
which is $1$ for $r=1$ and $0$ otherwise.
\end{proof}

\subsection{The square-root estimate}
\label{subsec:square-root}

For $u>0$ put
\begin{equation}\label{eq:A-RVM}
 \mathcal A_{\RVM}(u):=\sum_{k\ge1}\frac{e^{iu\tau_k}}{\tau_k},
\end{equation}
summed in the order of $k$, and write $\ell(t)=\log(t/2\pi)$, so that $\NR'(t)=\ell(t)/2\pi$ and $\ell(\tau_k)=\ell_k$.

\begin{lemma}[The weighted clock sum]\label{lem:A-RVM}
For every $u\ge\log2$ the series \CGeqref{eq:A-RVM} converges, and
\begin{equation}\label{eq:A-RVM-bound}
 |\mathcal A_{\RVM}(u)|\le\frac{\log^2(2+2u)}{4\pi}+\frac{\log(2+2u)}{2\pi}+C_{\mathcal A},
\end{equation}
with an absolute constant $C_{\mathcal A}$.
\end{lemma}

\begin{proof}
Put $Y_j=2^j\cdot4\pi(1+u)$ for $j\ge0$.

\emph{Head $\tau_k\le Y_0$.}  Let $n(t)=\#\{k:\tau_k\le t\}$.  By \CGeqref{eq:midpoint-condition}, $\tau_k\le t$ if and only if $k\le\NR(t)+\tfrac12$, so $n(t)\le\NR(t)+\tfrac12=\frac t{2\pi}(\ell(t)-1)+\frac{11}8\le\frac{t\,\ell(t)}{2\pi}$ for $t\ge\tau_1>\frac{11\pi}4$.  Stieltjes integration gives
\begin{equation}\label{eq:A-head}
 \sum_{\tau_k\le Y_0}\frac1{\tau_k}=\frac{n(Y_0)}{Y_0}+\int_{\tau_1}^{Y_0}\frac{n(t)}{t^2}\,dt\le\frac{\ell(Y_0)}{2\pi}+\int_{\tau_1}^{Y_0}\frac{\ell(t)}{2\pi t}\,dt\le\frac{\ell(Y_0)}{2\pi}+\frac{\ell(Y_0)^2}{4\pi},
\end{equation}
and $\ell(Y_0)=\log(2+2u)$.

\emph{Blocks.}  Let $\mathcal B_j=\{k:Y_j<\tau_k\le Y_{j+1}\}$, let $\mathcal T$ be the inverse clock of \CGeqref{eq:inverse-clock}, and put $\phi(y)=u\,\mathcal T(y-\tfrac12)/2\pi$, so that $e^{2\pi i\phi(k)}=e^{iu\tau_k}$.  By \CGeqref{eq:inverse-clock},
\begin{equation}\label{eq:phi-second}
 \phi''(y)=-\frac{2\pi u}{\mathcal T\,\ell(\mathcal T)^3},\qquad \mathcal T=\mathcal T(y-\tfrac12).
\end{equation}
On $\mathcal B_j$ put $\kappa_j=2\pi u/(Y_{j+1}\ell(Y_{j+1})^3)$.  Since $Y_{j+1}=2Y_j$ and $\ell(Y_{j+1})\le2\ell(Y_j)$ (because $\ell(Y_j)\ge\ell(Y_0)\ge\log2$), we have $\kappa_j\le-\phi''\le16\kappa_j$ there.  The number of indices in $\mathcal B_j$ is at most $\NR(Y_{j+1})-\NR(Y_j)+1\le Y_j\ell(Y_{j+1})/2\pi+1$.  The second-derivative test for exponential sums \cite[Thm.~5.9]{Titchmarsh1986} gives, for every run $a<k\le b$ of consecutive indices of $\mathcal B_j$,
\begin{equation}\label{eq:block-vdc}
 \Bigl|\sum_{a<k\le b}e^{iu\tau_k}\Bigr|\le A\Bigl(16\Bigl(\frac{Y_j\ell(Y_{j+1})}{2\pi}+1\Bigr)\kappa_j^{1/2}+\kappa_j^{-1/2}\Bigr)=:P_j,
\end{equation}
with an absolute constant $A$; a run of length less than $1$ has at most one term, and $1\le\kappa_j^{-1/2}$ because $\kappa_j\le u/(4(1+u)\log^32)<1$.  The weights $1/\tau_k$ are positive and decreasing, and $1/\tau_k<1/Y_j$ on $\mathcal B_j$, so Abel's inequality bounds the contribution of any run in $\mathcal B_j$ by $P_j/Y_j$:
\begin{equation}\label{eq:block-weighted}
 \frac{P_j}{Y_j}\le A\left(\frac{16}{2\pi}\Bigl(\frac{\pi u}{Y_j\ell(Y_{j+1})}\Bigr)^{1/2}+\frac{16}{Y_j}+\Bigl(\frac{\ell(Y_{j+1})^3}{\pi uY_j}\Bigr)^{1/2}\right).
\end{equation}
With $Y_j=2^j\cdot4\pi(1+u)$ and $\ell(Y_{j+1})=\log(2^{j+2}(1+u))\le j+2+\log(1+u)$, each of the three terms is at most an absolute constant times $(j+2)^{3/2}2^{-j/2}$, uniformly in $u\ge\log2$, because $\sup_{u\ge\log2}(1+\log(1+u))^3/(u(1+u))<\infty$.  Hence $\sum_{j\ge0}P_j/Y_j\le C_{\mathcal A}$ with $C_{\mathcal A}$ absolute.

\emph{Convergence and bound.}  If $\tau_K\in\mathcal B_J$, the terms with $k>K$ form a run of $\mathcal B_J$ followed by the full blocks $\mathcal B_j$, $j>J$, so $|\sum_{k>K'}^{K''}e^{iu\tau_k}/\tau_k|\le\sum_{j\ge J}P_j/Y_j\to0$ for $K''>K'\ge K$, $K\to\infty$.  The series converges, and \CGeqref{eq:A-head} together with $\sum_jP_j/Y_j\le C_{\mathcal A}$ gives \CGeqref{eq:A-RVM-bound}.
\end{proof}

\begin{theorem}[Square-root estimate for $J_{\RVM}$ and $\pi_{\RVM}$]\label{thm:Lambert-estimate}
For $x\ge8$,
\begin{equation}\label{eq:J-baseline}
 J_{\RVM}(x)-\li(x)=O\!\left(\frac{\sqrt x(\log\log x)^2}{\log x}\right),
\end{equation}
\begin{equation}\label{eq:pi-baseline}
 \pi_{\RVM}(x)-\li(x)=O\!\left(\frac{\sqrt x(\log\log x)^2}{\log x}\right),
\end{equation}
with absolute implied constants.  No hypothesis on $\zeta$ is used.
\end{theorem}

\begin{proof}
Let $u=\log x\ge\log8$ and $\rho_k=\tfrac12+i\tau_k$.  Lemma~\CGref{lem:Ei} with $w=\rho_ku$, $a=u/2$, $b=\tau_ku$ gives
\begin{equation}\label{eq:Ei-split}
 \Ei(\rho_ku)=\frac{e^{\rho_ku}}{\rho_ku}+\varrho_k,\qquad|\varrho_k|\le\frac{e^{u/2}}{\tau_k^2u^2}.
\end{equation}
Moreover $\frac1{\rho_k}=-\frac i{\tau_k}+r_k$ with $r_k=\frac{i/2}{\tau_k\rho_k}$, so $|r_k|\le\frac1{2\tau_k^2}$.  Since $\sum_k\tau_k^{-2}<\infty$, Lemma~\CGref{lem:A-RVM} shows that the series in \CGeqref{eq:J-RVM} converges in the order of $k$ and
\begin{equation}\label{eq:Ei-sum-A}
 \sum_{k\ge1}\Ei(\rho_ku)=\frac{e^{u/2}}u\Bigl(-i\,\mathcal A_{\RVM}(u)+\sum_kr_ke^{iu\tau_k}\Bigr)+\sum_k\varrho_k
 =\frac{e^{u/2}}u\bigl(-i\,\mathcal A_{\RVM}(u)+O(1)\bigr).
\end{equation}
For $u\ge2$ one has $\log(2+2u)\le\log(2u)+\log2\le3\log u$, so \CGeqref{eq:A-RVM-bound} gives $\mathcal A_{\RVM}(u)=O(\log^2u)$.  The remaining terms $-\log2+\int_x^\infty dy/(y(y^2-1)\log y)$ of \CGeqref{eq:J-RVM} are bounded for $x\ge2$.  With $e^{u/2}=\sqrt x$ and $\log u=\log\log x$ this proves \CGeqref{eq:J-baseline}.

For \CGeqref{eq:pi-baseline}, \CGeqref{eq:pi-RVM} gives
\begin{equation}\label{eq:pi-minus-li}
 \pi_{\RVM}(x)-\li(x)=\bigl(J_{\RVM}(x)-\li(x)\bigr)+\sum_{2\le m\le\log_2x}\frac{\mu(m)}m\,J_{\RVM}^{*}(x^{1/m}).
\end{equation}
By \CGeqref{eq:J-baseline}, $\li(y)\ll y/\log y$, and the continuity of $J_{\RVM}$ on $[2,8]$ (Theorem~\CGref{thm:J-exists}), $|J^*_{\RVM}(y)|\le C\,y/\log y$ for all $y\ge2$.  The term $m=2$ is therefore $O(\sqrt x/\log x)$, and the terms $3\le m\le\log_2x$ add up to at most $C\,x^{1/3}\log_2x/\log x=O(x^{1/3})$.  Both are absorbed in \CGeqref{eq:pi-baseline}.
\end{proof}

\begin{corollary}[The three limits]\label{cor:three-limits}
As $x\to\infty$,
\begin{equation}\label{eq:three-limits}
 \frac{\li(x)}{\pi_{\RVM}(x)}\to1,\qquad
 \frac{\pi(x)}{\li(x)}\to1,\qquad
 \frac{\pi(x)}{\pi_{\RVM}(x)}\to1 .
\end{equation}
More precisely, for some absolute $c>0$,
\begin{equation}\label{eq:pi-minus-piRVM}
 \pi(x)-\pi_{\RVM}(x)=O\!\left(x\,e^{-c\sqrt{\log x}}\right).
\end{equation}
\end{corollary}

\begin{proof}
Since $\li(x)\sim x/\log x$, \CGeqref{eq:pi-baseline} gives $\pi_{\RVM}(x)=\li(x)\bigl(1+O(x^{-1/2}(\log\log x)^2)\bigr)$; in particular $\pi_{\RVM}(x)>0$ for large $x$, and the first limit follows.  The second is the prime number theorem \cite{Hadamard1896,delaValleePoussin1896}, in the form $\pi(x)=\li(x)+O(x\,e^{-c\sqrt{\log x}})$ \cite[Ch.~18]{Davenport1980}.  Subtracting \CGeqref{eq:pi-baseline} gives \CGeqref{eq:pi-minus-piRVM}, and dividing by $\pi_{\RVM}(x)\sim x/\log x$ gives the third limit.
\end{proof}

The first limit and \CGeqref{eq:pi-baseline} concern only the synthetic system; the second and third compare it with the primes and use the prime number theorem.  Figures~\CGref{fig:prime-counts} and~\CGref{fig:pi-ratios} show them numerically.

\section{Integer weights and the continuous-to-atomic problem}
\label{sec:integer-weights}

Define
\begin{equation}\label{eq:cn}
 c_n^{\RVM}:=\pi_{\RVM}(n)-\pi_{\RVM}(n-1),\qquad n\ge2,
\end{equation}
and, initially for $\Re s>1$,
\begin{equation}\label{eq:weighted-zeta}
 \zeta_{\Z,\RVM}(s):=\prod_{n\ge2}(1-n^{-s})^{-c_n^{\RVM}}.
\end{equation}
No positivity or $0/1$ property is built into \CGeqref{eq:cn}.  Define the integer-step prime-power count and the right-continuous extension
\begin{equation}\label{eq:JZ}
 J_{\Z,\RVM}(x):=\sum_{1\le m\le\log_2x}\frac1m\,\pi_{\RVM}\bigl(\lfloor x^{1/m}\rfloor\bigr),
\end{equation}
\begin{equation}\label{eq:J-star-Z}
 J_{\RVM}^{*,\Z}(x):=J_{\RVM}^{*}(\lfloor x\rfloor).
\end{equation}

\begin{theorem}[Lipschitz control of $J_{\RVM}$]\label{thm:unit-cell}
There is an absolute constant $C_J$ such that
\begin{equation}\label{eq:unit-cell}
 |J_{\RVM}(x+\theta)-J_{\RVM}(x)|\le C_J\,\theta\qquad(x\ge2,\ 0<\theta\le x).
\end{equation}
\end{theorem}

\begin{proof}
Put $u_1=\log x$, $u_2=\log(x+\theta)$, so $0<u_2-u_1=\log(1+\theta/x)\le\theta/x$ and $u_2-u_1\le\log2$.  The terms $\li$ and $\int_x^\infty dy/(y(y^2-1)\log y)$ of \CGeqref{eq:J-RVM} have derivatives bounded by $1/\log2$ and $3/\log2$ on $[2,\infty)$, so they change by $O(\theta)$.  It remains to bound $\sum_kD_k$, where, by \CGeqref{eq:Ei-def} and $\frac{d}{dv}\Ei(\rho v)=e^{\rho v}/v$,
\begin{equation}\label{eq:Dk-def}
 D_k:=\Ei(\rho_k^{\RVM}u_2)-\Ei(\rho_k^{\RVM}u_1)=\int_{u_1}^{u_2}\frac{e^{v/2}e^{iv\tau_k}}{v}\,dv .
\end{equation}
The series $\sum_kD_k$ converges (Theorem~\CGref{thm:J-exists} and its proof).  By Abel's theorem for the increasing exponents $\tau_k$, $\sum_kD_k=\lim_{\eps\downarrow0}\sum_ke^{-\eps\tau_k}D_k$, and for $\eps>0$ Fubini gives
\begin{equation}\label{eq:damped-sum}
 \sum_ke^{-\eps\tau_k}D_k=\int_{u_1}^{u_2}\frac{e^{v/2}}v\,S_\eps(v)\,dv,\qquad S_\eps(v):=\sum_{k\ge1}g(\tau_k),\quad g(\tau):=e^{(iv-\eps)\tau}.
\end{equation}
\emph{Clock decomposition.}  With $\tau_0=4\pi$ and the sawtooth count \CGeqref{eq:sawtooth-count}, Stieltjes integration exactly as in \CGeqref{eq:transition-split} ($g(\infty)=0$, $g'\in L^1$) gives
\begin{equation}\label{eq:S-eps-split}
 S_\eps(v)=\int_{\tau_0}^\infty g\,d\NR+g(\tau_0)\NR(\tau_0)+\int_{\tau_0}^\infty\psi\bigl(\NR(\tau)+\tfrac12\bigr)g'(\tau)\,d\tau .
\end{equation}
The middle term is at most $\tfrac12$.  For the first, two integrations by parts with $\NR'(\tau)=\frac1{2\pi}\log\frac\tau{2\pi}$, $\NR''(\tau)=\frac1{2\pi\tau}$ and $v\ge\log2$ give $|\int_{\tau_0}^\infty g\,d\NR|\le C/v$, uniformly in $\eps\in(0,1]$.

\emph{Harmonics.}  The partial sums of $\psi(y+\frac12)=-\sum_{m\ge1}(-1)^m\sin(2\pi my)/(\pi m)$ are uniformly bounded and converge off a null set, and $g'\in L^1$ for $\eps>0$, so termwise integration is allowed:
\begin{equation}\label{eq:K-harmonics}
 \int_{\tau_0}^\infty\psi\bigl(\NR+\tfrac12\bigr)g'\,d\tau=-\sum_{m\ge1}\frac{(-1)^m}{2\pi i\,m}\bigl(K_m^+-K_m^-\bigr),\qquad
 K_m^\pm:=\int_{\tau_0}^\infty(iv-\eps)e^{-\eps\tau}e^{i\Phi_m^\pm(\tau)}\,d\tau,
\end{equation}
with $\Phi_m^\pm(\tau)=v\tau\pm2\pi m\NR(\tau)$.  The weight $e^{-\eps\tau}$ is positive and decreasing, so by the second mean-value theorem each integral is at most $2(v+1)\sup_{\xi}|\int_a^\xi e^{i\Phi}d\tau|$ over the relevant interval $[a,b]$.  We use the first- and second-derivative tests \cite[Lemmas~4.2 and~4.4]{Titchmarsh1986}: $|\int_a^b e^{i\Phi}|\le4/\mu$ if $\Phi'$ is monotone with $|\Phi'|\ge\mu$, and $\le8/\sqrt\lambda$ if $|\Phi''|\ge\lambda$.

For $K_m^+$: $(\Phi_m^+)'=v+m\log\frac\tau{2\pi}\ge v+m\log2$ is increasing, so $|K_m^+|\le8(v+1)/(v+m\log2)$.

For $K_m^-$: $(\Phi_m^-)'=v-m\log\frac\tau{2\pi}$ is decreasing and vanishes at $\tau_m:=2\pi e^{v/m}>2\pi$, and $(\Phi_m^-)''=-m/\tau$.  On $[\tau_0,\tau_m/2]$ (if nonempty) $(\Phi_m^-)'\ge m\log2$; on $[2\tau_m,\infty)$, $(\Phi_m^-)'\le-m\log2$; on $[\max(\tau_0,\tau_m/2),2\tau_m]$, $|(\Phi_m^-)''|\ge m/(2\tau_m)$.  Hence
\begin{equation}\label{eq:K-minus-bound}
 |K_m^-|\le2(v+1)\Bigl(\frac{8}{m\log2}+8\sqrt{\frac{2\tau_m}m}\Bigr)
 =2(v+1)\Bigl(\frac{8}{m\log2}+16\sqrt\pi\,\frac{e^{v/(2m)}}{\sqrt m}\Bigr).
\end{equation}
Summing with the weights $1/(2\pi m)$ and using $\sum_mm^{-3/2}e^{v/(2m)}\le e^{v/2}+\zeta(\tfrac32)e^{v/4}$,
\begin{equation}\label{eq:S-eps-bound}
 |S_\eps(v)|\le C\,(v+1)\,e^{v/2}\qquad(v\ge\log2,\ 0<\eps\le1),
\end{equation}
with an absolute constant $C$.

\emph{Conclusion.}  Inserting \CGeqref{eq:S-eps-bound} into \CGeqref{eq:damped-sum}, for $v\in[u_1,u_2]$ we have $e^{v/2}\le\sqrt{2x}$ and $(v+1)/v\le1+1/\log2$, so
\begin{equation}\label{eq:Dk-sum-bound}
 \Bigl|\sum_ke^{-\eps\tau_k}D_k\Bigr|\le C'(u_2-u_1)\,x\le C'\theta,
\end{equation}
uniformly in $\eps$.  Letting $\eps\downarrow0$ and taking $2\Re$ proves \CGeqref{eq:unit-cell}.
\end{proof}

The $m=1$ harmonic, with stationary point $\tau_1=2\pi x$, produces the unit-periodic ripple of $J_{\RVM}$ (Figure~\CGref{fig:prime-counts}); the harmonics $m\ge2$ have stationary points $2\pi x^{1/m}$ and contribute $O(x^{-1/2+1/(2m)}m^{-3/2})$ to a unit increment.

\begin{corollary}[Bounded cell masses]\label{cor:cn-bounded}
\begin{equation}\label{eq:unit-cell-consequence}
 \pi_{\RVM}(y)-\pi_{\RVM}(\lfloor y\rfloor)=O(1),\qquad c_n^{\RVM}=O(1).
\end{equation}
\end{corollary}

\begin{proof}
By \CGeqref{eq:pi-RVM}, $c_n^{\RVM}=\sum_{m\le\log_2n}\frac{\mu(m)}m\bigl(J_{\RVM}^{*}(n^{1/m})-J_{\RVM}^{*}((n-1)^{1/m})\bigr)$.  The term $m=1$ is $O(1)$ by \CGeqref{eq:unit-cell}.  For $m\ge2$ and $(n-1)^{1/m}\ge2$, \CGeqref{eq:unit-cell} with $\theta=n^{1/m}-(n-1)^{1/m}\le n^{1/m-1}$ bounds the term by $C_Jn^{1/m-1}\le C_Jn^{-1/2}$, and there are at most $\log_2n$ such terms.  The jump of $J_{\RVM}^{*}$ at $2$ enters only for $n=2^m$, once, with size $|J_{\RVM}(2)|/m$.  The same argument applies to $\pi_{\RVM}(y)-\pi_{\RVM}(\lfloor y\rfloor)$.
\end{proof}

\begin{corollary}[The integer-weighted zeta is a Dirichlet series]\label{cor:ZRVM-Dirichlet}
The product \CGeqref{eq:weighted-zeta} converges absolutely for $\Re s>1$ and
\begin{equation}\label{eq:ZRVM-Dirichlet}
 \zeta_{\Z,\RVM}(s)=\sum_{N\ge1}\frac{a_N^{\Z}}{N^s},\qquad a_1^{\Z}=1,\qquad \sum_{N\ge1}\frac{|a_N^{\Z}|}{N^\sigma}<\infty\quad(\sigma>1).
\end{equation}
The inverse is also an absolutely convergent Dirichlet series on the whole half-plane:
\begin{equation}\label{eq:ZRVM-inverse-Dirichlet}
 \frac1{\zeta_{\Z,\RVM}(s)}=\sum_{N\ge1}\frac{\bar a_N^{\Z}}{N^s},\qquad \bar a_1^{\Z}=1,\qquad \sum_{N\ge1}\frac{|\bar a_N^{\Z}|}{N^\sigma}<\infty\quad(\sigma>1).
\end{equation}
\end{corollary}

\begin{proof}
By Corollary~\CGref{cor:cn-bounded},
\begin{equation}\label{eq:log-ZRVM}
 \log\zeta_{\Z,\RVM}(s)=\sum_{n\ge2}c_n^{\RVM}\sum_{m\ge1}\frac{n^{-ms}}m=\sum_{N\ge2}\frac{b_N}{N^s},\qquad b_N:=\sum_{n^m=N}\frac{c_n^{\RVM}}m,
\end{equation}
with $\sum_N|b_N|N^{-\sigma}\le C\sum_{n\ge2}\sum_m n^{-m\sigma}/m<\infty$ for $\sigma>1$.  The exponential of an absolutely convergent Dirichlet series without constant term is an absolutely convergent Dirichlet series with $a_1^{\Z}=1$, and $\sum_N|a_N^{\Z}|N^{-\sigma}\le\exp\bigl(\sum_N|b_N|N^{-\sigma}\bigr)$.  The same applies to $1/\zeta_{\Z,\RVM}=\exp\bigl(-\sum_Nb_NN^{-s}\bigr)$, which gives \CGeqref{eq:ZRVM-inverse-Dirichlet}.
\end{proof}

Numerically the cell masses lie in $[-0.38,1.04]$ for $n\le1000$ (Experiment~\CGref{exp:primes}).

\begin{theorem}[Mellin--Stieltjes identification]\label{thm:MS}
There is a function $J_{\RVM}^{\rm M}$ on $(1,\infty)$ with $J_{\RVM}^{\rm M}(x)x^{-c}\in L^2((1,\infty),dx/x)$ for every $c>1$ such that
\begin{equation}\label{eq:MS-RVM}
 \log\zeta_{\RVM}(s)=s\int_1^\infty J_{\RVM}^{\rm M}(x)\,x^{-s-1}\,dx\qquad(\Re s>1),
\end{equation}
and $J_{\RVM}^{\rm M}-J_{\RVM}$ is constant on $(1,\infty)$.  In particular $J_{\RVM}^{\rm M}=J_{\RVM}^{*}+C_{\RVM}$ on $[2,\infty)$ for a constant $C_{\RVM}$.
\end{theorem}

The proof is given in Section~\CGref{sec:euler-products} (Theorems~\CGref{thm:mellin-counts} and~\CGref{thm:distributional-EF}); it uses the half-plane bound of Theorem~\CGref{thm:halfplane-RVM}.

\begin{corollary}[Weighted factorization]\label{cor:weighted-factorization}
\begin{equation}\label{eq:JZ-diff}
 J_{\Z,\RVM}(x)-J_{\RVM}^{*}(x)=O(\log\log x),
\end{equation}
and
\begin{equation}\label{eq:weighted-factorization}
 \zeta_{\Z,\RVM}(s)=\zeta_{\RVM}(s)\,e^{h_{\RVM}(s)},
\end{equation}
where
\begin{equation}\label{eq:hRVM}
 h_{\RVM}(s):=s\int_1^\infty\bigl(J_{\Z,\RVM}(x)-J_{\RVM}^{\rm M}(x)\bigr)x^{-s-1}\,dx
\end{equation}
is holomorphic for $\Re s>0$.  The two functions have the same zeros in $\Re s>0$; compare the modified Euler products of Grosswald and Schnitzer \cite{GrosswaldSchnitzer1978}.
\end{corollary}

\begin{proof}
By \CGeqref{eq:unit-cell-consequence} replacing $x^{1/m}$ by its integer part costs $O(1/m)$ in the $m$th summand; summing $m^{-1}$ up to $\log_2x$ gives \CGeqref{eq:JZ-diff}.  The product \CGeqref{eq:weighted-zeta} gives $\log\zeta_{\Z,\RVM}(s)=s\int_1^\infty J_{\Z,\RVM}(x)x^{-s-1}dx$ for $\Re s>1$, because $\sum_{n\le y}c_n^{\RVM}=\pi_{\RVM}(\lfloor y\rfloor)$.  Subtracting \CGeqref{eq:MS-RVM} gives \CGeqref{eq:weighted-factorization} on $\Re s>1$.  On $[2,\infty)$, $J_{\Z,\RVM}-J_{\RVM}^{\rm M}=(J_{\Z,\RVM}-J_{\RVM}^{*})-C_{\RVM}=O(\log\log x)$ by \CGeqref{eq:JZ-diff}; on $(1,2)$, $J_{\Z,\RVM}=0$ and $J_{\RVM}^{\rm M}$ is square-integrable.  Hence \CGeqref{eq:hRVM} converges locally normally for $\Re s>0$, and exponentiation introduces no zeros.
\end{proof}

\begin{lemma}[Exact first-harmonic cell cancellation]\label{lem:cell-cancel}
For every integer $n$,
\begin{equation}\label{eq:cell-cancel}
 \int_{n-1}^{n}2\cos(2\pi x)\,dx=0 .
\end{equation}
\end{lemma}

\begin{hypothesis}[Uniform cell remainder]\label{hyp:cell-remainder}
\begin{equation}\label{eq:cell-remainder}
 c_n^{\RVM}-\int_{n-1}^{n}\left(\frac1{\log x}+2\cos(2\pi x)\right)dx=o\!\left(\frac1{\log n}\right).
\end{equation}
\end{hypothesis}

\begin{proposition}[Conditional eventual positivity]\label{prop:eventual-positive}
Hypothesis~\CGref{hyp:cell-remainder} implies
\begin{equation}\label{eq:eventual-positive}
 c_n^{\RVM}\sim\frac1{\log n},
\end{equation}
hence $0<c_n^{\RVM}<1$ for all large $n$.
\end{proposition}

\begin{proof}
The main-density integral is $1/\log n+O(1/(n\log^2n))$; the first harmonic vanishes by Lemma~\CGref{lem:cell-cancel}.
\end{proof}

\section{The critical RVM Mertens transform}
\label{sec:mertens-transform}

Define
\begin{equation}\label{eq:F-RVM}
 F_{\RVM}(s)=\widehat M_{\RVM}^{\crit}(s):=\frac1{\zeta_{\RVM}(s)}=\frac{G(s)}{P_{\RVM}^{\crit}(s)} .
\end{equation}
Formally the associated summatory object is the inverse Mellin transform
\begin{equation}\label{eq:Mcrit-formal}
 M_{\RVM}^{\crit}(x)=\frac1{2\pi i}\int_{(c)}F_{\RVM}(s)\frac{x^s}{s}\,ds ;
\end{equation}
its existence as a symmetric limit, for every $c>\tfrac12$, is proved in Theorem~\CGref{thm:unsmoothed-RVM}.

\begin{proposition}[Synthetic pole abscissa]\label{prop:Mcrit-poles}
$F_{\RVM}$ is holomorphic in $\Re s>\tfrac12$.  Its poles are the simple poles
\begin{equation}\label{eq:Mcrit-poles}
 \frac12\pm i\tau_k\quad(k\ge1)\qquad\text{and}\qquad s=-2,-4,-6,\dots
\end{equation}
and $s=1$ is a simple zero of $F_{\RVM}$.
\end{proposition}

\begin{proof}
This is Theorem~\CGref{thm:canonical} applied to $F_{\RVM}=G/P_{\RVM}^{\crit}$: the zeros of $P_{\RVM}^{\crit}$ are \CGeqref{eq:synthetic-divisor}, and the poles of $G$ are those of $\Gamma(s/2)$ except $s=0$, where $s\Gamma(s/2)\to2$.
\end{proof}

\begin{proposition}[Gaussian square-root bound under vertical growth]\label{prop:Gaussian}
Fix $c>1$ and $0<\eps<c-\tfrac12$.  Assume, uniformly for $\tfrac12+\eps\le\sigma\le c$,
\begin{equation}\label{eq:subgaussian}
 F_{\RVM}(\sigma+it)=\exp(o(t^2))\qquad(|t|\to\infty).
\end{equation}
Define
\begin{equation}\label{eq:MG}
 M_{\RVM,G}^{\crit}(x):=\frac1{2\pi i}\int_{(c)}e^{s^2}F_{\RVM}(s)\frac{x^s}{s}\,ds .
\end{equation}
Then
\begin{equation}\label{eq:MG-bound}
 M_{\RVM,G}^{\crit}(x)=O_\eps(x^{1/2+\eps}).
\end{equation}
\end{proposition}

\begin{proof}
On $s=\sigma+it$, $|e^{s^2}|=e^{\sigma^2-t^2}$, so \CGeqref{eq:subgaussian} gives absolute convergence on vertical lines and makes the horizontal integrals tend to $0$.  The strip $\tfrac12+\eps\le\Re s\le c$ contains no pole of $F_{\RVM}$ (Proposition~\CGref{prop:Mcrit-poles}) and does not contain $s=0$, so Cauchy's theorem moves the line to $\Re s=\tfrac12+\eps$, where $|x^s|=x^{1/2+\eps}$.
\end{proof}

Theorem~\CGref{thm:unsmoothed-RVM} below removes both the assumption \CGeqref{eq:subgaussian} and the smoothing.

\section{The transition theorem and the unsmoothed synthetic Mertens bound}
\label{sec:transition}

Let
\begin{equation}\label{eq:chi}
 \chi(s):=2^s\pi^{s-1}\sin\!\Bigl(\frac{\pi s}2\Bigr)\Gamma(1-s)=\pi^{s-1/2}\frac{\Gamma((1-s)/2)}{\Gamma(s/2)},
\end{equation}
the factor of the functional equation $\zeta(s)=\chi(s)\zeta(1-s)$ \cite[\S2.1]{Titchmarsh1986}.  Write
\begin{equation}\label{eq:transition-zw}
 s=\tfrac12+z,\quad z=\sigma'+it,\qquad
 \tau_*:=-iz=t-i\sigma',\qquad
 w(s):=\exp\bigl\{-2\pi i\,\NR(\tau_*)\bigr\}.
\end{equation}
Stirling's formula gives, uniformly on compact $\sigma$-sets,
\begin{equation}\label{eq:chi-size}
 |\chi(\sigma+it)|=\Bigl(\frac{|t|}{2\pi}\Bigr)^{1/2-\sigma}\bigl(1+O(|t|^{-1})\bigr).
\end{equation}

\begin{theorem}[Transition evaluation of the midpoint product]\label{thm:transition}
There is a constant $\varpi_{\RVM}>0$ with the following properties.
\begin{enumerate}[label=\textup{(\alph*)}]
\item For every $B>0$ there are $t_0(B)$ and $C_B$ such that, for $0<\sigma'\le B$ and $t\ge t_0(B)$,
\begin{equation}\label{eq:transition-w}
 \zeta_{\RVM}(s)=\varpi_{\RVM}\bigl(1+w(s)\bigr)\bigl(1+\mathcal E_0(s)\bigr),\qquad|\mathcal E_0(s)|\le\frac{C_B}t .
\end{equation}
\item For every $0<\sigma'_1<\sigma'_2<\tfrac12$ there are $t_0$ and $C$ such that, for $\sigma'_1\le\sigma'\le\sigma'_2$ and $t\ge t_0$,
\begin{equation}\label{eq:transition-chi}
 \zeta_{\RVM}(s)=\varpi_{\RVM}\bigl(1+\chi(s)\bigr)\bigl(1+\mathcal E(s)\bigr),\qquad|\mathcal E(s)|\le\frac Ct .
\end{equation}
\end{enumerate}
For $t<0$ the same statements hold by $\zeta_{\RVM}(\bar s)=\overline{\zeta_{\RVM}(s)}$.
\end{theorem}

The phase identity $w=\chi(s)(1+O(1/t))$ alone does not give \CGeqref{eq:transition-w}; the product has to be evaluated across the transition region $\tau_k\asymp t$.  $\Log$ denotes the principal logarithm and $\Log\Gamma$ the branch of $\log\Gamma$ on $\Re v>0$ that is real on $(0,\infty)$.  Nothing about $\zeta$ is used.

\subsection{Proof of Theorem~\ref{thm:transition}}
\label{subsec:transition-proof}

Throughout, $t>0$, $0<\sigma'\le B$, $\tau_0:=4\pi$, $t^2\ge2(\tau_0^2+B^2)$, and implied constants depend only on $B$ (and on $\sigma'_1,\sigma'_2$ in part (b)).

\subsubsection*{Step 0: reduction to a sum over the clock}
Since $q(s)=z^2-\tfrac14$, each factor of \CGeqref{eq:Pcrit} equals $(\tau_k^2+z^2)/(\tau_k^2+\tfrac14)$, so
\begin{equation}\label{eq:transition-product}
 \frac{P_{\RVM}^{\crit}(s)}{P_{\RVM}^{\crit}(\tfrac12)}=\prod_{k\ge1}\Bigl(1+\frac{z^2}{\tau_k^2}\Bigr)=\exp\Lambda(z),\qquad
 \Lambda(z):=\sum_{k\ge1}f(\tau_k),\quad f(\tau):=\Log\Bigl(1+\frac{z^2}{\tau^2}\Bigr).
\end{equation}
For $\tau>0$, $\Im(1+z^2/\tau^2)=2\sigma't/\tau^2>0$, so $f$ is smooth on $(0,\infty)$, $f(\tau)=O(|z|^2\tau^{-2})$, the series converges absolutely, and
\begin{equation}\label{eq:fprime}
 f'(\tau)=-\frac{2z^2}{\tau(\tau^2+z^2)}=-\frac2\tau+\frac1{\tau-iz}+\frac1{\tau+iz}.
\end{equation}

\subsubsection*{Step 1: the midpoint sawtooth}
The clock increases on $[2\pi,\infty)$, $\NR(\tau_0)=2\log\frac2e+\frac78\in(0,\tfrac12)$, and $\tau_1>\tau_0$.  With $\psi(y):=y-\lfloor y\rfloor-\tfrac12$, the count \CGeqref{eq:n-RVM} satisfies, for $\tau\ge\tau_0$,
\begin{equation}\label{eq:sawtooth-count}
 n_{\RVM}(\tau)=\bigl\lfloor\NR(\tau)+\tfrac12\bigr\rfloor=\NR(\tau)-\psi\bigl(\NR(\tau)+\tfrac12\bigr),\qquad n_{\RVM}(\tau_0)=0 .
\end{equation}
Stieltjes integration of $f$ against \CGeqref{eq:sawtooth-count}, followed by integration by parts in the $\psi$-part ($f(\infty)=0$, $|\psi|\le\tfrac12$, $\psi(\NR(\tau_0)+\tfrac12)=\NR(\tau_0)$), gives
\begin{equation}\label{eq:transition-split}
\begin{split}
 &\Lambda(z)=\Lambda_{\rm sm}(z)+\Lambda_{\rm saw}(z),\qquad
 \Lambda_{\rm sm}(z):=\int_{\tau_0}^\infty f\,d\NR+f(\tau_0)\NR(\tau_0),\\
 &\Lambda_{\rm saw}(z):=\int_{\tau_0}^{\infty}\psi\bigl(\NR(\tau)+\tfrac12\bigr)f'(\tau)\,d\tau .
\end{split}
\end{equation}

\begin{lemma}[Smooth part]\label{lem:TR-smooth}
With $\log G(s):=\Log(z^2-\tfrac14)+\Log\Gamma(\tfrac s2)-\tfrac s2\log\pi$,
\begin{equation}\label{eq:smooth-difference}
 \Lambda_{\rm sm}(z)-\log G(s)=\nu_3+O(t^{-1}),\qquad\nu_3\in\R\ \text{independent of }z .
\end{equation}
\end{lemma}

\begin{proof}
A second integration by parts ($f\NR\to0$ at $\infty$) cancels the boundary term:
\begin{equation}\label{eq:smooth-J-ibp}
 \Lambda_{\rm sm}(z)=-\int_{\tau_0}^\infty\NR f'\,d\tau=\int_{\tau_0}^\infty\NR(\tau)\Bigl(\frac2\tau-\frac{2\tau}{\tau^2+z^2}\Bigr)d\tau .
\end{equation}
The constant $\tfrac78$ of the clock contributes $\tfrac78\Log(1+z^2/\tau_0^2)$, because $\Log\frac{\tau^2}{\tau^2+z^2}$ is an antiderivative of the bracket on $(0,\infty)$ tending to $0$ at $\infty$.  As $0<\arg z<\frac\pi2$, $\arg(1+\tau_0^2/z^2)=O(t^{-2})$ and
\begin{equation}\label{eq:J2}
 \tfrac78\Log\Bigl(1+\frac{z^2}{\tau_0^2}\Bigr)=\tfrac74\Log z-\tfrac74\log\tau_0+O(t^{-2}).
\end{equation}
The term $\frac\tau{2\pi}\log\frac\tau{2\pi e}$ contributes $\frac1\pi\int_{\tau_0}^\infty\log\frac{\tau}{2\pi e}\frac{z^2}{\tau^2+z^2}\,d\tau$.  Put
\begin{equation}\label{eq:I-def}
 \mathcal I(z):=\frac1\pi\int_0^\infty\log\frac{\tau}{2\pi e}\,\frac{z^2}{\tau^2+z^2}\,d\tau .
\end{equation}
The integrand is holomorphic in $z$ on $\Re z>0$ (its singularities are at $z=\pm i\tau$), $O(|z|^2\tau^{-2}\log\tau)$ at $\infty$ and $O(|\log\tau|)$ at $0$, locally uniformly; so $\mathcal I$ is holomorphic on $\Re z>0$.  For real $z=x>0$ the substitution $\tau=xv$ and $\int_0^\infty\frac{dv}{1+v^2}=\frac\pi2$, $\int_0^\infty\frac{\log v}{1+v^2}\,dv=0$ give $\mathcal I(x)=\frac x2\log\frac x{2\pi e}$; by the identity theorem $\mathcal I(z)=\frac z2\Log\frac z{2\pi e}$ on $\Re z>0$.  On $[0,\tau_0]$, $|\tau^2+z^2|\ge t^2/2$, so $\frac{z^2}{\tau^2+z^2}=1+O(t^{-2})$ there.  Hence
\begin{equation}\label{eq:J-asymp}
 \Lambda_{\rm sm}(z)=\frac z2\Log\frac z{2\pi e}+\frac74\Log z+\nu_1+O(t^{-2}),\qquad
 \nu_1=-\tfrac74\log\tau_0-\tfrac1\pi\int_0^{\tau_0}\log\tfrac\tau{2\pi e}\,d\tau .
\end{equation}
Stirling's formula $\Log\Gamma(v)=(v-\tfrac12)\Log v-v+\tfrac12\log2\pi+O(|v|^{-1})$ \cite[\S5.11]{DLMF2024} with $v=\frac s2=\frac{z+1/2}2$ and $\Log v=\Log z-\log2+\frac1{2z}+O(t^{-2})$ gives
\begin{equation}\label{eq:stirling-loggamma}
 \Log\Gamma(\tfrac s2)=\tfrac z2\Log\tfrac z2-\tfrac14\Log z+\tfrac14\log2-\tfrac z2+\tfrac12\log2\pi+O(t^{-1}),
\end{equation}
and with $\Log(z^2-\tfrac14)=2\Log z+O(t^{-2})$, $-\frac s2\log\pi=-\frac z2\log\pi-\frac14\log\pi$,
\begin{equation}\label{eq:G-asymp}
 \log G(s)=\frac z2\Log\frac z{2\pi e}+\frac74\Log z+\nu_2+O(t^{-1}),\qquad
 \nu_2=\tfrac14\log2+\tfrac12\log2\pi-\tfrac14\log\pi .
\end{equation}
Subtracting gives \CGeqref{eq:smooth-difference} with $\nu_3=\nu_1-\nu_2$.
\end{proof}

\begin{lemma}[The clock as a holomorphic phase]\label{lem:TR-clock}
$\NR(\tau)=\frac\tau{2\pi}\bigl(\Log\frac\tau{2\pi}-1\bigr)+\frac78$ is holomorphic on $\Re\tau>0$, $\NR(\bar\tau)=\overline{\NR(\tau)}$, and for $\Re\tau\ge\tau_0$,
\begin{equation}\label{eq:clock-halfplane-bound}
 2\pi\,\Im\NR(\tau)\ge(\log2)\,\Im\tau\quad(\Im\tau\ge0),\qquad
 2\pi\,\Im\NR(\tau)\le(\log2)\,\Im\tau\quad(\Im\tau\le0).
\end{equation}
\end{lemma}

\begin{proof}
For $\tau=X+iY$, $X\ge\tau_0$, $Y>0$, put $v=Y/X$.  Then
\begin{equation}\label{eq:clock-im-identity}
 2\pi\Im\NR(\tau)=Y\log\frac{|\tau|}{2\pi e}+X\arg\tau
 =Y\Bigl[\log\frac{X}{2\pi e}+\tfrac12\log(1+v^2)+\frac{\arctan v}v\Bigr]\ge Y\bigl[\log2+\varphi(v)\bigr],
\end{equation}
using $X\ge4\pi$, where $\varphi(v)=\frac12\log(1+v^2)+\frac{\arctan v}v-1$.  Since $\varphi(0^+)=0$ and $\varphi'(v)=(v-\arctan v)/v^2\ge0$, $\varphi\ge0$.  The case $Y<0$ follows by conjugation.
\end{proof}

In particular $|w(s)|=e^{2\pi\Im\NR(\tau_*)}\le2^{-\sigma'}<1$, and on the vertical rays $|e^{\pm2\pi im\NR(\tau_0\pm iy)}|\le2^{-my}$.

\begin{lemma}[Abel-regularized harmonic sum]\label{lem:TR-harmonic}
\begin{equation}\label{eq:Psi-asymp}
 \Lambda_{\rm saw}(z)=\Log\bigl(1+w(s)\bigr)+\nu_4+O(t^{-2}),\qquad\nu_4\in\R\ \text{independent of }z .
\end{equation}
\end{lemma}

\begin{proof}
\emph{Regularization.}  For $0<r<1$ let
\begin{equation}\label{eq:psi-r}
 \psi_r(y):=-\sum_{m\ge1}r^m\frac{\sin2\pi my}{\pi m}=\frac1\pi\Im\Log\bigl(1-re^{2\pi iy}\bigr).
\end{equation}
Since $\Re(1-re^{2\pi iy})>0$, $|\psi_r|\le\tfrac12$, and $\psi_r(y)\to\psi(y)$ as $r\uparrow1$ for $y\notin\Z$ (Abel summation of the convergent Fourier series of $\psi$).  The set $\{\tau:\NR(\tau)+\tfrac12\in\Z\}$ is discrete and $f'\in L^1(\tau_0,\infty)$, so by dominated convergence
\begin{equation}\label{eq:Psi-r-limit}
 \Lambda_{\rm saw}(z)=\lim_{r\uparrow1}\Lambda_r(z),\qquad
 \Lambda_r(z):=\int_{\tau_0}^\infty\psi_r\bigl(\NR+\tfrac12\bigr)f'\,d\tau .
\end{equation}
\emph{Termwise integration.}  For fixed $r<1$ the series for $\psi_r(\NR+\tfrac12)$, with $\sin2\pi m(y+\tfrac12)=(-1)^m\sin2\pi my$, is dominated by $\sum_mr^m/m$; by Fubini,
\begin{equation}\label{eq:Psi-r-series}
 \Lambda_r(z)=-\sum_{m\ge1}\frac{(-r)^m}{2\pi i\,m}\bigl(I_m^+-I_m^-\bigr),\qquad
 I_m^\pm:=\int_{\tau_0}^\infty e^{\pm2\pi im\NR(\tau)}f'(\tau)\,d\tau .
\end{equation}
\emph{Contours.}  Let $U^\pm=\{\Re\tau\ge\tau_0,\ \pm\Im\tau\ge0\}$.  On $U^\pm$, $|e^{\pm2\pi im\NR}|\le1$ by Lemma~\CGref{lem:TR-clock}, and $|f'(\tau)|\le4|z|^2|\tau|^{-3}$ for $|\tau|\ge2|z|$, so the quarter-circle arcs of radius $R$ contribute $O(|z|^2R^{-2})\to0$.  The poles of $f'$ are $0$, $iz=-t+i\sigma'\notin U^\pm$, and $\tau_*=-iz=t-i\sigma'\in U^-$, the last with residue $1$ (from $\frac1{\tau+iz}$).  Define the ray integrals
\begin{equation}\label{eq:ray-integrals}
 R_m^+:=i\int_0^\infty e^{2\pi im\NR(\tau_0+iy)}f'(\tau_0+iy)\,dy,\qquad
 R_m^-:=-i\int_0^\infty e^{-2\pi im\NR(\tau_0-iy)}f'(\tau_0-iy)\,dy .
\end{equation}
On $U^+$ Cauchy's theorem gives $I_m^+=R_m^+$.  On $U^-$ the positively oriented boundary of the quarter disc runs down the ray, along the arc, and back along $[\tau_0,\tau_0+R]$ from right to left; letting $R\to\infty$,
\begin{equation}\label{eq:lower-contour}
 R_m^--I_m^-=2\pi i\,e^{-2\pi im\NR(\tau_*)}=2\pi i\,w^m .
\end{equation}
Therefore, using $|rw|<1$,
\begin{equation}\label{eq:Psi-r-closed}
 \Lambda_r(z)=-\sum_{m\ge1}\frac{(-rw)^m}m+\Phi_r(z)=\Log(1+rw)+\Phi_r(z),\qquad
 \Phi_r(z):=-\sum_{m\ge1}\frac{(-r)^m}{2\pi i\,m}\bigl(R_m^+-R_m^-\bigr).
\end{equation}
\emph{The rays.}  Split $f'(\tau)=-\frac2\tau+g_z(\tau)$, $g_z(\tau)=\frac{2\tau}{\tau^2+z^2}$, and correspondingly $R_m^\pm=R_m^{\pm,(0)}+R_m^{\pm,(z)}$, $\Phi_r=\Phi_r^{(0)}+\Phi_r^{(z)}$.  The ray integrands carry the factor $2^{-my}$.  Since $|2/\tau|\le2/\tau_0$ on the rays, $|R^{\pm,(0)}_m|\le2/(\tau_0m\log2)$, so the series for $\Phi_r^{(0)}$ is dominated by $\sum_mC/m^2$ uniformly in $r$, and $\Phi_r^{(0)}\to\nu_4:=\Phi_1^{(0)}$, a constant depending only on $\tau_0$.  It is real: $\NR(\bar\tau)=\overline{\NR(\tau)}$ gives $R_m^{-,(0)}=\overline{R_m^{+,(0)}}$, so each term equals $-(-r)^m\Im R_m^{+,(0)}/(\pi m)$.  For $\tau=\tau_0\pm iy$, $\Re(\tau^2+z^2)=\tau_0^2+\sigma'^2-y^2-t^2\le-\frac12(y^2+t^2)$ because $t^2\ge2(\tau_0^2+B^2)$; hence $|g_z(\tau)|\le4(\tau_0+y)t^{-2}$ and
\begin{equation}\label{eq:ray-bound}
 |R^{\pm,(z)}_m|\le\frac4{t^2}\int_0^\infty2^{-my}(\tau_0+y)\,dy=\frac4{t^2}\Bigl(\frac{\tau_0}{m\log2}+\frac1{(m\log2)^2}\Bigr).
\end{equation}
So $|\Phi_r^{(z)}(z)|\le C\,t^{-2}\sum_mm^{-2}$, uniformly in $0<r<1$ and $0<\sigma'\le B$, and $\Phi_r^{(z)}$ converges as $r\uparrow1$ by dominated convergence of the series.  Since $\Log(1+rw)\to\Log(1+w)$, letting $r\uparrow1$ in \CGeqref{eq:Psi-r-closed} and \CGeqref{eq:Psi-r-limit} gives \CGeqref{eq:Psi-asymp}.
\end{proof}

\subsubsection*{Step 2: assembly}
By \CGeqref{eq:transition-product}, \CGeqref{eq:transition-split} and Lemmas~\CGref{lem:TR-smooth} and~\CGref{lem:TR-harmonic},
\begin{equation}\label{eq:assembly}
\begin{split}
 \zeta_{\RVM}(s)&=\frac{P_{\RVM}^{\crit}(s)}{G(s)}=P_{\RVM}^{\crit}(\tfrac12)\exp\bigl(\Lambda_{\rm sm}(z)-\log G(s)+\Lambda_{\rm saw}(z)\bigr)\\
 &=P_{\RVM}^{\crit}(\tfrac12)\,e^{\nu_3+\nu_4}(1+w)\bigl(1+O(t^{-1})\bigr),
\end{split}
\end{equation}
uniformly for $0<\sigma'\le B$.  This is \CGeqref{eq:transition-w} with
\begin{equation}\label{eq:varpi-def}
 \varpi_{\RVM}:=P_{\RVM}^{\crit}(\tfrac12)\,e^{\nu_3+\nu_4}>0,
\end{equation}
which depends only on the clock; its numerical value $0.9619415776\ldots$ plays no role in the proof.

\subsubsection*{Step 3: from the clock phase to \texorpdfstring{$\chi$}{chi}}
Let
\begin{equation}\label{eq:theta-def}
 \vartheta(v):=\frac1{2i}\Bigl(\Log\Gamma\bigl(\tfrac14+\tfrac{iv}2\bigr)-\Log\Gamma\bigl(\tfrac14-\tfrac{iv}2\bigr)\Bigr)-\frac v2\log\pi,
\end{equation}
holomorphic on $|\Im v|<\tfrac12$.  The identity $\chi(\tfrac12+iv)=e^{-2i\vartheta(v)}$ holds for real $v$, hence on the strip, and $s=\tfrac12+i\tau_*$ with $|\Im\tau_*|=\sigma'<\tfrac12$.  Stirling's formula in the strip gives $\vartheta(v)=\frac v2\Log\frac v{2\pi}-\frac v2-\frac\pi8+O(|v|^{-1})$, i.e.
\begin{equation}\label{eq:clock-theta}
 \NR(v)=\frac{\vartheta(v)}\pi+1+O(|v|^{-1}).
\end{equation}
Therefore
\begin{equation}\label{eq:w-chi}
 w=e^{-2i\vartheta(\tau_*)-2\pi i+O(1/t)}=\chi(s)\bigl(1+O(t^{-1})\bigr).
\end{equation}
For $\sigma'\ge\sigma'_1$, \CGeqref{eq:chi-size} gives $|\chi(s)|\le\tfrac12$ when $t\ge t_0(\sigma'_1)$, so $1+w=(1+\chi)\bigl(1+(w-\chi)/(1+\chi)\bigr)=(1+\chi)(1+O(t^{-1}))$, which with \CGeqref{eq:transition-w} gives \CGeqref{eq:transition-chi}.  This completes the proof of Theorem~\CGref{thm:transition}. \qed

\begin{lemma}[Extension to the critical line]\label{lem:TR-boundary}
The function $\mathcal E_0=\zeta_{\RVM}/(\varpi_{\RVM}(1+w))-1$ extends holomorphically to $\{|\sigma'|<\tfrac12,\ t>t_0\}$, and $|\mathcal E_0|\le C/t$ for $0\le\sigma'\le\tfrac14$.
\end{lemma}

\begin{proof}
By Lemma~\CGref{lem:TR-clock}, $\Im\NR(\tau_*)\ne0$ for $\sigma'\ne0$, so $1+w$ vanishes in the strip only on $\sigma'=0$, exactly at $t=\tau_k$; these zeros are simple because $\frac{d}{dt}w=-\log(\tau_*/2\pi)\,w\ne0$.  The zeros of $\zeta_{\RVM}$ in the strip are the same simple points.  Hence $\mathcal E_0$ has only removable singularities, and the bound extends to $\sigma'=0$ by continuity.
\end{proof}

\begin{corollary}[Derivative at the zeros]\label{cor:derivative-zeros}
For $k$ with $\tau_k>t_0$,
\begin{equation}\label{eq:derivative-zeros}
 \zeta_{\RVM}'(\rho_k^{\RVM})=\varpi_{\RVM}\log\frac{\tau_k}{2\pi}\Bigl(1+O\Bigl(\frac1{\tau_k}\Bigr)\Bigr).
\end{equation}
\end{corollary}

\begin{proof}
By Lemma~\CGref{lem:TR-boundary}, $\zeta_{\RVM}=\varpi_{\RVM}(1+w)(1+\mathcal E_0)$ holds holomorphically near $\sigma'=0$.  At $s=\rho_k^{\RVM}$, $1+w=0$, so $\zeta_{\RVM}'(\rho_k^{\RVM})=\varpi_{\RVM}w'(\rho_k^{\RVM})(1+\mathcal E_0(\rho_k^{\RVM}))$.  From \CGeqref{eq:transition-zw}, $w'(s)=-2\pi i\,\NR'(\tau_*)\cdot(-i)\,w(s)=-\log(\tau_*/2\pi)\,w(s)$, and at the zero $w=-1$, $\tau_*=\tau_k$.  Finally $|\mathcal E_0|\le C/\tau_k$.
\end{proof}

\begin{remark}[Structure of the proof]\label{rem:TR-audit}
The delicate points are handled as follows: the smooth part is computed with principal branches by the identity theorem (Lemma~\CGref{lem:TR-smooth}); the clock estimate is quantified (Lemma~\CGref{lem:TR-clock}); the sawtooth series is replaced by its Abel means, termwise integration is justified for $r<1$, and $r\uparrow1$ is taken after the contour argument uniformly in $r$ and $z$ (Lemma~\CGref{lem:TR-harmonic}); the orientation of the lower contour and the sign of the residue are written out in \CGeqref{eq:lower-contour}; the $z$-dependence of the rays is isolated in $g_z$ and is $O(t^{-2})$, the rest is the explicit real constant $\nu_4$; Stirling is used uniformly for $0<\sigma'\le B$, and the only place where $\sigma'<\tfrac12$ enters is Step~3.
\end{remark}

\begin{corollary}[Vertical constant limit]\label{cor:vertical-limit}
For every fixed $\sigma>\tfrac12$,
\begin{equation}\label{eq:vertical-limit}
 \zeta_{\RVM}(\sigma+it)\longrightarrow\varpi_{\RVM},\qquad
 F_{\RVM}(\sigma+it)\longrightarrow L:=\frac1{\varpi_{\RVM}}\qquad(|t|\to\infty).
\end{equation}
More precisely, with $\sigma'=\sigma-\tfrac12$,
\begin{equation}\label{eq:w-decay}
 |w(\sigma+it)|=O_\sigma(|t|^{-\sigma'}),
\end{equation}
and, with $\alpha_c:=\min(c-\tfrac12,1)$ for $c>\tfrac12$,
\begin{equation}\label{eq:FminusL-all}
 F_{\RVM}(c+it)-L=O_c(|t|^{-\alpha_c}).
\end{equation}
\end{corollary}

\begin{proof}
Expanding the clock at $\tau_*=t-i\sigma'$ gives $\Im\NR(\tau_*)=-\frac{\sigma'}{2\pi}\log\frac t{2\pi}+O_\sigma(t^{-1})$, hence $|w|=e^{2\pi\Im\NR(\tau_*)}=O_\sigma(t^{-\sigma'})$.  Inverting \CGeqref{eq:transition-w} for large $t$ gives $F_{\RVM}=L(1+O(|w|)+O(t^{-1}))$.
\end{proof}

\subsection{Bounds on the whole half-plane}
\label{subsec:halfplane}

The constants of Theorem~\CGref{thm:transition} depend on the width $B$ of the band.  The following estimate holds on the whole half-plane and does not use contours.  Put
\begin{equation}\label{eq:saw-primitive}
 \mathfrak s_1(\tau):=\int_{\tau_0}^{\tau}\psi\bigl(\NR(v)+\tfrac12\bigr)\,dv\qquad(\tau\ge\tau_0).
\end{equation}

\begin{lemma}[Bounded primitive of the midpoint sawtooth]\label{lem:saw-primitive}
$|\mathfrak s_1(\tau)|\le\dfrac{\pi}{4\log2}$ for all $\tau\ge\tau_0$.
\end{lemma}

\begin{proof}
Substitute $y=\NR(v)$, $dv=\mathcal T'(y)\,dy$ with $\mathcal T'(y)=2\pi/\log(\mathcal T(y)/2\pi)$ from \CGeqref{eq:inverse-clock}.  The weight $\mathcal T'$ is positive and decreasing, with $\mathcal T'\le2\pi/\log2$ on $y\ge\NR(\tau_0)$ because $\log(\tau_0/2\pi)=\log2$.  The primitive of $y\mapsto\psi(y+\tfrac12)$ is periodic with values in $[-\tfrac18,0]$, so every integral of $\psi(y+\frac12)$ over an interval has modulus at most $\tfrac18$.  The second mean-value theorem gives $|\mathfrak s_1(\tau)|\le\frac{2\pi}{\log2}\cdot\frac18$.
\end{proof}

\begin{theorem}[Two-sided bound and limit at $+\infty$]\label{thm:halfplane-RVM}
There is an absolute constant $R_0$ such that, for $s=\tfrac12+z$, $z=\sigma'+it$, $\sigma'>\tfrac12$ and $|z|\ge R_0$,
\begin{equation}\label{eq:halfplane-RVM}
 \zeta_{\RVM}(s)=\varpi_{\RVM}\,e^{\theta_{\RVM}(s)},\qquad |\theta_{\RVM}(s)|\le\frac{\pi^2}{2\log2}\cdot\frac1{\sigma'}+\frac{C}{|z|}.
\end{equation}
Consequently, for every $\sigma_0>1$ there is $C(\sigma_0)$ with
\begin{equation}\label{eq:halfplane-two-sided}
 C(\sigma_0)^{-1}\le|\zeta_{\RVM}(s)|\le C(\sigma_0)\qquad(\Re s\ge\sigma_0),
\end{equation}
and $\zeta_{\RVM}(s)\to\varpi_{\RVM}$ uniformly in $t$ as $\Re s\to+\infty$.  Moreover
\begin{equation}\label{eq:nu4-closed}
 \nu_4=-2\int_{\tau_0}^\infty\frac{\mathfrak s_1(\tau)}{\tau^2}\,d\tau .
\end{equation}
\end{theorem}

\begin{proof}
\emph{Smooth part.}  For $\Re z>0$ the computation in the proof of Lemma~\CGref{lem:TR-smooth} is exact:
\begin{equation}\label{eq:smooth-exact}
 \Lambda_{\rm sm}(z)=\frac z2\Log\frac z{2\pi e}+\frac78\Log\Bigl(1+\frac{z^2}{\tau_0^2}\Bigr)-\frac1\pi\int_0^{\tau_0}\log\frac\tau{2\pi e}\,\frac{z^2}{\tau^2+z^2}\,d\tau .
\end{equation}
For $|z|\ge2\tau_0$ one has $|\tau^2+z^2|\ge\frac34|z|^2$ on $[0,\tau_0]$, and Stirling's formula is uniform in $|\arg(s/2)|<\frac\pi2$; the expansions \CGeqref{eq:J-asymp} and \CGeqref{eq:G-asymp} therefore hold with $O(|z|^{-1})$ in place of $O(t^{-1})$, uniformly on $\Re z>\tfrac12$, and $\Lambda_{\rm sm}(z)-\log G(s)=\nu_3+O(|z|^{-1})$.

\emph{Sawtooth part.}  Integrating \CGeqref{eq:transition-split} by parts with \CGeqref{eq:saw-primitive} ($\mathfrak s_1(\tau_0)=0$, $f'(\tau)=O(|z|^2\tau^{-3})$),
\begin{equation}\label{eq:saw-ibp}
 \Lambda_{\rm saw}(z)=-\int_{\tau_0}^\infty\mathfrak s_1f''\,d\tau,\qquad f''(\tau)=\frac2{\tau^2}-\frac1{(\tau-iz)^2}-\frac1{(\tau+iz)^2}.
\end{equation}
Since $|\tau\mp iz|^2=(\tau\pm t)^2+\sigma'^2$ and $\int_{\R}\frac{d\tau}{(\tau-a)^2+\sigma'^2}=\frac\pi{\sigma'}$, Lemma~\CGref{lem:saw-primitive} gives
\begin{equation}\label{eq:saw-bound}
 \Bigl|\Lambda_{\rm saw}(z)+2\int_{\tau_0}^\infty\frac{\mathfrak s_1(\tau)}{\tau^2}\,d\tau\Bigr|\le\frac{\pi}{4\log2}\cdot\frac{2\pi}{\sigma'} .
\end{equation}

\emph{Assembly.}  By \CGeqref{eq:assembly}, $\log\zeta_{\RVM}(s)=\log\varpi_\infty+\theta_{\RVM}(s)$ with $\varpi_\infty:=P_{\RVM}^{\crit}(\tfrac12)\exp(\nu_3-2\int_{\tau_0}^\infty\mathfrak s_1\tau^{-2}d\tau)>0$ and $\theta_{\RVM}$ bounded as in \CGeqref{eq:halfplane-RVM}.  For fixed $\sigma'$ and $t\to\infty$, Corollary~\CGref{cor:vertical-limit} gives $|\log\varpi_{\RVM}-\log\varpi_\infty|\le\pi^2/(2\sigma'\log2)$; letting $\sigma'\to\infty$, $\varpi_\infty=\varpi_{\RVM}$, which is \CGeqref{eq:nu4-closed} by \CGeqref{eq:varpi-def}.  On the compact set $\{\Re s\ge\sigma_0,\ |z|\le R_0\}$ the function $\zeta_{\RVM}$ is continuous and has no zeros or poles, which gives \CGeqref{eq:halfplane-two-sided}.
\end{proof}

Theorem~\CGref{thm:halfplane-RVM} extends the two-sided bound \CGeqref{eq:zetaRVM-two-sided} of Corollary~\CGref{cor:Q-control} from bands to the whole half-plane $\Re s\ge\sigma_0>1$, and identifies $\varpi_{\RVM}$ also as the limit of $\zeta_{\RVM}(\sigma)$ as $\sigma\to+\infty$.

\begin{theorem}[No absolutely convergent Dirichlet expansion]\label{thm:no-Dirichlet-RVM}
Neither $\zeta_{\RVM}$ nor $F_{\RVM}=1/\zeta_{\RVM}$ is represented by a Dirichlet series $\sum_na_nn^{-s}$ converging absolutely in some half-plane $\Re s>\sigma_0$.
\end{theorem}

\begin{proof}
Suppose $F_{\RVM}(s)=\sum_{n\ge1}a_nn^{-s}$ absolutely for $\Re s>\sigma_0$.  Fix $\sigma>\max(\sigma_0,\tfrac12)$ and put $\tilde a_1=a_1-L$, $\tilde a_n=a_n$ ($n\ge2$), so $F_{\RVM}(\sigma+it)-L=\sum_n\tilde a_nn^{-\sigma-it}$ with $\sum_n|\tilde a_n|n^{-\sigma}<\infty$.  Then
\begin{equation}\label{eq:mean-square}
 \frac1{2T}\int_{-T}^{T}\bigl|F_{\RVM}(\sigma+it)-L\bigr|^2dt
 =\sum_{m,n}\tilde a_m\overline{\tilde a_n}(mn)^{-\sigma}\,\frac1{2T}\int_{-T}^{T}\Bigl(\frac nm\Bigr)^{it}dt ,
\end{equation}
the interchange being justified by absolute convergence of the double series, uniformly in $t$.  The inner average is $1$ for $m=n$ and is bounded by $\min\bigl(1,1/(T|\log(n/m)|)\bigr)\to0$ for $m\ne n$; each term is bounded by $|\tilde a_m\tilde a_n|(mn)^{-\sigma}$, which is summable, so dominated convergence gives
\begin{equation}\label{eq:mean-square-limit}
 \lim_{T\to\infty}\frac1{2T}\int_{-T}^{T}\bigl|F_{\RVM}(\sigma+it)-L\bigr|^2dt=\sum_{n\ge1}|\tilde a_n|^2n^{-2\sigma}.
\end{equation}
By \CGeqref{eq:vertical-limit} the integrand tends to $0$ and is bounded on the line (continuous, with a finite limit), so the left side of \CGeqref{eq:mean-square-limit} is $0$.  Hence every $\tilde a_n=0$, $F_{\RVM}\equiv L$ in the half-plane and, by analytic continuation, everywhere, which contradicts $F_{\RVM}(1)=0$.  The same argument applied to $\zeta_{\RVM}-\varpi_{\RVM}$ forces an absolutely convergent Dirichlet series for $\zeta_{\RVM}$ to be the constant $\varpi_{\RVM}$, which contradicts its pole at $s=1$.
\end{proof}

\begin{theorem}[Unsmoothed synthetic Mertens bound]\label{thm:unsmoothed-RVM}
For every $x>1$ and every $c>\tfrac12$ the symmetric Perron integral
\begin{equation}\label{eq:Mcrit-symmetric}
 M_{\RVM}^{\crit}(x):=\lim_{U\to\infty}\frac1{2\pi i}\int_{c-iU}^{c+iU}F_{\RVM}(s)\frac{x^s}{s}\,ds
\end{equation}
exists, is independent of $c>\tfrac12$, and equals the absolutely convergent expression
\begin{equation}\label{eq:Mcrit-renorm}
 M_{\RVM}^{\crit}(x)=L+\frac1{2\pi i}\int_{(c)}\bigl(F_{\RVM}(s)-L\bigr)\frac{x^s}{s}\,ds .
\end{equation}
For every $\eps\in(0,\tfrac12)$,
\begin{equation}\label{eq:Mcrit-unsmoothed-bound}
 M_{\RVM}^{\crit}(x)=L+O_\eps(x^{1/2+\eps}).
\end{equation}
\end{theorem}

\begin{proof}
Fix $c>\tfrac12$.  By \CGeqref{eq:FminusL-all} and continuity of $F_{\RVM}$ on $\Re s=c$ (no poles there by Proposition~\CGref{prop:Mcrit-poles}), the integral in \CGeqref{eq:Mcrit-renorm} converges absolutely.  The Perron kernel gives, for $x>1$ and $c>0$ \cite[Ch.~17]{Davenport1980},
\begin{equation}\label{eq:bromwich}
 \lim_{U\to\infty}\frac1{2\pi i}\int_{c-iU}^{c+iU}\frac{x^s}{s}\,ds=1 .
\end{equation}
Splitting $F_{\RVM}=L+(F_{\RVM}-L)$ proves existence of \CGeqref{eq:Mcrit-symmetric} and \CGeqref{eq:Mcrit-renorm}.  If $\tfrac12<c<c'$, the integrand of \CGeqref{eq:Mcrit-symmetric} is holomorphic on the closed strip $c\le\Re s\le c'$ (Proposition~\CGref{prop:Mcrit-poles}); by Theorem~\CGref{thm:transition}(a) with $B=c'-\tfrac12$, $F_{\RVM}$ is bounded on the horizontal segments $\Im s=\pm U$, $U\ge t_0(B)$, so these contribute $O(x^{c'}U^{-1})\to0$, and Cauchy's theorem gives independence of $c$.  Take $c=\tfrac12+\eps$ in \CGeqref{eq:Mcrit-renorm}:
\begin{equation}\label{eq:Mcrit-final-estimate}
 |M_{\RVM}^{\crit}(x)-L|\le\frac{x^{1/2+\eps}}{2\pi}\int_{-\infty}^{\infty}\frac{|F_{\RVM}(\tfrac12+\eps+it)-L|}{|\tfrac12+\eps+it|}\,dt,
\end{equation}
and the integral is finite by \CGeqref{eq:FminusL-all}, the integrand being $O(|t|^{-1-\alpha_c})$.
\end{proof}

\begin{remark}[Scope]\label{rem:unsmoothed-scope}
Theorem~\CGref{thm:unsmoothed-RVM} concerns $1/\zeta_{\RVM}$.  The transfer to $M$ goes through the endpoint of the $\lambda$-family and is Theorem~\CGref{thm:step3}.
\end{remark}

\section{The exact Hadamard product and the comparison with \texorpdfstring{$\zeta$}{zeta}}
\label{sec:comparison}

Put
\begin{equation}\label{eq:Pzeta}
 P_\zeta(s):=G(s)\zeta(s)=2\xi(s),\qquad \xi(s)=\tfrac12s(s-1)\pi^{-s/2}\Gamma(s/2)\zeta(s).
\end{equation}

\begin{theorem}[Exact Hadamard product]\label{thm:Hadamard}
Let $\rho$ run over the nontrivial zeros of $\zeta$ with $\Im\rho>0$, with multiplicity, ordered by $\Im\rho$.  The product
\begin{equation}\label{eq:Hadamard-exact}
 P_\zeta(s)=\prod_{\Im\rho>0}\left(1+\frac{q(s)}{\rho(1-\rho)}\right)
\end{equation}
converges absolutely and locally uniformly in $s\in\C$, and the identity holds with no exponential or constant factor.
\end{theorem}

\begin{proof}
$\xi$ is entire of order one with $\xi(0)=\tfrac12$, so Hadamard's factorization \cite{Hadamard1896}, \cite[\S2.12]{Titchmarsh1986}, \cite[Ch.~12]{Davenport1980} gives
\begin{equation}\label{eq:hadamard-general}
 2\xi(s)=e^{Bs}\prod_{\rho}\Bigl(1-\frac s\rho\Bigr)e^{s/\rho},
\end{equation}
the product running over all nontrivial zeros and converging absolutely because $\sum_\rho|\rho|^{-2}<\infty$.  The functional equation $\xi(s)=\xi(1-s)$ makes $\rho\mapsto1-\rho$ an involution of the zero multiset; no nontrivial zero is real, since $(1-2^{1-\sigma})\zeta(\sigma)=\sum_{n\ge1}(-1)^{n-1}n^{-\sigma}>0$ for $0<\sigma<1$; so this involution exchanges $\Im\rho>0$ and $\Im\rho<0$.  Grouping $\rho$ with $1-\rho$, which absolute convergence permits,
\begin{equation}\label{eq:pair-factor}
 \Bigl(1-\frac s\rho\Bigr)\Bigl(1-\frac s{1-\rho}\Bigr)e^{s/\rho+s/(1-\rho)}
 =\Bigl(1+\frac{q(s)}{\rho(1-\rho)}\Bigr)e^{s/(\rho(1-\rho))}.
\end{equation}
Since $|\rho(1-\rho)|\ge(\Im\rho)^2$ and $\sum_{\Im\rho>0}(\Im\rho)^{-2}<\infty$, the numbers $1/(\rho(1-\rho))$ are absolutely summable, so the exponentials collect into a single factor:
\begin{equation}\label{eq:hadamard-exp}
 2\xi(s)=e^{B's}\,\Pi_\zeta(s),\qquad \Pi_\zeta(s):=\prod_{\Im\rho>0}\Bigl(1+\frac{q(s)}{\rho(1-\rho)}\Bigr),\qquad B'=B+\sum_{\Im\rho>0}\frac1{\rho(1-\rho)} .
\end{equation}
$\Pi_\zeta$ depends on $s$ only through $q(s)$, hence $\Pi_\zeta(1-s)=\Pi_\zeta(s)$.  From $\xi(1-s)=\xi(s)$ we get $e^{B's}=e^{B'(1-s)}$ on the open set where $\Pi_\zeta\ne0$, so $B'=0$.
\end{proof}

Define the exact quotient
\begin{equation}\label{eq:Q-RVM}
 \mathcal Q_{\RVM}(s):=\frac{\zeta(s)}{\zeta_{\RVM}(s)}=\frac{P_\zeta(s)}{P_{\RVM}^{\crit}(s)}.
\end{equation}

\begin{proposition}[Exact transform relation]\label{prop:M-transform-relation}
For $\Re s>1$, and by meromorphic continuation,
\begin{equation}\label{eq:M-transform-relation}
 F_{\RVM}(s)=\mathcal Q_{\RVM}(s)\,\frac1{\zeta(s)} .
\end{equation}
If the inverse Mellin transform $\mathcal Q_{\RVM}^{\vee}$ of $\mathcal Q_{\RVM}$ and the inverse defining $M_{\RVM}^{\crit}$ exist in a common distribution space with Mellin convolution, then
\begin{equation}\label{eq:M-convolution}
 M_{\RVM}^{\crit}=\mathcal Q_{\RVM}^{\vee}\Mell M .
\end{equation}
\end{proposition}

\begin{proof}
\CGeqref{eq:M-transform-relation} is the identity $1/\zeta_{\RVM}=(\zeta/\zeta_{\RVM})(1/\zeta)$; \CGeqref{eq:M-convolution} is the Mellin convolution theorem under the stated assumptions.
\end{proof}

\begin{corollary}[All growth of $\zeta$ to the right of $\frac12$ lies in $\mathcal Q_{\RVM}$]\label{cor:Q-control}
Let $0<\eps<c'-\tfrac12$.  There are $t_0$ and $0<c_-\le c_+$ such that
\begin{equation}\label{eq:zetaRVM-two-sided}
 c_-\le|\zeta_{\RVM}(\sigma+it)|\le c_+\qquad(\tfrac12+\eps\le\sigma\le c',\ |t|\ge t_0),
\end{equation}
and therefore, in the same region,
\begin{equation}\label{eq:zeta-asymp-Q}
 c_-|\mathcal Q_{\RVM}(s)|\le|\zeta(s)|\le c_+|\mathcal Q_{\RVM}(s)|,\qquad
 \frac1{c_+|\mathcal Q_{\RVM}(s)|}\le\Bigl|\frac1{\zeta(s)}\Bigr|\le\frac1{c_-|\mathcal Q_{\RVM}(s)|}.
\end{equation}
Consequently:
\begin{enumerate}[label=\textup{(\roman*)}]
\item the Lindel\"of hypothesis holds if and only if $\mathcal Q_{\RVM}(\sigma+it)=O_{\sigma,\eps}(|t|^{\eps})$ for every $\sigma>\tfrac12$;
\item every bound for $\zeta$ or $1/\zeta$ on $\tfrac12+\eps\le\sigma\le c'$, $|t|\ge t_0$, is equivalent, up to the constants $c_\pm$, to the same bound for $\mathcal Q_{\RVM}$ or $1/\mathcal Q_{\RVM}$.
\end{enumerate}
\end{corollary}

\begin{proof}
By Theorem~\CGref{thm:transition}(a) with $B=c'-\tfrac12$, $\zeta_{\RVM}=\varpi_{\RVM}(1+w)(1+\mathcal E_0)$ with $|\mathcal E_0|\le C_B/|t|$ and, by \CGeqref{eq:w-decay}, $|w|\le C(|t|/2\pi)^{-\eps}$; for $|t|\ge t_0$ this gives $\tfrac12\le|1+w|\le2$ and $\tfrac12\le|1+\mathcal E_0|\le2$, hence \CGeqref{eq:zetaRVM-two-sided} with $c_-=\varpi_{\RVM}/4$, $c_+=4\varpi_{\RVM}$.  Then \CGeqref{eq:zeta-asymp-Q} follows from $\zeta=\mathcal Q_{\RVM}\zeta_{\RVM}$.  For (i): the Lindel\"of hypothesis is equivalent to $\zeta(\sigma+it)=O(|t|^\eps)$ for every $\sigma>\tfrac12$ \cite[Ch.~XIII]{Titchmarsh1986}, and \CGeqref{eq:zeta-asymp-Q} transfers this to $\mathcal Q_{\RVM}$.  (ii) is \CGeqref{eq:zeta-asymp-Q}.
\end{proof}

On the prime side define
\begin{equation}\label{eq:prime-defect}
 \mathfrak d_n^{\RVM}:=\mathbf 1_{\{n\ \mathrm{prime}\}}-c_n^{\RVM}.
\end{equation}
Since $\pi_{\RVM}(1)=0$, for every integer $N\ge2$
\begin{equation}\label{eq:prime-defect-cumulative}
 \sum_{2\le n\le N}\mathfrak d_n^{\RVM}=\pi(N)-\pi_{\RVM}(N).
\end{equation}

\section{The \texorpdfstring{$\lambda$}{lambda}-family from \texorpdfstring{$\zeta_{\RVM}$}{zeta\_RVM} to \texorpdfstring{$\zeta_{Z1}$}{zeta\_Z1}}
\label{sec:lambda-family}

For $\lambda\in[0,1]$ define the moving ordinates, the family product and the family zeta function
\begin{equation}\label{eq:gamma-lambda}
 \gamma_k(\lambda):=(1-\lambda)\gamma_k+\lambda\tau_k=\gamma_k+\lambda\eta_k,
\end{equation}
\begin{equation}\label{eq:P-lambda}
 P_\lambda(s):=\prod_{k\ge1}\left(1+\frac{q(s)}{\tfrac14+\gamma_k(\lambda)^2}\right),
\end{equation}
\begin{equation}\label{eq:zeta-lambda}
 \zeta_\lambda(s):=\frac{P_\lambda(s)}{G(s)},\qquad F_\lambda(s):=\frac1{\zeta_\lambda(s)},
\end{equation}
and the endpoint
\begin{equation}\label{eq:family-endpoints}
 P_{Z1}:=P_0,\qquad \zeta_{Z1}:=\zeta_0,\qquad \zeta_1=\zeta_{\RVM}.
\end{equation}

\begin{theorem}[The $\lambda$-family]\label{thm:path}
The product \CGeqref{eq:P-lambda} converges absolutely and locally uniformly jointly on $[0,1]\times\C$.  For every $\lambda$, $P_\lambda$ is entire of order exactly one, $P_\lambda(1-s)=P_\lambda(s)$, $P_\lambda(\bar s)=\overline{P_\lambda(s)}$, $P_\lambda(0)=P_\lambda(1)=1$, and its zeros, with multiplicity, are exactly
\begin{equation}\label{eq:family-divisor}
 \frac12\pm i\gamma_k(\lambda),\qquad k\ge1 .
\end{equation}
For $0<\lambda\le1$ the ordinates $\gamma_k(\lambda)$ are strictly increasing in $k$ and all zeros of $P_\lambda$ are simple; at $\lambda=0$ the multiplicities are those of the ordinate list $(\gamma_k)$.  For every $j\ge0$, $P_\lambda^{(j)}\to P_{\lambda_0}^{(j)}$ locally uniformly as $\lambda\to\lambda_0$.  In particular no member $\zeta_\lambda$ of the family, $\zeta_{Z1}=\zeta_0$ included, can have a zero off the critical line other than $-2,-4,\dots$.
\end{theorem}

\begin{proof}
By Lemma~\CGref{lem:bounded-displacement}, $|\gamma_k(\lambda)-\gamma_k|\le C_\eta$, and $\gamma_k\asymp k/\log k$; so $\gamma_k(\lambda)\ge\tfrac12\gamma_k$ for $k\ge k_0$, uniformly in $\lambda$.  On a compact $K\subset\C$,
\begin{equation}\label{eq:lambda-normal}
 \sum_{k\ge1}\sup_{s\in K,\ 0\le\lambda\le1}\left|\frac{q(s)}{\tfrac14+\gamma_k(\lambda)^2}\right|<\infty,
\end{equation}
which gives joint normal convergence and, by Cauchy's formula, convergence of all $s$-derivatives.  For $|s|\le R_0$, uniformly in $\lambda$,
\begin{equation}\label{eq:order-one}
 \log\max_{|s|\le R_0}|P_\lambda(s)|\le\sum_{k\ge1}\log\Bigl(1+\frac{R_0^2+R_0}{\tfrac14+\gamma_k(\lambda)^2}\Bigr)=O(R_0\log R_0),
\end{equation}
so the order is at most one; the zero count is $\asymp T\log T$, so the exponent of convergence is one and the order is exactly one.  Symmetries, normalization and divisor follow factor by factor as in Theorem~\CGref{thm:canonical}.  For $0<\lambda\le1$,
\begin{equation}\label{eq:ordinate-gaps}
 \gamma_{k+1}(\lambda)-\gamma_k(\lambda)=(1-\lambda)(\gamma_{k+1}-\gamma_k)+\lambda(\tau_{k+1}-\tau_k)\ge\lambda(\tau_{k+1}-\tau_k)>0,
\end{equation}
so distinct factors vanish at distinct points and each zero is simple; at $\lambda=0$ equal ordinates give equal factors.
\end{proof}

\begin{remark}[All differences lie in $\gamma_k-\tau_k$]\label{rem:only-gamma-tau}
For $\lambda,\lambda'\in[0,1]$ and $s$ off the divisors,
\begin{equation}\label{eq:ratio-members}
 \frac{\zeta_\lambda(s)}{\zeta_{\lambda'}(s)}=\prod_{k\ge1}\frac{\bigl(\tfrac14+(\gamma_k+\lambda'\eta_k)^2\bigr)\bigl(\tfrac14+(\gamma_k+\lambda\eta_k)^2+q(s)\bigr)}{\bigl(\tfrac14+(\gamma_k+\lambda\eta_k)^2\bigr)\bigl(\tfrac14+(\gamma_k+\lambda'\eta_k)^2+q(s)\bigr)} .
\end{equation}
The gamma factor cancels; the members differ only through $\eta_k=\tau_k-\gamma_k$, and all of them coincide if every $\eta_k$ vanishes.
\end{remark}

\begin{theorem}[Infinitesimal generator]\label{thm:generator}
Put $A_{k,\lambda}:=\tfrac14+\gamma_k(\lambda)^2$.  Off the divisor of $P_\lambda$,
\begin{equation}\label{eq:generator}
 \mathcal V_\lambda(s):=\partial_\lambda\log P_\lambda(s)=-2q(s)\sum_{k\ge1}\frac{\gamma_k(\lambda)\,\eta_k}{A_{k,\lambda}\bigl(A_{k,\lambda}+q(s)\bigr)},
\end{equation}
the series converging locally uniformly on compact sets separated from the divisor.  Since $G$ does not depend on $\lambda$,
\begin{equation}\label{eq:generator-zeta}
 \partial_\lambda\log\zeta_\lambda=\mathcal V_\lambda,\qquad \partial_\lambda\log F_\lambda=-\mathcal V_\lambda .
\end{equation}
\end{theorem}

\begin{proof}
Since $\partial_\lambda A_{k,\lambda}=2\gamma_k(\lambda)\eta_k$,
\begin{equation}\label{eq:mode-derivative}
 \partial_\lambda\log\Bigl(1+\frac q{A_{k,\lambda}}\Bigr)=-\frac{2q\,\gamma_k(\lambda)\eta_k}{A_{k,\lambda}(A_{k,\lambda}+q)} .
\end{equation}
On compact sets separated from the divisor the summand is $O_K(|\eta_k|\gamma_k^{-3})$, and $\sum_k\gamma_k^{-3}<\infty$, so termwise differentiation is justified.
\end{proof}

\begin{proposition}[Integrated multiplier]\label{prop:integrated-multiplier}
On every component avoiding the moving divisor,
\begin{equation}\label{eq:integrated-zeta}
 \zeta_\lambda(s)=\zeta_{Z1}(s)\exp\Bigl(\int_0^\lambda\mathcal V_{\lambda'}(s)\,d\lambda'\Bigr),
\end{equation}
\begin{equation}\label{eq:integrated-inverse}
 F_\lambda(s)=F_{Z1}(s)\exp\Bigl(-\int_0^\lambda\mathcal V_{\lambda'}(s)\,d\lambda'\Bigr),
\end{equation}
and the total multiplier is
\begin{equation}\label{eq:transport-multiplier}
 \Pi_{\rm tr}(s):=\exp\Bigl(-\int_0^1\mathcal V_{\lambda'}(s)\,d\lambda'\Bigr)=\frac{P_{Z1}(s)}{P_{\RVM}^{\crit}(s)},\qquad \zeta_{Z1}=\zeta_{\RVM}\,\Pi_{\rm tr} .
\end{equation}
The identities extend meromorphically to $\C$ through the product form.
\end{proposition}

\begin{proof}
\CGeqref{eq:integrated-zeta} and \CGeqref{eq:integrated-inverse} are the fundamental theorem of calculus applied to \CGeqref{eq:generator-zeta}.  For \CGeqref{eq:transport-multiplier}, integrate each mode with $r=\gamma_k(\lambda')$, $dr=\eta_k\,d\lambda'$:
\begin{equation}\label{eq:mode-integral}
 \int_0^1\partial_{\lambda'}\log\Bigl(1+\frac q{A_{k,\lambda'}}\Bigr)d\lambda'\equiv\log\frac{1+q/(\tfrac14+\tau_k^2)}{1+q/(\tfrac14+\gamma_k^2)}\pmod{2\pi i}.
\end{equation}
Normal convergence permits summation, and exponentiation removes the branch ambiguity.
\end{proof}

\begin{corollary}[Expansion at the endpoint]\label{cor:endpoint-expansion}
Let $K\subset\C$ be compact and disjoint from the zeros of $P_{Z1}$.  Uniformly for $s\in K$, as $\lambda\downarrow0$,
\begin{equation}\label{eq:endpoint-zeta-expansion}
 \zeta_\lambda(s)=\zeta_{Z1}(s)\bigl(1+\lambda\mathcal V_0(s)+O_K(\lambda^2)\bigr),
\end{equation}
\begin{equation}\label{eq:endpoint-inverse-expansion}
 F_\lambda(s)=F_{Z1}(s)\bigl(1-\lambda\mathcal V_0(s)+O_K(\lambda^2)\bigr).
\end{equation}
\end{corollary}

\begin{proof}
By Theorem~\CGref{thm:path} the zeros of $P_\lambda$ converge to those of $P_{Z1}$, so $K$ avoids the moving divisor for small $\lambda$.  The series for $\mathcal V_\lambda$ and its $\lambda$-derivative converge normally on $K$ (the differentiated summands are $O_K(\gamma_k^{-3})$), so $\mathcal V_\lambda=\mathcal V_0+O_K(\lambda)$, and \CGeqref{eq:integrated-zeta} gives $\log(\zeta_\lambda/\zeta_{Z1})=\lambda\mathcal V_0+O_K(\lambda^2)$.
\end{proof}

\begin{proposition}[Residues along the family]\label{prop:residues}
Suppose $\gamma_k$ is a simple ordinate (it occurs once in the list).  Put $\rho_{k,\lambda}=\tfrac12+i\gamma_k(\lambda)$ and, for $0<\lambda\le1$,
\begin{equation}\label{eq:residue-transport}
 \mathfrak r_{k,\lambda}:=\Res_{s=\rho_{k,\lambda}}\frac1{s\,\zeta_\lambda(s)}=\frac{G(\rho_{k,\lambda})}{\rho_{k,\lambda}P_\lambda'(\rho_{k,\lambda})}.
\end{equation}
Then
\begin{equation}\label{eq:residue-limit}
 \mathfrak r_{k,\lambda}\longrightarrow\frac{G(\tfrac12+i\gamma_k)}{(\tfrac12+i\gamma_k)P_{Z1}'(\tfrac12+i\gamma_k)}\qquad(\lambda\downarrow0).
\end{equation}
\end{proposition}

\begin{proof}
Theorem~\CGref{thm:path} gives convergence of the zero locations and local uniform convergence of $P_\lambda'$; the limit denominator is nonzero because the endpoint zero is simple, and $G$ is continuous.
\end{proof}

\begin{remark}[Multiplicity]\label{rem:multiplicity}
For $0<\lambda\le1$ every zero of $\zeta_\lambda$ is simple.  At $\lambda=0$ a repeated ordinate of $\zeta$ appears in $\zeta_{Z1}$ with the same multiplicity, and the residues of Proposition~\CGref{prop:residues} are replaced by the full principal part.  Multiplicity plays no role in Step~3: Theorem~\CGref{thm:step3} uses only that $\zeta$ does not vanish in $\Re s>\tfrac12$.
\end{remark}

\section{The endpoint \texorpdfstring{$\zeta_{Z1}$}{zeta\_Z1}, Hamburger's theorem, and Step 3}
\label{sec:Z1-step3}

\subsection{The analytic class of the endpoint}
\label{subsec:Z1-class}

\begin{proposition}[Analytic class of $\zeta_{Z1}$]\label{prop:Z1-class}
The endpoint
\begin{equation}\label{eq:Z1-product}
 P_{Z1}(s)=\prod_{k\ge1}\left(1+\frac{q(s)}{\tfrac14+\gamma_k^2}\right),\qquad \zeta_{Z1}(s)=\frac{P_{Z1}(s)}{G(s)},
\end{equation}
has the following properties.
\begin{enumerate}[label=\textup{(\roman*)}]
\item $P_{Z1}$ is entire of order exactly one, $P_{Z1}(1-s)=P_{Z1}(s)$, $P_{Z1}(0)=P_{Z1}(1)=1$;
\item the nontrivial zeros of $\zeta_{Z1}$ are exactly
\begin{equation}\label{eq:Z1-divisor}
 \frac12\pm i\gamma_k,\qquad k\ge1,
\end{equation}
with the multiplicities of the ordinate list; its other zeros are $s=-2,-4,\dots$.  It is impossible for $\zeta_{Z1}$ to vanish at any point off the critical line other than $-2,-4,\dots$: every factor of $P_{Z1}$ vanishes only at $\tfrac12\pm i\gamma_k$;
\item $\zeta_{Z1}$ has a simple pole at $s=1$ with residue $1$;
\item $\zeta_{Z1}$ satisfies the functional equation of $\zeta$:
\begin{equation}\label{eq:Z1-FE}
 \pi^{-s/2}\Gamma(s/2)\zeta_{Z1}(s)=\pi^{-(1-s)/2}\Gamma((1-s)/2)\zeta_{Z1}(1-s);
\end{equation}
\item $(s-1)\zeta_{Z1}(s)$ is entire of order exactly one.
\end{enumerate}
\end{proposition}

\begin{proof}
(i) and the divisor in (ii) are the case $\lambda=0$ of Theorem~\CGref{thm:path}; the trivial zeros and the absence of other zeros come from $1/G$ as in Theorem~\CGref{thm:canonical}.  For (iii), $P_{Z1}(1)=1$ and
\begin{equation}\label{eq:G-near-one}
 G(s)=(s-1)+O\bigl((s-1)^2\bigr)\qquad(s\to1).
\end{equation}
For (iv),
\begin{equation}\label{eq:Z1-FE-proof}
 \pi^{-s/2}\Gamma(s/2)\zeta_{Z1}(s)=\frac{P_{Z1}(s)}{q(s)},
\end{equation}
and both $P_{Z1}$ and $q$ are invariant under $s\mapsto1-s$.  For (v),
\begin{equation}\label{eq:Z1-entire-removal}
 (s-1)\zeta_{Z1}(s)=P_{Z1}(s)\,\frac{\pi^{s/2}}{s\,\Gamma(s/2)},
\end{equation}
where $s\Gamma(s/2)\to2$ at $s=0$ and $1/\Gamma$ is entire of order one; the product has order at most one, and the zero count $\asymp T\log T$ gives order at least one.
\end{proof}

Define the ratio
\begin{equation}\label{eq:Z1-defect-ratio}
 \mathscr R_{Z1}(s):=\frac{\zeta_{Z1}(s)}{\zeta(s)}=\frac{P_{Z1}(s)}{P_\zeta(s)},
\end{equation}
a meromorphic function; $\mathscr R_{Z1}\equiv1$ is the identity $\zeta_{Z1}=\zeta$.

\subsection{Hamburger's theorem at the endpoint}
\label{subsec:Hamburger}

\begin{theorem}[Hamburger closure]\label{thm:Hamburger-Z1}
Consider the condition
\begin{equation}\label{eq:Z1-D}
 \text{(D)}\qquad\zeta_{Z1}(s)=\sum_{n\ge1}\frac{a_n}{n^s}\qquad\text{absolutely for }\Re s>1 .
\end{equation}
Then
\begin{equation}\label{eq:Hamburger-equivalence}
 (\mathrm D)\quad\Longleftrightarrow\quad\mathscr R_{Z1}\equiv1\quad\Longleftrightarrow\quad\zeta_{Z1}=\zeta\quad\Longleftrightarrow\quad\mathrm{RH}.
\end{equation}
\end{theorem}

\begin{proof}
Assume (D).  Then $\zeta_{Z1}(1-s)=\sum_na_nn^{-(1-s)}$ converges absolutely for $\Re s<0$, $(s-1)\zeta_{Z1}$ is entire of finite order (Proposition~\CGref{prop:Z1-class}(v)), and \CGeqref{eq:Z1-FE} holds.  These are the hypotheses of Hamburger's theorem \cite{Hamburger1921}, \cite[\S2.13]{Titchmarsh1986}, which gives
\begin{equation}\label{eq:Hamburger-conclusion}
 \zeta_{Z1}=C_H\,\zeta
\end{equation}
for a constant $C_H$.  Both functions have residue $1$ at $s=1$, so $C_H=1$.  If $\zeta_{Z1}=\zeta$, the nontrivial zeros of $\zeta$ are those of $\zeta_{Z1}$, namely \CGeqref{eq:Z1-divisor}, which is RH.  If RH holds, every $\rho$ with $\Im\rho>0$ is $\tfrac12+i\gamma_k$, so $\rho(1-\rho)=\tfrac14+\gamma_k^2$, and Theorem~\CGref{thm:Hadamard} gives
\begin{equation}\label{eq:RH-product}
 P_\zeta(s)=\prod_{k\ge1}\Bigl(1+\frac{q(s)}{\tfrac14+\gamma_k^2}\Bigr)=P_{Z1}(s)
\end{equation}
exactly, hence $\zeta=\zeta_{Z1}$.  Finally $\zeta=\sum_nn^{-s}$ converges absolutely for $\Re s>1$, so $\zeta_{Z1}=\zeta$ gives (D).
\end{proof}

\subsection{Condition (D) and the integer-weighted zeta}
\label{subsec:D-multiplier}

By Corollary~\CGref{cor:ZRVM-Dirichlet}, $\zeta_{\Z,\RVM}$ is an absolutely convergent Dirichlet series in $\Re s>1$ with first coefficient $1$, unconditionally.  Define the multiplier
\begin{equation}\label{eq:K-multiplier}
 \mathcal K(s):=\frac{\zeta_{Z1}(s)}{\zeta_{\Z,\RVM}(s)} .
\end{equation}

\begin{theorem}[Reduction of (D) to the multiplier]\label{thm:D-multiplier}
\begin{enumerate}[label=\textup{(\roman*)}]
\item If $\mathcal K$ has an expansion $\mathcal K(s)=\sum_{n\ge1}k_nn^{-s}$ that converges absolutely in $\Re s>1$, then \CGeqref{eq:Z1-D} holds.
\item If \CGeqref{eq:Z1-D} holds, then $\mathcal K(s)=\sum_{n\ge1}k_nn^{-s}$ absolutely on the whole half-plane $\Re s>1$.
\end{enumerate}
Thus (D) holds if and only if $\mathcal K$ is an absolutely convergent Dirichlet series in $\Re s>1$.
Hence, by Theorem~\CGref{thm:Hamburger-Z1}, an absolutely convergent Dirichlet expansion of $\mathcal K$ in $\Re s>1$ gives $\zeta_{Z1}=\zeta$, RH, and, by Theorem~\CGref{thm:step3}, $M(x)=O_\eps(x^{1/2+\eps})$.
\end{theorem}

\begin{proof}
(i) $\zeta_{Z1}=\zeta_{\Z,\RVM}\,\mathcal K$ is the Dirichlet convolution of two absolutely convergent series in $\Re s>1$, and $\sum_N|(a^{\Z}*k)_N|N^{-\sigma}\le\sum_N|a_N^{\Z}|N^{-\sigma}\sum_n|k_n|n^{-\sigma}<\infty$.  (ii) By \CGeqref{eq:ZRVM-inverse-Dirichlet}, $1/\zeta_{\Z,\RVM}$ is an absolutely convergent Dirichlet series in $\Re s>1$, and so is its product $\mathcal K$ with \CGeqref{eq:Z1-D}.
\end{proof}

By \CGeqref{eq:transport-multiplier} and \CGeqref{eq:weighted-factorization} (Theorem~\CGref{thm:MS}),
\begin{equation}\label{eq:K-explicit}
 \mathcal K(s)=\Pi_{\rm tr}(s)\,e^{-h_{\RVM}(s)}=\frac{P_{Z1}(s)}{P_{\RVM}^{\crit}(s)}\,e^{-h_{\RVM}(s)} ,
\end{equation}
an explicit function of the displacements $\gamma_k-\tau_k$ and of the cell masses $c_n^{\RVM}$.  Condition (D) is thereby located in a single object built from the two constructions of this paper: the $\lambda$-family and the integer weights.

\begin{remark}[Continuous-to-atomic meaning of (D)]\label{rem:D-atomic}
The prime-power count $J_{Z1}^{\rm M}$ of $\zeta_{Z1}$ (Theorem~\CGref{thm:mellin-counts}) is a continuous object built from the zeros, and (D) makes its measure the atomic measure $\sum_{p,m}m^{-1}\delta_{p^m}$.  The final step of the method is therefore the passage of $\zeta_{Z1}$ into the discrete Dirichlet category.  Theorem~\CGref{thm:cell-form} states it exactly: (D) holds if and only if $J_{Z1}^{\rm M}$ is constant between consecutive integers, and Corollary~\CGref{cor:product-identity} reduces it to the single identity \CGeqref{eq:product-identity} between the displacements $\gamma_k-\tau_k$ and $\varpi_{\RVM}$.
\end{remark}

\subsection{Step 3: the Mertens bound from \texorpdfstring{$\zeta_{Z1}=\zeta$}{zeta\_Z1 = zeta}}
\label{subsec:step3}

\begin{lemma}[Vertical growth of $\zeta$]\label{lem:zeta-growth}
For $\tfrac14\le\sigma\le4$ and $|t|\ge2$,
\begin{equation}\label{eq:zeta-growth}
 |\zeta(\sigma+it)|\le15|t| .
\end{equation}
\end{lemma}

\begin{proof}
For $\sigma>0$ \cite[(2.1.5)]{Titchmarsh1986},
\begin{equation}\label{eq:zeta-integral-rep}
 \zeta(s)=\frac s{s-1}-s\int_1^\infty\{y\}y^{-s-1}\,dy .
\end{equation}
For $|t|\ge2$ and $\sigma\le4$ one has $|s|\le3|t|$, $|s/(s-1)|\le|s|/|t|\le3$, and the integral is at most $1/\sigma\le4$; hence $|\zeta(s)|\le3+12|t|\le15|t|$.
\end{proof}

\begin{lemma}[Logarithm of $\zeta$ off the critical line]\label{lem:log-zeta}
Assume $\zeta(s)\ne0$ for $\Re s>\tfrac12$.  Let $0<\delta\le\tfrac14$ and
\begin{equation}\label{eq:three-circles-exponent}
 \omega(\delta):=\frac{\log(3-2\delta)}{\log(3-\delta)}<1 .
\end{equation}
Then, with the branch of $\log\zeta$ equal to $\sum_{n\ge2}\Lambda(n)n^{-s}/\log n$ on $\Re s>1$,
\begin{equation}\label{eq:log-zeta-bound}
 |\log\zeta(\sigma+it)|\le\Bigl(\frac{48}\delta\log|t|\Bigr)^{\omega(\delta)}\qquad(\sigma\ge\tfrac12+\delta,\ |t|\ge4).
\end{equation}
Consequently, for every $\eps>0$ there is $t_2=t_2(\delta,\eps)$ with
\begin{equation}\label{eq:inv-zeta-bound}
 \Bigl|\frac1{\zeta(\sigma+it)}\Bigr|\le|t|^{\eps}\qquad(\sigma\ge\tfrac12+\delta,\ |t|\ge t_2).
\end{equation}
\end{lemma}

\begin{proof}
Fix $t$ with $|t|\ge4$ and put $s_0=2+it$, $R=\tfrac32-\tfrac\delta4$, $r_3=\tfrac32-\tfrac\delta2$.  The closed disc $|s-s_0|\le R$ lies in $\tfrac12+\tfrac\delta4\le\Re s\le\tfrac72$, $|\Im s-t|\le\tfrac32$; it avoids $s=1$ and, by hypothesis, every zero of $\zeta$.  Hence $\log\zeta$ has a holomorphic branch on a neighbourhood of it, equal to the Dirichlet series on its intersection with $\Re s>1$, which is convex.  On the disc $|\Im s|\ge|t|/2\ge2$ and $|\Im s|\le2|t|$, so Lemma~\CGref{lem:zeta-growth} gives $\Re\log\zeta=\log|\zeta|\le\log(30|t|)$ there, while $|\log\zeta(s_0)|\le\log\zeta(2)$.  The Borel--Carath\'eodory lemma \cite[\S5.5]{Titchmarsh1939}, applied with radii $r_3<R$,
\begin{equation}\label{eq:borel-caratheodory}
 \max_{|s-s_0|\le r_3}|\log\zeta(s)|\le\frac{2r_3}{R-r_3}\max_{|s-s_0|=R}\Re\log\zeta(s)+\frac{R+r_3}{R-r_3}|\log\zeta(s_0)|,
\end{equation}
together with $R-r_3=\delta/4$, $2r_3\le3$, $R+r_3\le3$ and $\log30+\log\zeta(2)<4\le3\log|t|$, gives
\begin{equation}\label{eq:BC-bound}
 \max_{|s-s_0|\le r_3}|\log\zeta(s)|\le\frac{12}\delta\bigl(\log(30|t|)+\log\zeta(2)\bigr)\le\frac{48}\delta\log|t| .
\end{equation}
Now let $\tfrac12+\delta\le\sigma<\tfrac32$ and set $r_1=\tfrac12$, $r_2=2-\sigma\in(\tfrac12,\tfrac32-\delta]$.  On $|s-s_0|=r_1$ one has $\Re s\ge\tfrac32$, so $|\log\zeta(s)|\le\log\zeta(\tfrac32)<1$.  Hadamard's three-circle theorem \cite[\S5.3]{Titchmarsh1939} gives, with $M_j=\max_{|s-s_0|=r_j}|\log\zeta|$,
\begin{equation}\label{eq:three-circles}
 M_2\le M_1^{1-\omega}M_3^{\omega},\qquad \omega=\frac{\log(r_2/r_1)}{\log(r_3/r_1)}\le\frac{\log(3-2\delta)}{\log(3-\delta)}=\omega(\delta).
\end{equation}
Since $M_1<1$ and $M_3\le\frac{48}\delta\log|t|$ with $\frac{48}\delta\log|t|>1$, \CGeqref{eq:three-circles} gives \CGeqref{eq:log-zeta-bound} at $s=\sigma+it$, which lies on the circle of radius $r_2$.  For $\sigma\ge\tfrac32$, $|\log\zeta|\le\log\zeta(\tfrac32)<1$.  Finally $|1/\zeta|=\exp(-\Re\log\zeta)\le\exp|\log\zeta|$, and $(\frac{48}\delta\log|t|)^{\omega(\delta)}\le\eps\log|t|$ for $|t|\ge t_2(\delta,\eps)$ because $\omega(\delta)<1$.
\end{proof}

\begin{lemma}[Perron kernel]\label{lem:perron-kernel}
For $y>0$, $y\ne1$, $c>0$, $T>0$ \cite[Ch.~17]{Davenport1980},
\begin{equation}\label{eq:perron-kernel}
 \left|\frac1{2\pi i}\int_{c-iT}^{c+iT}\frac{y^s}{s}\,ds-\mathbf 1_{\{y>1\}}\right|\le y^c\min\Bigl(1,\frac1{T|\log y|}\Bigr).
\end{equation}
\end{lemma}

\begin{theorem}[Step 3]\label{thm:step3}
If $\zeta_{Z1}=\zeta$, then for every $\eps>0$
\begin{equation}\label{eq:step3-bound}
 M(x)=O_\eps\bigl(x^{1/2+\eps}\bigr)\qquad(x\ge1).
\end{equation}
\end{theorem}

\begin{proof}
\emph{Non-vanishing.}  By Proposition~\CGref{prop:Z1-class}(ii), $\zeta_{Z1}$ has no zeros in $\Re s>\tfrac12$; since $\zeta=\zeta_{Z1}$, neither has $\zeta$, and Lemma~\CGref{lem:log-zeta} applies.

\emph{Truncation.}  Let $x=N+\tfrac12$ with $N\ge3$ an integer, $c=1+1/\log x$ and $T=x$.  Applying \CGeqref{eq:perron-kernel} with $y=x/n$ to each term of $\sum_n\mu(n)n^{-s}=1/\zeta(s)$, which converges absolutely on $\Re s=c$, gives
\begin{equation}\label{eq:perron-truncated}
 M(x)=\frac1{2\pi i}\int_{c-iT}^{c+iT}\frac{x^s}{s\,\zeta(s)}\,ds+O\Bigl(\sum_{n\ge1}\Bigl(\frac xn\Bigr)^c\min\Bigl(1,\frac1{T|\log(x/n)|}\Bigr)\Bigr).
\end{equation}
For $n\le x/2$ or $n\ge2x$, $|\log(x/n)|\ge\log2$, and these terms contribute at most $x^c\zeta(c)/(T\log2)\le e\,x(1+\log x)/(T\log2)$.  For $x/2<n<2x$, $(x/n)^c\le4$ and $|\log(x/n)|\ge|n-x|/(2x)$; as $|n-x|\in\{j-\tfrac12:j\ge1\}$, each value taken at most twice, these terms contribute at most $8\sum_{1\le j\le2x}\min(1,4x/(T(2j-1)))\ll1+(x/T)\log x$.  Hence
\begin{equation}\label{eq:perron-error}
 M(x)=\frac1{2\pi i}\int_{c-iT}^{c+iT}\frac{x^s}{s\,\zeta(s)}\,ds+O(\log x).
\end{equation}
\emph{Contour shift.}  Fix $0<\delta\le\tfrac14$.  The integrand is holomorphic on the closed rectangle with vertices $c\pm iT$, $\tfrac12+\delta\pm iT$: $\zeta$ has no zeros there, $s=1$ is a zero of $1/\zeta$, and $s=0$ lies outside.  By Cauchy's theorem
\begin{equation}\label{eq:contour-shift}
 \int_{c-iT}^{c+iT}=\int_{\frac12+\delta-iT}^{\frac12+\delta+iT}+\int_{\frac12+\delta+iT}^{c+iT}-\int_{\frac12+\delta-iT}^{c-iT}.
\end{equation}
For $x\ge t_2(\delta,\eps)$, \CGeqref{eq:inv-zeta-bound} on $|t|=T=x$ bounds the horizontal integrals by
\begin{equation}\label{eq:horizontal-bound}
 \frac1{2\pi}\int_{\frac12+\delta}^{c}\frac{x^\sigma T^{\eps}}{T}\,d\sigma\le x^cT^{\eps-1}=e\,x^{\eps}.
\end{equation}
On the left side, $1/\zeta$ is continuous on the compact segment $\{\tfrac12+\delta+it:|t|\le t_2\}$, bounded by a constant $C(\delta,\eps)$, and \CGeqref{eq:inv-zeta-bound} holds for $|t|\ge t_2$; hence
\begin{equation}\label{eq:left-bound}
 \frac1{2\pi}\int_{-T}^{T}\frac{x^{1/2+\delta}}{|\tfrac12+\delta+it|\,|\zeta(\tfrac12+\delta+it)|}\,dt\le x^{1/2+\delta}\Bigl(C'(\delta,\eps)+\int_{t_2}^{T}t^{\eps-1}dt\Bigr)\ll_{\delta,\eps}x^{1/2+\delta+\eps}.
\end{equation}
\emph{Conclusion.}  \CGeqref{eq:perron-error}--\CGeqref{eq:left-bound} give $M(x)\ll_{\delta,\eps}x^{1/2+\delta+\eps}$ for $x=N+\tfrac12$; taking $\delta=\eps=\eps_0/2$ (with $\eps_0\le\tfrac12$) gives $O_{\eps_0}(x^{1/2+\eps_0})$.  For arbitrary $x\ge3$, $M(x)=M(\lfloor x\rfloor+\tfrac12)$ and $\lfloor x\rfloor+\tfrac12\le2x$; for $1\le x<3$, $|M(x)|\le2$.
\end{proof}

\begin{theorem}[Transfer of the synthetic Mertens bound to $M$]\label{thm:transfer}
Assume $\zeta_{Z1}=\zeta$ and let $0<\delta\le\tfrac14$, $\eps>0$.  Then:
\begin{enumerate}[label=\textup{(\roman*)}]
\item on $\Re s>\tfrac12$,
\begin{equation}\label{eq:transfer-factorization}
 \frac1{\zeta(s)}=F_{\RVM}(s)\,\Pi_{\rm tr}(s)^{-1},
 \qquad
 \Pi_{\rm tr}=\frac{P_{Z1}}{P_{\RVM}^{\crit}}=\frac{\zeta}{\zeta_{\RVM}},
\end{equation}
where $\Pi_{\rm tr}$ is holomorphic and zero-free in $\Re s>\tfrac12$ and $\Pi_{\rm tr}(1)=1$;
\item there is $t_3=t_3(\delta,\eps)$ such that, for $\tfrac12+\delta\le\sigma\le2$ and $|t|\ge t_3$,
\begin{equation}\label{eq:transfer-multiplier-bound}
 |F_{\RVM}(\sigma+it)|\le2L,\qquad |\Pi_{\rm tr}(\sigma+it)^{\pm1}|\le|t|^{\eps};
\end{equation}
\item the synthetic bound of Theorem~\CGref{thm:unsmoothed-RVM}, $M_{\RVM}^{\crit}(x)=L+O_\eps(x^{1/2+\eps})$, carries over to the Mertens function:
\begin{equation}\label{eq:transfer-conclusion}
 M(x)=O_\eps\bigl(x^{1/2+\eps}\bigr).
\end{equation}
\end{enumerate}
\end{theorem}

\begin{proof}
(i) By \CGeqref{eq:transport-multiplier}, $\zeta_{Z1}=\zeta_{\RVM}\Pi_{\rm tr}$; with $\zeta_{Z1}=\zeta$ this is \CGeqref{eq:transfer-factorization}.  In $\Re s>\tfrac12$, $\zeta=\zeta_{Z1}$ has no zeros (Proposition~\CGref{prop:Z1-class}(ii)) and $\zeta_{\RVM}$ has none (Theorem~\CGref{thm:canonical}); both have a simple pole of residue $1$ at $s=1$, so $\Pi_{\rm tr}$ is holomorphic there with $\Pi_{\rm tr}(1)=1$, and it is zero-free in $\Re s>\tfrac12$.

(ii) By Theorem~\CGref{thm:transition}(a) with $B=\tfrac32$, for $\delta\le\sigma'\le\tfrac32$ and $|t|\ge t_0$ we have $\zeta_{\RVM}=\varpi_{\RVM}(1+w)(1+\mathcal E_0)$ with $|w|\le C|t|^{-\delta}\le\tfrac12$ and $|\mathcal E_0|\le\tfrac12$ for $|t|$ large; hence $|F_{\RVM}|\le2L$ and
\begin{equation}\label{eq:log-zetaRVM-bounded}
 \log\zeta_{\RVM}=\log\varpi_{\RVM}+\Log(1+w)+\Log(1+\mathcal E_0)=O(1).
\end{equation}
By Lemma~\CGref{lem:log-zeta}, which applies because $\zeta=\zeta_{Z1}$ has no zeros in $\Re s>\tfrac12$, $|\log\zeta(\sigma+it)|\le(\frac{48}\delta\log|t|)^{\omega(\delta)}$ with $\omega(\delta)<1$.  Therefore
\begin{equation}\label{eq:log-Pi-tr}
 |\log\Pi_{\rm tr}(\sigma+it)|=|\log\zeta-\log\zeta_{\RVM}|\le\Bigl(\frac{48}\delta\log|t|\Bigr)^{\omega(\delta)}+O(1)\le\eps\log|t|
\end{equation}
for $|t|\ge t_3(\delta,\eps)$, which gives \CGeqref{eq:transfer-multiplier-bound}.

(iii) The proof of Theorem~\CGref{thm:unsmoothed-RVM} moves the Perron line of $F_{\RVM}x^s/s$ to $\Re s=\tfrac12+\delta$ using only the holomorphy of $F_{\RVM}$ in $\Re s>\tfrac12$ and the vertical control of $F_{\RVM}$.  By \CGeqref{eq:transfer-factorization} the integrand for $M$ is the same integrand multiplied by $\Pi_{\rm tr}^{-1}$, which is holomorphic in $\Re s>\tfrac12$ by (i) and at most $|t|^{\eps}$ by (ii).  Hence $|1/\zeta(\sigma+it)|\le2L|t|^{\eps}$ for $\tfrac12+\delta\le\sigma\le2$, $|t|\ge t_3$, and the truncated Perron formula, contour shift and estimates \CGeqref{eq:perron-error}--\CGeqref{eq:left-bound} of Theorem~\CGref{thm:step3}, applied with this bound, give $M(x)\ll_{\delta,\eps}x^{1/2+\delta+\eps}$.  Choosing $\delta=\eps=\eps_0/2$ gives \CGeqref{eq:transfer-conclusion}.
\end{proof}

\begin{remark}[The transfer in Mellin form]\label{rem:transfer-mellin}
Formally, \CGeqref{eq:transfer-factorization} reads $M=(\Pi_{\rm tr}^{-1})^{\vee}\Mell M_{\RVM}^{\crit}$: the Mertens function is the synthetic Mertens function convolved with the inverse Mellin transform of the inverse transport multiplier.  The synthetic factor $F_{\RVM}$ contributes the exponent $\tfrac12$ and is bounded (Theorem~\CGref{thm:transition}); the multiplier contributes at most $|t|^{\eps}$.  All the information of the $\lambda$-family enters through $\Pi_{\rm tr}=P_{Z1}/P_{\RVM}^{\crit}$, that is, through the displacements $\gamma_k-\tau_k$.
\end{remark}

\begin{corollary}[Step 3 through (D)]\label{cor:step3-D}
Condition \CGeqref{eq:Z1-D} implies $M(x)=O_\eps(x^{1/2+\eps})$ for every $\eps>0$.
\end{corollary}

\begin{proof}
Theorem~\CGref{thm:Hamburger-Z1} and Theorem~\CGref{thm:step3}.
\end{proof}

\begin{remark}[Status of Step 3]\label{rem:Step3-status}
The implication $\zeta_{Z1}=\zeta\Rightarrow M(x)=O_\eps(x^{1/2+\eps})$ is proved completely above, and it requires nothing beyond the identity: no separate statement about the zeros of $\zeta$ enters, because $\zeta=\zeta_{Z1}$ already carries the divisor \CGeqref{eq:Z1-divisor}.  The identity $\zeta_{Z1}=\zeta$ itself, equivalently (D), equivalently an absolutely convergent Dirichlet expansion of the multiplier $\mathcal K$ of \CGeqref{eq:K-multiplier} in $\Re s>1$ (Theorem~\CGref{thm:D-multiplier}), is not proved in this paper.  It is the open step of the method.
\end{remark}

\subsection{Reciprocal endpoint and finite-height Perron integrals}
\label{subsec:reciprocal}

Put
\begin{equation}\label{eq:FZ1}
 F_{Z1}(s):=\frac1{\zeta_{Z1}(s)}=F_{\RVM}(s)\,\Pi_{\rm tr}(s)^{-1},
\end{equation}
the second identity by \CGeqref{eq:transport-multiplier}.  For $c>1$ and finite $U$ define
\begin{equation}\label{eq:truncated-Perron}
 M_{\lambda,U}(x):=\frac1{2\pi i}\int_{c-iU}^{c+iU}F_\lambda(s)\frac{x^s}{s}\,ds,\qquad M_{Z1,U}:=M_{0,U}.
\end{equation}

\begin{proposition}[Finite-height Perron integrals along the family]\label{prop:finite-Perron}
For fixed $c>1$, $U<\infty$ and $x>0$,
\begin{equation}\label{eq:finite-Perron-limit}
 M_{\lambda,U}(x)\longrightarrow M_{Z1,U}(x)\qquad(\lambda\downarrow0).
\end{equation}
\end{proposition}

\begin{proof}
The segment $\{c+it:|t|\le U\}$ is compact and disjoint from every divisor \CGeqref{eq:family-divisor}, and $P_\lambda\to P_{Z1}$ uniformly on it, so $F_\lambda\to F_{Z1}$ uniformly there.
\end{proof}

\begin{theorem}[Reciprocal-coefficient criterion]\label{thm:coefficient-criterion}
Assume that in $\Re s>1$
\begin{equation}\label{eq:Z1-reciprocal-Dirichlet}
 F_{Z1}(s)=\sum_{n\ge1}\frac{\mu_{Z1}(n)}{n^s}
\end{equation}
absolutely, and let $M_{Z1}(x)=\sum_{n\le x}\mu_{Z1}(n)$, so that
\begin{equation}\label{eq:Z1-increments}
 \mu_{Z1}(n)=M_{Z1}(n)-M_{Z1}(n-1).
\end{equation}
Then the following are equivalent: (i) $M_{Z1}(n)=M(n)$ for every $n\ge1$; (ii) $\mu_{Z1}=\mu$; (iii) $F_{Z1}=1/\zeta$ in $\Re s>1$; (iv) $\zeta_{Z1}=\zeta$; (v) RH.
\end{theorem}

\begin{proof}
(i)$\Leftrightarrow$(ii) by \CGeqref{eq:Z1-increments}; (ii)$\Leftrightarrow$(iii) by uniqueness of absolutely convergent Dirichlet coefficients; (iii)$\Leftrightarrow$(iv) by analytic continuation, both functions being nonzero in $\Re s>1$; (iv)$\Leftrightarrow$(v) by Theorem~\CGref{thm:Hamburger-Z1}.
\end{proof}

\begin{remark}[A route that does not close Step 3]\label{rem:negative-convolution}
The following approach does not work.  One may try to expand both factors of $F_{Z1}=F_{\RVM}\,\Pi_{\rm tr}^{-1}$ in \CGeqref{eq:FZ1} as absolutely convergent Dirichlet series and identify their Dirichlet convolution with $\mu$.  By Theorem~\CGref{thm:no-Dirichlet-RVM} the expansion of $F_{\RVM}$ does not exist in any half-plane, so there are no coefficients to convolve.  The endpoint criterion concerns $\zeta_{Z1}$ itself, as in Theorem~\CGref{thm:Hamburger-Z1}.
\end{remark}

\section{The endpoint at infinity and the Euler products}
\label{sec:euler-products}

This section applies the clock decomposition of Theorem~\CGref{thm:halfplane-RVM} to the ordinates $\gamma_k$ of the endpoint, and constructs the prime-power counts of $\zeta_{\RVM}$ and $\zeta_{Z1}$ as Mellin transforms of their logarithms.  No hypothesis is used.

\subsection{The value of the endpoint at infinity}
\label{subsec:value-infinity}

With the zero count $N_\zeta$ of \CGeqref{eq:RvM-full} put
\begin{equation}\label{eq:frak-e}
 \mathfrak e(\tau):=N_\zeta(\tau)-\NR(\tau),\qquad \mathfrak e_1(\tau):=\int_{\tau_0}^{\tau}\mathfrak e(v)\,dv\qquad(\tau\ge\tau_0=4\pi).
\end{equation}
By \CGeqref{eq:RvM-full}, $\mathfrak e(\tau)=S(\tau)+O(1/\tau)$, and Littlewood's bound $\int_0^TS(t)\,dt=O(\log T)$ \cite[Thm.~9.9]{Titchmarsh1986} gives an absolute constant $C_{\mathfrak e}$ with
\begin{equation}\label{eq:frak-e1-bound}
 |\mathfrak e_1(\tau)|\le C_{\mathfrak e}\log\tau\qquad(\tau\ge\tau_0).
\end{equation}
For the clock the corresponding primitive is $-\mathfrak s_1$, bounded by Lemma~\CGref{lem:saw-primitive}.

\begin{theorem}[Clock decomposition of the endpoint]\label{thm:Z1-clock}
Let $s=\tfrac12+z$, $z=\sigma'+it$, $\sigma'>0$.  Then
\begin{equation}\label{eq:Z1-clock}
 \log\zeta_{Z1}(s)=\log\varpi_{Z1}-\int_{\tau_0}^\infty\mathfrak e_1(\tau)\Bigl(\frac1{(\tau-iz)^2}+\frac1{(\tau+iz)^2}\Bigr)d\tau+O\Bigl(\frac1{|z|}\Bigr)\qquad(|z|\ge R_0),
\end{equation}
where
\begin{equation}\label{eq:varpi-Z1}
 \varpi_{Z1}:=P_{Z1}(\tfrac12)\exp\Bigl(\nu_3+2\int_{\tau_0}^\infty\frac{\mathfrak e_1(\tau)}{\tau^2}\,d\tau\Bigr)>0 .
\end{equation}
Consequently:
\begin{enumerate}[label=\textup{(\roman*)}]
\item for $\sigma'>0$ and $|z|\ge R_0$,
\begin{equation}\label{eq:Z1-growth}
 \bigl|\log\zeta_{Z1}(s)-\log\varpi_{Z1}\bigr|\le\frac{2(\pi+2)C_{\mathfrak e}}{\sigma'}\Bigl(1+\log\bigl(2|t|+2\sigma'+\tau_0\bigr)\Bigr)+\frac C{|z|};
\end{equation}
\item $\zeta_{Z1}(\sigma)\to\varpi_{Z1}$ as $\sigma\to+\infty$;
\item the value at infinity is the ordered product
\begin{equation}\label{eq:varpi-Z1-product}
 \varpi_{Z1}=\varpi_{\RVM}\prod_{k\ge1}\frac{\tfrac14+\tau_k^2}{\tfrac14+\gamma_k^2}:=\varpi_{\RVM}\lim_{K\to\infty}\prod_{k\le K}\frac{\tfrac14+\tau_k^2}{\tfrac14+\gamma_k^2}.
\end{equation}
\end{enumerate}
\end{theorem}

\begin{proof}
As in \CGeqref{eq:transition-product}, $\log P_{Z1}(s)-\log P_{Z1}(\frac12)=\sum_kf(\gamma_k)$ with $f(\tau)=\Log(1+z^2/\tau^2)$.  Since $\gamma_1>\tau_0$, $N_\zeta(\tau_0)=0$ and $\mathfrak e(\tau_0)=-\NR(\tau_0)$.  Stieltjes integration against $N_\zeta=\NR+\mathfrak e$ and two integrations by parts ($f\mathfrak e\to0$ and $\mathfrak e_1f'\to0$ at $\infty$, $\mathfrak e_1(\tau_0)=0$) give
\begin{equation}\label{eq:Z1-split}
 \sum_kf(\gamma_k)=\Lambda_{\rm sm}(z)+\int_{\tau_0}^\infty\mathfrak e_1f''\,d\tau,
\end{equation}
with the same $\Lambda_{\rm sm}$ as in \CGeqref{eq:transition-split}.  The proof of Theorem~\CGref{thm:halfplane-RVM} gives $\Lambda_{\rm sm}-\log G=\nu_3+O(|z|^{-1})$, and inserting $f''$ from \CGeqref{eq:saw-ibp} gives \CGeqref{eq:Z1-clock}.

For (i), $|\tau\mp iz|^2=(\tau\pm t)^2+\sigma'^2$.  For real $a$, splitting at $X=2|a|+2\sigma'+\tau_0$ and using $(\tau-a)^2\ge\tau^2/4$ for $\tau\ge X$,
\begin{equation}\label{eq:log-kernel}
 \int_{\tau_0}^\infty\frac{\log\tau}{(\tau-a)^2+\sigma'^2}\,d\tau\le\frac\pi{\sigma'}\log X+4\,\frac{1+\log X}X\le\frac{\pi+2}{\sigma'}(1+\log X).
\end{equation}
With \CGeqref{eq:frak-e1-bound} this gives \CGeqref{eq:Z1-growth}.  For (ii) take $t=0$: the right side of \CGeqref{eq:Z1-growth} tends to $0$ as $\sigma'\to\infty$.

For (iii), Theorem~\CGref{thm:halfplane-RVM} gives $\varpi_{\RVM}=P_{\RVM}^{\crit}(\frac12)\exp(\nu_3-2\int\mathfrak s_1\tau^{-2}d\tau)$, so
\begin{equation}\label{eq:varpi-ratio}
 \frac{\varpi_{Z1}}{\varpi_{\RVM}}=\frac{P_{Z1}(\frac12)}{P_{\RVM}^{\crit}(\frac12)}\exp\Bigl(2\int_{\tau_0}^\infty\frac{\mathfrak e_1(\tau)+\mathfrak s_1(\tau)}{\tau^2}\,d\tau\Bigr).
\end{equation}
Since $\mathfrak e+\psi(\NR+\frac12)=N_\zeta-n_{\RVM}$ and $(\mathfrak e_1+\mathfrak s_1)(X)/X\to0$,
\begin{equation}\label{eq:count-difference-integral}
 \int_{\tau_0}^\infty\frac{\mathfrak e_1+\mathfrak s_1}{\tau^2}\,d\tau=\lim_{X\to\infty}\int_{\tau_0}^X\frac{N_\zeta(\tau)-n_{\RVM}(\tau)}\tau\,d\tau=\lim_{K\to\infty}\sum_{k\le K}\log\frac{\tau_k}{\gamma_k}.
\end{equation}
In the last step, $\int_{\tau_0}^X(N_\zeta-n_{\RVM})\tau^{-1}d\tau=\sum_{\gamma_k\le X}\log\frac X{\gamma_k}-\sum_{\tau_k\le X}\log\frac X{\tau_k}$; the indices counted by only one of the two sums are $O(\log X)$ in number and, by Lemma~\CGref{lem:bounded-displacement}, have ordinates in $[X-C_\eta,X]$, so they contribute $O(\log X/X)$.  Finally $P_{Z1}(\frac12)/P_{\RVM}^{\crit}(\frac12)=\prod_k\frac{\gamma_k^2}{\frac14+\gamma_k^2}\cdot\frac{\frac14+\tau_k^2}{\tau_k^2}$, an absolutely convergent product, and multiplying by $\prod_{k\le K}\tau_k^2/\gamma_k^2$ gives \CGeqref{eq:varpi-Z1-product}.
\end{proof}

The ordinates enter $\varpi_{Z1}$ only through the displacements $\gamma_k-\tau_k$, in agreement with Remark~\CGref{rem:only-gamma-tau}.

\begin{theorem}[The value at infinity decides the endpoint identity]\label{thm:varpi-Z1}
$\varpi_{Z1}\ge1$, and $\varpi_{Z1}=1$ if and only if $\zeta_{Z1}=\zeta$.
\end{theorem}

\begin{proof}
Write $\zeta_{Z1}=\zeta\,\mathscr R_{Z1}$ with $\mathscr R_{Z1}=P_{Z1}/P_\zeta$ as in \CGeqref{eq:Z1-defect-ratio}, and compare the products factor by factor, with the zeros $\rho_k=\beta_k+i\gamma_k$ listed with the ordinates (Theorem~\CGref{thm:Hadamard}).  Factors with $\rho_k(1-\rho_k)=\frac14+\gamma_k^2$ cancel.  The remaining factors of $P_\zeta$ come in pairs $\{\rho,1-\bar\rho\}$ with the same ordinate $\gamma$; with $\delta=\Re\rho-\frac12$ and $W=(s-\frac12)^2+\gamma^2$, the corresponding factor of $\mathscr R_{Z1}$ is
\begin{equation}\label{eq:pair-ratio}
 \frac{W^2}{(W-\delta^2)^2+4\delta^2\gamma^2}\cdot c_\gamma,\qquad
 c_\gamma=\frac{|\rho(1-\rho)|^2}{(\frac14+\gamma^2)^2}=1+\frac{2\delta^2(\gamma^2-\frac14)+\delta^4}{(\frac14+\gamma^2)^2}>1 .
\end{equation}
For real $\sigma>1$, $W\ge\gamma^2\ge\gamma_1^2$, and the logarithm of the first factor is at most $14\delta^2/W\le14\delta^2/\gamma^2$ in modulus; since $\delta\le\frac12$ and $\sum\gamma^{-2}<\infty$, dominated convergence and $\zeta(\sigma)\to1$ give $\varpi_{Z1}=\lim_{\sigma\to\infty}\zeta_{Z1}(\sigma)=\prod c_\gamma\ge1$, with equality exactly when there are no such pairs, that is, when $\mathscr R_{Z1}\equiv1$.
\end{proof}

\begin{corollary}[The endpoint identity as an identity of the displacements]\label{cor:product-identity}
The following are equivalent: (D); $\zeta_{Z1}=\zeta$; $\varpi_{Z1}=1$;
\begin{equation}\label{eq:product-identity}
 \prod_{k\ge1}\frac{\tfrac14+\gamma_k^2}{\tfrac14+\tau_k^2}=\varpi_{\RVM}
\end{equation}
(ordered product); and $\lim_{\sigma\to+\infty}\Pi_{\rm tr}(\sigma)=L=1/\varpi_{\RVM}$.
\end{corollary}

\begin{proof}
Theorems~\CGref{thm:Hamburger-Z1}, \CGref{thm:Z1-clock} and~\CGref{thm:varpi-Z1}, and $\Pi_{\rm tr}=\zeta_{Z1}/\zeta_{\RVM}\to\varpi_{Z1}/\varpi_{\RVM}$ by Theorems~\CGref{thm:halfplane-RVM} and~\CGref{thm:Z1-clock}(ii).
\end{proof}

\begin{proposition}[Boundedness of the endpoint]\label{prop:Z1-bounded}
Let $\sigma_0>1$.  Then $|\zeta_{Z1}|$ is bounded above and below on $\Re s\ge\sigma_0$ if and only if
\begin{equation}\label{eq:conjugate-poisson}
 \sup_{\sigma'\ge\sigma_0-\frac12,\ t\ge0}\Bigl|\Re\int_{\tau_0}^\infty\frac{\mathfrak e_1(\tau)+\mathfrak s_1(\tau)}{(\tau-t+i\sigma')^2}\,d\tau\Bigr|<\infty .
\end{equation}
\end{proposition}

\begin{proof}
$\tau+iz=\tau-t+i\sigma'$.  For $t\ge0$ the term with $(\tau-iz)^{-2}$ in \CGeqref{eq:Z1-clock} is bounded by \CGeqref{eq:frak-e1-bound}, since $|\tau-iz|\ge\tau+t$; the term with $\mathfrak s_1$ is bounded by Lemma~\CGref{lem:saw-primitive}.  By \CGeqref{eq:Z1-clock}, $\log|\zeta_{Z1}|$ is bounded on $\Re s\ge\sigma_0$, $t\ge0$, exactly when \CGeqref{eq:conjugate-poisson} holds, and $\zeta_{Z1}(\bar s)=\overline{\zeta_{Z1}(s)}$ covers $t<0$.
\end{proof}

The function in \CGeqref{eq:conjugate-poisson} is, after one integration by parts, the conjugate Poisson integral at distance $\sigma'$ of the count difference $N_\zeta-n_{\RVM}$; it depends only on the displacements $\gamma_k-\tau_k$.

\subsection{The Euler products of \texorpdfstring{$\zeta_{\RVM}$}{zeta\_RVM} and \texorpdfstring{$\zeta_{Z1}$}{zeta\_Z1}}
\label{subsec:euler-products}

\begin{theorem}[Prime-power counts as Mellin transforms]\label{thm:mellin-counts}
Let $X\in\{\RVM,Z1\}$.  There is a unique function $J_X^{\rm M}$ on $(1,\infty)$ with $J_X^{\rm M}(x)x^{-c}\in L^2((1,\infty),dx/x)$ for every $c>1$ and
\begin{equation}\label{eq:mellin-count}
 \log\zeta_X(s)=s\int_1^\infty J_X^{\rm M}(x)\,x^{-s-1}\,dx\qquad(\Re s>1),
\end{equation}
the branch of $\log\zeta_X$ being real on $(1,\infty)$.  Equivalently $\zeta_X(s)=\exp\bigl(\int_{[1,\infty)}x^{-s}\,dJ_X^{\rm M}(x)\bigr)$.
\end{theorem}

\begin{proof}
$\zeta_X$ has no zeros or poles in $\Re s>1$, so $\log\zeta_X$ is holomorphic there.  For $c>1$, Theorem~\CGref{thm:halfplane-RVM} ($X=\RVM$) and Theorem~\CGref{thm:Z1-clock} ($X=Z1$), together with continuity on the compact part, give $|\log\zeta_X(s)|\le C_c(1+\log(|t|+2))$ on $\Re s\ge c$.  Hence $F_\sigma(t):=\log\zeta_X(\sigma+it)/(\sigma+it)$ lies in $L^2(\R)$ with $\|F_\sigma\|_2\le B_c$ for all $\sigma\ge c$.  In the variable $u=\log x$, define $J_X^{\rm M}(e^u)e^{-\sigma u}$ as the inverse Fourier--Plancherel transform of $F_\sigma$; Cauchy's theorem and the decay of $F$ on horizontal segments show that $J_X^{\rm M}$ does not depend on $\sigma$.  Plancherel gives $\int_{\R}|J_X^{\rm M}(e^u)|^2e^{-2\sigma u}du=\|F_\sigma\|_2^2/2\pi\le B_c^2/2\pi$ for all $\sigma\ge c$; letting $\sigma\to\infty$ forces $J_X^{\rm M}(e^u)=0$ for almost every $u<0$.  Plancherel again gives \CGeqref{eq:mellin-count}, and uniqueness follows from the injectivity of the Fourier transform.
\end{proof}

\begin{theorem}[Distributional explicit formula]\label{thm:distributional-EF}
Let $X\in\{\RVM,Z1\}$, with ordinates $\kappa_k=\tau_k$ for $X=\RVM$ and $\kappa_k=\gamma_k$ for $X=Z1$.  In $\mathcal D'(1,\infty)$,
\begin{equation}\label{eq:distributional-EF}
 (\log x)\,\frac{dJ_X^{\rm M}}{dx}=1-\frac2{\sqrt x}\,\Re\sum_{k\ge1}x^{i\kappa_k}-\frac1{x(x^2-1)},
\end{equation}
the series converging in $\mathcal D'(1,\infty)$.  The explicit-formula count
\begin{equation}\label{eq:J-X}
 J_X(x):=\li(x)-2\Re\sum_{k\ge1}\Ei\Bigl(\bigl(\tfrac12+i\kappa_k\bigr)\log x\Bigr)-\log2+\int_x^\infty\frac{dy}{y(y^2-1)\log y}
\end{equation}
converges in $\mathcal D'(1,\infty)$ ($J_{\RVM}$ is the function \CGeqref{eq:J-RVM}), and $J_X^{\rm M}-J_X$ is a constant $C_X$ on $(1,\infty)$.
\end{theorem}

\begin{proof}
With $\rho_k=\frac12+i\kappa_k$, the factorization $\zeta_X=P_X/G$, $P_X'/P_X=\sum_k\bigl(\frac1{s-\rho_k}+\frac1{s-\bar\rho_k}\bigr)$, and the digamma expansion $\frac12\frac{\Gamma'}{\Gamma}(\frac s2)+\frac1s=-\frac{\gamma_E}2+\sum_{n\ge1}\bigl(\frac1{2n}-\frac1{s+2n}\bigr)$ give, for $\Re s>1$,
\begin{equation}\label{eq:log-derivative-X}
 -\frac{\zeta_X'}{\zeta_X}(s)=\frac1{s-1}-\sum_{k\ge1}\Bigl(\frac1{s-\rho_k}+\frac1{s-\bar\rho_k}\Bigr)+\sum_{n\ge1}\Bigl(\frac1{2n}-\frac1{s+2n}\Bigr)-\frac{\gamma_E}2-\frac12\log\pi .
\end{equation}
Each term is a Mellin--Stieltjes transform $\int_{[1,\infty)}x^{-s}d\mu$: $dx\mapsto\frac1{s-1}$, $x^{\rho-1}dx\mapsto\frac1{s-\rho}$, $x^{-2n-1}dx\mapsto\frac1{s+2n}$, $\delta_1\mapsto1$.  After multiplication by $x^{-c}$ and the change $u=\log x$, the Fourier transforms of the partial sums over $k$ and over $n$ converge locally uniformly with polynomial bounds, so the series converge in $\mathcal S'$, and on $(1,\infty)$ the limit measure is the right side of \CGeqref{eq:distributional-EF} times $dx$.  Differentiating \CGeqref{eq:mellin-count} in $s$ shows that $-\zeta_X'/\zeta_X$ is the Mellin--Stieltjes transform of $(\log x)\,dJ_X^{\rm M}$; injectivity gives \CGeqref{eq:distributional-EF}.  The series over $k$ converges in $\mathcal D'(1,\infty)$ because $\int\varphi(x)x^{-1/2+i\kappa}dx$ is rapidly decreasing in $\kappa$ for $\varphi\in C_c^\infty(1,\infty)$ and $\kappa_k\asymp k/\log k$.  By Lemma~\CGref{lem:Ei}, $\Ei(\rho_ku)=e^{\rho_ku}/(\rho_ku)+\varrho_k$ with $\sum|\varrho_k|$ locally uniformly convergent, so \CGeqref{eq:J-X} converges in $\mathcal D'(1,\infty)$ for the same reason; differentiating termwise with $\frac d{dx}\Ei(\rho\log x)=x^{\rho-1}/\log x$, $(\log x)J_X'$ equals the right side of \CGeqref{eq:distributional-EF}.  A distribution on an interval with zero derivative is constant.
\end{proof}

\begin{proof}[Proof of Theorem~\CGref{thm:MS}]
Take $X=\RVM$ in Theorems~\CGref{thm:mellin-counts} and~\CGref{thm:distributional-EF}.  By Theorem~\CGref{thm:J-exists}, $J_{\RVM}$ is continuous on $(1,\infty)$, so $J_{\RVM}^{\rm M}=J_{\RVM}+C_{\RVM}$ pointwise, and $J_{\RVM}^{*}=J_{\RVM}$ on $[2,\infty)$.
\end{proof}

\begin{theorem}[Cell form of the endpoint identity]\label{thm:cell-form}
The following are equivalent.
\begin{enumerate}[label=\textup{(\roman*)}]
\item (D), that is, $\zeta_{Z1}=\zeta$.
\item $J_{Z1}^{\rm M}$ is constant on every interval $(n,n+1)$, $n\ge1$, and its jumps $\mathsf b_n:=v_n-v_{n-1}$, where $v_n$ is its value on $(n,n+1)$ and $v_0:=0$, satisfy $\sum_{n\le X}|\mathsf b_n|=O_\eps(X^{1+\eps})$ for every $\eps>0$.
\item For every $n\ge1$, in $\mathcal D'(n,n+1)$,
\begin{equation}\label{eq:cell-identity}
 2\,\Re\sum_{k\ge1}\bigl(x^{i\gamma_k}-x^{i\tau_k}\bigr)=\sqrt x\,\log x\;J_{\RVM}'(x),
\end{equation}
and the jumps of $J_{Z1}^{\rm M}$ satisfy the bound in (ii).
\end{enumerate}
In this case $J_{Z1}^{\rm M}$ is Riemann's prime-power count $J(x)=\sum_{p^m\le x}1/m$.
\end{theorem}

\begin{proof}
(ii)$\Rightarrow$(i): $|v_n|\le\sum_{m\le n}|\mathsf b_m|\ll n^{1+\eps}$, so for $\Re s>1$ summation by parts in \CGeqref{eq:mellin-count} gives $\log\zeta_{Z1}(s)=\sum_nv_n(n^{-s}-(n+1)^{-s})=\sum_{n\ge1}\mathsf b_nn^{-s}$, absolutely convergent by partial summation.  The exponential of an absolutely convergent Dirichlet series is one, so $\zeta_{Z1}$ satisfies \CGeqref{eq:Z1-D}.

(i)$\Rightarrow$(ii): by Theorem~\CGref{thm:Hamburger-Z1}, $\zeta_{Z1}=\zeta$, and $\log\zeta(s)=s\int_1^\infty J(x)x^{-s-1}dx$ for $\Re s>1$; by uniqueness in Theorem~\CGref{thm:mellin-counts}, $J_{Z1}^{\rm M}=J$, which is constant between consecutive integers, with jumps $1/m$ at $p^m$ and $\sum_{n\le X}|\mathsf b_n|=J(X)\ll X$.

(ii)$\Leftrightarrow$(iii): Theorem~\CGref{thm:distributional-EF} for $X=Z1$ and $X=\RVM$, with $dJ_{\RVM}^{\rm M}/dx=J_{\RVM}'$, gives in $\mathcal D'(1,\infty)$
\begin{equation}\label{eq:cell-difference}
 (\log x)\Bigl(\frac{dJ_{Z1}^{\rm M}}{dx}-J_{\RVM}'(x)\Bigr)=-\frac2{\sqrt x}\,\Re\sum_{k\ge1}\bigl(x^{i\gamma_k}-x^{i\tau_k}\bigr),
\end{equation}
so $dJ_{Z1}^{\rm M}/dx=0$ on $(n,n+1)$ is \CGeqref{eq:cell-identity}.
\end{proof}

\begin{remark}[The final step in cell form]\label{rem:cell-form}
Theorem~\CGref{thm:cell-form} is the exact form of the continuous-to-atomic passage of Remark~\CGref{rem:D-atomic}.  Both sides of \CGeqref{eq:cell-identity} are built from the midpoint construction: the right side from the clock alone ($J_{\RVM}'$ is bounded on $[2,\infty)$ by Theorem~\CGref{thm:unit-cell}, equal to $1/\log x$ plus the unit-periodic first harmonic and the harmonics with stationary points $2\pi x^{1/m}$, Section~\CGref{sec:integer-weights}), the left side from the displacements $\gamma_k-\tau_k$.  The final step of the method is the identity \CGeqref{eq:cell-identity} on every cell, equivalently \CGeqref{eq:product-identity}.  Numerically, the explicit formula \CGeqref{eq:J-X} with the $\gamma_k$ is constant between consecutive integers up to the truncation of the list (Experiment~\CGref{exp:primes}, Figure~\CGref{fig:prime-counts}), and the partial products of \CGeqref{eq:product-identity} approach $\varpi_{\RVM}$ (Experiment~\CGref{exp:product-identity}).
\end{remark}

\section{The atomized family and the programme for the final step}
\label{sec:atomized}

The continuous family $\zeta_\lambda$ keeps, for every $\lambda$, the functional equation and the zeros on the critical line.  Its atomization $\tilde\zeta_{\Z,\lambda}$, built from the cell masses of the prime coordinate, keeps the Euler product and the Dirichlet series.  The two are joined by a factor $e^{\tilde h_\lambda}$, and both families have a limit at $\lambda=0$.  Hamburger's theorem needs the functional equation and the Dirichlet series on one and the same function; at $\lambda=0$ this is the vanishing of the joining factor.  This section constructs the atomized family for every $\lambda$, proves its analytic properties, writes the joining factor cell by cell, and lists in Subsection~\CGref{subsec:programme} the steps that remain, in the order in which they will be followed.  The passage from a continuous prime system to a discrete one is classical for Beurling generalized primes \cite{Beurling1937,DiamondMontgomeryVorhauer2006,BrouckeVindas2024,DiamondZhang2016}; the atomization used here places the masses at the integers, which produces ordinary Dirichlet series.

Throughout, $\eta_k=\tau_k-\gamma_k$, $\gamma_k(\lambda)=(1-\lambda)\gamma_k+\lambda\tau_k$, $\rho_k(\lambda)=\frac12+i\gamma_k(\lambda)$, and $\lambda'$ denotes an integration variable in $[0,1]$.

\subsection{The family at infinity}
\label{subsec:family-infinity}

Put $N_\lambda(\tau)=\#\{k:\gamma_k(\lambda)\le\tau\}$, $\mathfrak e^{(\lambda)}=N_\lambda-\NR$ and $\mathfrak e_1^{(\lambda)}(\tau)=\int_{\tau_0}^\tau\mathfrak e^{(\lambda)}$, so that $\mathfrak e_1^{(0)}=\mathfrak e_1$ and $\mathfrak e_1^{(1)}=-\mathfrak s_1$.

\begin{proposition}[Clock decomposition along the family]\label{prop:family-clock}
There is an absolute constant $C$ with $|\mathfrak e_1^{(\lambda)}(\tau)|\le C\log\tau$ for $\tau\ge\tau_0$ and all $\lambda\in[0,1]$.  Consequently Theorem~\CGref{thm:Z1-clock} holds for every $\zeta_\lambda$ with $\mathfrak e_1$ replaced by $\mathfrak e_1^{(\lambda)}$, uniformly in $\lambda$; in particular $\varpi_\lambda:=\lim_{\sigma\to+\infty}\zeta_\lambda(\sigma)$ exists and
\begin{equation}\label{eq:varpi-lambda}
 \varpi_\lambda=\varpi_{\RVM}\prod_{k\ge1}\frac{\frac14+\tau_k^2}{\frac14+\gamma_k(\lambda)^2},\qquad \varpi_1=\varpi_{\RVM},\quad\varpi_0=\varpi_{Z1}.
\end{equation}
\end{proposition}

\begin{proof}
$\gamma_k(\lambda)$ lies between $\gamma_k$ and $\tau_k$, so $\gamma_1(\lambda)>\tau_0$.  For $X\ge\tau_0$,
\begin{equation}\label{eq:family-count-primitive}
 \int_{\tau_0}^X(N_\lambda-n_{\RVM})=\sum_k\bigl((X-\gamma_k(\lambda))^+-(X-\tau_k)^+\bigr)=(1-\lambda)\sum_{\gamma_k(\lambda),\tau_k\le X}\eta_k+O(\log X),
\end{equation}
the indices counted by only one of the two sums being $O(\log X)$ in number with ordinates in $[X-C_\eta,X]$ (Lemma~\CGref{lem:bounded-displacement}).  The case $\lambda=0$ of \CGeqref{eq:family-count-primitive} identifies the last sum with $\mathfrak e_1+\mathfrak s_1+O(\log X)=O(\log X)$, by \CGeqref{eq:frak-e1-bound} and Lemma~\CGref{lem:saw-primitive}.  Since $\mathfrak e_1^{(\lambda)}=\int(N_\lambda-n_{\RVM})-\mathfrak s_1$, the bound follows, and the proof of Theorem~\CGref{thm:Z1-clock} applies verbatim.
\end{proof}

Define the \emph{flow at infinity}
\begin{equation}\label{eq:Upsilon}
 \Upsilon(\lambda):=\sum_{k\ge1}\frac{2\gamma_k(\lambda)\,\eta_k}{\frac14+\gamma_k(\lambda)^2}\qquad(\text{summed in the order of }k).
\end{equation}

\begin{proposition}[The value at infinity along the family]\label{prop:family-flow}
The series \CGeqref{eq:Upsilon} converges uniformly on $[0,1]$, $\lambda\mapsto\log\varpi_\lambda$ is $C^1$ on $[0,1]$, and
\begin{equation}\label{eq:varpi-derivative}
 \frac{d}{d\lambda}\log\varpi_\lambda=-\Upsilon(\lambda)=\lim_{\sigma\to+\infty}\mathcal V_\lambda(\sigma),\qquad
 \Upsilon'(\lambda)=\sum_{k\ge1}\frac{2\eta_k^2\bigl(\frac14-\gamma_k(\lambda)^2\bigr)}{\bigl(\frac14+\gamma_k(\lambda)^2\bigr)^2}<0 ,
\end{equation}
with $\mathcal V_\lambda$ the generator of Theorem~\CGref{thm:generator}.  Moreover
\begin{equation}\label{eq:flow-integral}
 \int_{0^+}^1\Upsilon(\lambda)\,d\lambda=\sum_{k\ge1}\log\frac{\frac14+\tau_k^2}{\frac14+\gamma_k^2}=\log L+\log\varpi_{Z1},
\end{equation}
so that (D) holds if and only if $\int_{0^+}^1\Upsilon(\lambda)\,d\lambda=\log L$, and $\Upsilon(1)<\log L+\log\varpi_{Z1}<\Upsilon(0)$.
\end{proposition}

\begin{proof}
Put $g(u)=2u/(\frac14+u^2)$.  The series $\Upsilon(1)=\sum_kg(\tau_k)\eta_k$ converges by summation by parts against $N_\zeta-n_{\RVM}$, exactly as in \CGeqref{eq:count-difference-integral}.  Since $|g(\gamma_k(\lambda))-g(\tau_k)|\le C|\eta_k|\gamma_k^{-2}$ and $\sum_k\eta_k^2\gamma_k^{-2}<\infty$, the series $\Upsilon(\lambda)-\Upsilon(1)$ converges absolutely and uniformly; the same holds for the termwise derivative, which is the series in \CGeqref{eq:varpi-derivative} (each term is negative because $\gamma_k(\lambda)>\frac12$).  Differentiating \CGeqref{eq:varpi-lambda} termwise gives $\frac d{d\lambda}\log\varpi_\lambda=-\Upsilon(\lambda)$, and letting $\sigma\to+\infty$ in the generator \CGeqref{eq:generator} gives the same series.  Integrating from $0$ to $1$ and using $\varpi_1=\varpi_{\RVM}$, $\varpi_0=\varpi_{Z1}$ gives \CGeqref{eq:flow-integral}; the last statement follows from Theorem~\CGref{thm:varpi-Z1} and the monotonicity of $\Upsilon$.
\end{proof}

\begin{proposition}[The value at infinity and the verified zeros]\label{prop:varpi-verified}
If every zero $\beta+i\gamma$ with $0<\gamma\le T_0$ has $\beta=\frac12$, then
\begin{equation}\label{eq:varpi-verified}
 0\le\log\varpi_{Z1}\le0.26\sum_{\gamma>T_0}\frac1{\gamma^2}.
\end{equation}
With $T_0=3\cdot10^{12}$ \cite{PlattTrudgian2021}, the right side is $\approx\frac{0.26}{2\pi T_0}\bigl(\log\frac{T_0}{2\pi}+1\bigr)<4\cdot10^{-13}$.
\end{proposition}

\begin{proof}
By the proof of Theorem~\CGref{thm:varpi-Z1}, $\log\varpi_{Z1}=\sum\log c_\gamma$ over the pairs $\{\rho,1-\bar\rho\}$ with $\rho(1-\rho)\ne\frac14+\gamma^2$, and all of them have $\gamma>T_0$.  For $\delta\le\frac12$ and $\gamma\ge T_0$, \CGeqref{eq:pair-ratio} gives $\log c_\gamma\le(2\delta^2\gamma^2+\delta^4)\gamma^{-4}\le0.51\,\gamma^{-2}$; a pair consists of two zeros, which gives \CGeqref{eq:varpi-verified}.  The numerical value follows from $N_\zeta(T)=\frac T{2\pi}\log\frac T{2\pi e}+O(\log T)$.
\end{proof}

\subsection{Atomization by smooth cells}
\label{subsec:smooth-cells}

Fix $\varphi_n\in C_c^\infty(n-1,n+1)$, $n\ge2$, with $\sum_{n\ge2}\varphi_n=1$ on $[2,\infty)$ and $\varphi_n(m)=\delta_{nm}$ for integers $m$.  For $\lambda\in[0,1]$ let $J_\lambda$ be the explicit-formula count \CGeqref{eq:J-X} with the ordinates $\gamma_k(\lambda)$ (a distribution on $(1,\infty)$ by Theorem~\CGref{thm:distributional-EF}, which applies to every $\zeta_\lambda$ by Proposition~\CGref{prop:family-clock}), $J^*_\lambda=J_\lambda\mathbf 1_{[2,\infty)}$, and $\pi_\lambda=\sum_{j\ge1}\frac{\mu(j)}jJ^*_\lambda(x^{1/j})$.

\begin{definition}[Atomized family]\label{def:atomized}
The cell masses, the atomized zeta function and its coefficients are
\begin{equation}\label{eq:atomized-family}
 \tilde c_n(\lambda):=\langle d\pi_\lambda,\varphi_n\rangle,\qquad
 \tilde\zeta_{\Z,\lambda}(s):=\prod_{n\ge2}\bigl(1-n^{-s}\bigr)^{-\tilde c_n(\lambda)}=\sum_{N\ge1}\frac{\tilde a_N(\lambda)}{N^s},
\end{equation}
\begin{equation}\label{eq:atomized-coefficients}
 \tilde a_N(\lambda)=\sum_{\prod_nn^{r_n}=N}\ \prod_{n:\,r_n\ge1}\binom{\tilde c_n(\lambda)+r_n-1}{r_n},\qquad \tilde a_1(\lambda)=1 .
\end{equation}
\end{definition}

Unlike the discretizations of Beurling systems, where the atoms are placed at chosen real points \cite{DiamondMontgomeryVorhauer2006,BrouckeVindas2024,DiamondZhang2016}, the masses of $d\pi_\lambda$ are placed at the integers with real exponents; measures with continuous part and non-integer multiplicities of this kind appear also in \cite[Ch.~3]{LapidusvanFrankenhuysen2000}.  Thus $\tilde a_p(\lambda)=\tilde c_p(\lambda)$ for primes $p$, $\tilde a_4=\tilde c_4+\frac12\tilde c_2(\tilde c_2+1)$, $\tilde a_6=\tilde c_6+\tilde c_2\tilde c_3$.  The pairing with a fixed smooth partition defines $\tilde c_n(\lambda)$ for every $\lambda$ without point values of $\pi_\lambda$, and $\lambda\mapsto\tilde c_n(\lambda)$ is continuous.  If $d\pi_0=\sum_p\delta_p$, then $\tilde c_n(0)=\mathbf 1_{\rm prime}(n)$ and $\tilde a_N(0)=1$ for all $N$.  Along the family, in $\mathcal D'(1,\infty)$,
\begin{equation}\label{eq:J-lambda-derivative}
 \frac{\partial J_\lambda}{\partial\lambda}(x)=2\sum_{k\ge1}\eta_k\,\Im\frac{x^{\rho_k(\lambda)}}{\rho_k(\lambda)},\qquad
 \frac{\partial}{\partial\lambda}\Bigl((\log x)\frac{dJ_\lambda}{dx}\Bigr)=\frac{2\log x}{\sqrt x}\sum_{k\ge1}\eta_k\sin\bigl(\gamma_k(\lambda)\log x\bigr).
\end{equation}

\subsection{The joining factor cell by cell}
\label{subsec:joining-factor}

Define the \emph{atomization kernel} and the \emph{cell defects}
\begin{equation}\label{eq:atomization-kernel}
 \mathfrak a_s(x):=\sum_{n\ge2}\varphi_n(x)\,n^{-s}-x^{-s},\qquad
 \mathfrak h_n(\lambda;s):=\bigl\langle d\pi_\lambda,\ \varphi_n(x)\bigl(n^{-s}-x^{-s}\bigr)\bigr\rangle .
\end{equation}
$\mathfrak a_s$ vanishes at every integer $\ge2$, and $\mathfrak h_n(\lambda;s)$ measures how far the mass of the cell $n$ is from being concentrated at $n$.

The explicit-formula count $J_\lambda$ need not vanish on $(1,2)$, whereas $\pi_\lambda$ is built from $J^*_\lambda=J_\lambda\mathbf 1_{[2,\infty)}$.  The first cell $(1,2)$ enters the joining factor through
\begin{equation}\label{eq:first-cell}
 \mathfrak u_\lambda(s):=C_\lambda\,2^{-s}+s\int_1^2J^{\rm M}_\lambda(x)\,x^{-s-1}\,dx,\qquad C_\lambda:=J^{\rm M}_\lambda-J_\lambda ,
\end{equation}
where $J^{\rm M}_\lambda$ is the Mellin count of $\zeta_\lambda$ (Theorem~\CGref{thm:mellin-counts}, which applies to every $\zeta_\lambda$ by Proposition~\CGref{prop:family-clock}) and $C_\lambda$ is the constant of Theorem~\CGref{thm:distributional-EF}.

\begin{proposition}[The first cell]\label{prop:first-cell}
Let $\lambda\in[0,1]$.
\begin{enumerate}[label=\textup{(\roman*)}]
\item $\mathfrak u_\lambda$ is entire, and for $\Re s>1$
\begin{equation}\label{eq:first-cell-split}
 \log\zeta_\lambda(s)=\mathfrak u_\lambda(s)+\langle dJ^*_\lambda,x^{-s}\rangle,\qquad \langle dJ^*_\lambda,x^{-s}\rangle:=s\int_2^\infty J_\lambda(x)\,x^{-s-1}\,dx .
\end{equation}
\item $\mathfrak u_\lambda(\sigma+it)=o(|t|)$ as $|t|\to\infty$, uniformly for $\sigma$ in compact sets.
\item $\mathfrak u_\lambda$ is constant if and only if $J_\lambda=0$ on $(1,2)$; in that case $\mathfrak u_\lambda\equiv C_\lambda$.
\end{enumerate}
For Riemann's count $J(x)=\sum_{p^m\le x}1/m$, which vanishes on $(1,2)$, the corresponding term is $0$.
\end{proposition}

\begin{proof}
(i) $J^{\rm M}_\lambda(x)x^{-c}\in L^2((1,\infty),dx/x)$ for $c>1$, so $J^{\rm M}_\lambda\in L^1(1,2)$ and the integral in \CGeqref{eq:first-cell} is entire in $s$.  By \CGeqref{eq:mellin-count}, $\log\zeta_\lambda(s)=s\int_1^2J^{\rm M}_\lambda x^{-s-1}dx+s\int_2^\infty(J_\lambda+C_\lambda)x^{-s-1}dx$, and $s\int_2^\infty C_\lambda x^{-s-1}dx=C_\lambda2^{-s}$.  The second integral is the pairing of $x^{-s}$ with the distributional derivative of $J^*_\lambda$, which vanishes on $(0,2)$.
(ii) $C_\lambda2^{-s}$ is bounded on vertical strips, and $\int_1^2J^{\rm M}_\lambda(x)x^{-s-1}dx=\int_0^{\log2}J^{\rm M}_\lambda(e^u)e^{-\sigma u}e^{-itu}du\to0$ as $|t|\to\infty$ by the Riemann--Lebesgue lemma, uniformly for $\sigma$ in compact sets.
(iii) For a constant $c$ one has $c=c\,2^{-s}+s\int_1^2c\,x^{-s-1}dx$, so $\mathfrak u_\lambda(s)-c=s\int_1^\infty H(x)x^{-s-1}dx$ with $H=(J^{\rm M}_\lambda-c)\mathbf 1_{(1,2)}+(C_\lambda-c)\mathbf 1_{[2,\infty)}$.  By the injectivity used in Theorem~\CGref{thm:mellin-counts}, $\mathfrak u_\lambda\equiv c$ if and only if $H=0$ almost everywhere, that is, $C_\lambda=c$ and $J^{\rm M}_\lambda=C_\lambda$ on $(1,2)$, which is $J_\lambda=0$ on $(1,2)$.
\end{proof}

\begin{theorem}[The joining factor]\label{thm:joining-factor}
Let $\tilde h_\lambda:=\log\tilde\zeta_{\Z,\lambda}-\log\zeta_\lambda$ on the half-plane where $\tilde\zeta_{\Z,\lambda}$ converges absolutely (Theorem~\CGref{thm:atomized-continuation}).  Then
\begin{equation}\label{eq:joining-factor}
 \tilde\zeta_{\Z,\lambda}=\zeta_\lambda\,e^{\tilde h_\lambda},\qquad
 \tilde h_\lambda(s)=-\mathfrak u_\lambda(s)+\sum_{m\ge1}\frac1m\sum_{n\ge2}\mathfrak h_n(\lambda;ms),
\end{equation}
and, explicitly,
\begin{equation}\label{eq:joining-explicit}
 \tilde h_\lambda(s)=-\mathfrak u_\lambda(s)+\sum_{m\ge1}\frac1m\sum_{j\ge1}\frac{\mu(j)}j\int_2^\infty\mathfrak a_{ms}(y^j)\,\frac1{\log y}\Bigl(1-\frac2{\sqrt y}\Re\sum_{k\ge1}y^{i\gamma_k(\lambda)}-\frac1{y(y^2-1)}\Bigr)dy,
\end{equation}
the series over $k$ converging absolutely after integration.  Its derivative along the family is
\begin{equation}\label{eq:joining-derivative}
 \partial_\lambda\tilde h_\lambda(s)=-\partial_\lambda\mathfrak u_\lambda(s)+\sum_{m\ge1}\frac1m\sum_{j\ge1}\frac{\mu(j)}j\int_2^\infty\mathfrak a_{ms}(y^j)\,\frac2{\sqrt y}\sum_{k\ge1}\eta_k\sin\bigl(\gamma_k(\lambda)\log y\bigr)\,dy .
\end{equation}
Here $\partial_\lambda\mathfrak u_\lambda(s)=\mathcal V_\lambda(s)-\partial_\lambda\langle dJ^*_\lambda,x^{-s}\rangle$ by \CGeqref{eq:first-cell-split} and \CGeqref{eq:generator-zeta}.  Moreover $\mathfrak u_\lambda(\sigma)\to\log\varpi_\lambda$ as $\sigma\to+\infty$, and $\mathfrak u_\lambda\equiv\log\varpi_\lambda$ if $J_\lambda$ vanishes on $(1,2)$.
\end{theorem}

\begin{proof}
By the Euler product, $\log\tilde\zeta_{\Z,\lambda}(s)=\sum_{m\ge1}\frac1m\sum_n\tilde c_n(\lambda)n^{-ms}$.  By Proposition~\CGref{prop:first-cell}, $\log\zeta_\lambda(s)=\mathfrak u_\lambda(s)+\langle dJ^*_\lambda,x^{-s}\rangle$; finite inversion (Proposition~\CGref{prop:finite-inversion}, whose proof applies to $J^*_\lambda$) gives $dJ^*_\lambda=\sum_m\frac1m\,d[\pi_\lambda(x^{1/m})]$, and $\langle d[\pi_\lambda(x^{1/m})],x^{-s}\rangle=\langle d\pi_\lambda,y^{-ms}\rangle$.  Inserting $1=\sum_n\varphi_n$ on $[2,\infty)$ gives $\langle d\pi_\lambda,y^{-ms}\rangle=\sum_n\bigl(\tilde c_n(\lambda)n^{-ms}-\mathfrak h_n(\lambda;ms)\bigr)$, hence \CGeqref{eq:joining-factor}.  Letting $\sigma\to+\infty$, $\log\tilde\zeta_{\Z,\lambda}\to0$ (because $\tilde a_1=1$), $\log\zeta_\lambda\to\log\varpi_\lambda$ (Proposition~\CGref{prop:family-clock}) and the double series tends to $0$ by the bounds of Theorem~\CGref{thm:atomized-continuation}; so $\mathfrak u_\lambda(\sigma)\to\log\varpi_\lambda$, and a constant $\mathfrak u_\lambda$ (Proposition~\CGref{prop:first-cell}(iii)) equals $\log\varpi_\lambda$.  Atoms of $d\pi_\lambda$ at integers contribute nothing because the test functions vanish there.  Formula \CGeqref{eq:joining-explicit} is \CGeqref{eq:joining-factor} with $d\pi_\lambda$ written by \CGeqref{eq:distributional-EF} and the substitution $x=y^j$; the integrated terms in $k$ are $O(\gamma_k(\lambda)^{-2})$ because $\mathfrak a_{ms}(2)=0$.  Differentiating in $\lambda$ with \CGeqref{eq:J-lambda-derivative} gives \CGeqref{eq:joining-derivative}; $\mathfrak u_\lambda(s)$ is differentiable in $\lambda$ because $J^{\rm M}_\lambda$ is the Fourier--Plancherel transform of $\log\zeta_\lambda(\sigma+it)/(\sigma+it)$, whose $\lambda$-derivative $\mathcal V_\lambda(\sigma+it)/(\sigma+it)$ is bounded in $L^2(\R)$ locally uniformly in $\lambda$.
\end{proof}

\begin{proposition}[Explicit cell defects]\label{prop:cell-explicit}
For every $n\ge2$, $\lambda\in[0,1]$ and $\Re s>1$,
\begin{equation}\label{eq:cell-explicit}
 \mathfrak h_n(\lambda;s)=\sum_{j\ge1}\frac{\mu(j)}j\int_2^\infty\varphi_n(y^j)\bigl(n^{-s}-y^{-js}\bigr)\,\frac1{\log y}\Bigl(1-\frac2{\sqrt y}\Re\sum_{k\ge1}y^{i\gamma_k(\lambda)}-\frac1{y(y^2-1)}\Bigr)dy,
\end{equation}
\begin{equation}\label{eq:cell-explicit-derivative}
 \partial_\lambda\mathfrak h_n(\lambda;s)=\sum_{j\ge1}\frac{\mu(j)}j\int_2^\infty\varphi_n(y^j)\bigl(n^{-s}-y^{-js}\bigr)\,\frac2{\sqrt y}\sum_{k\ge1}\eta_k\sin\bigl(\gamma_k(\lambda)\log y\bigr)\,dy ,
\end{equation}
the sums over $j$ being finite and the sums over $k$ converging absolutely after integration, because the test functions are smooth with compact support.  At $\lambda=1$ the cell defect $\mathfrak h_n(1;\cdot)$ is the dispersion of the cell $n$ for $\zeta_{\RVM}$, computed from the clock alone; $d\pi_{\RVM}$ has a bounded density there (Corollary~\CGref{cor:cn-bounded}), so $\mathfrak h_n(1;\cdot)\equiv0$ holds only if this density vanishes on the support of $\varphi_n$.
\end{proposition}

\begin{proof}
Insert \CGeqref{eq:distributional-EF} for $d\pi_\lambda=\sum_j\frac{\mu(j)}j\,d[J^*_\lambda(x^{1/j})]$ into the definition \CGeqref{eq:atomization-kernel} and substitute $x=y^j$; the atoms of $d\pi_\lambda$ at integers, for instance the jump of $J^*_\lambda$ at $2$, contribute nothing because $n^{-s}-x^{-s}$ vanishes at $x=n$ and $\varphi_n$ vanishes at the other integers.  Differentiation in $\lambda$ uses \CGeqref{eq:J-lambda-derivative}; the factor $\log y$ cancels.  For the last statement, $\mathfrak h_n(1;\cdot)\equiv0$ forces $\varphi_n\,d\pi_{\RVM}=\tilde c_n(1)\delta_n$ by the injectivity used in Proposition~\CGref{prop:cell-limit}, and an absolutely continuous measure equals a multiple of $\delta_n$ only if both vanish.
\end{proof}

In the form of the path, the final step requires that in every cell the flow of the displacements cancel exactly the dispersion of the clock, \CGeqref{eq:closing-cells}.

\begin{proposition}[The limit $\lambda\to0$ and the meaning of a vanishing cell]\label{prop:cell-limit}
For every $n\ge2$ and $\Re s>1$, $\mathfrak h_n(\lambda;s)\to\mathfrak h_n(0;s)$ and $\tilde c_n(\lambda)\to\tilde c_n(0)$ as $\lambda\to0^+$, and
\begin{equation}\label{eq:cell-flow}
 \mathfrak h_n(0;s)=\mathfrak h_n(1;s)-\int_{0^+}^1\partial_\lambda\mathfrak h_n(\lambda;s)\,d\lambda .
\end{equation}
Moreover $\mathfrak h_n(0;s)=0$ for all $\Re s>1$ if and only if $\varphi_n\,d\pi_{Z1}=\tilde c_n(0)\,\delta_n$.
\end{proposition}

\begin{proof}
The test function $\varphi_n(x)(n^{-s}-x^{-s})$ is smooth with compact support, and the integrated terms over $k$ are dominated uniformly in $\lambda$ because $|\gamma_k(\lambda)-\gamma_k|\le C_\eta$; dominated convergence gives the limits and \CGeqref{eq:cell-flow}.  Since $\mathfrak h_n(0;s)=\tilde c_n(0)n^{-s}-\langle\varphi_n\,d\pi_{Z1},x^{-s}\rangle$, its vanishing for all $s$ says that the compactly supported distribution $\varphi_n\,d\pi_{Z1}$ has the Mellin transform of $\tilde c_n(0)\delta_n$; the Mellin transform is injective on such distributions.
\end{proof}

\subsection{Mean square of the cell masses}
\label{subsec:mean-square}

From now on the cells are translates of one function: $\varphi_n(x)=\varphi(x-n)$, where $\varphi=\chi/\sum_{m\in\Z}\chi(\cdot-m)$ for a fixed $\chi\in C_c^\infty(-1,1)$ positive on $[-\frac12,\frac12]$.  Then $\sum_{n\ge2}\varphi_n=1$ on $[2,\infty)$, $\varphi_n(m)=\delta_{nm}$, and $\|\varphi_n^{(A)}\|_\infty=\|\varphi^{(A)}\|_\infty=:C_A$ for every $A\ge0$, uniformly in $n$.  For $n\ge2$ put
\begin{equation}\label{eq:VW}
 V_n(\gamma):=\sum_{j\ge1}\frac{\mu(j)}j\int_2^\infty\varphi_n(y^j)\,y^{-1/2}e^{i\gamma\log y}\,dy,\qquad
 W_n(\gamma):=\sum_{j\ge1}\frac{\mu(j)}j\int_2^\infty\varphi_n(y^j)\,\frac{y^{-1/2}e^{i\gamma\log y}}{\log y}\,dy ,
\end{equation}
so that $W_n'=iV_n$; both are rapidly decreasing in $\gamma$ when $n$ is not a power of $2$ (for $n=2^j$ the term $j$ carries the cut of $J^*_\lambda$ at $y=2$ and is treated in the proof of Theorem~\CGref{thm:mean-square-cells}).  The density of $d\pi_\lambda$ used in \CGeqref{eq:cell-explicit} and the derivative \CGeqref{eq:J-lambda-derivative} give, with $\eta_k=\tau_k-\gamma_k$,
\begin{equation}\label{eq:cell-mass-W}
 \tilde c_n(\lambda)=\mathfrak m_n-2\,\Re\sum_{k\ge1}W_n\bigl(\gamma_k(\lambda)\bigr),\qquad
 \partial_\lambda\tilde c_n(\lambda)=2\sum_{k\ge1}\eta_k\,\Im V_n\bigl(\gamma_k(\lambda)\bigr),
\end{equation}
where $\mathfrak m_n$ collects the terms that do not involve the ordinates and is $O(1)$.

\begin{lemma}[Mean value over the moving ordinates]\label{lem:mean-value-ordinates}
There is $C$ such that for all $\lambda\in[0,1]$, $T\ge2$, $u_0\in\R$ and complex $a_k$ with $|a_k|\le1$,
\begin{equation}\label{eq:mean-value-ordinates}
 \int_{u_0-1}^{u_0+1}\Bigl|\sum_{\gamma_k(\lambda)\le T}a_k\,e^{i\gamma_k(\lambda)u}\Bigr|^2du\le C\,T\log^2T .
\end{equation}
\end{lemma}

\begin{proof}
Let $w(u)=C_0\bigl(\sin(\frac{u-u_0}2)/\frac{u-u_0}2\bigr)^2$ with $C_0$ such that $w\ge1$ on $|u-u_0|\le1$; its Fourier transform is bounded and supported in $[-1,1]$.  Expanding the square,
\begin{align}
\int w(u)\Bigl|\sum_ka_ke^{i\gamma_k(\lambda)u}\Bigr|^2du&=\sum_{k,k'}a_k\bar a_{k'}\,e^{i(\gamma_k(\lambda)-\gamma_{k'}(\lambda))u_0}\,\hat w\bigl(\gamma_k(\lambda)-\gamma_{k'}(\lambda)\bigr)\nonumber\\
&\ll\#\{(k,k'):\gamma_k(\lambda)\le T,\ |\gamma_k(\lambda)-\gamma_{k'}(\lambda)|\le1\}.\label{eq:fejer-pairs}
\end{align}
Since $|\gamma_k(\lambda)-\tau_k|\le C_\eta$, every interval $[t,t+1]$ contains $O(\log(t+2))$ of the $\gamma_k(\lambda)$, uniformly in $\lambda$, because it contains that many $\tau_k$ after widening by $C_\eta$.  The number of pairs is therefore $O(T\log^2T)$.
\end{proof}

\begin{theorem}[Mean square of the cell masses along the family]\label{thm:mean-square-cells}
For every $\eps>0$ there is $C_\eps$ such that for all $X\ge2$ and all $\lambda\in[0,1]$
\begin{equation}\label{eq:mean-square-cells}
 \sum_{X<n\le2X}|\tilde c_n(\lambda)|^2\le C_\eps X^{1+\eps},\qquad
 \sum_{X<n\le2X}|\partial_\lambda\tilde c_n(\lambda)|^2\le C_\eps X^{1+\eps}.
\end{equation}
In particular, letting $\lambda\to0^+$ (each $\tilde c_n$ is continuous in $\lambda$ and the constant does not depend on $\lambda$), $\sum_{X<n\le2X}|\tilde c_n(0)|^2\le C_\eps X^{1+\eps}$.
\end{theorem}

\begin{proof}
Fix $\eps$, $X$, $\lambda$ and put $T=X^{1+\eps}$.  By \CGeqref{eq:cell-mass-W}, $\tilde c_n(\lambda)-\mathfrak m_n=-2\Re\sum_{j}\frac{\mu(j)}jA_{n,j}$, with
\begin{equation}\label{eq:A-nj}
 A_{n,j}=\sum_k\int_2^\infty\varphi_n(y^j)\,\frac{y^{-1/2+i\gamma_k(\lambda)}}{\log y}\,dy ,\qquad 1\le j\le J:=\Bigl\lfloor\frac{\log(2X+1)}{\log2}\Bigr\rfloor .
\end{equation}
\emph{Large ordinates.}  In $u=\log y$ the function $\varphi_n(e^{ju})e^{u/2}/u$ has support of length $O((jn)^{-1})$, size $O(n^{1/(2j)})$ and $r$-th derivative $O_r((jn)^rn^{1/(2j)})$: by Fa\`a di Bruno's formula the $r$-th derivative of $\varphi_n(e^{ju})$ is a sum of terms $\varphi_n^{(m)}(e^{ju})(je^{ju})^mj^{r-m}$, $m\le r$, and $\|\varphi_n^{(m)}\|_\infty=C_m$ does not depend on $n$; the derivatives of $e^{u/2}/u$ are $O_r(e^{u/2})$ because $u\ge\log2$.  Integrating by parts $A$ times, each term with $\gamma_k(\lambda)>T$ is $O_A\bigl((jn/\gamma_k(\lambda))^An^{1/(2j)-1}\bigr)$, and their sum over $k$ is $O(X^{-1})$ once $A>(3+\eps)/\eps$.
\emph{The cell $n=2^j$.}  Each block contains exactly one $n=2^j$; for it the term $j$ carries the cut of $J^*_\lambda$ at $y=2$ and is treated directly.  It equals $-\frac{\mu(j)}j\int_2^3J_\lambda(y)\,\frac d{dy}\varphi_n(y^j)\,dy$; the weight has $L^2$ norm $O((jn)^{1/2})$, and $\|J_\lambda\|_{L^2(2,3)}$ is bounded uniformly in $\lambda$ (Theorem~\CGref{thm:mellin-counts} with the uniform bound of Proposition~\CGref{prop:family-clock}).  This term is therefore $O((n/j)^{1/2})$ and contributes $O(X)$ to the sum of squares; the same holds for $\partial_\lambda\tilde c_n$, with $\partial_\lambda J_\lambda$ in place of $J_\lambda$.
\emph{Small ordinates.}  Let $S(y)=y^{-1/2}(\log y)^{-1}\sum_{\gamma_k(\lambda)\le T}y^{i\gamma_k(\lambda)}$ and let $I_{n,j}$ be the support of $y\mapsto\varphi_n(y^j)$, of length $O(X^{1/j-1})$.  By Cauchy--Schwarz in the cell,
\begin{equation}\label{eq:cell-CS}
 \Bigl|\int\varphi_n(y^j)S(y)\,dy\Bigr|^2\le\|\varphi\|_\infty^2\,|I_{n,j}|\int_{I_{n,j}}|S|^2dy .
\end{equation}
For fixed $j$ the intervals $I_{n,j}$, $X<n\le2X$, cover each point at most twice and lie in $[(X-1)^{1/j},(2X+1)^{1/j}]$, an interval of length at most $2$ in $u=\log y$.  With $dy/y=du$ and Lemma~\CGref{lem:mean-value-ordinates} ($a_k=1$),
\begin{equation}\label{eq:cell-block}
 \sum_{X<n\le2X}\Bigl|\int\varphi_n(y^j)S\,dy\Bigr|^2\ll X^{1/j-1}\Bigl(\frac{j}{\log X}\Bigr)^2\int_{u_0-1}^{u_0+1}\Bigl|\sum_{\gamma_k(\lambda)\le T}e^{i\gamma_k(\lambda)u}\Bigr|^2du\ll j^2X^{1/j+2\eps}.
\end{equation}
\emph{Assembly.}  By Minkowski's inequality over $j$, the $\ell^2$ norm over $X<n\le2X$ of $\tilde c_n(\lambda)-\mathfrak m_n$ is $\ll\sum_{j\le J}\bigl(X^{1/(2j)+\eps}+X^{-1/2}\bigr)\ll X^{1/2+\eps}$, and $\sum_{X<n\le2X}\mathfrak m_n^2\ll X$.  This is the first bound with $3\eps$ in place of $\eps$.  For $\partial_\lambda\tilde c_n$ the argument is the same, by \CGeqref{eq:cell-mass-W}, with coefficients $a_k=\eta_k/C_\eta$ in Lemma~\CGref{lem:mean-value-ordinates} and without the factor $(\log y)^{-1}$.
\end{proof}

\begin{corollary}[Abscissa $1$ for every member]\label{cor:abscissa-one}
For every $\lambda\in[0,1]$ and $\sigma>1$,
\begin{equation}\label{eq:abscissa-one}
 \sum_{n\ge2}|\tilde c_n(\lambda)|\,n^{-\sigma}<\infty ;
\end{equation}
hence $\tilde\zeta_{\Z,\lambda}$ and $1/\tilde\zeta_{\Z,\lambda}$ are absolutely convergent Dirichlet series in $\Re s>1$, for every $\lambda$, including the endpoint $\lambda=0$.  Moreover, uniformly in $\lambda$, $\#\{n\le X:|\tilde c_n(\lambda)|>n^{\delta}\}\ll_\eps X^{1-2\delta+\eps}$ for every $\delta>0$.
\end{corollary}

\begin{proof}
By Cauchy--Schwarz on dyadic blocks, $\sum_{X<n\le2X}|\tilde c_n(\lambda)|\le(X\cdot C_\eps X^{1+\eps})^{1/2}$, and summing over $X=2^i$ gives \CGeqref{eq:abscissa-one} for $\sigma>1+\eps/2$, with $\eps$ arbitrary.  Then $\log\tilde\zeta_{\Z,\lambda}=\sum_m\frac1m\sum_n\tilde c_n(\lambda)n^{-ms}$ converges absolutely in $\Re s>1$, and the exponentials of $\pm$ this series are absolutely convergent Dirichlet series there.  The last statement is Chebyshev's inequality on dyadic blocks.
\end{proof}

The abscissa $\frac32$ obtained below from the primitive $\pi_\lambda-\mathfrak L$ (Theorem~\CGref{thm:atomized-continuation}(ii)) is not the true one: a bound on the size of $\pi_\lambda-\mathfrak L$ does not see the cancellation of the density of $d\pi_\lambda$ in mean square, which Theorem~\CGref{thm:mean-square-cells} captures.  Absolute convergence in $\Re s>1$ is exactly what Hamburger's theorem requires, and Knopp's obstruction (Subsection~\CGref{subsec:hamburger-table}) does not arise.

\subsection{Analytic continuation of the atomized family}
\label{subsec:atomized-continuation}

Let $\mathfrak L(x)=\int_1^x(1-y^{-1})\,dy/\log y$; then $\log\frac s{s-1}=\int_1^\infty x^{-s}\,d\mathfrak L(x)$ for $\Re s>1$ (both sides have derivative $\frac1s-\frac1{s-1}$ and tend to $0$ as $s\to+\infty$), and $\mathfrak L(x)=\li(x)-\log\log x-\gamma_E$.

\begin{theorem}[Continuation of the atomized family]\label{thm:atomized-continuation}
Let $\lambda\in[0,1]$.
\begin{enumerate}[label=\textup{(\roman*)}]
\item For every $c>\frac12$, $(\pi_\lambda-\mathfrak L)\,x^{-c}\in L^2\bigl((2,\infty),dx/x\bigr)$.
\item $|\tilde c_n(\lambda)|\le C(1+I_n^{1/2})$ with $I_n$ as below, and $\sum_n|\tilde c_n(\lambda)|n^{-\sigma}<\infty$ for $\sigma>\frac32$; by Corollary~\CGref{cor:abscissa-one} this holds for $\sigma>1$, and $\tilde\zeta_{\Z,\lambda}$ is an absolutely convergent Dirichlet series in $\Re s>1$.
\item $\tilde h_\lambda$ extends holomorphically to $\Re s>\frac12$.  Hence $\tilde\zeta_{\Z,\lambda}=\zeta_\lambda e^{\tilde h_\lambda}$ is meromorphic in $\Re s>\frac12$, with a single pole at $s=1$ and no zeros.
\item For $\lambda=1$, $\tilde h_1$ extends holomorphically to $\Re s>0$, and the zeros of $\tilde\zeta_{\Z,1}$ in $\Re s>0$ are exactly $\frac12\pm i\tau_k$.
\end{enumerate}
\end{theorem}

\begin{proof}
(i) By Proposition~\CGref{prop:family-clock}, for $c>\frac12$ there is $C_c$ with $|\log\zeta_\lambda(s)-\log\frac s{s-1}|\le C_c(1+\log(|t|+2))$ on $\Re s\ge c$; the function $\log\zeta_\lambda-\log\frac s{s-1}$ is holomorphic on $\Re s>\frac12$, where $\zeta_\lambda$ has no zeros, and its pole at $1$ is removed.  It equals $s\int_1^\infty(J_\lambda^{\rm M}-\mathfrak L)x^{-s-1}dx$ for $\Re s>1$.  The argument of Theorem~\CGref{thm:mellin-counts}, applied on the lines $\Re s=c>\frac12$, gives $(J^{\rm M}_\lambda-\mathfrak L)x^{-c}\in L^2$.  The constant $C_\lambda=J^{\rm M}_\lambda-J_\lambda$, the truncation at $2$, the terms $j\ge2$ of the M\"obius inversion (of size $O(\sqrt x)$ in mean) and $\mathfrak L(x)-\mathfrak L(x^{1/j})$ preserve the property for $c>\frac12$.

(ii) Put $E=\pi_\lambda-\mathfrak L$ and $I_n=\int_{n-1}^{n+1}|E|^2dx$.  Since $\int\varphi_n'=0$, $\tilde c_n(\lambda)=\int\varphi_n\,d\mathfrak L-\int E\varphi_n'\,dx$, so $|\tilde c_n(\lambda)|\le C(1+I_n^{1/2})$.  By (i), $\sum_nI_nn^{-2c-1}<\infty$, and Cauchy--Schwarz gives $\sum_nI_n^{1/2}n^{-\sigma}<\infty$ for $\sigma>c+1$.

(iii) With $\psi_n=\varphi_n(x)(n^{-s}-x^{-s})$ one has $|\psi_n'|\le C|s|n^{-\sigma-1}$ and $\psi_n$ vanishes at the ends of the cell, so
\begin{equation}\label{eq:cell-defect-bound}
 |\mathfrak h_n(\lambda;s)|\le C|s|\,n^{-\sigma-1}\bigl(1+I_n^{1/2}\bigr),\qquad \sum_nI_n^{1/2}n^{-\sigma-1}<\infty\quad(\sigma>c).
\end{equation}
Choosing $c\in(\frac12,\sigma)$, the series \CGeqref{eq:joining-factor} converges locally uniformly on $\Re s>\frac12$; the terms $m\ge2$ are evaluated at $\Re(ms)\ge2\sigma$.  Since $\mathfrak u_\lambda$ is entire (Proposition~\CGref{prop:first-cell}), $\tilde h_\lambda$ is holomorphic on $\Re s>\frac12$, and $\zeta_\lambda$ has no zeros there.

(iv) By Corollary~\CGref{cor:cn-bounded}, $\pi_{\RVM}$ oscillates by $O(1)$ on every unit cell; integrating by parts, $|\mathfrak h_n(1;s)|\le C|s|n^{-\sigma-1}$, and the series converges for $\sigma>0$.  The factor $e^{\tilde h_1}$ has no zeros, so the zeros of $\tilde\zeta_{\Z,1}$ in $\Re s>0$ are those of $\zeta_{\RVM}$.
\end{proof}

\begin{theorem}[Closing criterion; via Hamburger's theorem \cite{Hamburger1921}]\label{thm:closing}
The following are equivalent.
\begin{enumerate}[label=\textup{(\roman*)}]
\item (D), that is, $\zeta_{Z1}=\zeta$.
\item $\tilde h_0\equiv0$.
\item $J_{Z1}=0$ on $(1,2)$, and for every $n\ge2$ and every $\Re s>1$,
\begin{equation}\label{eq:closing-cells}
 \int_{0^+}^1\partial_\lambda\mathfrak h_n(\lambda;s)\,d\lambda=\mathfrak h_n(1;s).
\end{equation}
\end{enumerate}
In this case $\tilde c_n(0)=\mathbf 1_{\rm prime}(n)$, $\tilde a_N(0)=1$ for all $N$, and $\varpi_{Z1}=1$.
\end{theorem}

\begin{proof}
(ii)$\Rightarrow$(i): $\tilde h_0\equiv0$ gives $\zeta_{Z1}=\tilde\zeta_{\Z,0}=\exp\bigl(\sum_m\frac1m\sum_n\tilde c_n(0)n^{-ms}\bigr)$, an absolutely convergent Dirichlet series in $\Re s>1$ by Corollary~\CGref{cor:abscissa-one}.  (iii)$\Rightarrow$(ii): by \CGeqref{eq:cell-flow} every $\mathfrak h_n(0;\cdot)$ vanishes, and $\mathfrak u_0\equiv\log\varpi_{Z1}$ (Theorem~\CGref{thm:joining-factor}), so $\tilde h_0\equiv-\log\varpi_{Z1}$ is constant; then $\zeta_{Z1}=\varpi_{Z1}\tilde\zeta_{\Z,0}$ is an absolutely convergent Dirichlet series in $\Re s>1$ (Corollary~\CGref{cor:abscissa-one}), Theorem~\CGref{thm:Hamburger-Z1} gives $\zeta_{Z1}=\zeta$, and comparing the first coefficients, $\varpi_{Z1}=1$.  (i)$\Rightarrow$(iii): if $\zeta_{Z1}=\zeta$, Theorem~\CGref{thm:cell-form} gives $J^{\rm M}_{Z1}=J$, and $J_{Z1}=J$ by von Mangoldt's explicit formula \cite{Edwards1974}, so $J_{Z1}$ vanishes on $(1,2)$, $d\pi_{Z1}=\sum_p\delta_p$, $\tilde c_n(0)=\mathbf 1_{\rm prime}(n)$, every $\mathfrak h_n(0;\cdot)$ vanishes, and \CGeqref{eq:cell-flow} gives \CGeqref{eq:closing-cells}.
\end{proof}

\begin{remark}[The pointwise size of the cell masses]\label{rem:A6-scale}
Theorem~\CGref{thm:closing} needs no pointwise bound on $\tilde c_n(0)$: its proof uses the cell masses only through \CGeqref{eq:abscissa-one}.  The bound $\tilde c_n(0)=O(n^{\eps})$ for every $n$ is not proved directly; Corollary~\CGref{cor:abscissa-one} restricts the exceptions to a set of $n\le X$ of size $O(X^{1-2\delta+\eps})$ for the threshold $n^{\delta}$, uniformly in $\lambda$, and under the closure every cell mass is bounded, $\tilde c_n(0)=\mathbf 1_{\rm prime}(n)$.  In the ordinates, \CGeqref{eq:cell-mass-W} gives the trivial pointwise bound $\tilde c_n(\lambda)=O(n^{1/2+\eps})$; the mean square \CGeqref{eq:mean-square-cells} shows that the typical cell mass is $O(n^{\eps})$.
\end{remark}

\subsection{The two endpoints of the atomized family}
\label{subsec:atomized-endpoints}

\begin{proposition}[Atomization at the clock endpoint]\label{prop:atomized-clock}
At $\lambda=1$:
\begin{enumerate}[label=\textup{(\roman*)}]
\item $\tilde c_n(1)=O(1)$, and $\tilde\zeta_{\Z,1}$ is an absolutely convergent Dirichlet series in $\Re s>1$;
\item $\tilde h_1(\sigma)\to-\log\varpi_{\RVM}=\log L\ne0$ as $\sigma\to+\infty$;
\item $\tilde h_1$ is not constant: for every $\sigma>1$, $\zeta_{\RVM}(\sigma+it)\to\varpi_{\RVM}$ as $|t|\to\infty$, whereas $\tilde\zeta_{\Z,1}(\sigma+it)$ has no limit as $|t|\to\infty$.
\end{enumerate}
\end{proposition}

\begin{proof}
(i) By Corollary~\CGref{cor:cn-bounded}, $\pi_{\RVM}$ oscillates by $O(1)$ on every unit cell, and $\tilde c_n(1)=-\int\pi_{\RVM}\varphi_n'\,dx$ with $\int\varphi_n'=0$.  (ii) is \CGeqref{eq:joining-factor} with $\sigma\to+\infty$, since $\varpi_{\RVM}<1$.  (iii) The first limit is \CGeqref{eq:vertical-limit}.  If $\tilde\zeta_{\Z,1}(\sigma+it)$ had a limit $\ell$ as $|t|\to\infty$, the mean-square argument \CGeqref{eq:mean-square}--\CGeqref{eq:mean-square-limit} applied to $\tilde\zeta_{\Z,1}-\ell$ on the line $\Re s=\sigma$ would give $\tilde a_N(1)=0$ for $N\ge2$, so $\tilde\zeta_{\Z,1}\equiv1$, which contradicts its pole at $s=1$ (Theorem~\CGref{thm:atomized-continuation}(iii)).  If $\tilde h_1$ were constant, $\zeta_{\RVM}$ would be a constant multiple of an absolutely convergent Dirichlet series, which Theorem~\CGref{thm:no-Dirichlet-RVM} excludes.
\end{proof}

The two families thus keep complementary structures and are joined by a zero-free factor:
\begin{equation}\label{eq:duality}
 \underbrace{\zeta_\lambda}_{\substack{\text{functional equation}\\ \text{zeros }\frac12\pm i\gamma_k(\lambda)}}\cdot\ e^{\tilde h_\lambda}
 \ =\ \underbrace{\tilde\zeta_{\Z,\lambda}}_{\substack{\text{Dirichlet series}\\ \text{Euler product}}},\qquad
 \tilde h_1\ \text{not constant},\qquad (D)\iff\tilde h_0\equiv0 .
\end{equation}
At $\lambda=1$ the joining factor carries the whole difference between the clock and an Euler product; Theorem~\CGref{thm:closing} states that this difference disappears at $\lambda=0$.

The reciprocal of the atomized family is again an Euler product:
\begin{equation}\label{eq:atomized-mobius}
 \frac1{\tilde\zeta_{\Z,\lambda}(s)}=\prod_{n\ge2}\bigl(1-n^{-s}\bigr)^{\tilde c_n(\lambda)}=\sum_{N\ge1}\frac{\tilde\mu_\lambda(N)}{N^s},\qquad
 \tilde\mu_\lambda(N)=\sum_{\prod_nn^{r_n}=N}\ \prod_{n:\,r_n\ge1}(-1)^{r_n}\binom{\tilde c_n(\lambda)}{r_n},
\end{equation}
absolutely convergent in $\Re s>1$ for every $\lambda$ (Corollary~\CGref{cor:abscissa-one}), because it is the exponential of the absolutely convergent series $-\sum_m\frac1m\sum_n\tilde c_n(\lambda)n^{-ms}$.  Thus $\tilde\mu_\lambda(p)=-\tilde c_p(\lambda)$, and the atomized Mertens functions
\begin{equation}\label{eq:atomized-mertens}
 \tilde M_\lambda(x):=\sum_{N\le x}\tilde\mu_\lambda(N),\qquad \frac1{\tilde\zeta_{\Z,\lambda}}=e^{-\tilde h_\lambda}F_\lambda ,
\end{equation}
exist for every $\lambda\in[0,1]$.  Under the conditions of Theorem~\CGref{thm:closing}, $\tilde c_n(0)=\mathbf 1_{\rm prime}(n)$, so $\tilde\mu_0=\mu$ and $\tilde M_0=M$.

\begin{proposition}[Mellin form of the endpoint Mertens function]\label{prop:Mellin-Mertens}
There are $c_1\ge2$ and a unique measurable $B_{Z1}$ on $(0,\infty)$, vanishing on $(0,1)$, with $x^{-c}B_{Z1}\in L^2\bigl((1,\infty),dx/x\bigr)$ and
\begin{equation}\label{eq:Mellin-Mertens}
 F_{Z1}(s)=s\int_1^\infty B_{Z1}(x)\,x^{-s-1}\,dx=\exp\Bigl(-\int_{[1,\infty)}x^{-s}\,dJ^{\rm M}_{Z1}(x)\Bigr)\qquad(\Re s>c_1),
\end{equation}
for every $c>c_1$.  Moreover $(D)$ holds if and only if $B_{Z1}=M$ almost everywhere; and if $F_{Z1}$ has an absolutely convergent Dirichlet expansion \CGeqref{eq:Z1-reciprocal-Dirichlet}, then $B_{Z1}=M_{Z1}$.
\end{proposition}

\begin{proof}
Put $A=2(\pi+2)C_{\mathfrak e}$.  By \CGeqref{eq:Z1-growth}, for $\sigma'\ge\max(4A,R_0)$, $|F_{Z1}(s)|\le C(|t|+\sigma)^{1/4}$; take $c_1=\max(2,\frac12+4A,\frac12+R_0)$.  Then $\int_{\R}|F_{Z1}(\sigma+it)|^2|\sigma+it|^{-2}dt\le C'\sigma^{-1/2}$ for $\sigma>c_1$, and the argument of Theorem~\CGref{thm:mellin-counts} gives $B_{Z1}$, its vanishing on $(0,1)$, its independence of $c$ and its uniqueness; the second equality in \CGeqref{eq:Mellin-Mertens} is Theorem~\CGref{thm:mellin-counts}.  If $\zeta_{Z1}=\zeta$, then $F_{Z1}(s)=1/\zeta(s)=s\int_1^\infty M(x)x^{-s-1}dx$ and uniqueness gives $B_{Z1}=M$.  Conversely, $B_{Z1}=M$ gives $F_{Z1}=1/\zeta$ on $\Re s>c_1$, hence $\zeta_{Z1}=\zeta$ by analytic continuation.  Under \CGeqref{eq:Z1-reciprocal-Dirichlet}, partial summation gives $F_{Z1}(s)=s\int_1^\infty M_{Z1}(x)x^{-s-1}dx$, and uniqueness gives $B_{Z1}=M_{Z1}$.
\end{proof}

\subsection{Hamburger's hypotheses at the two endpoints}
\label{subsec:hamburger-table}

Hamburger's theorem \cite{Hamburger1921}, in the form of \cite[\S2.13]{Titchmarsh1986} (see also \cite{Siegel1922,BochnerChandrasekharan1956}), has four hypotheses: (H0) $f=G/P$ with $G$ entire of finite order and $P$ a polynomial; (H1) $f(s)=\sum_na_nn^{-s}$ absolutely for $\sigma>1$, with arbitrary complex $a_n$; (H2) $f(s)\Gamma(\frac s2)\pi^{-s/2}=g(1-s)\Gamma(\frac{1-s}2)\pi^{-(1-s)/2}$; (H3) $g(1-s)=\sum_nb_nn^{s-1}$ absolutely for $\sigma<-\alpha<0$.  It contains no hypothesis on the location of zeros, and its conclusion is $f=C\zeta$.  Absolute convergence for $\sigma>1$ in (H1) cannot be weakened to some half-plane: Knopp constructed infinitely many linearly independent solutions of (H0), (H2), (H3) that are Dirichlet series absolutely convergent only in some half-plane \cite{Knopp1994,Nakamura2020}.  Within the Selberg class, the functions of degree $1$ are classified in \cite{KaczorowskiPerelli1999}.  The table records the status at the two endpoints and is updated as the programme advances.

\begin{center}
\small
\begin{tabular}{p{0.06\textwidth}p{0.40\textwidth}p{0.44\textwidth}}
\toprule
 & $\zeta_{Z1}$ (continuous limit) & $\tilde\zeta_{\Z,0}$ (atomized limit)\\
\midrule
H0 & proved: $(s-1)\zeta_{Z1}$ entire of order $1$ (Proposition~\CGref{prop:Z1-class}) & open: meromorphic in $\Re s>\frac12$ only (Theorem~\CGref{thm:atomized-continuation})\\
H1 & open: this is (D) & proved in $\Re s>1$ (Corollary~\CGref{cor:abscissa-one})\\
H2 & proved, with $g=f$ (Proposition~\CGref{prop:Z1-class}) & open\\
H3 & follows from H1 & open\\
zeros & no hypothesis & no hypothesis\\
\bottomrule
\end{tabular}
\end{center}

\noindent The two columns meet when $\tilde h_0\equiv0$: then $\zeta_{Z1}=\tilde\zeta_{\Z,0}$ has H0--H3, and Hamburger's theorem gives $\zeta_{Z1}=\zeta$ (Theorem~\CGref{thm:closing}).  It suffices that the joining factor be symmetric (Theorem~\CGref{thm:pair-criterion}), or that $\zeta_{Z1}$ have a positive representing measure (Theorem~\CGref{thm:positivity}).

The parallel between the two families, topic by topic of the paper, is the following.

\begingroup\small
\begin{longtable}{>{\raggedright\arraybackslash}p{0.20\textwidth}>{\raggedright\arraybackslash}p{0.355\textwidth}>{\raggedright\arraybackslash}p{0.355\textwidth}}
\toprule
Topic of the paper & continuous family $\zeta_\lambda$ & atomized family $\tilde\zeta_{\Z,\lambda}$\\
\midrule
\endfirsthead
\toprule
Topic of the paper & continuous family $\zeta_\lambda$ & atomized family $\tilde\zeta_{\Z,\lambda}$\\
\midrule
\endhead
\bottomrule
\endfoot
\addlinespace[2pt]
clock, $\tau_k$, bounded displacement & same & same\\
definition & Hadamard product over $\frac12\pm i\gamma_k(\lambda)$ & Euler product over $n\ge2$ with masses $\tilde c_n(\lambda)=\langle d\pi_\lambda,\varphi_n\rangle$\\
zeros & exactly $\frac12\pm i\gamma_k(\lambda)$; impossible elsewhere (Theorem~\CGref{thm:path}) & none in $\Re s>\frac12$ for every $\lambda$, since $\tilde h_\lambda$ is holomorphic there (Theorem~\CGref{thm:atomized-continuation}(iii)); for $\lambda=1$, exactly $\frac12\pm i\tau_k$ in $\Re s>0$ (Theorem~\CGref{thm:atomized-continuation}(iv))\\
functional equation & proved for every $\lambda$ & open; the factor $e^{\tilde h_\lambda}$ carries no functional equation\\
Dirichlet series & not for $\lambda=1$ in any half-plane (Theorem~\CGref{thm:no-Dirichlet-RVM}); not in $\Re s>1$ for $0<\lambda<1$ (Proposition~\CGref{prop:nonconstant}); at $\lambda=0$ this is (D) & proved in $\Re s>1$ for every $\lambda$, including $\lambda=0$ (Corollary~\CGref{cor:abscissa-one})\\
first coefficient ($\sigma\to+\infty$) & $\varpi_\lambda=\varpi_{\RVM}\prod_k\frac{1/4+\tau_k^2}{1/4+\gamma_k(\lambda)^2}$ \CGeqref{eq:varpi-lambda} & $\tilde a_1(\lambda)=1$ for every $\lambda$\\
vertical behaviour & $\lambda=1$: $\zeta_{\RVM}(\sigma+it)\to\varpi_{\RVM}$ as $|t|\to\infty$ \CGeqref{eq:vertical-limit} & almost periodic on every line of absolute convergence; $\lambda=1$: no limit as $|t|\to\infty$ (Proposition~\CGref{prop:atomized-clock})\\
half-plane bounds & $\lambda=1$: $C^{-1}\le|\zeta_{\RVM}|\le C$ on $\Re s\ge\sigma_0>1$ (Theorem~\CGref{thm:halfplane-RVM}); every $\lambda$: $|\log\zeta_\lambda-\log\frac s{s-1}|\ll_c1+\log(|t|+2)$ on $\Re s\ge c>\frac12$ & $|\tilde\zeta_{\Z,\lambda}|\le\sum_N|\tilde a_N(\lambda)|N^{-\sigma_0}$ and $|\tilde\zeta_{\Z,\lambda}^{-1}|\le\sum_N|\tilde\mu_\lambda(N)|N^{-\sigma_0}$ on $\Re s\ge\sigma_0>1$, for every $\lambda$\\
prime coordinate & counts $J_\lambda$, $\pi_\lambda$ (distributions, \CGeqref{eq:J-X}) & masses $\tilde c_n(\lambda)$ at the integers; $\tilde c_n(1)=O(1)$; mean square $X^{1+\eps}$ on $(X,2X]$ for every $\lambda$ (Theorem~\CGref{thm:mean-square-cells})\\
reciprocal & $F_\lambda=1/\zeta_\lambda$; $F_{Z1}$ is the Mellin transform of $B_{Z1}$ (Proposition~\CGref{prop:Mellin-Mertens}) & Dirichlet series with coefficients $\tilde\mu_\lambda(N)$ \CGeqref{eq:atomized-mobius}\\
Mertens & $M^{\crit}_{\RVM}=L+O_\eps(x^{1/2+\eps})$ (Theorem~\CGref{thm:unsmoothed-RVM}); $B_{Z1}=M\iff(D)$ & $\tilde M_\lambda=\sum_{N\le x}\tilde\mu_\lambda(N)$ exists for every $\lambda$ \CGeqref{eq:atomized-mertens}; the bound $x^{1/2+\eps}$ is open (it needs control of $e^{-\tilde h_\lambda}$)\\
flow in $\lambda$ & generator $\mathcal V_\lambda$ on the Hadamard product; $\Upsilon(\lambda)$ at $+\infty$ (Proposition~\CGref{prop:family-flow}) & $\log\tilde\zeta_{\Z,\lambda}=\sum_n\tilde c_n(\lambda)\bigl(-\log(1-n^{-s})\bigr)$, each mass differentiable in $\lambda$ by \CGeqref{eq:J-lambda-derivative}; joining factor \CGeqref{eq:joining-derivative}\\
limit $\lambda\to0^+$ & $\zeta_\lambda\to\zeta_{Z1}$ & $\tilde c_n(\lambda)\to\tilde c_n(0)$, $\mathfrak h_n(\lambda;\cdot)\to\mathfrak h_n(0;\cdot)$ (Proposition~\CGref{prop:cell-limit})\\
endpoint: Hamburger's hypotheses & H0 and H2 proved; H1 is (D), open & H1 proved in $\Re s>1$ (Corollary~\CGref{cor:abscissa-one}); H0 and H2 open\\
relation between the endpoints & $\tilde\zeta_{\Z,0}=\zeta_{Z1}\,e^{\tilde h_0}$ & same\\
what closes the route & $\zeta_{Z1}$ satisfies H1 & $\tilde\zeta_{\Z,0}$ satisfies H0 and H2\\
common form of the closure & $\tilde h_0\equiv0$ & $\tilde h_0\equiv0$\\
weaker forms of the closure & $\zeta_{Z1}$ has a positive representing measure (Theorem~\CGref{thm:positivity}) & $e^{\tilde h_0}$ of finite order and symmetric under $s\mapsto1-s$ (Theorem~\CGref{thm:pair-criterion})\\
consequence of the closure & $\zeta_{Z1}=\zeta$; Step~3 gives $M(x)=O_\eps(x^{1/2+\eps})$ (Theorem~\CGref{thm:step3}) & $\tilde a_N(0)=1$, $\tilde c_n(0)=\mathbf 1_{\rm prime}(n)$, $\tilde\mu_0=\mu$, $\tilde M_0=M$; $\tilde\zeta_{\Z,0}=\zeta=\zeta_{Z1}$\\
\end{longtable}
\endgroup

\noindent\emph{How the two routes complement each other.}  Along $\lambda$ each family keeps exactly the half of Hamburger's hypotheses that the other one lacks:
\begin{equation}\label{eq:complement}
 \underbrace{\zeta_\lambda}_{\substack{\text{keeps the functional equation}\\ \text{and the zeros on }\Re s=\frac12}}\qquad\longleftrightarrow\qquad\underbrace{\tilde\zeta_{\Z,\lambda}}_{\substack{\text{keeps the Dirichlet series}\\ \text{and the Euler product}}},\qquad \tilde\zeta_{\Z,\lambda}=\zeta_\lambda\,e^{\tilde h_\lambda}.
\end{equation}
The closure is the same in both columns, the vanishing of the joining factor at the endpoint: by \CGeqref{eq:joining-factor} and Theorem~\CGref{thm:closing},
\begin{equation}\label{eq:closure-common}
 \tilde h_0\equiv0\iff J_{Z1}=0\text{ on }(1,2),\ \ \mathfrak h_n(0;s)=0\ \ \text{for all }n\ge2,\ \Re s>1,\ \text{and}\ \varpi_{Z1}=1 .
\end{equation}
No further condition is needed in the atomized column: the cell masses have mean square $X^{1+\eps}$ uniformly in $\lambda$ (Theorem~\CGref{thm:mean-square-cells}), so $\tilde\zeta_{\Z,0}$ satisfies H1 (Corollary~\CGref{cor:abscissa-one}).

\subsection{Programme}
\label{subsec:programme}

The steps are numbered in the order in which they are followed.  The status column is updated with each step.

\begin{center}
\small
\begin{tabular}{p{0.05\textwidth}p{0.55\textwidth}p{0.30\textwidth}}
\toprule
Step & Statement & Status\\
\midrule
A1 & Atomization by smooth cells, coefficients $\tilde a_N(\lambda)$ & proved (Definition~\CGref{def:atomized})\\
A2 & Joining factor and explicit cell defects, \CGeqref{eq:joining-factor}--\CGeqref{eq:cell-explicit-derivative} & proved (Thm.~\CGref{thm:joining-factor}, Prop.~\CGref{prop:cell-explicit})\\
A3 & $L^2$ square-root estimate for every $\lambda$; Dirichlet series in $\Re s>\frac32$ & proved (Theorem~\CGref{thm:atomized-continuation}); superseded by A6\\
A4 & $\tilde h_\lambda$ holomorphic in $\Re s>\frac12$; in $\Re s>0$ for $\lambda=1$ & proved (Theorem~\CGref{thm:atomized-continuation})\\
A5 & Limits $\lambda\to0^+$ of $\tilde c_n$, $\mathfrak h_n$, $\varpi_\lambda$; flow at infinity $\Upsilon$ & proved (Propositions~\CGref{prop:family-flow}, \CGref{prop:cell-limit})\\
A5$'$ & Endpoints: $\tilde h_1$ not constant; atomized M\"obius coefficients $\tilde\mu_\lambda$; Mellin form of $M_{Z1}$ & proved (Propositions~\CGref{prop:atomized-clock}, \CGref{prop:Mellin-Mertens})\\
A6 & Mean square of the cell masses uniformly in $\lambda$; Dirichlet series in $\Re s>1$ for every $\lambda$, including $\lambda=0$ & proved (Theorem~\CGref{thm:mean-square-cells}, Corollary~\CGref{cor:abscissa-one})\\
A7 & Cell identity \CGeqref{eq:closing-cells} for every $n\ge2$, and $J_{Z1}=0$ on the first cell $(1,2)$ & open\\
A8 & Equivalent scalar form: $\int_{0^+}^1\Upsilon\,d\lambda=\log L$, i.e.\ \CGeqref{eq:product-identity} & open; follows from A7\\
A9 & $\zeta_{Z1}=\zeta$ by Hamburger; $M(x)=O_\eps(x^{1/2+\eps})$ by Step~3 & proved as implication from A7 (Thms.~\CGref{thm:closing}, \CGref{thm:step3})\\
A10 & Continuation of $\tilde h_\lambda$ to $0<\Re s\le\frac12$ for $\lambda<1$ & open; not needed for A9\\
A11 & $\zeta_\lambda$ is not a Dirichlet series and $\tilde h_\lambda$ is not constant for $0<\lambda\le1$; convexity of $\log\zeta_\lambda(\sigma)$ in $\lambda$ & proved (Propositions~\CGref{prop:nonconstant}, \CGref{prop:convexity})\\
A12 & Hamburger's theorem for the pair: A9 follows if $e^{\tilde h_0}$ is of finite order and symmetric under $s\mapsto1-s$ & proved as implication (Theorem~\CGref{thm:pair-criterion})\\
A13 & Positivity: A9 follows if $\zeta_{Z1}$ has a positive representing measure; then $\varpi_{Z1}\le1$ & proved as implication (Theorem~\CGref{thm:positivity}, Corollary~\CGref{cor:positivity})\\
A14 & Clock end: density of $d\pi_{\RVM}$, cell masses $\frac1{\log n}+O(n^{-1/4})$, defects $O(|s|n^{-\sigma-4/3})$ & proved (Theorem~\CGref{thm:clock-density}, Corollary~\CGref{cor:clock-cells})\\
A15 & Height kernel of the cell masses; flow as count difference; Gaussian defect and the density condition $\tilde h_0(1)=0$ & proved (Propositions~\CGref{prop:height-kernels}, \CGref{prop:flow-count}, \CGref{prop:gaussian-defect})\\
\bottomrule
\end{tabular}
\end{center}

\begin{remark}[Input for A7]\label{rem:input-A7}
Steps A1--A6 use only two properties of an ordinate list: bounded displacement from the clock and a primitive of the count difference of size $O(\log\tau)$.  The list $(\tau_k)$ has both properties, and for it the corresponding function $\zeta_{\RVM}$ is not a Dirichlet series (Theorem~\CGref{thm:no-Dirichlet-RVM}).  Step A7 must therefore use a property of the ordinates $\gamma_k$ that the clock ordinates do not have; the natural candidate is the arithmetic of $\zeta$ carried by the $\gamma_k$, for instance Landau's formula $\sum_{0<\gamma\le T}x^\rho=-\frac T{2\pi}\Lambda(x)+O(\log T)$, in which the prime powers appear as resonances of the ordinates.
\end{remark}

\section{Criteria for the final step and the clock end of the family}
\label{sec:criteria}

This section collects the consequences of Corollary~\CGref{cor:abscissa-one} for the final step.  Along the open path $0<\lambda\le1$ no member $\zeta_\lambda$ is a Dirichlet series (Proposition~\CGref{prop:nonconstant}); the family is convex in $\lambda$ on the real axis (Proposition~\CGref{prop:convexity}); Hamburger's theorem applies to the pair $(\zeta_\lambda,\tilde\zeta_{\Z,\lambda})$ as soon as the joining factor is symmetric (Theorem~\CGref{thm:pair-criterion}); positivity of the measure representing $\zeta_{Z1}$ gives $\varpi_{Z1}\le1$ and closes the route (Theorem~\CGref{thm:positivity}).  At the clock end the density of $d\pi_{\RVM}$ is computed explicitly (Theorem~\CGref{thm:clock-density}), and the flow of the cell masses is written in the heights of the ordinates (Propositions~\CGref{prop:height-kernels} and~\CGref{prop:flow-count}).

Throughout this section the function $\chi$ of Subsection~\CGref{subsec:mean-square} is even, so that $\varphi$ is even, and
\begin{equation}\label{eq:cell-moments}
 \hat\varphi(\xi):=\int_{\R}\varphi(t)e^{-2\pi i\xi t}\,dt,\qquad \mathsf M_r:=\int_{\R}t^r\varphi(t)\,dt,\qquad \mathsf C_r:=\int_{\R}t^r\varphi(t)\cos(2\pi t)\,dt .
\end{equation}
$\hat\varphi$ is real and even, and $\mathsf M_1=\mathsf C_1=0$.  Since $\sum_{m\in\Z}\varphi(t-m)=1$,
\begin{equation}\label{eq:partition-fourier}
 \hat\varphi(q)=\int_0^1\sum_{m\in\Z}\varphi(t-m)\,e^{-2\pi iqt}\,dt=\delta_{q0}\qquad(q\in\Z),
\end{equation}
so $\mathsf M_0=1$ and $\mathsf C_0=\hat\varphi(1)=0$.

\subsection{The open path}
\label{subsec:open-path}

\begin{proposition}[No member of the open path is a Dirichlet series]\label{prop:nonconstant}
Let $0<\lambda\le1$.  Then $\zeta_\lambda$ is not a constant multiple of a Dirichlet series $\sum_na_nn^{-s}$ converging absolutely in $\Re s>1$.  In particular the joining factor $\tilde h_\lambda$ is not constant.
\end{proposition}

\begin{proof}
Suppose $\zeta_\lambda(s)=\sum_na_nn^{-s}$ absolutely for $\Re s>1$.  By Theorem~\CGref{thm:path}, exactly as in Proposition~\CGref{prop:Z1-class}, $(s-1)\zeta_\lambda(s)$ is entire of order one and $\zeta_\lambda$ satisfies \CGeqref{eq:Z1-FE}; these are the hypotheses of Hamburger's theorem, so $\zeta_\lambda=C\zeta$ with $C\ne0$.  The nontrivial zeros of $\zeta$ with positive ordinate are then $\frac12+i\gamma_k(\lambda)$, $k\ge1$, with multiplicity.  The ordinates of $\zeta$, listed in non-decreasing order, are the $\gamma_k$, and the $\gamma_k(\lambda)$ are increasing; hence $\gamma_k(\lambda)=\gamma_k$, that is $\lambda\eta_k=0$, for every $k$.  Since $\lambda>0$, every $\eta_k$ vanishes, all members of the family coincide (Remark~\CGref{rem:only-gamma-tau}), and $\zeta_{\RVM}=\zeta_\lambda$ is an absolutely convergent Dirichlet series, against Theorem~\CGref{thm:no-Dirichlet-RVM}.  If $\tilde h_\lambda$ were a constant $c$, then $\zeta_\lambda=e^{-c}\tilde\zeta_{\Z,\lambda}$ would be such a series by Corollary~\CGref{cor:abscissa-one}.
\end{proof}

The two columns of \CGeqref{eq:complement} exist for every $\lambda$, and for every $\lambda\in(0,1]$ they are different functions.  The joining factor converges, $\tilde h_\lambda\to\tilde h_0$ as $\lambda\to0^+$, through non-constant functions; (D) is the statement that the limit is constant.  A property of the displacements that yields (D) holds for the sequence $(\eta_k)$ and fails for every sequence $(c\,\eta_k)$ with $0\le c<1$, because the ordinates $\gamma_k(\lambda)=\tau_k-(1-\lambda)\eta_k$ are obtained from the clock by the displacements $(1-\lambda)\eta_k$.

\subsection{Convexity along the family}
\label{subsec:convexity}

\begin{proposition}[Convexity on the real axis]\label{prop:convexity}
For every real $\sigma>1$ the function $\lambda\mapsto\log\zeta_\lambda(\sigma)$ is strictly convex on $[0,1]$:
\begin{equation}\label{eq:convexity-second}
 \partial_\lambda^2\log\zeta_\lambda(\sigma)=\sum_{k\ge1}\eta_k^2\Bigl(\mathcal C\bigl(\sigma-\tfrac12,\gamma_k(\lambda)\bigr)-\mathcal C\bigl(\tfrac12,\gamma_k(\lambda)\bigr)\Bigr)>0,\qquad \mathcal C(a,\gamma):=\frac{2(a^2-\gamma^2)}{(a^2+\gamma^2)^2}.
\end{equation}
Consequently, with the generator $\mathcal V_\lambda$ of Theorem~\CGref{thm:generator},
\begin{equation}\label{eq:convexity-sandwich}
 \log\zeta_{\RVM}(\sigma)-\mathcal V_1(\sigma)<\log\zeta_{Z1}(\sigma)<\log\zeta_{\RVM}(\sigma)-\mathcal V_0(\sigma)\qquad(\sigma>1).
\end{equation}
\end{proposition}

\begin{proof}
With $a=\sigma-\frac12>\frac12$ one has $q(\sigma)=a^2-\frac14$, so by \CGeqref{eq:P-lambda} $\log\zeta_\lambda(\sigma)=-\log G(\sigma)+\sum_k\bigl(\log(a^2+\gamma_k(\lambda)^2)-\log(\frac14+\gamma_k(\lambda)^2)\bigr)$.  Since $\partial_\gamma^2\log(a^2+\gamma^2)=\mathcal C(a,\gamma)$ and $\partial_\lambda\gamma_k(\lambda)=\eta_k$, termwise differentiation gives the series in \CGeqref{eq:convexity-second}; its terms are $O(\eta_k^2\gamma_k^{-2})$, so it converges uniformly on $[0,1]$.  Fix $\gamma>\frac12$.  Then $\mathcal C(\frac12,\gamma)<0$; $\partial_a\mathcal C(a,\gamma)=4a(3\gamma^2-a^2)(a^2+\gamma^2)^{-3}>0$ for $0<a<\sqrt3\,\gamma$; and $\mathcal C(a,\gamma)>0$ for $a>\gamma$.  Hence $\mathcal C(a,\gamma)>\mathcal C(\frac12,\gamma)$ for every $a>\frac12$.  All ordinates satisfy $\gamma_k(\lambda)>\tau_0>\frac12$ and $\eta_1=\tau_1-\gamma_1\ne0$ ($\tau_1=14.52\ldots$ by \CGeqref{eq:lambert}, $\gamma_1=14.13\ldots$), so the sum is positive.  A strictly convex function $F$ on $[0,1]$ satisfies $F(1)-F'(1)<F(0)<F(1)-F'(0)$; with $F(\lambda)=\log\zeta_\lambda(\sigma)$ and $F'=\mathcal V_\lambda(\sigma)$ \CGeqref{eq:generator-zeta} this is \CGeqref{eq:convexity-sandwich}.
\end{proof}

As $\sigma\to+\infty$, \CGeqref{eq:convexity-second} becomes $-\Upsilon'(\lambda)>0$ and \CGeqref{eq:convexity-sandwich} becomes the last statement of Proposition~\CGref{prop:family-flow}.  The sign in \CGeqref{eq:convexity-second} comes from the test function $x^{-\sigma}$, whose transform is not oscillatory in the ordinate; for the cell functions $\varphi_n$ the corresponding second derivative $\partial_\lambda^2\tilde c_n(\lambda)=-2\Re\sum_k\eta_k^2W_n''(\gamma_k(\lambda))$ has terms of both signs.

\subsection{Hamburger's theorem for the pair}
\label{subsec:pair}

\begin{theorem}[Hamburger's theorem for the pair $(\zeta_\lambda,\tilde\zeta_{\Z,\lambda})$]\label{thm:pair-criterion}
Let $\lambda\in[0,1]$.  Suppose that the joining factor $e^{\tilde h_\lambda}$ extends to a meromorphic function on $\C$ of finite order with finitely many poles, and that
\begin{equation}\label{eq:pair-symmetry}
 e^{\tilde h_\lambda(s)}=e^{\tilde h_\lambda(1-s)} .
\end{equation}
Then $\tilde\zeta_{\Z,\lambda}=\zeta$.  For $\lambda=0$, moreover, $\zeta_{Z1}=\zeta$ and $\tilde h_0\equiv0$.
\end{theorem}

\begin{proof}
Put $f=\tilde\zeta_{\Z,\lambda}=\zeta_\lambda e^{\tilde h_\lambda}$.  By Theorem~\CGref{thm:path}, $(s-1)\zeta_\lambda$ is entire of order one, so $f=\mathsf G/\mathsf P$ with $\mathsf G$ entire of finite order and $\mathsf P$ a polynomial: this is (H0).  By \CGeqref{eq:Z1-FE} for $\zeta_\lambda$ and \CGeqref{eq:pair-symmetry}, $f(s)\Gamma(\frac s2)\pi^{-s/2}=f(1-s)\Gamma(\frac{1-s}2)\pi^{-(1-s)/2}$: this is (H2) with $g=f$.  (H1) is Corollary~\CGref{cor:abscissa-one}, and it gives (H3) for $g=f$.  Hamburger's theorem gives $f=C\zeta$, and $C=1$ because $\tilde a_1(\lambda)=1$.  By Theorem~\CGref{thm:atomized-continuation}(iii) and the functional equation, the nontrivial zeros of $f=\zeta$ lie on $\Re s=\frac12$; they are the points $\frac12\pm i\gamma_k$.  For $\lambda=0$, Theorem~\CGref{thm:Hamburger-Z1} then gives $\zeta_{Z1}=\zeta=\tilde\zeta_{\Z,0}$, and $e^{\tilde h_0}=1$.
\end{proof}

Theorem~\CGref{thm:pair-criterion} uses the functional equation of $\zeta_\lambda$ and the Dirichlet series of $\tilde\zeta_{\Z,\lambda}$ as they stand; the only hypothesis is on the factor that joins them, and it asks for symmetry instead of constancy.  At $\lambda=1$ the two hypotheses do not hold together: the conclusion $\tilde\zeta_{\Z,1}=\zeta$ would place the zeros of $\zeta$ at $\frac12\pm i\tau_k$ (Theorem~\CGref{thm:atomized-continuation}(iv)), so $\zeta=\zeta_{\RVM}$ by Theorem~\CGref{thm:Hadamard}, against Theorem~\CGref{thm:no-Dirichlet-RVM}.

\subsection{Positivity and the value at infinity}
\label{subsec:positivity}

For an even Schwartz function $\mathsf g$ on $\R$ put $\hat{\mathsf g}(\xi)=\int_{\R}\mathsf g(u)e^{-2\pi i\xi u}du$ and $\mathsf g^{\sharp}(s)=\int_0^\infty\mathsf g(u)u^{s-1}du$.

\begin{lemma}[Summation formula for a representing measure]\label{lem:summation-positive}
Let $\lambda\in[0,1]$ and let $\mathsf A$ be a positive measure on $[1,\infty)$ with
\begin{equation}\label{eq:representing-measure}
 \zeta_\lambda(s)=\int_{[1,\infty)}u^{-s}\,d\mathsf A(u)\qquad(\Re s>1).
\end{equation}
Then for every even Schwartz function $\mathsf g$
\begin{equation}\label{eq:summation-positive}
 \mathsf g(0)+2\int_{[1,\infty)}\mathsf g\,d\mathsf A=\hat{\mathsf g}(0)+2\int_{[1,\infty)}\hat{\mathsf g}\,d\mathsf A .
\end{equation}
\end{lemma}

\begin{proof}
$\mathsf g^{\sharp}$ is meromorphic with simple poles at $s=0,-2,-4,\dots$, residue $\mathsf g(0)$ at $s=0$, and decays faster than any power of $|t|$ on vertical strips; $\mathsf g^{\sharp}(1)=\frac12\hat{\mathsf g}(0)$; and
\begin{equation}\label{eq:tate-local}
 \pi^{-s/2}\Gamma(\tfrac s2)\,\hat{\mathsf g}^{\sharp}(1-s)=\pi^{-(1-s)/2}\Gamma(\tfrac{1-s}2)\,\mathsf g^{\sharp}(s)
\end{equation}
for $0<\Re s<1$, and then by meromorphic continuation: the Fourier transform of the even homogeneous distribution $|\xi|^{-s}$ is a multiple of $|u|^{s-1}$, and the multiple is fixed by $\mathsf g(u)=e^{-\pi u^2}$, for which $\hat{\mathsf g}=\mathsf g$ and $\mathsf g^{\sharp}(s)=\frac12\pi^{-s/2}\Gamma(\frac s2)$.  Fix $1<c<2$.  By Mellin inversion and Fubini, $\int\mathsf g\,d\mathsf A=\frac1{2\pi i}\int_{(c)}\zeta_\lambda(s)\mathsf g^{\sharp}(s)\,ds$, since $\int u^{-c}d\mathsf A<\infty$.  On $\Re s\ge c$, $|\zeta_\lambda(s)|\le\zeta_\lambda(c)$; with the functional equation \CGeqref{eq:Z1-FE}, valid for every $\zeta_\lambda$ by Theorem~\CGref{thm:path}, and the finite order of $(s-1)\zeta_\lambda$, the Phragm\'en--Lindel\"of principle bounds $\zeta_\lambda$ by a power of $|t|$ on $1-c\le\Re s\le c$, $|t|\ge1$.  Moving the line to $\Re s=1-c$ crosses $s=1$, where $\zeta_\lambda$ has residue $1$, and $s=0$, where $\zeta_\lambda(0)=P_\lambda(0)/G(0)=-\frac12$:
\begin{equation}\label{eq:summation-shift}
 \int\mathsf g\,d\mathsf A=\tfrac12\hat{\mathsf g}(0)-\tfrac12\mathsf g(0)+\frac1{2\pi i}\int_{(1-c)}\zeta_\lambda(s)\mathsf g^{\sharp}(s)\,ds .
\end{equation}
By \CGeqref{eq:Z1-FE} and \CGeqref{eq:tate-local}, $\zeta_\lambda(s)\mathsf g^{\sharp}(s)=\zeta_\lambda(1-s)\hat{\mathsf g}^{\sharp}(1-s)$, and the last integral is $\frac1{2\pi i}\int_{(c)}\zeta_\lambda(s)\hat{\mathsf g}^{\sharp}(s)\,ds=\int\hat{\mathsf g}\,d\mathsf A$.
\end{proof}

\begin{theorem}[Positivity criterion]\label{thm:positivity}
Let $\lambda\in[0,1]$.  If $\zeta_\lambda$ has a representation \CGeqref{eq:representing-measure} with a positive measure $\mathsf A$, then
\begin{equation}\label{eq:positivity-bound}
 \varpi_\lambda=\mathsf A(\{1\})\le1 .
\end{equation}
In particular, if $\zeta_{Z1}$ has such a representation, then $\varpi_{Z1}=1$ and $\zeta_{Z1}=\zeta$.
\end{theorem}

\begin{proof}
By monotone convergence $\zeta_\lambda(\sigma)=\int u^{-\sigma}d\mathsf A\to\mathsf A(\{1\})$ as $\sigma\to+\infty$, and the limit is $\varpi_\lambda$ (Proposition~\CGref{prop:family-clock}).  Let $\mathsf T(\xi)=(1-|\xi|)_+$ and $\mathsf B=\mathsf T*\mathsf T$, so that $\mathsf B(\xi)=\frac23-\xi^2+\frac12|\xi|^3$ for $|\xi|\le1$, $\mathsf B(\xi)=\frac16(2-|\xi|)^3$ for $1\le|\xi|\le2$, $\mathsf B=0$ elsewhere, and $\mathsf B''(\xi)=\mathsf T(\xi+1)-2\mathsf T(\xi)+\mathsf T(\xi-1)$.  Put
\begin{equation}\label{eq:extremal-function}
 \mathsf g(u):=\Bigl(\frac{\sin\pi u}{\pi u}\Bigr)^4(1-u^2),\qquad \hat{\mathsf g}(\xi)=\mathsf B(\xi)+\frac1{4\pi^2}\mathsf B''(\xi) .
\end{equation}
Then $\mathsf g(0)=1$; $\mathsf g(u)\le0$ for $|u|\ge1$; $\hat{\mathsf g}(0)=\frac23-\frac1{2\pi^2}$ and $\hat{\mathsf g}(1)=\frac16+\frac1{4\pi^2}$; and $\hat{\mathsf g}\ge0$ on $\R$: for $0\le\xi\le1$, $\hat{\mathsf g}(\xi)=\frac23-\xi^2+\frac12\xi^3+\frac{3\xi-2}{4\pi^2}$ is positive at $\xi=0$ and $\xi=1$ and its derivative $-2\xi+\frac32\xi^2+\frac3{4\pi^2}$ changes sign once, from positive to negative, and for $1\le\xi\le2$, $\hat{\mathsf g}(\xi)=\frac16(2-\xi)^3+\frac{2-\xi}{4\pi^2}\ge0$.  For $\epsilon>0$ let $\mathsf g_\epsilon(u)=\mathsf g(u)e^{-\epsilon\pi u^2}$, an even Schwartz function with $\mathsf g_\epsilon(0)=1$, $\mathsf g_\epsilon\le0$ on $|u|\ge1$, and $\hat{\mathsf g}_\epsilon=\hat{\mathsf g}*(\epsilon^{-1/2}e^{-\pi\xi^2/\epsilon})\ge0$.  By Lemma~\CGref{lem:summation-positive}, since $\mathsf A\ge0$ is carried by $[1,\infty)$ and has mass $\varpi_\lambda$ at $u=1$,
\begin{equation}\label{eq:lp-inequality}
 1\ \ge\ \mathsf g_\epsilon(0)+2\int\mathsf g_\epsilon\,d\mathsf A=\hat{\mathsf g}_\epsilon(0)+2\int\hat{\mathsf g}_\epsilon\,d\mathsf A\ \ge\ \hat{\mathsf g}_\epsilon(0)+2\varpi_\lambda\,\hat{\mathsf g}_\epsilon(1) .
\end{equation}
Letting $\epsilon\to0$, $\hat{\mathsf g}_\epsilon\to\hat{\mathsf g}$ pointwise because $\hat{\mathsf g}$ is continuous, and
\begin{equation}\label{eq:lp-value}
 \varpi_\lambda\le\frac{1-\hat{\mathsf g}(0)}{2\hat{\mathsf g}(1)}=\frac{\frac13+\frac1{2\pi^2}}{\frac13+\frac1{2\pi^2}}=1 .
\end{equation}
The last statement follows from Theorem~\CGref{thm:varpi-Z1}.
\end{proof}

\begin{corollary}[Positive form of the final step]\label{cor:positivity}
The following are equivalent: \textup{(i)} (D); \textup{(ii)} $\zeta_{Z1}$ is the Mellin--Stieltjes transform \CGeqref{eq:representing-measure} of a positive measure on $[1,\infty)$; \textup{(iii)} $J^{\rm M}_{Z1}$, extended by $0$ to $(0,1]$, is non-decreasing on $(0,\infty)$.
\end{corollary}

\begin{proof}
(i)$\Rightarrow$(iii): then $J^{\rm M}_{Z1}=J$ (Theorem~\CGref{thm:cell-form}).  (iii)$\Rightarrow$(ii): $J^{\rm M}_{Z1}(x)x^{-c}\in L^2((1,\infty),dx/x)$ for every $c>1$, so a non-decreasing $J^{\rm M}_{Z1}$ is $O_c(x^{c})$ and $dJ^{\rm M}_{Z1}$ is a positive measure on $[1,\infty)$ with $\int x^{-\sigma}dJ^{\rm M}_{Z1}<\infty$ for $\sigma>1$.  By Theorem~\CGref{thm:mellin-counts}, $\zeta_{Z1}(s)=\exp\int x^{-s}dJ^{\rm M}_{Z1}=\int u^{-s}d\mathsf A(u)$, where $\mathsf A=\sum_{r\ge0}\frac1{r!}(dJ^{\rm M}_{Z1})^{*r}$ (multiplicative convolution powers) is a positive measure on $[1,\infty)$.  (ii)$\Rightarrow$(i) is Theorem~\CGref{thm:positivity}.
\end{proof}

In Theorem~\CGref{thm:positivity} the Dirichlet series is not used: (H1) is replaced by the positivity of the representing measure, the functional equation gives $\varpi_{Z1}\le1$ through \CGeqref{eq:lp-inequality}, and the Hadamard product gives $\varpi_{Z1}\ge1$ (Theorem~\CGref{thm:varpi-Z1}).  The function \CGeqref{eq:extremal-function} is extremal for the integers: for $\mathsf A=\sum_{n\ge1}\delta_n$ both inequalities in \CGeqref{eq:lp-inequality} are equalities in the limit, since $\mathsf g(n)=0$ for $n\ge1$ and $\hat{\mathsf g}(n)=0$ for $n\ge2$.  At the clock end $\varpi_{\RVM}<1$, and the density of $dJ_{\RVM}$ takes negative values in every cell $(n,n+1)$ with $n$ large (Theorem~\CGref{thm:clock-density}).

\subsection{The density of the clock}
\label{subsec:clock-density}

Let $\Psi(t):=t\log t-t$ for $t>0$, and for $M\ge1$, $0<t_0\le\frac12$ let
\begin{equation}\label{eq:U-def}
 \mathcal U(M,t_0):=\int_{t_0}^{\infty}e^{-iM\Psi(t)}\,dt:=\lim_{T\to\infty}\int_{t_0}^{T}e^{-iM\Psi(t)}\,dt .
\end{equation}

\begin{lemma}[Stationary phase for $\Psi$]\label{lem:stationary-U}
The limit \CGeqref{eq:U-def} exists, and there is an absolute constant $C_{\mathcal U}$ with
\begin{equation}\label{eq:U-asymptotic}
 \Bigl|\mathcal U(M,t_0)-\sqrt{\frac{2\pi}M}\,e^{i(M-\pi/4)}\Bigr|\le\frac{C_{\mathcal U}}M\qquad(M\ge1,\ 0<t_0\le\tfrac12).
\end{equation}
\end{lemma}

\begin{proof}
$\Psi'(t)=\log t$ is increasing and unbounded, so the limit exists by the first-derivative test.  $\Psi$ is convex with minimum $\Psi(1)=-1$.  Put $\mathsf w:=\operatorname{sgn}(t-1)\sqrt{2(\Psi(t)+1)}$.  Near $t=1$, $\Psi(t)+1=\frac12(t-1)^2(1+O(t-1))$ with an analytic factor, so $t\mapsto\mathsf w$ is analytic with derivative $1$ at $t=1$; elsewhere $\mathsf w\,d\mathsf w=\log t\,dt$.  Hence $t\mapsto\mathsf w$ is an increasing analytic bijection of $(0,\infty)$ onto $(-\sqrt2,\infty)$, with
\begin{equation}\label{eq:U-substitution}
 \mathsf p(\mathsf w):=\frac{dt}{d\mathsf w}=\frac{\mathsf w}{\log t},\qquad \mathsf p(0)=1,\qquad \mathsf q(\mathsf w):=\frac{\mathsf p(\mathsf w)-1}{\mathsf w}=\frac1{\log t}-\frac1{\mathsf w}.
\end{equation}
$\mathsf q$ is analytic on $(-\sqrt2,\infty)$, bounded, tends to $0$ as $\mathsf w\to\infty$, and $\mathsf q'=-\mathsf w/(t\log^3t)+\mathsf w^{-2}$.  Away from $\mathsf w=0$ one has $\int|\mathsf w|\,|t\log^3t|^{-1}d\mathsf w=\int dt/(t\log^2t)$, which converges at $t=0$ and at $t=\infty$, and $\int\mathsf w^{-2}d\mathsf w$ converges at $\infty$; so $\int_{-\sqrt2}^\infty|\mathsf q'|\,d\mathsf w<\infty$.  With $\mathsf w_0=\mathsf w(t_0)\le\mathsf w(\frac12)=-\sqrt{1-\log2}<-\frac12$ and $e^{-iM\Psi}=e^{iM}e^{-iM\mathsf w^2/2}$,
\begin{equation}\label{eq:U-split}
 e^{-iM}\,\mathcal U(M,t_0)=\int_{\mathsf w_0}^\infty e^{-iM\mathsf w^2/2}\,d\mathsf w+\int_{\mathsf w_0}^\infty e^{-iM\mathsf w^2/2}\,\mathsf w\,\mathsf q(\mathsf w)\,d\mathsf w .
\end{equation}
The first integral is $\int_{\R}e^{-iM\mathsf w^2/2}d\mathsf w=\sqrt{2\pi/M}\,e^{-i\pi/4}$ minus $\int_{-\infty}^{\mathsf w_0}$, and the latter is at most $4/(M|\mathsf w_0|)\le8/M$ by the first-derivative test \cite[Lemma~4.2]{Titchmarsh1986}.  Since $\mathsf w\,e^{-iM\mathsf w^2/2}=\frac iM\frac d{d\mathsf w}e^{-iM\mathsf w^2/2}$, integration by parts bounds the second integral by $M^{-1}\bigl(\sup|\mathsf q|+\int|\mathsf q'|\bigr)$.
\end{proof}

For $m\ge1$ and $y>1$ put
\begin{equation}\label{eq:clock-harmonic}
 \mathcal R_m(y):=-2(-1)^m\,m^{-3/2}\,y^{\frac1{2m}-\frac12}\cos\Bigl(2\pi m\,y^{1/m}-\frac{(7m+1)\pi}4\Bigr),
\end{equation}
so that $\mathcal R_1(y)=2\cos(2\pi y)$ and $\mathcal R_2(y)=-2^{-1/2}y^{-1/4}\cos(4\pi\sqrt y+\frac\pi4)$.

\begin{theorem}[The density of the clock]\label{thm:clock-density}
$J_{\RVM}$ is continuously differentiable on $(2,\infty)$, and for $y\ge16$
\begin{equation}\label{eq:clock-density}
 J_{\RVM}'(y)=\frac1{\log y}+2\cos(2\pi y)-\frac{y^{-1/4}}{\sqrt2}\cos\Bigl(4\pi\sqrt y+\frac\pi4\Bigr)+O\bigl(y^{-1/3}\bigr).
\end{equation}
On $(2,\infty)$ the measure $d\pi_{\RVM}$ is the sum of the atoms $\frac{\mu(j)}jJ_{\RVM}(2)\,\delta_{2^j}$, $j\ge1$, and of a density that is continuous off the points $2^j$ and equals $J_{\RVM}'(y)+O(y^{-1/2})$.
\end{theorem}

\begin{proof}
\emph{An exact series.}  We use the notation of the proof of Theorem~\CGref{thm:unit-cell}: $v=\log y$, $S_\eps(v)=\sum_ke^{(iv-\eps)\tau_k}$, the decomposition \CGeqref{eq:S-eps-split} and the harmonics $K_m^\pm$ of \CGeqref{eq:K-harmonics}, with phases $\Phi_m^\pm(\tau)=v\tau\pm2\pi m\NR(\tau)$.  Put
\begin{equation}\label{eq:H-harmonic}
 \mathcal H_m^\pm(v):=\int_{\tau_0}^\infty e^{i\Phi_m^\pm(\tau)}\,d\tau ,
\end{equation}
improper integrals that converge, locally uniformly in $v\ge\log2$, by the first-derivative test, because $(\Phi_m^\pm)'$ is monotone and unbounded.  As $\eps\downarrow0$, $K_m^\pm\to iv\,\mathcal H_m^\pm(v)$, and the first two terms of \CGeqref{eq:S-eps-split} tend to a continuous function $S_{\rm sm}(v)$ with $|S_{\rm sm}(v)|\le C/v+\frac12$.  The bounds for $K_m^\pm$ in that proof are uniform in $\eps$ and summable against the weights $1/(2\pi m)$, so by dominated convergence in \CGeqref{eq:damped-sum} and \CGeqref{eq:K-harmonics}
\begin{equation}\label{eq:S-zero}
 \sum_kD_k=\int_{u_1}^{u_2}\frac{e^{v/2}}v\,S_0(v)\,dv,\qquad S_0(v):=S_{\rm sm}(v)+\sum_{m\ge1}\frac{(-1)^m\,v}{2\pi m}\bigl(\mathcal H_m^-(v)-\mathcal H_m^+(v)\bigr),
\end{equation}
the series converging uniformly on compact sets.  Thus $S_0$ is continuous, $J_{\RVM}$ is continuously differentiable on $(2,\infty)$, and by \CGeqref{eq:J-RVM}
\begin{equation}\label{eq:J-prime-exact}
 J_{\RVM}'(y)=\frac1{\log y}\Bigl(1-\frac1{y(y^2-1)}-\frac2{\sqrt y}\,\Re S_0(\log y)\Bigr).
\end{equation}

\emph{The terms without a stationary point.}  $(\Phi_m^+)'\ge v+m\log2$, so $|\mathcal H_m^+|\le4/(v+m\log2)$ and $\sum_m\frac v{2\pi m}|\mathcal H_m^+|\ll1+\log v$.  For $m>v/\log2$, $(\Phi_m^-)'\le v-m\log2<0$ on $[\tau_0,\infty)$; the first-derivative test, and the second-derivative test on $[\tau_0,2\tau_0]$ with $|(\Phi_m^-)''|=m/\tau$, give $|\mathcal H_m^-|\ll\min\bigl((m\log2-v)^{-1},m^{-1/2}\bigr)$, and $\sum_{m>v/\log2}\frac v{2\pi m}|\mathcal H_m^-|\ll1+\log v$.

\emph{The terms $3\le m\le v/\log2$.}  By the argument for \CGeqref{eq:K-minus-bound}, $|\mathcal H_m^-|\le\frac8{m\log2}+16\sqrt\pi\,e^{v/(2m)}m^{-1/2}$, so these terms contribute to $S_0$ at most $C\,v\sum_{m\ge3}\bigl(m^{-2}+m^{-3/2}e^{v/(2m)}\bigr)\ll v\,e^{v/6}$.

\emph{The terms $m=1,2$.}  The stationary point of $\Phi_m^-$ is $\mathsf t_m:=2\pi e^{v/m}=2\pi y^{1/m}$.  With $\tau=\mathsf t_mt$ and \CGeqref{eq:clock}, $\Phi_m^-(\mathsf t_mt)=-m\mathsf t_m\Psi(t)-\frac{7\pi m}4$, hence
\begin{equation}\label{eq:H-by-U}
 \mathcal H_m^-(v)=\mathsf t_m\,e^{-7\pi im/4}\,\mathcal U\bigl(m\mathsf t_m,\tau_0/\mathsf t_m\bigr),\qquad \tau_0/\mathsf t_m=2e^{-v/m}\le\tfrac12\quad(y\ge4^m).
\end{equation}
By Lemma~\CGref{lem:stationary-U} with $M=m\mathsf t_m$,
\begin{equation}\label{eq:H-asymptotic}
 \frac{(-1)^mv}{2\pi m}\,\mathcal H_m^-(v)=(-1)^m\,\frac v{m^{3/2}}\,y^{\frac1{2m}}\exp\Bigl(i\Bigl(2\pi m\,y^{1/m}-\frac{(7m+1)\pi}4\Bigr)\Bigr)+O\Bigl(\frac v{m^2}\Bigr).
\end{equation}
Inserting the four estimates into \CGeqref{eq:J-prime-exact} gives $J_{\RVM}'(y)=\frac1{\log y}+\mathcal R_1(y)+\mathcal R_2(y)+O(y^{-1/3})$, which is \CGeqref{eq:clock-density}.

\emph{The prime coordinate.}  By \CGeqref{eq:pi-RVM}, on each interval $(2^r,2^{r+1})$ the function $\pi_{\RVM}$ is the finite sum of the continuously differentiable functions $\frac{\mu(j)}jJ_{\RVM}(x^{1/j})$, $j\le r$; the terms $j\ge2$ have derivative $\frac{\mu(j)}{j^2}x^{1/j-1}J_{\RVM}'(x^{1/j})$, and $|J_{\RVM}'|\le C_J$ by \CGeqref{eq:unit-cell}, so their sum is $O(x^{-1/2})$.  At $x=2^j$ the term $j$ jumps by $\frac{\mu(j)}jJ_{\RVM}(2)$.
\end{proof}

\begin{corollary}[Cell masses and cell defects at the clock end]\label{cor:clock-cells}
For $n\ge17$, with $\mathbf 1_{2^j}(n)=1$ if $n=2^j$ and $0$ otherwise,
\begin{equation}\label{eq:clock-cell-sharp}
 c_n^{\RVM}=\int_{n-1}^n\Bigl(\frac1{\log x}+2\cos(2\pi x)\Bigr)dx+\mathbf 1_{2^j}(n)\,\frac{\mu(j)}j\,J_{\RVM}(2)+O\bigl(n^{-1/4}\bigr),
\end{equation}
\begin{equation}\label{eq:clock-cell-smooth}
 \tilde c_n(1)=\frac1{\log n}+\mathbf 1_{2^j}(n)\,\frac{\mu(j)}j\,J_{\RVM}(2)+O\bigl(n^{-1/4}\bigr),
\end{equation}
and for $0<\Re s\le\sigma_1$
\begin{equation}\label{eq:clock-defect}
 |\mathfrak h_n(1;s)|\le C_{\sigma_1}\bigl(|s|\,n^{-\sigma-4/3}+|s|(1+|s|)\,n^{-\sigma-2}\bigr).
\end{equation}
In particular Hypothesis~\CGref{hyp:cell-remainder} holds along the integers $n$ that are not of the form $2^j$ with $\mu(j)\ne0$, and $c_n^{\RVM}\sim1/\log n$, $0<c_n^{\RVM}<1$ for all large such $n$; at $n=2^j$, $c_n^{\RVM}=\bigl(1+\mu(j)J_{\RVM}(2)\log2\bigr)/\log n+O(n^{-1/4})$.
\end{corollary}

\begin{proof}
\CGeqref{eq:clock-cell-sharp}: $c_n^{\RVM}=\pi_{\RVM}(n)-\pi_{\RVM}(n-1)$ is the integral of the density over $(n-1,n)$ plus the atom at $n$, and $\int_{n-1}^n(|\mathcal R_2|+O(x^{-1/3}))dx=O(n^{-1/4})$.  \CGeqref{eq:clock-cell-smooth}: $\varphi_n(2^j)=\mathbf 1_{2^j}(n)$; $\int\varphi(t)\,dt/\log(n+t)=\frac1{\log n}+O(\frac1{n\log^2n})$ because $\mathsf M_0=1$; and $\int\varphi(t)\,2\cos(2\pi(n+t))\,dt=2\mathsf C_0=0$ by \CGeqref{eq:partition-fourier}.  \CGeqref{eq:clock-defect}: the test function $\varphi_n(x)(n^{-s}-x^{-s})$ vanishes at the atoms, and $n^{-s}-(n+t)^{-s}=s\,n^{-s-1}t+O_{\sigma_1}(|s|(1+|s|)n^{-\sigma-2})$ for $|t|\le1$.  The first moment of the density is $\int t\varphi(t)\pi_{\RVM}'(n+t)\,dt=\frac{\mathsf M_1}{\log n}+2\mathsf C_1+O(n^{-1/3})=O(n^{-1/3})$, because $\mathsf M_1=\mathsf C_1=0$, $\mathcal R_2(n+t)=\mathcal R_2(n)+O(n^{-3/4})$ and the density is bounded.  The last statements follow from \CGeqref{eq:clock-cell-sharp} and Lemma~\CGref{lem:cell-cancel}.
\end{proof}

By \CGeqref{eq:clock-density}, $\frac{d\pi_{\RVM}}{dy}=\bigl(\frac1{\log y}-1\bigr)+\bigl(1+2\cos(2\pi y)\bigr)+O(y^{-1/4})$: the clock carries the first harmonic of the comb $\sum_n\delta_n=1+2\sum_{q\ge1}\cos(2\pi qy)$ of all integers, with weight $1$, on a smooth background.  This harmonic comes from the ordinates $\tau_k$ near the height $2\pi y$, the stationary point of $\Phi_1^-$; by \CGeqref{eq:partition-fourier} it does not enter the cell masses \CGeqref{eq:clock-cell-smooth}, which are $1/\log n$ to leading order, and it enters the cell defects only through the moments $\mathsf C_r$, $r\ge2$.

\subsection{The cell masses in the heights of the ordinates}
\label{subsec:height-kernels}

\begin{proposition}[Height kernel of the cell masses]\label{prop:height-kernels}
For every $\eps>0$, $\lambda\in[0,1]$ and $n\ge3$ that is not a power of $2$,
\begin{equation}\label{eq:height-kernel}
 \tilde c_n(\lambda)=\mathfrak m_n-\frac2{\sqrt n\,\log n}\sum_{k\ge1}\hat\varphi\Bigl(\frac{\gamma_k(\lambda)}{2\pi n}\Bigr)\cos\bigl(\gamma_k(\lambda)\log n\bigr)+E_n(\lambda),
\end{equation}
where $\sum_{X<n\le2X}|E_n(\lambda)|^2\le C_\eps X^{1/2+\eps}$ uniformly in $\lambda$ (sum over the $n$ that are not powers of $2$).  The kernel $\hat\varphi(\gamma/2\pi n)$ vanishes at the heights $\gamma=2\pi qn$, $q=1,2,\dots$
\end{proposition}

\begin{proof}
The term $j=1$ of $W_n$ in \CGeqref{eq:VW} is $\int\varphi(t)(n+t)^{-1/2+i\gamma}\,dt/\log(n+t)$.  For $0<\gamma\le n^{1+\eps}$, $(1+t/n)^{-1/2+i\gamma}=e^{i\gamma t/n}\bigl(1+O((1+\gamma/n)/n)\bigr)$ and $\frac1{\log(n+t)}=\frac1{\log n}(1+O(\frac1n))$ on $|t|\le1$, so this term equals $\frac{n^{-1/2+i\gamma}}{\log n}\hat\varphi(\frac\gamma{2\pi n})+O\bigl(n^{-3/2}(1+\gamma/n)\bigr)$, because $\int\varphi(t)e^{i\gamma t/n}dt=\hat\varphi(\gamma/2\pi n)$ for even $\varphi$.  Summed over the $O(n^{1+\eps}\log n)$ ordinates $\gamma_k(\lambda)\le n^{1+\eps}$ the error is $O(n^{-1/2+3\eps})$.  For $\gamma>n^{1+\eps}$ both the term and $\hat\varphi(\gamma/2\pi n)$ are $O_A((n/\gamma)^A)$, as in the proof of Theorem~\CGref{thm:mean-square-cells}.  The terms $j\ge2$ have sum of squares $O(X^{1/2+3\eps})$ over $X<n\le2X$ by \CGeqref{eq:cell-block}.  The zeros of the kernel are \CGeqref{eq:partition-fourier}.
\end{proof}

Since $\gamma_k(\lambda)\log n=\tau_k\log n-(1-\lambda)\eta_k\log n$, the passage along the family rotates each clock phase $n^{i\tau_k}$ by the angle $-(1-\lambda)\eta_k\log n$: up to the change of $\hat\varphi$ over a distance $O(1/n)$, the sum in \CGeqref{eq:height-kernel} is $\Re\,\Xi_n\bigl((1-\lambda)\log n\bigr)$ with
\begin{equation}\label{eq:Xi-rotation}
 \Xi_n(\Omega):=\sum_{k\ge1}\hat\varphi\Bigl(\frac{\tau_k}{2\pi n}\Bigr)n^{i\tau_k}e^{-i\Omega\eta_k}\qquad(0\le\Omega\le\log n).
\end{equation}
At $\Omega=0$ this is the clock sum, whose real part is $O(n^{1/4}\log n)$ by \CGeqref{eq:clock-cell-smooth} and Theorem~\CGref{thm:clock-density}, because the resonant heights $\tau_k\approx2\pi n$ are zeros of the kernel; the endpoint of the family is $\Omega=\log n$.

\begin{proposition}[The flow of a cell mass and the count difference]\label{prop:flow-count}
For every $n\ge3$ that is not a power of $2$,
\begin{equation}\label{eq:flow-count}
 \tilde c_n(0)-\tilde c_n(1)=-2\int_{0^+}^1\sum_{k\ge1}\eta_k\,\Im V_n\bigl(\gamma_k(\lambda)\bigr)\,d\lambda=-2\,\Im\int_{\tau_0}^\infty V_n(t)\bigl(N_\zeta(t)-n_{\RVM}(t)\bigr)\,dt .
\end{equation}
\end{proposition}

\begin{proof}
The first equality is \CGeqref{eq:cell-mass-W} and the continuity of $\tilde c_n$ at $\lambda=0$.  By \CGeqref{eq:cell-mass-W}, $\tilde c_n(0)-\tilde c_n(1)=2\Re\sum_k\bigl(W_n(\tau_k)-W_n(\gamma_k)\bigr)=2\Re\int_{\tau_0}^\infty W_n\,d(n_{\RVM}-N_\zeta)$.  Both counts vanish at $\tau_0$, their difference is $O(\log t)$ by Lemma~\CGref{lem:bounded-displacement}, and $W_n$ is rapidly decreasing; integration by parts with $W_n'=iV_n$ gives $2\Re\bigl(i\int V_n(N_\zeta-n_{\RVM})\,dt\bigr)$.
\end{proof}

\subsection{The Gaussian defect of the atomization}
\label{subsec:gaussian-defect}

\begin{lemma}[Vertical growth of the joining factor]\label{lem:vertical-growth}
For every $\lambda\in[0,1]$ and $\frac12<\sigma_0<\sigma_1$, $\tilde h_\lambda(\sigma+it)=o(|t|)$ as $|t|\to\infty$, uniformly for $\sigma_0\le\sigma\le\sigma_1$.
\end{lemma}

\begin{proof}
By \CGeqref{eq:joining-factor} and Proposition~\CGref{prop:first-cell}(ii) it suffices to treat the double series.  Fix $c\in(\frac12,\sigma_0)$.  By \CGeqref{eq:cell-defect-bound}, $\sum_{m\ge1}\frac1m\sum_{n>N}|\mathfrak h_n(\lambda;ms)|\le C|s|\sum_{n>N}n^{-\sigma_0-1}(1+I_n^{1/2})$ and $\sum_{m>N}\frac1m\sum_{n\ge2}|\mathfrak h_n(\lambda;ms)|\le C|s|\sum_{n\ge2}n^{-N\sigma_0-1}(1+I_n^{1/2})$; both are at most $\epsilon_N|s|$ with $\epsilon_N\to0$ as $N\to\infty$.  For the finitely many pairs $m,n\le N$, with $E=\pi_\lambda-\mathfrak L$ as in the proof of Theorem~\CGref{thm:atomized-continuation},
\begin{equation}\label{eq:cell-RL}
\begin{split}
 \mathfrak h_n(\lambda;ms)={}&\int\varphi_n(x)\bigl(n^{-ms}-x^{-ms}\bigr)d\mathfrak L(x)-\int E(x)\varphi_n'(x)\bigl(n^{-ms}-x^{-ms}\bigr)dx\\
 &-ms\int E(x)\varphi_n(x)\,x^{-ms-1}dx .
\end{split}
\end{equation}
The first two terms are bounded in $t$, and the last integral tends to $0$ as $|t|\to\infty$ by the Riemann--Lebesgue lemma, because $E\in L^1(n-1,n+1)$.  Hence $\limsup_{|t|\to\infty}|t|^{-1}\sup_\sigma|\tilde h_\lambda(\sigma+it)|\le2\epsilon_N$ for every $N$.
\end{proof}

For $y>0$ and $c>1$ put
\begin{equation}\label{eq:gaussian-defect}
 \Theta_\lambda(y):=\frac1{4\pi i}\int_{(c)}\zeta_\lambda(s)\,\Gamma(\tfrac s2)\,\pi^{-s/2}\,y^{-s/2}\,ds,\qquad
 \mathfrak G_\lambda(y):=\sum_{N\ge1}\tilde a_N(\lambda)\,e^{-\pi N^2y}-\Theta_\lambda(y).
\end{equation}
For $\zeta$ itself $\Theta(y)=\sum_{n\ge1}e^{-\pi n^2y}$, and if $\zeta_\lambda$ has a representation \CGeqref{eq:representing-measure}, then $\Theta_\lambda(y)=\int e^{-\pi u^2y}\,d\mathsf A(u)$.  $\mathfrak G_\lambda$ is the change produced by the atomization in the Gaussian count of the integers of the system.

\begin{proposition}[Gaussian defect]\label{prop:gaussian-defect}
Let $\lambda\in[0,1]$.
\begin{enumerate}[label=\textup{(\roman*)}]
\item $\Theta_\lambda(y)=y^{-1/2}\,\Theta_\lambda(1/y)+\frac12\bigl(y^{-1/2}-1\bigr)$ for $y>0$, and $\Theta_\lambda(y)=O_A(y^{-A})$ as $y\to\infty$ for every $A$.
\item The coefficients $\tilde a_N(\lambda)$ satisfy the theta relation of $\zeta$,
\begin{equation}\label{eq:theta-atomized}
 \sum_{N\ge1}\tilde a_N(\lambda)\,e^{-\pi N^2y}=y^{-1/2}\sum_{N\ge1}\tilde a_N(\lambda)\,e^{-\pi N^2/y}+\tfrac12\bigl(y^{-1/2}-1\bigr)\qquad(y>0),
\end{equation}
if and only if
\begin{equation}\label{eq:defect-selfdual}
 \mathfrak G_\lambda(y)=y^{-1/2}\,\mathfrak G_\lambda(1/y)\qquad(y>0).
\end{equation}
Under the hypotheses of Theorem~\CGref{thm:pair-criterion}, \CGeqref{eq:defect-selfdual} holds.
\item For every $\delta>0$, as $y\to0^+$,
\begin{equation}\label{eq:defect-expansion}
 \mathfrak G_\lambda(y)=\tfrac12\bigl(e^{\tilde h_\lambda(1)}-1\bigr)\,y^{-1/2}+O_\delta\bigl(y^{-1/4-\delta}\bigr),
\end{equation}
and $e^{\tilde h_\lambda(1)}=\Res_{s=1}\tilde\zeta_{\Z,\lambda}(s)$.
\item If \CGeqref{eq:defect-selfdual} holds, then $\tilde h_\lambda(1)=0$.
\end{enumerate}
\end{proposition}

\begin{proof}
(i) $\zeta_\lambda$ is bounded by a power of $|t|$ on vertical strips (proof of Lemma~\CGref{lem:summation-positive}; on $\Re s\ge c$ use Proposition~\CGref{prop:family-clock}), and $\Gamma(\frac s2)$ decays exponentially.  Moving the line to $\Re s=1-c$ crosses $s=1$, with residue $\Gamma(\frac12)\pi^{-1/2}y^{-1/2}=y^{-1/2}$, and $s=0$, with residue $2\zeta_\lambda(0)=-1$; on the new line \CGeqref{eq:Z1-FE} and $y^{-s/2}=y^{-1/2}(1/y)^{-(1-s)/2}$ give $y^{-1/2}\Theta_\lambda(1/y)$.  The decay follows by moving the line to $\Re s=2A$.
(ii) By $e^{-x}=\frac1{2\pi i}\int_{(c/2)}\Gamma(w)x^{-w}dw$ and Corollary~\CGref{cor:abscissa-one}, $\sum_N\tilde a_N(\lambda)e^{-\pi N^2y}=\frac1{4\pi i}\int_{(c)}\tilde\zeta_{\Z,\lambda}(s)\Gamma(\frac s2)\pi^{-s/2}y^{-s/2}ds$; subtracting (i) from the theta relation gives the equivalence.  Under the hypotheses of Theorem~\CGref{thm:pair-criterion}, $\tilde\zeta_{\Z,\lambda}=\zeta$.
(iii) By (ii), $2\mathfrak G_\lambda(y)=\frac1{2\pi i}\int_{(c)}\zeta_\lambda(s)\bigl(e^{\tilde h_\lambda(s)}-1\bigr)\Gamma(\frac s2)\pi^{-s/2}y^{-s/2}ds$.  By Lemma~\CGref{lem:vertical-growth}, $|e^{\tilde h_\lambda(s)}|\le e^{o(|t|)}$ on $\frac12+2\delta\le\Re s\le c$, so the line can be moved to $\Re s=\frac12+2\delta$.  The only pole crossed is $s=1$, where $\zeta_\lambda$ has residue $1$ and $\tilde h_\lambda$ is holomorphic; the residue is $(e^{\tilde h_\lambda(1)}-1)y^{-1/2}$, and the remaining integral is $O_\delta(y^{-1/4-\delta})$.
(iv) $\sum_N\tilde a_N(\lambda)e^{-\pi N^2y}=O(e^{-\pi y})$ as $y\to\infty$, so by (i) $\mathfrak G_\lambda(y)=O_A(y^{-A})$ there; then \CGeqref{eq:defect-selfdual} gives $\mathfrak G_\lambda(y)=O_A(y^{A})$ as $y\to0^+$ for every $A$, and \CGeqref{eq:defect-expansion} forces $e^{\tilde h_\lambda(1)}=1$; $\tilde h_\lambda$ is real on the real axis.
\end{proof}

\begin{corollary}[The density of the atomized integers]\label{cor:residue-condition}
$\sum_{N\le x}\tilde a_N(\lambda)$ has mean density $e^{\tilde h_\lambda(1)}$ in the sense of \CGeqref{eq:defect-expansion}, and (D) implies
\begin{equation}\label{eq:residue-condition}
 \tilde h_0(1)=0,\qquad\text{that is,}\qquad \mathfrak u_0(1)=\sum_{m\ge1}\frac1m\sum_{n\ge2}\mathfrak h_n(0;m).
\end{equation}
\end{corollary}

\begin{proof}
Under (D), $\tilde h_0\equiv0$ (Theorem~\CGref{thm:closing}); the second form is \CGeqref{eq:joining-factor} at $s=1$.
\end{proof}

The identity $\tilde h_0(\sigma)=0$ for real $\sigma$ has two distinguished members: $\sigma=+\infty$ is $\varpi_{Z1}=1$, the first coefficient (Corollary~\CGref{cor:product-identity}), and $\sigma=1$ is \CGeqref{eq:residue-condition}, the density of the integers.

\section{Numerical illustration with \texorpdfstring{$2\,001\,052$}{2001052} zeros}
\label{sec:numerics}

This section shows what is being constructed and checks that the theorems are consistent with computation.  It uses the ordinates $\gamma_1,\dots,\gamma_N$ of the first $N=2\,001\,052$ zeros of $\zeta$ (file \texttt{zeross.bin}, $\gamma_N=1\,132\,490.6587$) and the midpoint ordinates $\tau_1,\dots,\tau_N$ computed from \CGeqref{eq:lambert} (file \texttt{zLamb.bin}, $\tau_N=1\,132\,490.7226$).  The programs are listed in Appendix~\CGref{app:code} and their complete console output in Appendix~\CGref{app:output}; all figures have linear axes.  Each item is an \emph{experiment}: it reports a computation and does not replace a proof.

Products over the ordinates are evaluated as follows: the first $10^5$ factors exactly, the factors $10^5<k\le N$ through power sums of $(\tfrac14+o_k^2)^{-1}$, and the factors $k>N$ by the midpoint rule for the clock, which gives the tail $\frac1{2\pi}\bigl(q\frac{L_U+1}U-\frac{q^2}2\frac{3L_U+1}{9U^3}+\frac{q^3}3\frac{5L_U+1}{25U^5}\bigr)$ with $\NR(U)=N$, $L_U=\log(U/2\pi)$.

\subsection{The clock and the zeros}
\label{subsec:num-clock}

\begin{experiment}[Clock, midpoints and displacements]\label{exp:clock}
\emph{Object:} $\NR$, $N_\zeta$, $\tau_k$, $\eta_k=\tau_k-\gamma_k$.  \emph{Range:} $k\le N$.  \emph{Expected:} \CGeqref{eq:midpoint-condition} to rounding, and $|\eta_k|\le C_\eta$ (Lemma~\CGref{lem:bounded-displacement}).  \emph{Observed:} $\max_k|\NR(\tau_k)-(k-\tfrac12)|=6.985\cdot10^{-10}$; $\max_k|\eta_k|=0.9937$ (at $k=34$), mean $|\eta_k|=0.1512$, root mean square $0.1881$, mean $\eta_k=1.7\cdot10^{-7}$.  By blocks of $k$, $\max|\eta_k|$ is $0.840,0.994,0.915,0.834,0.786,0.763,0.717$ on $[1,10]$, $[11,10^2]$, $[10^2,10^3]$, \dots, $[10^6,N]$.
\end{experiment}

Figure~\CGref{fig:clock-zoom} shows how the clock places the midpoint ordinates.  The staircase $N_\zeta$ jumps by one at each $\gamma_k$; the smooth clock $\NR$ crosses the half-integer level $k-\tfrac12$ at $\tau_k$.  Since $N_\zeta-\NR=S(T)+O(1/T)$ (Figure~\CGref{fig:remainder}), the horizontal distance between the open square $(\gamma_k,k-\tfrac12)$ and the red point $(\tau_k,k-\tfrac12)$ is the remainder $S$ converted into ordinate units by the slope $\NR'=\frac1{2\pi}\log\frac T{2\pi}$.  The right panel repeats the picture at height $1.13\cdot10^6$, where the mean spacing is $0.519$: every $\tau_k$ still lies next to its $\gamma_k$.

\begin{figure}[htbp]
\centering
\includegraphics[width=0.86\textwidth]{fig00_clock_zoom.png}
\caption{The midpoint clock $\NR$ (red), the zero staircase $N_\zeta$ (black), the midpoint ordinates $\tau_k$ (red points) and the points $(\gamma_k,k-\frac12)$ (squares).  Left: the first nine zeros.  Right: the last twelve zeros of the file.}
\label{fig:clock-zoom}
\end{figure}

\begin{figure}[htbp]
\centering
\includegraphics[width=0.86\textwidth]{fig00b_remainder.png}
\caption{The remainder $N_\zeta(T)-\NR(T)=S(T)+O(1/T)$ of \CGeqref{eq:RvM-full} for $T\le200$.}
\label{fig:remainder}
\end{figure}

\subsection{The \texorpdfstring{$\lambda$}{lambda}-family and the endpoint}
\label{subsec:num-family}

Figure~\CGref{fig:family} shows the core of the method.  The lower left panel draws the moving ordinates $\gamma_k(\lambda)$ in clock units, $\NR(\gamma_k(\lambda))-(k-\frac12)$: at $\lambda=1$ all lines meet at $0$ (the $\tau_k$), at $\lambda=0$ they sit at the ordinates of $\zeta$; for $\lambda>0$ the ordinates $\gamma_k(\lambda)$ remain strictly increasing in $k$ (Theorem~\CGref{thm:path}).  The lower right panel draws the real functions $Z_\lambda(t)=e^{i\vartheta(t)}\zeta_\lambda(\frac12+it)$ for $\lambda=0,0.25,0.5,0.75,1$ together with Hardy's function $Z(t)=e^{i\vartheta(t)}\zeta(\frac12+it)$: the zeros of every $Z_\lambda$ are real, they move continuously with $\lambda$, and the curve for $\lambda=0$ lies on top of $Z$.

\begin{figure}[htbp]
\centering
\includegraphics[width=0.80\textwidth]{fig01_clock_family.png}
\caption{Top: $\eta_k$ for every 50th $k\le N$ and its distribution.  Bottom: moving ordinates in clock units, $k\le30$, and the functions $Z_\lambda(t)$ on the critical line.}
\label{fig:family}
\end{figure}

\begin{experiment}[The endpoint $\zeta_{Z1}$ against $\zeta$]\label{exp:Z1}
\emph{Object:} $\zeta_{Z1}=P_{Z1}/G$ computed from the $N$ ordinates $\gamma_k$, and $\zeta$ computed by mpmath.  \emph{Range:} $\sigma\in\{2,1.25,0.75\}$, $1\le t\le2000$ ($499$ points per line), and the critical line $0.5\le t\le60$.  \emph{Expected:} if $\zeta_{Z1}=\zeta$, the ratio is $1$ up to the truncation of the product.  \emph{Observed:} $\max|\zeta_{Z1}/\zeta-1|=8.7\cdot10^{-10}$, $1.1\cdot10^{-9}$, $2.6\cdot10^{-9}$ on the three lines, and $\max|Z_{Z1}-Z|/\max(1,|Z|)=8.4\cdot10^{-10}$ on the critical line (Figure~\CGref{fig:Z1}).
\end{experiment}

\begin{figure}[htbp]
\centering
\includegraphics[width=0.86\textwidth]{fig02_zetaZ1_vs_zeta.png}
\caption{Left: $|\zeta_{Z1}/\zeta-1|$ on three vertical lines.  Right: $Z_{Z1}$ and Hardy's $Z$ on the critical line.}
\label{fig:Z1}
\end{figure}

Experiment~\CGref{exp:Z1} is consistent with the identity $\zeta_{Z1}=\zeta$ to nine digits in the range computed; it is a finite computation and not a proof of the identity.

\begin{experiment}[The value at infinity: the product identity]\label{exp:product-identity}
\emph{Object:} the ordered partial products $P_K=\prod_{k\le K}(\frac14+\gamma_k^2)/(\frac14+\tau_k^2)$ of \CGeqref{eq:product-identity}.  \emph{Range:} $K\le N$.  \emph{Expected:} by Corollary~\CGref{cor:product-identity}, $\zeta_{Z1}=\zeta$ is the statement $P_K\to\varpi_{\RVM}$.  \emph{Observed:} $P_K=0.96256$, $0.96248$, $0.96229$, $0.961951$, $0.961930$, $0.9619418189$, $0.9619415390$ for $K=10,10^2,\dots,10^6,N$, against $\varpi_{\RVM}=0.9619415807$ from Experiment~\CGref{exp:transition}; for $K\ge10^6$, $P_K$ stays in $[0.9619389,0.9619440]$ with mean $0.9619416633$ (Figure~\CGref{fig:product-identity}).  The flow at infinity \CGeqref{eq:Upsilon} is $\Upsilon(\lambda)=0.0436950$, $0.0412445$, $0.0387987$, $0.0363572$, $0.0339198$ for $\lambda=0,\frac14,\frac12,\frac34,1$, decreasing as in \CGeqref{eq:varpi-derivative} with slope close to $-\sum_k2\eta_k^2/\gamma_k^2=-0.0098211$, and $\int_0^1\Upsilon\,d\lambda=\sum_{k\le N}\log\frac{1/4+\tau_k^2}{1/4+\gamma_k^2}=0.0388016004$ against $\log L=0.0388015571$.
\end{experiment}

\begin{figure}[htbp]
\centering
\includegraphics[width=0.86\textwidth]{fig11_product_identity.png}
\caption{The partial products $P_K$ of \CGeqref{eq:product-identity} against $\varpi_{\RVM}$ (dashed).  Left: $K\le10^4$.  Right: $(P_K-\varpi_{\RVM})\cdot10^6$ for $10^5\le K\le N$.  The value of the endpoint at $+\infty$ is $\varpi_{Z1}=\varpi_{\RVM}/\lim_KP_K$ (Theorem~\CGref{thm:Z1-clock}).}
\label{fig:product-identity}
\end{figure}

\subsection{The transition theorem and \texorpdfstring{$\zeta_{\RVM}$}{zeta\_RVM} against \texorpdfstring{$\zeta$}{zeta}}
\label{subsec:num-transition}

\begin{experiment}[Transition theorem]\label{exp:transition}
\emph{Object:} $\zeta_{\RVM}$, $w(s)$, $\varpi_{\RVM}$, $L$.  \emph{Range:} $\sigma'\in\{0.25,0.5,1.5\}$, $10\le t\le10^4$.  \emph{Expected:} Theorem~\CGref{thm:transition} and Corollary~\CGref{cor:vertical-limit}: $t\,|\zeta_{\RVM}/(\varpi_{\RVM}(1+w))-1|$ bounded, $|w|=(t/2\pi)^{-\sigma'}(1+O(1/t))$.  \emph{Observed:} Richardson extrapolation on $\Re s=2$ gives $\varpi_{\RVM}=0.9619415807$ and $L=1.0395641690$; $\max_{t>100}t\,|\zeta_{\RVM}/(\varpi_{\RVM}(1+w))-1|=0.0617,0.0618,0.0620$ for the three values of $\sigma'$, and $0.0790$ for $\varpi_{\RVM}(1+\chi)$ at $\sigma'=0.25$; $|w|\,t^{\sigma'}=1.5832$, $2.5066$, $15.66$--$15.75$, i.e.\ $(2\pi)^{\sigma'}$.
\end{experiment}

\begin{figure}[htbp]
\centering
\includegraphics[width=0.80\textwidth]{fig03_transition.png}
\caption{Top: the transition error multiplied by $t$, and the rescaled vertical limit $t^{\min(\sigma',1)}|1/\zeta_{\RVM}-L|$.  Bottom: $|\zeta|$ and $|\zeta_{\RVM}|$ on $\sigma=2$ and $\sigma=0.75$.}
\label{fig:transition}
\end{figure}

The bottom panels of Figure~\CGref{fig:transition} compare $|\zeta|$ and $|\zeta_{\RVM}|$ on two vertical lines: $\zeta$ carries the Euler product and oscillates with the primes, whereas $\zeta_{\RVM}$ approaches the constant $\varpi_{\RVM}$ on $\sigma=2$ and $\varpi_{\RVM}|1+\chi|$ on $\sigma=0.75$.  This is the analytic content of Theorem~\CGref{thm:no-Dirichlet-RVM}.

\subsection{Exponential sums}
\label{subsec:num-expsums}

\begin{experiment}[Exponential sums over the two spectra]\label{exp:expsums}
\emph{Object:} $S(N;\log x)=\sum_{k\le N}x^{io_k}$ for $o_k=\tau_k$ and $o_k=\gamma_k$.  \emph{Range:} $N\le2\,001\,052$, $x\in\{2,3,4,6,10\}$.  \emph{Expected:} $|S_{\RVM}|\ll1+\log N$ (Theorem~\CGref{thm:exp-sum}); for $\gamma_k$, Landau's term $-\frac T{2\pi}\frac{\Lambda(x)}{\sqrt x}$.  \emph{Observed:} $\max_N|S_{\RVM}(N;\log x)|=3.57,3.98,3.95,4.53,9.70$; at $N=2\,001\,052$, $|\sum x^{i\gamma_k}|=88\,301$, $114\,268$, $62\,435$, $0.26$, $2.42$, against Landau's $88\,342$, $114\,324$, $62\,467$, $0$, $0$.
\end{experiment}

\begin{figure}[htbp]
\centering
\includegraphics[width=0.86\textwidth]{fig04_exponential_sums.png}
\caption{The midpoint spectrum carries no prime frequencies (left); the spectrum of $\zeta$ resonates at prime powers (right).}
\label{fig:expsums}
\end{figure}

\subsection{Prime counts: \texorpdfstring{$\pi$}{pi} and \texorpdfstring{$\pi_{\RVM}$}{pi\_RVM}}
\label{subsec:num-primes}

\begin{experiment}[Prime counts and RVM primes]\label{exp:primes}
\emph{Object:} $J_{\RVM}$, $\pi_{\RVM}$, $c_n^{\RVM}$, $p_j^{\RVM}$ against $J$, $\pi$, the primes.  \emph{Range:} $2\le x\le100$ with all $N$ ordinates of each list ($3000$ points); $2\le x\le1000$ with the first $K=17\,061$ ordinates and a smooth cut ($40\,000$ points).  \emph{Expected:} the explicit formula with the $\gamma_k$ reproduces $\pi(x)$; $\pi_{\RVM}$ is continuous and $\pi_{\RVM}(x)-\li(x)=O(\sqrt x(\log\log x)^2/\log x)$ (Theorem~\CGref{thm:Lambert-estimate}).  \emph{Observed:} away from the primes, $|\pi_{\gamma}-\pi|\le0.40$ on $[2,100]$; $\pi_{\RVM}(100)=26.08$, $\pi(100)=25$, $\operatorname{li}(100)=30.13$; $\pi_{\RVM}(1000)=169.98$, $\pi(1000)=168$, $\operatorname{li}(1000)=177.61$.  $\pi_{\RVM}$ is not monotone: it carries a ripple of unit scale, with minimum grid increment $-0.25$ on $[2,100]$.  The first crossings of the levels $j=1,\dots,8$ are $p_j^{\RVM}=2.21,3.16,5.19,7.16,11.14,14.02,17.26,19.18$ against $p_j=2,3,5,7,11,13,17,19$.
\end{experiment}

\begin{figure}[htbp]
\centering
\includegraphics[width=0.86\textwidth]{fig05_prime_counts.png}
\caption{$\pi(x)$ (black staircase), the explicit formula with the zeros of $\zeta$ (blue), $\pi_{\RVM}(x)$ (red) and $\operatorname{li}(x)-\operatorname{li}(2)$ (green).  The distance between the red and green curves is the quantity bounded in Theorem~\CGref{thm:Lambert-estimate}; the ratios are in Figure~\CGref{fig:pi-ratios}.}
\label{fig:prime-counts}
\end{figure}

\begin{figure}[htbp]
\centering
\includegraphics[width=0.86\textwidth]{fig06_rvm_primes.png}
\caption{Left: cell masses $c_n^{\RVM}$, the prime indicator and $1/\log n$.  Right: real RVM prime levels against the primes up to $1000$.}
\label{fig:rvm-primes}
\end{figure}

Figure~\CGref{fig:prime-counts} shows $\pi$ and $\pi_{\RVM}$ in the usual format.  The RVM count follows the prime staircase much more closely than $\operatorname{li}$: replacing each $\gamma_k$ by $\tau_k$ moves every zero by less than one unit (Experiment~\CGref{exp:clock}), and the low ordinates, which shape $\pi(x)$ for small $x$, are moved least in relative terms.  The steps are replaced by a continuous ripple, so the cell masses $c_n^{\RVM}$ of Figure~\CGref{fig:rvm-primes} are real numbers, occasionally negative, and are not a prime indicator.

\begin{experiment}[The three ratios and the square-root scale]\label{exp:pi-ratios}
\emph{Object:} $\li/\pi_{\RVM}$, $\pi/\li$, $\pi/\pi_{\RVM}$, and $(\pi_{\RVM}-\li)/B$, $(\pi-\li)/B$ with $B(x)=\sqrt x(\log\log x)^2/\log x$.  \emph{Range:} the grid of Experiment~\CGref{exp:primes} on $[8,1000]$ ($K=17\,061$ ordinates).  \emph{Expected:} the three ratios tend to $1$ (Corollary~\CGref{cor:three-limits}) and $(\pi_{\RVM}-\li)/B$ stays bounded (Theorem~\CGref{thm:Lambert-estimate}).  \emph{Observed:} at $x=100,500,1000$, $\li/\pi_{\RVM}=1.155,1.086,1.045$, $\pi/\li=0.830,0.933,0.946$, $\pi/\pi_{\RVM}=0.959,1.013,0.988$; $\max|\pi_{\RVM}-\li|/B=2.72$ on $[8,1000]$, $1.35$ on $[50,1000]$ and $0.72$ on $[500,1000]$, while $\max|\pi-\li|/B=2.45$ on $[8,1000]$ (Figure~\CGref{fig:pi-ratios}).
\end{experiment}

\begin{figure}[htbp]
\centering
\includegraphics[width=0.86\textwidth]{fig10_prime_ratios.png}
\caption{Left: the three ratios of Corollary~\CGref{cor:three-limits}.  Right: $\pi_{\RVM}-\li$ and $\pi-\li$ in units of the scale $B(x)$ of Theorem~\CGref{thm:Lambert-estimate}.  The ratio $\pi/\pi_{\RVM}$ (black) is the closest of the three to $1$; the ratios with $\li$ approach $1$ at the rate $\tfrac12\li(\sqrt x)/\li(x)$ of the bias $J-\pi$.}
\label{fig:pi-ratios}
\end{figure}

On $[2,1000]$ the two normalized errors in Figure~\CGref{fig:pi-ratios} have the same size and sign: both counts lie below $\li$ by about $\tfrac12\li(\sqrt x)$, the $m=2$ term of the M\"obius inversion \CGeqref{eq:pi-minus-li}.  This is why $\pi/\pi_{\RVM}$ is much closer to $1$ than either ratio with $\li$.

\subsection{The Mertens functions and the ternary division}
\label{subsec:num-mertens}

The Mertens function is the sum of the ternary sequence $\mu(n)\in\{-1,0,+1\}$.  Its values split into two independent pieces of information:
\begin{enumerate}[label=\textup{(\alph*)}]
\item \emph{the zero set} $\{n:\mu(n)=0\}$, the non-squarefree integers, of density $1-6/\pi^2=0.392073$;
\item \emph{the signs} on the squarefree integers, $\mu(n)=(-1)^{\omega(n)}$, each sign of density $3/\pi^2=0.303964$.
\end{enumerate}
The ternary walk $\mu_h$ of \CGeqref{eq:intro-muh}--\CGeqref{eq:intro-quantizer} keeps (a) exactly, $\mu_h(n)=0\Leftrightarrow\mu(n)=0$, and replaces (b): at a squarefree $n$ it steps up when $M_h(n-1)$ lies below the target $\mathcal M(n)=M_{\RVM}^{\crit}(n)$ and down otherwise.  Thus
\begin{equation}\label{eq:ternary-update}
 M_h(n)=M_h(n-1)+\mu(n)^2\,\varsigma_n,\qquad |M_h(n)-\mathcal M(n)|\ \text{small},
\end{equation}
so that $M_h$ is an integer walk with the zero set of $\mu$ that follows $M_{\RVM}^{\crit}$, in the same way as $M$ is the integer walk with the zero set of $\mu$ whose signs are $(-1)^{\omega(n)}$.

\begin{table}[htbp]
\centering
\small
\begin{tabular}{rrrrrcrrr}
\toprule
$n$ & $\mu(n)$ & $\mu(n)^2$ & $\mathcal M(n)$ & $M_h(n-1)$ & $\varsigma_n$ & $\mu_h(n)$ & $M_h(n)$ & $M(n)$\\
\midrule
1 & 1 & 1 & 0.2925 & 0 & $+1$ & 1 & 1 & 1\\
2 & $-1$ & 1 & 0.4563 & 1 & $-1$ & $-1$ & 0 & 0\\
3 & $-1$ & 1 & $-0.6523$ & 0 & $-1$ & $-1$ & $-1$ & $-1$\\
4 & 0 & 0 & $-0.8720$ & $-1$ & -- & 0 & $-1$ & $-1$\\
5 & $-1$ & 1 & $-1.7255$ & $-1$ & $-1$ & $-1$ & $-2$ & $-2$\\
6 & 1 & 1 & $-1.3063$ & $-2$ & $+1$ & 1 & $-1$ & $-1$\\
7 & $-1$ & 1 & $-1.5931$ & $-1$ & $-1$ & $-1$ & $-2$ & $-2$\\
8 & 0 & 0 & $-2.4029$ & $-2$ & -- & 0 & $-2$ & $-2$\\
9 & 0 & 0 & $-1.3292$ & $-2$ & -- & 0 & $-2$ & $-2$\\
10 & 1 & 1 & $-1.3280$ & $-2$ & $+1$ & 1 & $-1$ & $-1$\\
11 & $-1$ & 1 & $-2.3272$ & $-1$ & $-1$ & $-1$ & $-2$ & $-2$\\
12 & 0 & 0 & $-2.2067$ & $-2$ & -- & 0 & $-2$ & $-2$\\
13 & $-1$ & 1 & $-1.8468$ & $-2$ & $+1$ & 1 & $-1$ & $-3$\\
14 & 1 & 1 & $-1.9324$ & $-1$ & $-1$ & $-1$ & $-2$ & $-2$\\
15 & 1 & 1 & $-1.5435$ & $-2$ & $+1$ & 1 & $-1$ & $-1$\\
\bottomrule
\end{tabular}
\caption{The ternary division for $n\le15$ (the first $40$ values are in Appendix~\CGref{app:out-s05-mertens}).  The zero columns of $\mu$ and $\mu_h$ coincide; the signs of $\mu_h$ are chosen by comparing $M_h(n-1)$ with $\mathcal M(n)$.  For $n\le12$ the two walks agree; at $n=13$ the target $-1.85$ lies above $M_h(12)=-2$, so $\mu_h(13)=+1$ while $\mu(13)=-1$.}
\label{tab:ternary}
\end{table}

\begin{experiment}[Mertens functions and the ternary walk]\label{exp:mertens}
\emph{Object:} $M$, $M_{\RVM}^{\crit}$, $M_h$; the explicit formulas.  \emph{Range:} $n\le10^6$; explicit formulas on $[2,2000]$ with all $N$ ordinates of each list.  \emph{Expected:} the explicit formula with the zeros of $\zeta$ reproduces $M$; $M_{\RVM}^{\crit}(x)=O(x^{1/2+\eps})$ (Theorem~\CGref{thm:unsmoothed-RVM}); $M_h$ tracks $M_{\RVM}^{\crit}$ with the densities of $\mu$.  \emph{Observed:}
\begin{itemize}
\item explicit formula for $M$ with the $\gamma_k$: root mean square error $0.0005$ on $[2,2000]$ off the integers;
\item $\zeta_{\RVM}'(\rho_k^{\RVM})/\log(\tau_k/2\pi)=0.96194$ at $k=1000$, against $\varpi_{\RVM}=0.96194$ (Corollary~\CGref{cor:derivative-zeros}); the values $|\zeta'(\rho_k)|/\log(\gamma_k/2\pi)$ are spread over $(0,3)$ (Figure~\CGref{fig:explicit});
\item spectral sums $2\sum|a|^2$: $0.029034$ for $\zeta$, $0.027203$ for the midpoint system;
\item the integer target agrees with the full explicit formula for $M_{\RVM}^{\crit}$ to $0.019$ at $100$ integers $\le10^5$;
\item $\max|M_h-\mathcal M|=2.62$ on $[1,10^6]$ and $1.65$ on $[10^4,10^6]$;
\item densities of $(-1,0,+1)$ for $\mu$: $0.303857$, $0.392074$, $0.304069$;
\item densities of $(-1,0,+1)$ for $\mu_h$: $0.304003$, $0.392074$, $0.303923$;
\item $\mu_h(n)=\mu(n)$ at $50.03\%$ of the squarefree $n\le10^6$, and $\sum_{n\le10^6}\mu(n)\mu_h(n)/\#\{n\le10^6:\mu(n)\ne0\}=0.0007$: the sign sequences are uncorrelated although both walks have the same zero set and comparable size;
\item $\max_{10\le n\le10^6}|W(n)|/\sqrt n=0.832,0.754,0.688$ for $W=M,M_{\RVM}^{\crit},M_h$;
\item energies $\frac1N\sum_{n\le N}W(n)^2$ at $N=10^6$: $14\,851$, $12\,538$, $12\,538$.
\end{itemize}
\end{experiment}

\begin{figure}[htbp]
\centering
\includegraphics[width=0.86\textwidth]{fig07_mertens.png}
\caption{Left: $M$ (black), $M_{\RVM}^{\crit}$ (red) and $M_h$ (green) at every integer $n\le2000$.  Right: range of $M$ and mean of $M_{\RVM}^{\crit}$ in blocks of $500$ integers up to $10^6$, with $\pm\sqrt n$.}
\label{fig:mertens}
\end{figure}

\begin{figure}[htbp]
\centering
\includegraphics[width=0.86\textwidth]{fig08_ternary_walk.png}
\caption{Left: $\mu(n)$ and $\mu_h(n)$ for $n\le60$: same zeros, different signs.  Right: $M$, the target $M_{\RVM}^{\crit}$ and $M_h$ for $n\le300$.}
\label{fig:ternary}
\end{figure}

\begin{figure}[htbp]
\centering
\includegraphics[width=0.86\textwidth]{fig09_explicit_weights.png}
\caption{Left: $M(x)$, its explicit formula from the zeros of $\zeta$ and $M_{\RVM}^{\crit}(x)$ from the $\tau_k$ on $[2,300]$.  Right: normalized derivatives at the zeros: erratic for $\zeta$, constant $\varpi_{\RVM}$ for $\zeta_{\RVM}$.}
\label{fig:explicit}
\end{figure}

Figures~\CGref{fig:mertens}--\CGref{fig:explicit} summarize the Mertens side.  $M_{\RVM}^{\crit}$ has the size and the visual texture of $M$, it stays inside $\pm\sqrt n$ as Theorem~\CGref{thm:unsmoothed-RVM} requires, and its residues are deterministic; $M_h$ is its ternary integer version with the zero set of $\mu$.  The passage from these synthetic objects to $M$ itself is not numerical: it is the identity $\zeta_{Z1}=\zeta$ of Step~3, which Experiment~\CGref{exp:Z1} is consistent with.

\section{Conclusions and open problems}
\label{sec:conclusions}

The following statements are proved:
\begin{enumerate}[label=\textup{(\arabic*)}]
\item the midpoint ordinates $\tau_k$ are given by the Lambert formula \CGeqref{eq:lambert} and stay within a bounded distance of the ordinates $\gamma_k$ of $\zeta$ (Lemma~\CGref{lem:bounded-displacement});
\item $P_{\RVM}^{\crit}$ is entire and symmetric, and $\zeta_{\RVM}$ has exactly the simple zeros $\tfrac12\pm i\tau_k$, the trivial zeros and a residue-one pole at $1$ (Theorem~\CGref{thm:canonical});
\item $S_{\RVM}(N;u)\ll1+\log N$ uniformly on compact positive frequency intervals, so the point centroid and $M_H^{\RVM}(x)$ tend to $0$ (Section~\CGref{sec:exp-sums});
\item $J_{\RVM}$ exists compact-uniformly and is related to $\pi_{\RVM}$ by exact finite M\"obius inversion (Section~\CGref{sec:explicit-formula});
\item $J_{\RVM}(x)-\li(x)$ and $\pi_{\RVM}(x)-\li(x)$ are $O(\sqrt x(\log\log x)^2/\log x)$, so $\li/\pi_{\RVM}\to1$, and with the prime number theorem $\pi/\li\to1$ and $\pi/\pi_{\RVM}\to1$ (Theorem~\CGref{thm:Lambert-estimate}, Corollary~\CGref{cor:three-limits});
\item $J_{\RVM}$ is Lipschitz on $[2,\infty)$, $c_n^{\RVM}=O(1)$, and $\zeta_{\Z,\RVM}$ is an absolutely convergent Dirichlet series in $\Re s>1$ (Section~\CGref{sec:integer-weights});
\item $\zeta_{\RVM}=\varpi_{\RVM}(1+w)(1+O(1/|t|))$ on every band to the right of $\tfrac12$ and $\zeta_{\RVM}=\varpi_{\RVM}(1+\chi)(1+O(1/|t|))$ on closed substrips of $\tfrac12<\Re s<1$ (Theorem~\CGref{thm:transition});
\item $1/\zeta_{\RVM}\to L$ on every vertical line to the right of $\tfrac12$, $\zeta_{\RVM}^{\pm1}$ have no absolutely convergent Dirichlet expansion, and $M_{\RVM}^{\crit}(x)=L+O_\eps(x^{1/2+\eps})$ (Section~\CGref{sec:transition});
\item $2\xi$ equals its Hadamard product over $\rho(1-\rho)$ exactly (Theorem~\CGref{thm:Hadamard}), and $|\zeta|\asymp|\mathcal Q_{\RVM}|$ on $\frac12+\eps\le\sigma\le c'$, $|t|\ge t_0$ (Corollary~\CGref{cor:Q-control});
\item the $\lambda$-family from $\zeta_{\RVM}$ to $\zeta_{Z1}$ converges jointly, has simple zeros for $\lambda>0$, an explicit generator and total multiplier, and its members differ only through $\gamma_k-\tau_k$ (Section~\CGref{sec:lambda-family});
\item $\zeta_{Z1}$ has the functional equation, the pole, the residue and the order of $\zeta$, and $\zeta_{Z1}=\zeta\Leftrightarrow$(D)$\Leftrightarrow$RH (Theorem~\CGref{thm:Hamburger-Z1});
\item $\zeta_{\RVM}$ is bounded above and below on every half-plane $\Re s\ge\sigma_0>1$ and tends to $\varpi_{\RVM}$ as $\Re s\to+\infty$ (Theorem~\CGref{thm:halfplane-RVM}); $1/\zeta_{\Z,\RVM}$ is an absolutely convergent Dirichlet series in $\Re s>1$, so (D) is equivalent to an absolutely convergent Dirichlet expansion of $\mathcal K$ in $\Re s>1$ (Theorem~\CGref{thm:D-multiplier});
\item no member of the $\lambda$-family, $\zeta_{Z1}$ included, can have a zero off the critical line other than the trivial zeros (Theorem~\CGref{thm:path});
\item the endpoint has the value $\varpi_{Z1}=\varpi_{\RVM}\prod_k(\frac14+\tau_k^2)/(\frac14+\gamma_k^2)\ge1$ at $+\infty$, and $\zeta_{Z1}=\zeta$ if and only if $\varpi_{Z1}=1$ (Theorems~\CGref{thm:Z1-clock}, \CGref{thm:varpi-Z1}, Corollary~\CGref{cor:product-identity});
\item the prime-power counts of $\zeta_{\RVM}$ and $\zeta_{Z1}$ are the Mellin transforms of their logarithms and satisfy a distributional explicit formula; this proves the Mellin--Stieltjes identification and the weighted factorization $\zeta_{\Z,\RVM}=\zeta_{\RVM}e^{h_{\RVM}}$, and (D) holds if and only if the prime-power count of $\zeta_{Z1}$ is constant between consecutive integers (Theorems~\CGref{thm:mellin-counts}, \CGref{thm:distributional-EF}, \CGref{thm:MS}, \CGref{thm:cell-form});
\item along the family $\frac d{d\lambda}\log\varpi_\lambda=-\Upsilon(\lambda)$, and (D) is the identity $\int_{0^+}^1\Upsilon\,d\lambda=\log L$; with the verified zeros, $0\le\log\varpi_{Z1}<4\cdot10^{-13}$ (Propositions~\CGref{prop:family-flow}, \CGref{prop:varpi-verified});
\item the atomized family $\tilde\zeta_{\Z,\lambda}=\zeta_\lambda e^{\tilde h_\lambda}$ consists of absolutely convergent Dirichlet series in $\Re s>1$ for every $\lambda$, the endpoint included, because the cell masses have mean square $X^{1+\eps}$ on dyadic blocks uniformly in $\lambda$; the joining factor is explicit cell by cell and holomorphic in $\Re s>\frac12$ (in $\Re s>0$ for $\lambda=1$), and (D) is equivalent to the vanishing of every cell defect at $\lambda=0$ and of $J_{Z1}$ on the first cell $(1,2)$ (Proposition~\CGref{prop:first-cell}, Theorems~\CGref{thm:joining-factor}, \CGref{thm:mean-square-cells}, \CGref{thm:atomized-continuation}, \CGref{thm:closing});
\item at the clock endpoint the joining factor is not constant, $\tilde h_1(\sigma)\to\log L$; the atomized reciprocals have M\"obius coefficients $\tilde\mu_\lambda$, and the endpoint Mertens function has a Mellin form $B_{Z1}$ with $B_{Z1}=M$ if and only if (D) (Propositions~\CGref{prop:atomized-clock}, \CGref{prop:Mellin-Mertens});
\item Step~3: $\zeta_{Z1}=\zeta$ implies $M(x)=O_\eps(x^{1/2+\eps})$ (Theorem~\CGref{thm:step3}); equivalently, the synthetic bound $M_{\RVM}^{\crit}(x)=L+O_\eps(x^{1/2+\eps})$ transfers to $M$ through $1/\zeta=F_{\RVM}\Pi_{\rm tr}^{-1}$ (Theorem~\CGref{thm:transfer});
\item for $0<\lambda\le1$ no $\zeta_\lambda$ is an absolutely convergent Dirichlet series in $\Re s>1$ and no joining factor $\tilde h_\lambda$ is constant, and $\lambda\mapsto\log\zeta_\lambda(\sigma)$ is strictly convex for every real $\sigma>1$ (Propositions~\CGref{prop:nonconstant}, \CGref{prop:convexity});
\item if the joining factor $e^{\tilde h_\lambda}$ is meromorphic of finite order with finitely many poles and symmetric under $s\mapsto1-s$, Hamburger's theorem gives $\tilde\zeta_{\Z,\lambda}=\zeta$, and for $\lambda=0$ also $\zeta_{Z1}=\zeta$ (Theorem~\CGref{thm:pair-criterion});
\item if $\zeta_\lambda$ is the Mellin--Stieltjes transform of a positive measure, then $\varpi_\lambda\le1$; (D) holds if and only if $\zeta_{Z1}$ has such a representation, if and only if $J^{\rm M}_{Z1}$ is non-decreasing (Theorem~\CGref{thm:positivity}, Corollary~\CGref{cor:positivity});
\item $J_{\RVM}$ is continuously differentiable on $(2,\infty)$ with $J_{\RVM}'(y)=\frac1{\log y}+2\cos(2\pi y)-2^{-1/2}y^{-1/4}\cos(4\pi\sqrt y+\frac\pi4)+O(y^{-1/3})$; consequently $\tilde c_n(1)=\frac1{\log n}+O(n^{-1/4})$ off the powers of $2$, and Hypothesis~\CGref{hyp:cell-remainder} holds off the integers $2^j$ with $\mu(j)\ne0$ (Theorem~\CGref{thm:clock-density}, Corollary~\CGref{cor:clock-cells});
\item the cell masses are sums over the ordinates with the height kernel $\hat\varphi(\gamma/2\pi n)$, which vanishes at the heights $2\pi qn$; their flow is the pairing of $V_n$ with the count difference $N_\zeta-n_{\RVM}$; and the Gaussian defect of the atomization is $\frac12(e^{\tilde h_\lambda(1)}-1)y^{-1/2}+O(y^{-1/4-\delta})$ (Propositions~\CGref{prop:height-kernels}, \CGref{prop:flow-count}, \CGref{prop:gaussian-defect}).
\end{enumerate}
The computations of Section~\CGref{sec:numerics} with $2\,001\,052$ zeros are consistent with each of these statements.

\begin{remark}[Open problems]\label{rem:open-problems}
\begin{enumerate}[label=\textup{(\roman*)}]
\item Prove the identity $\zeta_{Z1}=\zeta$, equivalently condition (D), equivalently an absolutely convergent Dirichlet expansion of $\mathcal K=\zeta_{Z1}/\zeta_{\Z,\RVM}$ in $\Re s>1$, equivalently the product identity \CGeqref{eq:product-identity}, equivalently the cell identity \CGeqref{eq:cell-identity} with the jump bound of Theorem~\CGref{thm:cell-form}, equivalently step A7 of the programme in Subsection~\CGref{subsec:programme}, equivalently the positivity of the measure representing $\zeta_{Z1}$ (Corollary~\CGref{cor:positivity}).  This is the final step of the method.
\item Hypothesis~\CGref{hyp:cell-remainder} is proved for the integers $n$ that are not of the form $2^j$ with $\mu(j)\ne0$ (Corollary~\CGref{cor:clock-cells}); at $n=2^j$ the cell mass is $(1+\mu(j)J_{\RVM}(2)\log2)/\log n+O(n^{-1/4})$.  Determine the sign of $1\pm J_{\RVM}(2)\log2$.
\item Determine the monotone branches of $\pi_{\RVM}$, on which the real prime levels $p_j^{\RVM}$ are defined, and explain the observed proximity $p_j^{\RVM}\approx p_j$ for small $j$ (Experiment~\CGref{exp:primes}).
\end{enumerate}
\end{remark}

\appendix
\section{Python programs}
\label{app:code}

The programs below produce every number, table and figure of Section~\CGref{sec:numerics}.  They read the two data files \texttt{data/zeross.bin} (ordinates $\gamma_k$ of the zeros of $\zeta$) and \texttt{data/zLamb.bin} (midpoint ordinates $\tau_k$), both little-endian float64 with $2\,001\,052$ entries, and write their console output, CSV, MAT and binary files to \texttt{output/} and the figures to \texttt{figures/}.  They require Python~3 with NumPy, SciPy, mpmath and Matplotlib, and are run from the directory \texttt{code/} in the order \texttt{s01}--\texttt{s07}.

\subsection{\texttt{rvm\_common.py}}
\label{app:code-rvm-common}

\begin{lstlisting}[language=Python]
"""Common routines for the midpoint Riemann-von Mangoldt (RVM) zeta system.

Data files (little-endian float64, one ordinate per entry):
    data/zeross.bin  ordinates gamma_k of the nontrivial zeros of zeta (2,001,052 values)
    data/zLamb.bin   midpoint RVM ordinates tau_k, N_RVM(tau_k) = k - 1/2

Notation follows the paper: N_RVM is the clock, tau_k the midpoint ordinates,
gamma_k the ordinates of zeta, eta_k = tau_k - gamma_k, q(s) = s(s-1),
G(s) = s(s-1) Gamma(s/2) pi^(-s/2), zeta_lambda = P_lambda / G with ordinates
gamma_k(lambda) = (1-lambda) gamma_k + lambda tau_k.
"""
import os
import sys
import time
import numpy as np
from scipy import special

ROOT = os.path.dirname(os.path.dirname(os.path.abspath(__file__)))
DATA = os.path.join(ROOT, "data")
OUT = os.path.join(ROOT, "output")
FIG = os.path.join(ROOT, "figures")
os.makedirs(OUT, exist_ok=True)
os.makedirs(FIG, exist_ok=True)

TWO_PI = 2.0 * np.pi


class Tee:
    """Write console output both to the screen and to output/<name>_console.txt."""

    def __init__(self, name):
        self.f = open(os.path.join(OUT, name + "_console.txt"), "w")
        self.t0 = time.time()

    def __call__(self, fmt, *args):
        s = fmt % args if args else fmt
        sys.stdout.write(s + "\n")
        sys.stdout.flush()
        self.f.write(s + "\n")
        self.f.flush()

    def close(self):
        self("total run time %.1f s" % (time.time() - self.t0))
        self.f.close()


def read_bin(name):
    return np.fromfile(os.path.join(DATA, name), dtype="<f8")


def load_ordinates():
    gam = read_bin("zeross.bin")
    tau = read_bin("zLamb.bin")
    n = min(gam.size, tau.size)
    return gam[:n], tau[:n]


def clock(T):
    """N_RVM(T) = T/(2 pi) log(T/(2 pi e)) + 7/8 (real or complex T, principal log)."""
    return T / TWO_PI * (np.log(T / TWO_PI) - 1.0) + 0.875


def tail_coefficients(U):
    """Midpoint-rule tail of sum_{k>N} log(1 + q/tau_k^2), N = N_RVM(U):
    returns (c1, c2, c3) with tail(q) = c1 q - c2 q^2 + c3 q^3."""
    LU = np.log(U / TWO_PI)
    c1 = (LU + 1.0) / U / TWO_PI
    c2 = (3.0 * LU + 1.0) / (9.0 * U ** 3) / 2.0 / TWO_PI
    c3 = (5.0 * LU + 1.0) / (25.0 * U ** 5) / 3.0 / TWO_PI
    return c1, c2, c3


def clock_inverse(y):
    """Solve N_RVM(U) = y by Newton iteration."""
    U = TWO_PI * max(y, 2.0) / np.log(max(y, 3.0))
    for _ in range(80):
        U -= (clock(U) - y) / (np.log(U / TWO_PI) / TWO_PI)
    return U


class Product:
    """log P(s) = sum_k log(1 + q(s)/(1/4 + o_k^2)) for an ordinate list o_k.

    Exact summation for k <= K1, power sums for K1 < k <= N, and the clock tail
    for k > N.  Accurate for |q(s)| << o_{K1}^2."""

    def __init__(self, ords, K1=100000, nterms=6):
        self.o = np.asarray(ords, dtype=float)
        self.N = self.o.size
        self.K1 = min(K1, self.N)
        A = 0.25 + self.o[self.K1:] ** 2
        self.pw = [np.sum(A ** (-j)) for j in range(1, nterms + 1)]
        self.A1 = 0.25 + self.o[: self.K1] ** 2
        self.tail = tail_coefficients(clock_inverse(float(self.N)))

    def logP(self, s):
        s = np.atleast_1d(np.asarray(s, dtype=complex))
        q = s * (s - 1.0)
        out = np.empty(s.shape, dtype=complex)
        for i, qi in enumerate(q):
            v = np.sum(np.log1p(qi / self.A1))
            for j, p in enumerate(self.pw, start=1):
                v += (-1) ** (j + 1) * qi ** j / j * p
            c1, c2, c3 = self.tail
            v += c1 * qi - c2 * qi ** 2 + c3 * qi ** 3
            out[i] = v
        return out


def logG(s):
    s = np.asarray(s, dtype=complex)
    return np.log(s) + np.log(s - 1.0) + special.loggamma(s / 2.0) - s / 2.0 * np.log(np.pi)


def zeta_from_product(prod, s):
    s = np.atleast_1d(np.asarray(s, dtype=complex))
    return np.exp(prod.logP(s) - logG(s))


def chi(s):
    s = np.asarray(s, dtype=complex)
    return np.exp((s - 0.5) * np.log(np.pi) + special.loggamma((1 - s) / 2) - special.loggamma(s / 2))


def w_transition(s):
    """w(s) = exp(-2 pi i N_RVM(tau_*)), tau_* = -i (s - 1/2)."""
    s = np.asarray(s, dtype=complex)
    tstar = -1j * (s - 0.5)
    return np.exp(-2j * np.pi * clock(tstar))


def theta_rs(t):
    t = np.asarray(t, dtype=float)
    return t / 2 * np.log(t / TWO_PI) - t / 2 - np.pi / 8 + 1 / (48 * t) + 7 / (5760 * t ** 3)


def gamma_lambda(gam, tau, lam):
    return (1.0 - lam) * gam + lam * tau


def ei_upper(w):
    """Ei on Im w != 0 along the horizontal line (paper, eq. Ei-def): Ei(w) = -E1(-w)."""
    return -special.exp1(-np.asarray(w, dtype=complex))


def mobius_sieve(X):
    mu = np.ones(X + 1, dtype=np.int8)
    mu[0] = 0
    is_comp = np.zeros(X + 1, dtype=bool)
    for p in range(2, X + 1):
        if not is_comp[p]:
            is_comp[2 * p::p] = True
            mu[p::p] *= -1
            if p * p <= X:
                mu[p * p::p * p] = 0
    return mu


def primes_upto(X):
    s = np.ones(X + 1, dtype=bool)
    s[:2] = False
    for p in range(2, int(X ** 0.5) + 1):
        if s[p]:
            s[p * p::p] = False
    return np.nonzero(s)[0]


def write_csv(name, header, cols):
    A = np.column_stack(cols)
    np.savetxt(os.path.join(OUT, name), A, delimiter=",", header=",".join(header), comments="", fmt="%.12g")


def savefig(fig, name):
    fig.savefig(os.path.join(FIG, name), dpi=150, bbox_inches="tight")
\end{lstlisting}

\subsection{\texttt{s01\_clock\_and\_family.py}}
\label{app:code-s01-clock-and-family}

\begin{lstlisting}[language=Python]
"""Experiment 1: the midpoint clock, the displacements eta_k = tau_k - gamma_k,
the lambda-family zeta_lambda, and the numerical comparison zeta_Z1 versus zeta."""
import numpy as np
import mpmath as mp
import scipy.io as sio
import matplotlib
matplotlib.use("Agg")
import matplotlib.pyplot as plt
from rvm_common import (Tee, load_ordinates, clock, Product, zeta_from_product, gamma_lambda,
                        write_csv, savefig, OUT)
import os

say = Tee("s01_clock_and_family")
gam, tau = load_ordinates()
N = gam.size
k = np.arange(1, N + 1)
say("zeros loaded: N = %d, gamma_1 = %.9f, gamma_N = %.6f, tau_1 = %.9f, tau_N = %.6f",
    N, gam[0], gam[-1], tau[0], tau[-1])

# ---- clock check and displacements
clock_err = np.max(np.abs(clock(tau) - (k - 0.5)))
eta = tau - gam
say("max |N_RVM(tau_k) - (k - 1/2)| = %.3e", clock_err)
say("eta_k = tau_k - gamma_k:  max |eta| = %.6f (k = %d),  mean |eta| = %.6f,  rms = %.6f,  mean eta = %.3e",
    np.max(np.abs(eta)), int(np.argmax(np.abs(eta))) + 1, np.mean(np.abs(eta)),
    np.sqrt(np.mean(eta ** 2)), np.mean(eta))
blocks = [(1, 10), (11, 100), (101, 1000), (1001, 10 ** 4), (10 ** 4 + 1, 10 ** 5), (10 ** 5 + 1, 10 ** 6), (10 ** 6 + 1, N)]
rows = []
say("%12s %12s %12s %12s", "k from", "k to", "max|eta|", "rms eta")
for a, b in blocks:
    e = eta[a - 1:b]
    rows.append([a, b, np.max(np.abs(e)), np.sqrt(np.mean(e ** 2))])
    say("%12d %12d %12.6f %12.6f", a, b, rows[-1][2], rows[-1][3])
rows = np.array(rows)
write_csv("s01_eta_blocks.csv", ["k_from", "k_to", "max_abs_eta", "rms_eta"], [rows[:, i] for i in range(4)])
say("first ordinates:")
say("%4s %16s %16s %12s", "k", "gamma_k", "tau_k", "eta_k")
for i in range(12):
    say("%4d %16.9f %16.9f %12.6f", i + 1, gam[i], tau[i], eta[i])

# ---- zeta_Z1 versus zeta on vertical lines
PZ = Product(gam)
PR = Product(tau)
tline = np.concatenate([np.linspace(1.0, 100.0, 199), np.linspace(101.0, 2000.0, 300)])
mp.mp.dps = 20
cmp_rows = []
for sigma in (2.0, 1.25, 0.75):
    s = sigma + 1j * tline
    zz1 = zeta_from_product(PZ, s)
    ze = np.array([complex(mp.zeta(complex(x))) for x in s])
    rel = np.abs(zz1 / ze - 1.0)
    cmp_rows.append([sigma, np.max(rel), np.median(rel)])
    say("sigma = %.2f, t in [1,2000] (%d points): max |zeta_Z1/zeta - 1| = %.3e, median = %.3e",
        sigma, tline.size, np.max(rel), np.median(rel))
    write_csv("s01_zetaZ1_vs_zeta_sigma%.2f.csv" % sigma, ["t", "abs_zeta", "abs_zetaZ1", "rel_err"],
              [tline, np.abs(ze), np.abs(zz1), rel])
cmp_rows = np.array(cmp_rows)

# ---- Hardy-type functions Z_lambda(t) = exp(i theta(t)) zeta_lambda(1/2+it) on the critical line
def theta_exact(t):
    return np.array([float(mp.siegeltheta(x)) for x in t])

tc = np.linspace(0.5, 60.0, 1200)
th = theta_exact(tc)
lams = [0.0, 0.25, 0.5, 0.75, 1.0]
Zl = {}
for lam in lams:
    P = PZ if lam == 0.0 else (PR if lam == 1.0 else Product(gamma_lambda(gam, tau, lam)))
    Zl[lam] = np.real(np.exp(1j * th) * zeta_from_product(P, 0.5 + 1j * tc))
Zhardy = np.array([float(mp.siegelz(x)) for x in tc])
dZ = np.max(np.abs(Zl[0.0] - Zhardy) / np.maximum(1.0, np.abs(Zhardy)))
say("critical line, t in [0.5,60]: max |Z_Z1(t) - Z(t)| / max(1,|Z(t)|) = %.3e", dZ)
write_csv("s01_Z_lambda_critical_line.csv", ["t", "Z_hardy"] + ["Z_lambda_%.2f" % l for l in lams],
          [tc, Zhardy] + [Zl[l] for l in lams])

# ---- figures
dec = np.arange(0, N, 50)
fig, ax = plt.subplots(2, 2, figsize=(12, 8))
ax[0, 0].plot(k[dec], eta[dec], ".", ms=1.5, color="tab:blue")
ax[0, 0].axhline(0, color="k", lw=0.5)
ax[0, 0].set_xlabel("k")
ax[0, 0].set_ylabel(r"$\eta_k=\tau_k-\gamma_k$")
ax[0, 0].set_title(r"Bounded displacement, $N=%d$ zeros" % N)
ax[0, 1].hist(eta, bins=200, density=True, color="tab:blue")
ax[0, 1].set_xlabel(r"$\eta_k$")
ax[0, 1].set_title(r"Distribution of $\eta_k$ (all $k\leq N$)")
lam = np.linspace(0, 1, 101)
for kk in range(30):
    gl = (1 - lam) * gam[kk] + lam * tau[kk]
    ax[1, 0].plot(lam, clock(gl) - (kk + 0.5), "-", lw=1)
ax[1, 0].axhline(0, color="k", lw=0.5)
ax[1, 0].set_xlabel(r"$\lambda$")
ax[1, 0].set_ylabel(r"$\mathcal{N}_{\rm RVM}(\gamma_k(\lambda))-(k-\frac{1}{2})$")
ax[1, 0].set_title(r"Moving ordinates in clock units, $k\leq30$ ($\lambda=0$: $\zeta_{Z1}$, $\lambda=1$: $\zeta_{\rm RVM}$)", fontsize=9)
ax[1, 1].plot(tc, Zhardy, "k-", lw=2.2, label=r"Hardy $Z(t)$ of $\zeta$")
cols = ["tab:red", "tab:orange", "tab:green", "tab:cyan", "tab:blue"]
for lam, c in zip(lams, cols):
    ax[1, 1].plot(tc, Zl[lam], "-", color=c, lw=1.0, label=r"$Z_\lambda$, $\lambda=%.2f$" % lam)
ax[1, 1].axhline(0, color="k", lw=0.5)
ax[1, 1].set_xlabel("t")
ax[1, 1].set_title(r"$e^{i\vartheta(t)}\zeta_\lambda(\frac{1}{2}+it)$ along the family")
ax[1, 1].legend(fontsize=7, ncol=2)
fig.tight_layout()
savefig(fig, "fig01_clock_family.png")

fig, ax = plt.subplots(1, 2, figsize=(12, 4))
for sigma, c in zip((2.0, 1.25, 0.75), ("tab:blue", "tab:green", "tab:red")):
    d = np.loadtxt(os.path.join(OUT, "s01_zetaZ1_vs_zeta_sigma%.2f.csv" % sigma), delimiter=",", skiprows=1)
    ax[0].plot(d[:, 0], d[:, 3], ".", ms=2, color=c, label=r"$\sigma=%.2f$" % sigma)
ax[0].set_xlabel("t")
ax[0].set_ylabel(r"$|\zeta_{Z1}/\zeta-1|$")
ax[0].set_title(r"$\zeta_{Z1}$ from %d ordinates versus $\zeta$" % N)
ax[0].legend()
ax[1].plot(tc, Zhardy, "k-", lw=2.5, label=r"$Z(t)$ (mpmath)")
ax[1].plot(tc, Zl[0.0], "r--", lw=1.0, label=r"$Z_{Z1}(t)$ (product)")
ax[1].set_xlabel("t")
ax[1].legend()
ax[1].set_title(r"Critical line: $\lambda=0$ endpoint and Hardy $Z$")
fig.tight_layout()
savefig(fig, "fig02_zetaZ1_vs_zeta.png")

sio.savemat(os.path.join(OUT, "s01_results.mat"), {
    "N": N, "clock_err": clock_err, "eta_blocks": rows, "zetaZ1_cmp": cmp_rows,
    "t_critical": tc, "Z_hardy": Zhardy, "Z_lambda": np.array([Zl[l] for l in lams]), "lambdas": np.array(lams)})
say.close()

# ---- zoom: how the smooth clock places the midpoint ordinates next to the zeros
say2 = Tee("s01b_clock_zoom")
fig, ax = plt.subplots(1, 2, figsize=(12, 4))
T = np.linspace(9.0, 52.0, 4000)
ax[0].step(T, np.searchsorted(gam, T, side="right"), where="post", color="k", lw=1.5, label=r"$N_\zeta(T)$")
ax[0].plot(T, clock(T), "r-", lw=1.5, label=r"$\mathcal{N}_{\rm RVM}(T)$")
for i in range(9):
    ax[0].plot([gam[i], gam[i]], [0, i + 1], "k:", lw=0.8)
ax[0].plot(tau[:9], np.arange(9) + 0.5, "ro", ms=5, label=r"$\tau_k$: $\mathcal{N}_{\rm RVM}(\tau_k)=k-\frac{1}{2}$")
ax[0].plot(gam[:9], np.arange(9) + 0.5, "ks", ms=4, mfc="none", label=r"$(\gamma_k,\,k-\frac{1}{2})$")
ax[0].set_xlabel("T")
ax[0].set_title("Clock and staircase at the first zeros")
ax[0].legend(loc="upper left", fontsize=8)
kz = np.arange(N - 12, N)
T3 = np.linspace(gam[kz[0]] - 0.3, gam[kz[-1]] + 0.3, 6000)
ax[1].step(T3 - 1132480, np.searchsorted(gam, T3, side="right") - 2001000, where="post", color="k", lw=1.5, label=r"$N_\zeta(T)$")
ax[1].plot(T3 - 1132480, clock(T3) - 2001000, "r-", lw=1.5, label=r"$\mathcal{N}_{\rm RVM}(T)$")
ax[1].plot(tau[kz] - 1132480, kz + 0.5 - 2001000, "ro", ms=5, label=r"$\tau_k$")
ax[1].plot(gam[kz] - 1132480, kz + 0.5 - 2001000, "ks", ms=4, mfc="none", label=r"$(\gamma_k,\,k-\frac{1}{2})$")
ax[1].set_xlabel(r"$T-1132480$")
ax[1].set_ylabel(r"count $-\,2001000$")
ax[1].set_title(r"Zoom at the last 12 zeros of the file (mean spacing %.3f)" % (2 * np.pi / np.log(gam[-1] / (2 * np.pi))))
ax[1].legend(loc="upper left", fontsize=8)
fig.tight_layout()
savefig(fig, "fig00_clock_zoom.png")
fig, ax = plt.subplots(1, 1, figsize=(12, 3))
T2 = np.linspace(9.0, 200.0, 20000)
ax.plot(T2, np.searchsorted(gam, T2, side="right") - clock(T2), "k-", lw=0.8)
ax.axhline(0, color="r", lw=0.8)
ax.set_xlabel("T")
ax.set_ylabel(r"$S(T)+O(1/T)$")
ax.set_title(r"$N_\zeta(T)-\mathcal{N}_{\rm RVM}(T)$: the remainder of the Riemann-von Mangoldt formula (jump $+1$ at each $\gamma_k$)")
fig.tight_layout()
savefig(fig, "fig00b_remainder.png")
say2("zoom table near the top of the file:")
say2("%9s %18s %18s %10s %10s", "k", "gamma_k", "tau_k", "eta_k", "N*-N_RVM")
for i in kz:
    say2("%9d %18.6f %18.6f %10.6f %10.6f", i + 1, gam[i], tau[i], tau[i] - gam[i],
         np.searchsorted(gam, gam[i], side="right") - 0.5 - clock(gam[i]))
say2.close()
\end{lstlisting}

\subsection{\texttt{s02\_transition.py}}
\label{app:code-s02-transition}

\begin{lstlisting}[language=Python]
"""Experiment 2: the transition theorem zeta_RVM = varpi (1 + w)(1 + O(1/t)),
the constant varpi_RVM, the vertical limit of 1/zeta_RVM, and zeta_RVM versus zeta."""
import os
import numpy as np
import mpmath as mp
import scipy.io as sio
import matplotlib
matplotlib.use("Agg")
import matplotlib.pyplot as plt
from rvm_common import (Tee, load_ordinates, Product, zeta_from_product, chi, w_transition,
                        write_csv, savefig, OUT)

say = Tee("s02_transition")
gam, tau = load_ordinates()
PR = Product(tau)

# ---- varpi_RVM by Richardson extrapolation on Re s = 2
Ts = 1.0e4 * np.array([0.25, 0.5, 1.0])
v = np.array([zeta_from_product(PR, 2 + 1j * T)[0] / (1 + w_transition(2 + 1j * T)) for T in Ts])
A = np.column_stack([np.ones(3), 1 / Ts, 1 / Ts ** 2])
c = np.linalg.solve(A.astype(complex), v)
varpi = c[0].real
L = 1.0 / varpi
say("varpi_RVM (Richardson, sigma = 2, T = 2500, 5000, 10000) = %.10f %+.2e i", c[0].real, c[0].imag)
say("L = 1/varpi_RVM = %.10f", L)

# ---- relative error of the transition formula
tt = np.exp(np.linspace(np.log(10.0), np.log(1.0e4), 160))
rows = [tt]
hdr = ["t"]
fits = []
for sp in (0.25, 0.5, 1.5):
    s = 0.5 + sp + 1j * tt
    zr = zeta_from_product(PR, s)
    err_w = np.abs(zr / (varpi * (1 + w_transition(s))) - 1)
    rows.append(err_w)
    hdr.append("err_w_sigmap%.2f" % sp)
    sel = tt > 100
    slope = np.polyfit(np.log(tt[sel]), np.log(err_w[sel]), 1)[0]
    fits.append([sp, slope, np.max(err_w[sel] * tt[sel])])
    say("sigma' = %.2f: |zeta_RVM/(varpi(1+w)) - 1|: log-log slope on t>100 = %.3f, max t*err = %.4f",
        sp, slope, np.max(err_w[sel] * tt[sel]))
    if sp < 0.5:
        err_c = np.abs(zr / (varpi * (1 + chi(s))) - 1)
        rows.append(err_c)
        hdr.append("err_chi_sigmap%.2f" % sp)
        say("sigma' = %.2f: |zeta_RVM/(varpi(1+chi)) - 1|: max t*err (t>100) = %.4f",
            sp, np.max(err_c[sel] * tt[sel]))
    wmod = np.abs(w_transition(s))
    rows.append(wmod)
    hdr.append("abs_w_sigmap%.2f" % sp)
    Fv = 1 / zr
    rows.append(np.abs(Fv - L))
    hdr.append("abs_F_minus_L_sigmap%.2f" % sp)
    say("          |w| * t^sigma' in [%.4f, %.4f];  |1/zeta_RVM - L| at t = 1e4: %.3e",
        np.min(wmod * tt ** sp), np.max(wmod * tt ** sp), np.abs(Fv[-1] - L))
write_csv("s02_transition_errors.csv", hdr, rows)

# ---- zeta_RVM versus zeta on two vertical lines and on the critical line
mp.mp.dps = 20
tl = np.linspace(1.0, 120.0, 900)
cmp = {}
for sigma in (2.0, 0.75):
    s = sigma + 1j * tl
    zr = zeta_from_product(PR, s)
    ze = np.array([complex(mp.zeta(complex(x))) for x in s])
    cmp[sigma] = (zr, ze)
    write_csv("s02_zetaRVM_vs_zeta_sigma%.2f.csv" % sigma, ["t", "abs_zeta", "abs_zetaRVM", "abs_varpi_1_plus_w"],
              [tl, np.abs(ze), np.abs(zr), np.abs(varpi * (1 + w_transition(s)))])

fig, ax = plt.subplots(2, 2, figsize=(12, 8))
cols = {0.25: "tab:red", 0.5: "tab:green", 1.5: "tab:blue"}
j = 1
for sp in (0.25, 0.5, 1.5):
    ax[0, 0].plot(tt, tt * rows[j], "-", color=cols[sp], label=r"$\sigma'=%.2f$" % sp)
    j += 4 if sp < 0.5 else 3
ax[0, 0].set_xlabel("t")
ax[0, 0].set_ylabel(r"$t\,|\zeta_{\rm RVM}/(\varpi_{\rm RVM}(1+w))-1|$")
ax[0, 0].set_title(r"Transition theorem: error times $t$ stays bounded, $\varpi_{\rm RVM}=%.10f$" % varpi, fontsize=9)
ax[0, 0].legend()
for sp in (0.25, 0.5, 1.5):
    s = 0.5 + sp + 1j * tt
    ax[0, 1].plot(tt, np.abs(1 / zeta_from_product(PR, s) - L) * tt ** min(sp, 1.0), "-", color=cols[sp], label=r"$\sigma'=%.2f$" % sp)
ax[0, 1].set_xlabel("t")
ax[0, 1].set_ylabel(r"$t^{\min(\sigma',1)}\,|1/\zeta_{\rm RVM}(\sigma+it)-L|$")
ax[0, 1].set_title(r"Vertical limit $L=1/\varpi_{\rm RVM}$, rescaled error", fontsize=9)
ax[0, 1].legend()
for a, sigma in zip((ax[1, 0], ax[1, 1]), (2.0, 0.75)):
    zr, ze = cmp[sigma]
    a.plot(tl, np.abs(ze), "k-", lw=1.2, label=r"$|\zeta(\sigma+it)|$")
    a.plot(tl, np.abs(zr), "r-", lw=1.2, label=r"$|\zeta_{\rm RVM}(\sigma+it)|$")
    a.set_xlabel("t")
    a.set_title(r"$\sigma=%.2f$" % sigma)
    a.legend()
fig.tight_layout()
savefig(fig, "fig03_transition.png")

sio.savemat(os.path.join(OUT, "s02_results.mat"), {"varpi": varpi, "L": L, "richardson": c,
                                                   "slopes": np.array(fits), "t": tt})
say.close()
\end{lstlisting}

\subsection{\texttt{s03\_exponential\_sums.py}}
\label{app:code-s03-exponential-sums}

\begin{lstlisting}[language=Python]
"""Experiment 3: exponential sums S(N;u) = sum_{k<=N} exp(i u o_k) over the midpoint
ordinates tau_k and over the ordinates gamma_k of zeta, and the height transforms M_H(x)."""
import os
import numpy as np
import scipy.io as sio
import matplotlib
matplotlib.use("Agg")
import matplotlib.pyplot as plt
from rvm_common import Tee, load_ordinates, write_csv, savefig, OUT

say = Tee("s03_exponential_sums")
gam, tau = load_ordinates()
N = gam.size
step = 2000
Nk = np.arange(step, N + 1, step)
xs = [2.0, 3.0, 4.0, 6.0, 10.0]
Lam = {2.0: np.log(2), 3.0: np.log(3), 4.0: np.log(2), 6.0: 0.0, 10.0: 0.0}
res = {}
say("%6s %14s %14s %14s %16s %14s" % ("x", "max|S_RVM|", "max|S_RVM|/log", "|S_zeta(N)|",
                                     "Landau -T/2pi*L/sqrt(x)", "|M_H^RVM(x)|"))
for x in xs:
    u = np.log(x)
    SR = np.cumsum(np.exp(1j * u * tau))[Nk - 1]
    SZ = np.cumsum(np.exp(1j * u * gam))[Nk - 1]
    res[x] = (SR, SZ)
    landau = -gam[-1] / (2 * np.pi) * Lam[x] / np.sqrt(x)
    say("%6.1f %14.4f %14.4f %14.2f %16.2f %14.3e" % (x, np.max(np.abs(SR)), np.max(np.abs(SR) / np.log(Nk)),
                                                     np.abs(SZ[-1]), landau, np.abs(SR[-1]) / N))
    write_csv("s03_expsum_x%g.csv" % x, ["N", "re_S_RVM", "im_S_RVM", "re_S_zeta", "im_S_zeta"],
              [Nk, SR.real, SR.imag, SZ.real, SZ.imag])

fig, ax = plt.subplots(1, 2, figsize=(12, 4))
cols = {2.0: "tab:blue", 3.0: "tab:green", 4.0: "tab:purple", 6.0: "tab:orange", 10.0: "tab:red"}
for x in xs:
    ax[0].plot(Nk, np.abs(res[x][0]), "-", color=cols[x], lw=0.8, label=r"$x=%g$" % x)
ax[0].plot(Nk, np.log(Nk), "k--", lw=1, label=r"$\log N$")
ax[0].set_xlabel("N")
ax[0].set_ylabel(r"$|S_{\rm RVM}(N;\log x)|$")
ax[0].set_title(r"Midpoint ordinates: $|\sum_{k\leq N}x^{i\tau_k}|$ stays $O(\log N)$")
ax[0].legend(fontsize=8, ncol=2)
for x in xs:
    ax[1].plot(Nk, res[x][1].real, "-", color=cols[x], lw=0.8, label=r"$x=%g$" % x)
T = np.interp(Nk, np.arange(1, N + 1), gam)
for x in (2.0, 3.0, 4.0):
    ax[1].plot(Nk, -T / (2 * np.pi) * Lam[x] / np.sqrt(x), ":", color=cols[x], lw=1.5)
ax[1].set_xlabel("N")
ax[1].set_ylabel(r"$\mathrm{Re}\sum_{k\leq N}x^{i\gamma_k}$")
ax[1].set_title(r"Zeros of $\zeta$: Landau term $-\frac{T}{2\pi}\frac{\Lambda(x)}{\sqrt{x}}$ (dotted)")
ax[1].legend(fontsize=8, ncol=2)
fig.tight_layout()
savefig(fig, "fig04_exponential_sums.png")
sio.savemat(os.path.join(OUT, "s03_results.mat"), {"N": Nk, "x": np.array(xs),
                                                   "S_RVM": np.array([res[x][0] for x in xs]),
                                                   "S_zeta": np.array([res[x][1] for x in xs])})
say.close()
\end{lstlisting}

\subsection{\texttt{s04\_primes.py}}
\label{app:code-s04-primes}

\begin{lstlisting}[language=Python]
"""Experiment 4: the prime-power count J_RVM, the synthetic prime count pi_RVM, the cell
masses c_n^RVM and the real RVM prime levels p_j^RVM, compared with J, pi(x) and the primes.

J_o(x) = li(x) - 2 Re sum_k Ei((1/2 + i o_k) log x) - log 2 + int_x^inf dt/(t(t^2-1) log t)
for an ordinate list o_k (o = tau: J_RVM; o = gamma: Riemann's explicit formula for J),
pi_o(x) = sum_{m <= log2 x} mu(m)/m J_o(x^(1/m)).
The zero sums are evaluated on a uniform grid in u = log x by the recurrence
exp(i o_k (u + du)) = exp(i o_k u) exp(i o_k du); Ei is exact (scipy exp1) for k <= 300 and
uses Ei(w) = e^w (1/w + 1/w^2 + 2/w^3 + 6/w^4) for k > 300, where |w| > 400."""
import os
import numpy as np
import scipy.io as sio
from scipy import special, integrate
import matplotlib
matplotlib.use("Agg")
import matplotlib.pyplot as plt
from rvm_common import (Tee, load_ordinates, ei_upper, mobius_sieve, primes_upto, write_csv,
                        savefig, clock_inverse, OUT)

say = Tee("s04_primes")
gam, tau = load_ordinates()
KE = 300


def zero_sum_ei(ords, u, K=None, taper=None):
    """sum_k Ei((1/2 + i o_k) u) on a uniform grid u, k <= K, optional weights."""
    o = ords if K is None else ords[:K]
    wts = np.ones(o.size) if taper is None else taper
    rho = 0.5 + 1j * o
    out = np.zeros(u.size, dtype=complex)
    # exact part
    for i in range(0, KE):
        out += wts[i] * ei_upper(rho[i] * u)
    # asymptotic part by recurrence
    r = rho[KE:]
    b = [wts[KE:] / r ** m for m in (1, 2, 3, 4)]
    du = u[1] - u[0]
    stepf = np.exp(1j * o[KE:] * du)
    for j in range(u.size):
        if j % 200 == 0:
            E = np.exp(1j * o[KE:] * u[j])
        A = [E @ bm for bm in b]
        uj = u[j]
        out[j] += np.exp(uj / 2) * (A[0] / uj + A[1] / uj ** 2 + 2 * A[2] / uj ** 3 + 6 * A[3] / uj ** 4)
        E = E * stepf
    return out


def tail_integral(x):
    return integrate.quad(lambda t: 1.0 / (t * (t * t - 1.0) * np.log(t)), x, np.inf, limit=200)[0]


def J_from(ords, u, K=None, taper=None):
    x = np.exp(u)
    S = zero_sum_ei(ords, u, K, taper)
    ti = np.array([tail_integral(v) for v in x])
    return special.expi(u) - 2 * S.real - np.log(2.0) + ti


def pi_from_J(u, J):
    mu = mobius_sieve(64)
    out = np.zeros(u.size)
    for i, ui in enumerate(u):
        m = 1
        while ui / m >= np.log(2.0) - 1e-12:
            if mu[m] != 0:
                out[i] += mu[m] / m * np.interp(ui / m, u, J)
            m += 1
    return out


def exact_J_pi(x):
    P = primes_upto(int(x.max()) + 1)
    pi = np.searchsorted(P, x, side="right").astype(float)
    J = np.zeros(x.size)
    for m in range(1, 40):
        J += np.searchsorted(P, x ** (1.0 / m) + 1e-12, side="right") / m
    return J, pi


def run(xmax, npts, K, label):
    u = np.linspace(np.log(2.0), np.log(xmax), npts)
    if K is None:
        tg = tr = None
        Kg = Kr = None
    else:
        Kg = Kr = K
        v = np.arange(K) / K
        tg = tr = np.where(v < 0.6, 1.0, np.cos(np.pi / 2 * (v - 0.6) / 0.4) ** 2)
    JR = J_from(tau, u, Kr, tr)
    JG = J_from(gam, u, Kg, tg)
    PR = pi_from_J(u, JR)
    PG = pi_from_J(u, JG)
    x = np.exp(u)
    Je, Pe = exact_J_pi(x)
    write_csv("s04_prime_counts_%s.csv" % label, ["x", "J_exact", "J_from_zeta_zeros", "J_RVM",
                                                  "pi_exact", "pi_from_zeta_zeros", "pi_RVM", "li"],
              [x, Je, JG, JR, Pe, PG, PR, special.expi(u)])
    return x, Je, JG, JR, Pe, PG, PR


# x <= 100 with all zeros; x <= 1000 with a smooth cut at K zeros
x1, Je1, JG1, JR1, Pe1, PG1, PR1 = run(100.0, 3000, None, "x100_allzeros")
say("x <= 100, all %d zeros of each list, %d grid points:" % (gam.size, x1.size))
say("  max |pi_from_zeta_zeros - pi| away from primes (distance > 0.2): %.4f",
    np.max(np.abs(PG1 - Pe1)[np.min(np.abs(x1[:, None] - primes_upto(101)[None, :]), axis=1) > 0.2]))
say("  pi_RVM(100) = %.4f, pi(100) = %d, li(100) = %.4f, J_RVM(100) = %.4f, J(100) = %.4f",
    PR1[-1], Pe1[-1], special.expi(np.log(100.0)), JR1[-1], Je1[-1])
dPR = np.diff(PR1)
say("  pi_RVM monotone on [2,100]: %s (min increment on the grid %.3e)", bool(np.all(dPR > 0)), dPR.min())

tauK = 2 * np.pi * 1000.0 * 2.5
K = int(np.searchsorted(tau, tauK))
x2, Je2, JG2, JR2, Pe2, PG2, PR2 = run(1000.0, 40000, K, "x1000_K%d" % K)
say("x <= 1000, smooth cut at K = %d zeros (tau_K = %.1f), %d grid points:" % (K, tau[K - 1], x2.size))
say("  pi_RVM(1000) = %.4f, pi(1000) = %d, li(1000) = %.4f", PR2[-1], Pe2[-1], special.expi(np.log(1000.0)))

# cell masses and real prime levels from the fine x<=1000 run
n = np.arange(2, 1001)
pin = np.interp(n, x2, PR2)
c = np.diff(np.concatenate([[0.0], pin]))
P = primes_upto(1000)
isp = np.isin(n, P).astype(float)
say("cell masses c_n = pi_RVM(n) - pi_RVM(n-1):")
say("%5s %10s %10s %8s" % ("n", "c_n", "1/log n", "prime"))
for i in range(0, 40):
    say("%5d %10.5f %10.5f %8d" % (n[i], c[i], 1 / np.log(n[i]), isp[i]))
say("sum_{n<=1000} c_n = %.4f,  pi(1000) = %d", c.sum(), len(P))
levels = np.arange(1, int(np.floor(PR2.max())) + 1)
pj = np.array([x2[np.argmax(PR2 >= j)] for j in levels])
say("real RVM prime levels p_j^RVM (first crossing of level j) versus p_j:")
say("%5s %12s %8s %10s" % ("j", "p_j^RVM", "p_j", "diff"))
for j in range(min(30, levels.size)):
    say("%5d %12.4f %8d %10.4f" % (levels[j], pj[j], P[j], pj[j] - P[j]))
write_csv("s04_cell_masses.csv", ["n", "c_n_RVM", "inv_log_n", "is_prime"], [n, c, 1 / np.log(n), isp])
m = min(levels.size, P.size)
write_csv("s04_rvm_prime_levels.csv", ["j", "p_j_RVM", "p_j"], [levels[:m], pj[:m], P[:m]])

# ---- figures (linear axes)
fig, ax = plt.subplots(1, 2, figsize=(12, 4))
ax[0].step(x1, Pe1, where="post", color="k", lw=1.6, label=r"$\pi(x)$")
ax[0].plot(x1, PG1, "b-", lw=0.6, label=r"$\pi(x)$ from %d zeros of $\zeta$" % gam.size)
ax[0].plot(x1, PR1, "r-", lw=1.4, label=r"$\pi_{\rm RVM}(x)$ from %d $\tau_k$" % tau.size)
ax[0].plot(x1, special.expi(np.log(x1)) - special.expi(np.log(2.0)), "g--", lw=1, label=r"$\mathrm{li}(x)-\mathrm{li}(2)$")
ax[0].set_xlabel("x")
ax[0].set_title(r"Prime counts on $[2,100]$")
ax[0].legend(fontsize=7)
ax[1].step(x2, Pe2, where="post", color="k", lw=1.2, label=r"$\pi(x)$")
ax[1].plot(x2, PR2, "r-", lw=1.0, label=r"$\pi_{\rm RVM}(x)$")
ax[1].set_xlabel("x")
ax[1].set_title(r"Prime counts on $[2,1000]$ ($K=%d$ ordinates, smooth cut)" % K)
ax[1].legend(fontsize=8)
fig.tight_layout()
savefig(fig, "fig05_prime_counts.png")

fig, ax = plt.subplots(1, 2, figsize=(12, 4))
nn = n[:98]
ax[0].bar(nn, isp[:98], width=0.8, color="0.75", label="prime indicator")
ax[0].plot(nn, c[:98], "r.-", lw=0.8, ms=4, label=r"$c_n^{\rm RVM}$")
ax[0].plot(nn, 1 / np.log(nn), "g--", lw=1, label=r"$1/\log n$")
ax[0].set_xlabel("n")
ax[0].set_title(r"Cell masses $c_n^{\rm RVM}=\pi_{\rm RVM}(n)-\pi_{\rm RVM}(n-1)$")
ax[0].legend(fontsize=8)
ax[1].plot(P[:m], pj[:m], "r.", ms=4, label=r"$(p_j,\,p_j^{\rm RVM})$")
ax[1].plot([0, 1000], [0, 1000], "k-", lw=0.8, label="diagonal")
ax[1].set_xlabel(r"$p_j$")
ax[1].set_ylabel(r"$p_j^{\rm RVM}$")
ax[1].set_title(r"Real RVM prime levels $\pi_{\rm RVM}(p_j^{\rm RVM})=j$ versus primes")
ax[1].legend(fontsize=8)
fig.tight_layout()
savefig(fig, "fig06_rvm_primes.png")

sio.savemat(os.path.join(OUT, "s04_results.mat"), {"x100": x1, "pi100": Pe1, "piZeta100": PG1, "piRVM100": PR1,
                                                   "x1000": x2, "pi1000": Pe2, "piRVM1000": PR2,
                                                   "c_n": c, "p_RVM": pj, "K": K})
say.close()
\end{lstlisting}

\subsection{\texttt{s05\_mertens.py}}
\label{app:code-s05-mertens}

\begin{lstlisting}[language=Python]
"""Experiment 5: the Mertens function M, the synthetic Mertens function M_RVM^crit and the
ternary walk M_h with values in {-1,0,+1} on the squarefree support.

Explicit formulas (simple zeros, residues a = 1/(rho zeta'(rho))):
  M_zeta(x) = -2 + 2 Re sum_k a_k x^(rho_k) + trivial terms        (zeros of zeta)
  M_RVM(x)  = -2 + 2 Re sum_k b_k x^(1/2 + i tau_k) + trivial terms (b_k = 1/(rho_k zeta_RVM'(rho_k)))
Ternary walk: mu_h(n) = mu(n)^2 * sign, sign = +1 if M_h(n-1) < target(n) else -1."""
import os
import math
import numpy as np
import mpmath as mp
import scipy.io as sio
from scipy.interpolate import PchipInterpolator
import matplotlib
matplotlib.use("Agg")
import matplotlib.pyplot as plt
from rvm_common import (Tee, load_ordinates, Product, logG, clock_inverse, tail_coefficients,
                        theta_rs, mobius_sieve, write_csv, savefig, OUT)

say = Tee("s05_mertens")
gam, tau = load_ordinates()
Nz = gam.size
XMAX = 10 ** 6
K0 = 1000

# ---------------------------------------------------------------- RVM residues
c1, c2, c3 = tail_coefficients(clock_inverse(float(Nz)))
t2 = tau ** 2
zpR = np.zeros(K0, dtype=complex)
for k in range(K0):
    tk = tau[k]
    fac = (t2 - tk ** 2) / (t2 + 0.25)
    fac[k] = 1.0
    q = -0.25 - tk ** 2
    lmod = np.sum(np.log(np.abs(fac))) + c1 * q - c2 * q ** 2 + c3 * q ** 3
    rho = 0.5 + 1j * tk
    lP = np.log(2j * tk / (0.25 + tk ** 2)) + lmod
    zpR[k] = (-1) ** k * np.exp(lP - logG(rho))
ratio = zpR / np.log(tau[:K0] / (2 * np.pi))
PR = Product(tau)
Ts = 1.0e4 * np.array([0.25, 0.5, 1.0])
from rvm_common import zeta_from_product, w_transition
v = np.array([zeta_from_product(PR, 2 + 1j * T)[0] / (1 + w_transition(2 + 1j * T)) for T in Ts])
varpi = np.linalg.solve(np.column_stack([np.ones(3), 1 / Ts, 1 / Ts ** 2]).astype(complex), v)[0].real
sel = np.arange(K0 // 2, K0)
cfit = np.mean((ratio[sel] / varpi - 1) * tau[sel])
say("zeta_RVM'(rho_k) exact for k <= %d; zeta_RVM'(rho_k)/log(tau_k/2pi) at k = %d: %.10f %+.2e i",
    K0, K0, ratio[-1].real, ratio[-1].imag)
say("varpi_RVM = %.10f; fit zeta_RVM'(rho_k) = varpi log(tau_k/2pi)(1 + c/tau_k), c = %.4f %+.4f i",
    varpi, cfit.real, cfit.imag)
zR = varpi * np.log(tau / (2 * np.pi)) * (1 + cfit / tau)
zR[:K0] = zpR
bR = 1.0 / ((0.5 + 1j * tau) * zR)

# ---------------------------------------------------------------- residues of zeta
def rs_Z(t):
    a = np.sqrt(t / (2 * np.pi))
    Nn = np.floor(a).astype(int)
    p = a - Nn
    nmax = Nn.max()
    n = np.arange(1, nmax + 1)
    ph = theta_rs(t)[:, None] - t[:, None] * np.log(n)[None, :]
    mask = n[None, :] <= Nn[:, None]
    Z = 2 * np.sum(np.cos(ph) * mask / np.sqrt(n)[None, :], axis=1)
    cc = np.cos(2 * np.pi * p)
    C0 = np.cos(2 * np.pi * (p ** 2 - p - 1.0 / 16)) / cc
    return Z + (-1.0) ** (Nn - 1) * (t / (2 * np.pi)) ** (-0.25) * C0


mp.mp.dps = 20
KM = 1000
zG = np.zeros(Nz, dtype=complex)
for k in range(KM + 50):
    zG[k] = complex(mp.zeta(mp.mpc(0.5, gam[k]), 1, 1))
h = 1e-5
for i0 in range(KM, Nz, 4000):
    i1 = min(Nz, i0 + 4000)
    t = gam[i0:i1]
    Zp = (rs_Z(t + h) - rs_Z(t - h)) / (2 * h)
    val = -1j * np.exp(-1j * theta_rs(t)) * Zp
    if i0 == KM:
        chk = np.max(np.abs(val[:50] / zG[KM:KM + 50] - 1))
        say("Riemann-Siegel zeta'(rho) versus mpmath at k = %d..%d: max relative difference %.2e",
            KM + 1, KM + 50, chk)
        zG[KM + 50:i1] = val[50:]
    else:
        zG[i0:i1] = val
aG = 1.0 / ((0.5 + 1j * gam) * zG)
S2G = 2 * np.sum(np.abs(aG) ** 2)
S2R = 2 * np.sum(np.abs(bR) ** 2)
say("median |zeta'(rho)| = %.4f;  spectral sums 2 sum|a|^2: zeta %.6f, RVM %.6f", np.median(np.abs(zG)), S2G, S2R)


def trivial_zeta(x, n=6):
    v = np.zeros_like(x)
    for k in range(1, n + 1):
        v += (-1) ** (k - 1) * (2 * np.pi / x) ** (2 * k) / (math.factorial(2 * k) * k * float(mp.zeta(2 * k + 1)))
    return v


def trivial_rvm(x, n=6):
    v = np.zeros_like(x)
    for k in range(1, n + 1):
        s = -2.0 * k
        lP = np.real(PR.logP(s)[0])
        v -= x ** (-2 * k) * (2 * k + 1) * np.pi ** k * 2 * (-1) ** k / (math.factorial(k) * np.exp(lP))
    return v


def explicit_uniform(ords, wts, u):
    """2 Re sqrt(x) sum_k w_k exp(i o_k u) on a uniform u grid (recurrence)."""
    out = np.zeros(u.size)
    st = np.exp(1j * ords * (u[1] - u[0]))
    for j in range(u.size):
        if j % 200 == 0:
            E = np.exp(1j * ords * u[j])
        out[j] = 2 * np.exp(u[j] / 2) * np.real(E @ wts)
        E = E * st
    return out


# ---------------------------------------------------------------- exact M by sieve
mu = mobius_sieve(XMAX)
M = np.cumsum(mu.astype(np.int64))
sqf = mu != 0
say("M(10^6) = %d;  squarefree density Q(10^6)/10^6 = %.6f (6/pi^2 = %.6f)", M[XMAX], sqf[1:].mean(), 6 / np.pi ** 2)

# ---------------------------------------------------------------- explicit formulas on [2, 2000]
u = np.linspace(np.log(2.0), np.log(2000.0), 2500)
xg = np.exp(u)
Mz = -2 + explicit_uniform(gam, aG, u) + trivial_zeta(xg)
Mc = -2 + explicit_uniform(tau, bR, u) + trivial_rvm(xg)
Mfl = M[np.floor(xg).astype(int)]
off = np.abs(xg - np.round(xg)) > 0.05
say("explicit formula with %d zeros of zeta on [2,2000]: rms(M_zeta - M) = %.4f off integers (|x-n|>0.05)",
    Nz, np.sqrt(np.mean((Mz - Mfl)[off] ** 2)))
write_csv("s05_explicit_formulas.csv", ["x", "M_floor", "M_zeta_explicit", "M_RVM_explicit"], [xg, Mfl, Mz, Mc])

# ---------------------------------------------------------------- integer target for M_RVM^crit
def target_grid(X, div=24):
    x = [1.0]
    while x[-1] < X:
        x.append(min(X, x[-1] + max(1.0, np.floor(np.sqrt(x[-1]) / div))))
    return np.array(x)


def tapered(x, ords, a, C1=2.0, C2=4.0):
    S = np.zeros(x.size, dtype=complex)
    lx = np.log(x)
    i0 = 0
    while i0 < x.size:
        i1 = min(x.size, i0 + 200)
        K = max(1, int(np.searchsorted(ords, C2 * 2 * np.pi * np.sqrt(x[i1 - 1]))))
        uu = ords[None, :K] / (2 * np.pi * np.sqrt(x[i0:i1, None]))
        w = np.where(uu <= C1, 1.0, np.where(uu >= C2, 0.0, np.cos(np.pi / 2 * (uu - C1) / (C2 - C1)) ** 2))
        E = np.exp(1j * lx[i0:i1, None] * ords[None, :K])
        S[i0:i1] = np.sqrt(x[i0:i1]) * np.sum(E * w * a[None, :K], axis=1)
        i0 = i1
    return S


xT = target_grid(float(XMAX))
TT = -2 + 2 * np.real(tapered(xT, tau, bR)) + trivial_rvm(xT)
nn = np.arange(1, XMAX + 1)
T = np.concatenate([[0.0], PchipInterpolator(xT, TT)(nn)])
nchk = np.arange(1000, 100001, 997)
full = -2 + 2 * np.real(np.array([np.sqrt(n) * np.sum(np.exp(1j * np.log(n) * tau) * bR) for n in nchk])) + trivial_rvm(nchk.astype(float))
say("target on %d grid points; at %d integers <= 1e5: max |M_RVM(n) (all zeros) - target(n)| = %.4f",
    xT.size, nchk.size, np.max(np.abs(full - T[nchk])))

# ---------------------------------------------------------------- ternary walk
mu_h = np.zeros(XMAX + 1, dtype=np.int8)
Mh = np.zeros(XMAX + 1)
sgn = np.zeros(XMAX + 1, dtype=np.int8)
cur = 0.0
for n in range(1, XMAX + 1):
    if sqf[n]:
        sgn[n] = 1 if cur < T[n] else -1
        cur += sgn[n]
        mu_h[n] = sgn[n]
    Mh[n] = cur
err = Mh[1:] - T[1:]
say("ternary walk: max |M_h - target| on [1,1e6] = %.4f, on [1e4,1e6] = %.4f", np.max(np.abs(err)), np.max(np.abs(err[9999:])))
dens = lambda a: [np.mean(a[1:] == -1), np.mean(a[1:] == 0), np.mean(a[1:] == 1)]
say("densities of (-1, 0, +1):  mu   %.6f %.6f %.6f", *dens(mu))
say("                           mu_h %.6f %.6f %.6f   (3/pi^2 = %.6f, 1 - 6/pi^2 = %.6f)", *dens(mu_h), 3 / np.pi ** 2, 1 - 6 / np.pi ** 2)
say("sign agreement mu_h = mu on squarefrees: %.6f;  sum mu mu_h / Q = %.6f",
    np.mean(mu_h[sqf] == mu[sqf]), np.sum(mu_h.astype(int) * mu.astype(int)) / sqf[1:].sum())
say("max |M(n)|/sqrt(n) on [10,1e6] = %.4f;  max |M_RVM(n)|/sqrt(n) = %.4f;  max |M_h(n)|/sqrt(n) = %.4f",
    np.max(np.abs(M[10:]) / np.sqrt(nn[9:])), np.max(np.abs(T[10:]) / np.sqrt(nn[9:])), np.max(np.abs(Mh[10:]) / np.sqrt(nn[9:])))
say("\nNUMERICAL EXAMPLE OF THE TERNARY DIVISION (first 40 integers)")
say("%4s %4s %6s %10s %8s %6s %6s %6s %6s" % ("n", "mu", "mu^2", "target", "M_h(n-1)", "sign", "mu_h", "M_h", "M"))
rows = []
for n in range(1, 41):
    say("%4d %4d %6d %10.4f %8d %6s %6d %6d %6d" % (n, mu[n], abs(mu[n]), T[n], Mh[n - 1],
                                                     ("+1" if sgn[n] > 0 else ("-1" if sgn[n] < 0 else "-")),
                                                     mu_h[n], Mh[n], M[n]))
    rows.append([n, mu[n], abs(mu[n]), T[n], Mh[n - 1], sgn[n], mu_h[n], Mh[n], M[n]])
rows = np.array(rows, dtype=float)
write_csv("s05_ternary_example.csv", ["n", "mu", "mu2", "target", "Mh_prev", "sign", "mu_h", "Mh", "M"],
          [rows[:, i] for i in range(rows.shape[1])])

Ns = 10 ** np.arange(3, 7)
say("\nENERGY E_N = (1/N) sum_{n<=N} W(n)^2")
say("%10s %14s %14s %14s" % ("N", "E_N(M)", "E_N(M_RVM)", "E_N(M_h)"))
EN = []
for Nn in Ns:
    EN.append([Nn, np.mean(M[1:Nn + 1].astype(float) ** 2), np.mean(T[1:Nn + 1] ** 2), np.mean(Mh[1:Nn + 1] ** 2)])
    say("%10d %14.4f %14.4f %14.4f" % tuple(EN[-1]))
EN = np.array(EN)
write_csv("s05_energy.csv", ["N", "E_M", "E_MRVM", "E_Mh"], [EN[:, i] for i in range(4)])
nall = np.arange(1, XMAX + 1)
write_csv("s05_walks.csv", ["n", "M", "M_RVM_target", "M_h"], [nall, M[nall], T[nall], Mh[nall]])
mu_h[1:].astype(np.int8).tofile(os.path.join(OUT, "mu_h_int8_1e6.bin"))

# ---------------------------------------------------------------- figures (linear axes)
fig, ax = plt.subplots(1, 2, figsize=(12, 4))
w = np.arange(1, 2001)
ax[0].step(w, M[w], where="post", color="k", lw=1.0, label=r"$M(n)$")
ax[0].plot(w, T[w], "r-", lw=1.2, label=r"$M_{\rm RVM}^{\rm crit}(n)$")
ax[0].step(w, Mh[w], where="post", color="tab:green", lw=0.8, label=r"$M_h(n)$")
ax[0].set_xlabel("n")
ax[0].set_title(r"Mertens functions on $[1,2000]$, every integer")
ax[0].legend(fontsize=8)
nb = 2000
bw = XMAX // nb
xb = (np.arange(nb) + 0.5) * bw
Mb = M[1:].reshape(nb, bw)
Tb = T[1:].reshape(nb, bw)
ax[1].fill_between(xb, Mb.min(1), Mb.max(1), color="0.6", label=r"$M(n)$ (range per block)")
ax[1].plot(xb, Tb.mean(1), "r-", lw=1.0, label=r"$M_{\rm RVM}^{\rm crit}(n)$")
ax[1].plot(xb, np.sqrt(xb), "k--", lw=0.8, label=r"$\pm\sqrt{n}$")
ax[1].plot(xb, -np.sqrt(xb), "k--", lw=0.8)
ax[1].set_xlabel("n")
ax[1].set_title(r"$[1,10^6]$, blocks of %d integers" % bw)
ax[1].legend(fontsize=8)
fig.tight_layout()
savefig(fig, "fig07_mertens.png")

fig, ax = plt.subplots(1, 2, figsize=(12, 4))
w = np.arange(1, 61)
ax[0].vlines(w - 0.15, 0, mu[w], color="k", lw=2, label=r"$\mu(n)$")
ax[0].vlines(w + 0.15, 0, mu_h[w], color="tab:green", lw=2, label=r"$\mu_h(n)$")
ax[0].axhline(0, color="k", lw=0.5)
ax[0].set_xlabel("n")
ax[0].set_yticks([-1, 0, 1])
ax[0].set_title(r"Ternary values: same zero set $\{\mu(n)=0\}$, signs chosen by tracking")
ax[0].legend(fontsize=8)
w = np.arange(1, 301)
ax[1].step(w, M[w], where="post", color="k", lw=1.0, label=r"$M(n)$")
ax[1].plot(w, T[w], "r.-", lw=0.8, ms=3, label=r"target $M_{\rm RVM}^{\rm crit}(n)$")
ax[1].step(w, Mh[w], where="post", color="tab:green", lw=1.2, label=r"$M_h(n)$")
ax[1].set_xlabel("n")
ax[1].set_title(r"The walk $M_h$ follows $M_{\rm RVM}^{\rm crit}$ within $\pm1$ step")
ax[1].legend(fontsize=8)
fig.tight_layout()
savefig(fig, "fig08_ternary_walk.png")

fig, ax = plt.subplots(1, 2, figsize=(12, 4))
sel = xg <= 300
ax[0].step(np.arange(1, 301), M[1:301], where="post", color="k", lw=1.0, label=r"$M(x)$")
ax[0].plot(xg[sel], Mz[sel], "b-", lw=0.7, label=r"explicit formula, %d zeros of $\zeta$" % Nz)
ax[0].plot(xg[sel], Mc[sel], "r-", lw=1.0, label=r"$M_{\rm RVM}^{\rm crit}(x)$, %d $\tau_k$" % Nz)
ax[0].set_xlabel("x")
ax[0].set_title("Explicit formulas on [2,300]")
ax[0].legend(fontsize=7)
kk = np.arange(5000)
ax[1].plot(gam[kk], np.abs(zG[kk]) / np.log(gam[kk] / (2 * np.pi)), "b.", ms=1.5, label=r"$|\zeta'(\rho_k)|/\log(\gamma_k/2\pi)$")
ax[1].plot(tau[kk], np.abs(zR[kk]) / np.log(tau[kk] / (2 * np.pi)), "r.", ms=1.5, label=r"$|\zeta_{\rm RVM}'(\rho^{\rm RVM}_k)|/\log(\tau_k/2\pi)$")
ax[1].axhline(varpi, color="k", lw=0.8, ls="--", label=r"$\varpi_{\rm RVM}$")
ax[1].set_xlabel("ordinate")
ax[1].set_ylim(0, 4)
ax[1].set_title(r"Derivatives at the zeros, $k\leq5000$")
ax[1].legend(fontsize=7)
fig.tight_layout()
savefig(fig, "fig09_explicit_weights.png")

sio.savemat(os.path.join(OUT, "s05_results.mat"), {"varpi": varpi, "cfit": cfit, "S2G": S2G, "S2R": S2R,
                                                   "energy": EN, "ternary_example": rows,
                                                   "x_explicit": xg, "M_zeta_explicit": Mz, "M_RVM_explicit": Mc})
say.close()
\end{lstlisting}

\subsection{\texttt{s06\_prime\_ratios.py}}
\label{app:code-s06-prime-ratios}

\begin{lstlisting}[language=Python]
"""Experiment 6: the three ratios li/pi_RVM, pi/li, pi/pi_RVM and the square-root estimate
pi_RVM(x) - li(x) = O(sqrt(x) (log log x)^2 / log x), from the prime counts written by s04_primes.py.

Input: output/s04_prime_counts_x1000_K17061.csv (columns x, pi_exact, pi_RVM, li, ...).
B(x) = sqrt(x) (log log x)^2 / log x is the scale of the square-root estimate."""
import os
import numpy as np
import matplotlib
matplotlib.use("Agg")
import matplotlib.pyplot as plt
from rvm_common import Tee, savefig, write_csv, OUT

say = Tee("s06_prime_ratios")
D = np.genfromtxt(os.path.join(OUT, "s04_prime_counts_x1000_K17061.csv"), delimiter=",", names=True)
x, pe, pr, li = D["x"], D["pi_exact"], D["pi_RVM"], D["li"]
say("grid: %d points on [%.0f, %.0f]; pi_RVM from the first 17061 ordinates tau_k (smooth cut)"
    % (x.size, x[0], x[-1]))

sel = x >= 10
r1, r2, r3 = li[sel] / pr[sel], pe[sel] / li[sel], pe[sel] / pr[sel]
xs = x[sel]
say("")
say("ratios at selected x")
say("%8s %12s %12s %12s %12s %12s %12s" % ("x", "pi", "pi_RVM", "li", "li/pi_RVM", "pi/li", "pi/pi_RVM"))
for X in [10, 20, 50, 100, 200, 500, 1000]:
    i = int(np.argmin(np.abs(x - X)))
    say("%8.1f %12.4f %12.4f %12.4f %12.6f %12.6f %12.6f"
        % (x[i], pe[i], pr[i], li[i], li[i] / pr[i], pe[i] / li[i], pe[i] / pr[i]))

s8 = x >= 8
B = np.sqrt(x[s8]) * np.log(np.log(x[s8])) ** 2 / np.log(x[s8])
e_rvm = (pr[s8] - li[s8]) / B
e_pi = (pe[s8] - li[s8]) / B
say("")
say("square-root scale B(x) = sqrt(x) (log log x)^2 / log x on [8, 1000]")
say("max |pi_RVM - li| / B = %.4f   (at x = %.2f)" % (np.max(np.abs(e_rvm)), x[s8][np.argmax(np.abs(e_rvm))]))
say("max |pi - li| / B     = %.4f   (at x = %.2f)" % (np.max(np.abs(e_pi)), x[s8][np.argmax(np.abs(e_pi))]))
for X in [100, 1000]:
    m = x[s8] >= X / 2
    say("max |pi_RVM - li| / B on [%d, 1000] = %.4f" % (X // 2, np.max(np.abs(e_rvm[m]))))

# CSV files: full grid
write_csv("s06_prime_ratios.csv", ["x", "li_over_pi_RVM", "pi_over_li", "pi_over_pi_RVM"],
          [xs, r1, r2, r3])
write_csv("s06_square_root_scale.csv", ["x", "pi_RVM_minus_li_over_B", "pi_minus_li_over_B"],
          [x[s8], e_rvm, e_pi])

# ---- figure (linear axes)
fig, ax = plt.subplots(1, 2, figsize=(12, 4))
ax[0].plot(xs, r3, "k-", lw=0.8, label=r"$\pi(x)/\pi_{\rm RVM}(x)$")
ax[0].plot(xs, r2, "b-", lw=0.8, label=r"$\pi(x)/\mathrm{li}(x)$")
ax[0].plot(xs, r1, "r-", lw=1.4, label=r"$\mathrm{li}(x)/\pi_{\rm RVM}(x)$")
ax[0].axhline(1.0, color="0.5", lw=0.8, ls="--")
ax[0].set_xlabel("x")
ax[0].set_title(r"The three ratios on $[10,1000]$")
ax[0].legend(fontsize=8)
ax[1].plot(x[s8], e_pi, "b-", lw=0.6, label=r"$(\pi-\mathrm{li})/B$")
ax[1].plot(x[s8], e_rvm, "r-", lw=1.4, label=r"$(\pi_{\rm RVM}-\mathrm{li})/B$")
ax[1].axhline(0.0, color="0.5", lw=0.8, ls="--")
ax[1].set_xlabel("x")
ax[1].set_title(r"Square-root scale $B(x)=\sqrt{x}\,(\log\log x)^2/\log x$")
ax[1].legend(fontsize=8)
fig.tight_layout()
savefig(fig, "fig10_prime_ratios.png")
say.close()
\end{lstlisting}

\subsection{\texttt{s07\_product\_identity.py}}
\label{app:code-s07-product-identity}

\begin{lstlisting}[language=Python]
"""Experiment 7: the value at infinity of the endpoint.

varpi_Z1 = lim_{sigma -> +inf} zeta_Z1(sigma) = varpi_RVM * prod_k (1/4 + tau_k^2) / (1/4 + gamma_k^2),
so zeta_Z1 = zeta holds exactly when the ordered partial products
P_K = prod_{k <= K} (1/4 + gamma_k^2) / (1/4 + tau_k^2)
converge to varpi_RVM.  The partial products are computed from data/zeross.bin (gamma_k) and
data/zLamb.bin (tau_k) and compared with varpi_RVM from s02_transition.py."""
import numpy as np
import matplotlib
matplotlib.use("Agg")
import matplotlib.pyplot as plt
from rvm_common import Tee, load_ordinates, savefig, write_csv

say = Tee("s07_product_identity")
gam, tau = load_ordinates()
VARPI = 0.9619415807          # Richardson value of varpi_RVM (s02_transition.py)
N = gam.size
logP = np.cumsum(np.log1p(0.25 / gam ** 2) + 2 * np.log(gam) - np.log1p(0.25 / tau ** 2) - 2 * np.log(tau))
P = np.exp(logP)
say("N = %d ordinates; varpi_RVM = %.10f" % (N, VARPI))
say("%10s %16s %14s" % ("K", "P_K", "P_K - varpi"))
for K in [10, 100, 1000, 10 ** 4, 10 ** 5, 10 ** 6, N]:
    say("%10d %16.10f %14.3e" % (K, P[K - 1], P[K - 1] - VARPI))
for K0 in [10 ** 4, 10 ** 5, 10 ** 6]:
    seg = P[K0 - 1:]
    say("K >= %7d: min P_K = %.10f, max P_K = %.10f, mean = %.10f" % (K0, seg.min(), seg.max(), seg.mean()))
# the flow along the family: Upsilon(lam) = sum_k 2 gamma_k(lam) eta_k / (1/4 + gamma_k(lam)^2),
# eta_k = tau_k - gamma_k, gamma_k(lam) = (1 - lam) gamma_k + lam tau_k; d/dlam log varpi_lam = -Upsilon(lam)
eta = tau - gam
say("")
say("flow along the lambda-family")
for lam in [0.0, 0.25, 0.5, 0.75, 1.0]:
    gl = (1 - lam) * gam + lam * tau
    say("Upsilon(%.2f) = %.10f" % (lam, np.sum(2 * gl * eta / (0.25 + gl ** 2))))
say("sum_k 2 eta_k^2 / gamma_k^2 = %.10f" % np.sum(2 * eta ** 2 / gam ** 2))
total = np.sum(np.log((0.25 + tau ** 2) / (0.25 + gam ** 2)))
say("integral_0^1 Upsilon = sum_k log((1/4+tau_k^2)/(1/4+gamma_k^2)) = %.10f" % total)
say("log L = log(1/varpi_RVM)                               = %.10f" % np.log(1 / VARPI))
step = 50
idx = np.arange(step - 1, N, step)
write_csv("s07_partial_products.csv", ["K", "P_K"], [idx + 1, P[idx]])

fig, ax = plt.subplots(1, 2, figsize=(12, 4))
k1 = np.arange(1, 10 ** 4 + 1)
ax[0].plot(k1, P[:10 ** 4], "r-", lw=0.6, label=r"$P_K$")
ax[0].axhline(VARPI, color="k", lw=0.8, ls="--", label=r"$\varpi_{\rm RVM}$")
ax[0].set_xlabel("K")
ax[0].set_title(r"Partial products $P_K$, $K\leq 10^4$")
ax[0].legend(fontsize=8)
sel = idx >= 10 ** 5 - 1
ax[1].plot(idx[sel] + 1, (P[idx[sel]] - VARPI) * 1e6, "r-", lw=0.6)
ax[1].axhline(0.0, color="k", lw=0.8, ls="--")
ax[1].set_xlabel("K")
ax[1].set_ylabel(r"$(P_K-\varpi_{\rm RVM})\cdot 10^{6}$")
ax[1].set_title(r"$10^5\leq K\leq %d$" % N)
fig.tight_layout()
savefig(fig, "fig11_product_identity.png")
say.close()
\end{lstlisting}


\section{Console output of the programs}
\label{app:output}

The complete console output of each program, as written to \texttt{output/*\_console.txt}.

\subsection{\texttt{s01\_clock\_and\_family}}
\label{app:out-s01-clock-and-family}

\begin{lstlisting}[style=out]
zeros loaded: N = 2001052, gamma_1 = 14.134725142, gamma_N = 1132490.658714, tau_1 = 14.521346953, tau_N = 1132490.722588
max |N_RVM(tau_k) - (k - 1/2)| = 6.985e-10
eta_k = tau_k - gamma_k:  max |eta| = 0.993711 (k = 34),  mean |eta| = 0.151209,  rms = 0.188066,  mean eta = 1.731e-07
      k from         k to     max|eta|      rms eta
           1           10     0.840022     0.545768
          11          100     0.993711     0.453166
         101         1000     0.915450     0.345177
        1001        10000     0.833881     0.274177
       10001       100000     0.785835     0.226352
      100001      1000000     0.762816     0.192411
     1000001      2001052     0.717359     0.178900
first ordinates:
   k          gamma_k            tau_k        eta_k
   1     14.134725142     14.521346953     0.386622
   2     21.022039639     20.655740356    -0.366299
   3     25.010857580     25.492675432     0.481818
   4     30.424876126     29.739411632    -0.685464
   5     32.935061588     33.624531889     0.689470
   6     37.586178159     37.257370087    -0.328808
   7     40.918719012     40.700551604    -0.218167
   8     43.327073281     43.994013953     0.666941
   9     48.005150881     47.165128427    -0.840022
  10     49.773832478     50.233653391     0.459821
  11     52.970321478     53.214422874     0.244101
  12     56.446247697     56.118921948    -0.327326
sigma = 2.00, t in [1,2000] (499 points): max |zeta_Z1/zeta - 1| = 8.706e-10, median = 2.685e-10
sigma = 1.25, t in [1,2000] (499 points): max |zeta_Z1/zeta - 1| = 1.128e-09, median = 3.693e-10
sigma = 0.75, t in [1,2000] (499 points): max |zeta_Z1/zeta - 1| = 2.629e-09, median = 5.179e-10
critical line, t in [0.5,60]: max |Z_Z1(t) - Z(t)| / max(1,|Z(t)|) = 8.370e-10
total run time 30.6 s
\end{lstlisting}

\subsection{\texttt{s01b\_clock\_zoom}}
\label{app:out-s01b-clock-zoom}

\begin{lstlisting}[style=out]
zoom table near the top of the file:
        k            gamma_k              tau_k      eta_k   N*-N_RVM
  2001041     1132485.303923     1132485.011569  -0.292355  -0.563104
  2001042     1132485.662638     1132485.530752  -0.131886  -0.254025
  2001043     1132486.244063     1132486.049936  -0.194127  -0.373908
  2001044     1132486.392256     1132486.569120   0.176863   0.340656
  2001045     1132486.874881     1132487.088303   0.213422   0.411072
  2001046     1132487.571868     1132487.607487   0.035619   0.068605
  2001047     1132487.950694     1132488.126670   0.175977   0.338949
  2001048     1132488.905328     1132488.645854  -0.259474  -0.499774
  2001049     1132489.162780     1132489.165037   0.002257   0.004348
  2001050     1132489.585215     1132489.684221   0.099006   0.190696
  2001051     1132490.165305     1132490.203404   0.038099   0.073383
  2001052     1132490.658714     1132490.722588   0.063873   0.123026
total run time 0.6 s
\end{lstlisting}

\subsection{\texttt{s02\_transition}}
\label{app:out-s02-transition}

\begin{lstlisting}[style=out]
varpi_RVM (Richardson, sigma = 2, T = 2500, 5000, 10000) = 0.9619415807 -1.28e-10 i
L = 1/varpi_RVM = 1.0395641690
sigma' = 0.25: |zeta_RVM/(varpi(1+w)) - 1|: log-log slope on t>100 = -1.181, max t*err = 0.0617
sigma' = 0.25: |zeta_RVM/(varpi(1+chi)) - 1|: max t*err (t>100) = 0.0790
          |w| * t^sigma' in [1.5832, 1.5832];  |1/zeta_RVM - L| at t = 1e4: 1.425e-01
sigma' = 0.50: |zeta_RVM/(varpi(1+w)) - 1|: log-log slope on t>100 = -1.181, max t*err = 0.0618
          |w| * t^sigma' in [2.5061, 2.5066];  |1/zeta_RVM - L| at t = 1e4: 2.543e-02
sigma' = 1.50: |zeta_RVM/(varpi(1+w)) - 1|: log-log slope on t>100 = -1.182, max t*err = 0.0620
          |w| * t^sigma' in [15.6619, 15.7496];  |1/zeta_RVM - L| at t = 1e4: 1.598e-05
total run time 11.5 s
\end{lstlisting}

\subsection{\texttt{s03\_exponential\_sums}}
\label{app:out-s03-exponential-sums}

\begin{lstlisting}[style=out]
     x     max|S_RVM| max|S_RVM|/log    |S_zeta(N)| Landau -T/2pi*L/sqrt(x)   |M_H^RVM(x)|
   2.0         3.5745         0.2695       88300.77        -88341.59      1.722e-06
   3.0         3.9750         0.3759      114267.92       -114324.31      1.175e-06
   4.0         3.9523         0.3732       62435.29        -62466.94      8.004e-07
   6.0         4.5264         0.4547           0.26            -0.00      2.262e-06
  10.0         9.6989         1.1495           2.42            -0.00      4.031e-06
total run time 1.7 s
\end{lstlisting}

\subsection{\texttt{s04\_primes}}
\label{app:out-s04-primes}

\begin{lstlisting}[style=out]
x <= 100, all 2001052 zeros of each list, 3000 grid points:
  max |pi_from_zeta_zeros - pi| away from primes (distance > 0.2): 0.4035
  pi_RVM(100) = 26.0782, pi(100) = 25, li(100) = 30.1261, J_RVM(100) = 29.5889, J(100) = 28.5333
  pi_RVM monotone on [2,100]: False (min increment on the grid -2.528e-01)
x <= 1000, smooth cut at K = 17061 zeros (tau_K = 15707.6), 40000 grid points:
  pi_RVM(1000) = 169.9750, pi(1000) = 168, li(1000) = 177.6097
cell masses c_n = pi_RVM(n) - pi_RVM(n-1):
    n        c_n    1/log n    prime
    2    0.52944    1.44270        1
    3    1.04188    0.91024        1
    4    0.53593    0.72135        0
    5    0.46739    0.62133        1
    6    0.36089    0.55811        0
    7    0.63254    0.51390        1
    8    0.62393    0.48090        0
    9   -0.37669    0.45512        0
   10    0.17151    0.43429        0
   11    0.69779    0.41703        1
   12    0.32703    0.40243        0
   13    0.36227    0.38987        1
   14    0.57662    0.37892        0
   15    0.19671    0.36927        0
   16   -0.05284    0.36067        0
   17    0.39958    0.35296        1
   18    0.76865    0.34598        0
   19    0.43500    0.33962        1
   20   -0.03342    0.33381        0
   21    0.03009    0.32846        0
   22    0.37466    0.32352        0
   23    0.40277    0.31893        1
   24    0.11768    0.31466        0
   25    0.01006    0.31067        0
   26    0.27350    0.30693        0
   27    0.55508    0.30341        0
   28    0.47627    0.30010        0
   29    0.12577    0.29697        1
   30   -0.08367    0.29401        0
   31    0.07329    0.29121        1
   32    0.38040    0.28854        0
   33    0.43927    0.28600        0
   34    0.34619    0.28358        0
   35    0.06419    0.28127        0
   36   -0.02475    0.27906        0
   37    0.15448    0.27694        1
   38    0.40992    0.27491        0
   39    0.49628    0.27296        0
   40    0.35226    0.27109        0
   41    0.14030    0.26928        1
sum_{n<=1000} c_n = 169.9750,  pi(1000) = 168
real RVM prime levels p_j^RVM (first crossing of level j) versus p_j:
    j      p_j^RVM      p_j       diff
    1       2.2139        2     0.2139
    2       3.1550        3     0.1550
    3       5.1887        5     0.1887
    4       7.1593        7     0.1593
    5      11.1406       11     0.1406
    6      14.0208       13     1.0208
    7      17.2552       17     0.2552
    8      19.1781       19     0.1781
    9      26.0536       23     3.0536
   10      28.0447       29    -0.9553
   11      33.0707       31     2.0707
   12      38.0931       37     1.0931
   13      40.2282       41    -0.7718
   14      44.1861       43     1.1861
   15      46.1221       47    -0.8779
   16      57.1779       53     4.1779
   17      59.1201       59     0.1201
   18      61.0809       61     0.0809
   19      67.0384       67     0.0384
   20      74.0931       71     3.0931
   21      76.1235       73     3.1235
   22      83.0822       79     4.0822
   23      86.0513       83     3.0513
   24      92.0825       89     3.0825
   25      94.0781       97    -2.9219
   26      96.1319      101    -4.8681
   27     103.2219      103     0.2219
   28     112.1163      107     5.1163
   29     115.0277      109     6.0277
   30     118.1431      113     5.1431
total run time 107.7 s
\end{lstlisting}

\subsection{\texttt{s05\_mertens}}
\label{app:out-s05-mertens}

\begin{lstlisting}[style=out]
zeta_RVM'(rho_k) exact for k <= 1000; zeta_RVM'(rho_k)/log(tau_k/2pi) at k = 1000: 0.9619388429 -1.41e-05 i
varpi_RVM = 0.9619415807; fit zeta_RVM'(rho_k) = varpi log(tau_k/2pi)(1 + c/tau_k), c = -0.0052 -0.0208 i
Riemann-Siegel zeta'(rho) versus mpmath at k = 1001..1050: max relative difference 1.19e-05
median |zeta'(rho)| = 7.8899;  spectral sums 2 sum|a|^2: zeta 0.029034, RVM 0.027203
M(10^6) = 212;  squarefree density Q(10^6)/10^6 = 0.607926 (6/pi^2 = 0.607927)
explicit formula with 2001052 zeros of zeta on [2,2000]: rms(M_zeta - M) = 0.0005 off integers (|x-n|>0.05)
target on 49896 grid points; at 100 integers <= 1e5: max |M_RVM(n) (all zeros) - target(n)| = 0.0191
ternary walk: max |M_h - target| on [1,1e6] = 2.6155, on [1e4,1e6] = 1.6539
densities of (-1, 0, +1):  mu   0.303857 0.392074 0.304069
                           mu_h 0.304003 0.392074 0.303923   (3/pi^2 = 0.303964, 1 - 6/pi^2 = 0.392073)
sign agreement mu_h = mu on squarefrees: 0.500344;  sum mu mu_h / Q = 0.000688
max |M(n)|/sqrt(n) on [10,1e6] = 0.8321;  max |M_RVM(n)|/sqrt(n) = 0.7541;  max |M_h(n)|/sqrt(n) = 0.6882

NUMERICAL EXAMPLE OF THE TERNARY DIVISION (first 40 integers)
   n   mu   mu^2     target M_h(n-1)   sign   mu_h    M_h      M
   1    1      1     0.2925        0     +1      1      1      1
   2   -1      1     0.4563        1     -1     -1      0      0
   3   -1      1    -0.6523        0     -1     -1     -1     -1
   4    0      0    -0.8720       -1      -      0     -1     -1
   5   -1      1    -1.7255       -1     -1     -1     -2     -2
   6    1      1    -1.3063       -2     +1      1     -1     -1
   7   -1      1    -1.5931       -1     -1     -1     -2     -2
   8    0      0    -2.4029       -2      -      0     -2     -2
   9    0      0    -1.3292       -2      -      0     -2     -2
  10    1      1    -1.3280       -2     +1      1     -1     -1
  11   -1      1    -2.3272       -1     -1     -1     -2     -2
  12    0      0    -2.2067       -2      -      0     -2     -2
  13   -1      1    -1.8468       -2     +1      1     -1     -3
  14    1      1    -1.9324       -1     -1     -1     -2     -2
  15    1      1    -1.5435       -2     +1      1     -1     -1
  16    0      0    -1.0414       -1      -      0     -1     -1
  17   -1      1    -1.6970       -1     -1     -1     -2     -2
  18    0      0    -3.0095       -2      -      0     -2     -2
  19   -1      1    -3.2869       -2     -1     -1     -3     -3
  20    0      0    -2.2604       -3      -      0     -3     -3
  21    1      1    -1.2515       -3     +1      1     -2     -2
  22    1      1    -1.1120       -2     +1      1     -1     -1
  23   -1      1    -1.3607       -1     -1     -1     -2     -2
  24    0      0    -1.3669       -2      -      0     -2     -2
  25    0      0    -1.3810       -2      -      0     -2     -2
  26    1      1    -1.9805       -2     +1      1     -1     -1
  27    0      0    -3.0141       -1      -      0     -1     -1
  28    0      0    -3.6155       -1      -      0     -1     -1
  29   -1      1    -3.2066       -1     -1     -1     -2     -2
  30   -1      1    -2.1599       -2     -1     -1     -3     -3
  31   -1      1    -1.3649       -3     +1      1     -2     -4
  32    0      0    -1.2973       -2      -      0     -2     -4
  33    1      1    -1.6435       -2     +1      1     -1     -3
  34    1      1    -1.7783       -1     -1     -1     -2     -2
  35    1      1    -1.4711       -2     +1      1     -1     -1
  36    0      0    -1.0723       -1      -      0     -1     -1
  37   -1      1    -1.0774       -1     -1     -1     -2     -2
  38    1      1    -1.5897       -2     +1      1     -1     -1
  39    1      1    -2.2227       -1     -1     -1     -2      0
  40    0      0    -2.4881       -2      -      0     -2      0

ENERGY E_N = (1/N) sum_{n<=N} W(n)^2
         N         E_N(M)     E_N(M_RVM)       E_N(M_h)
      1000        14.6400        18.6280        18.5960
     10000       163.6747       195.6965       196.0683
    100000      1622.3189      1308.3041      1308.8186
   1000000     14851.0857     12537.9008     12538.1174
total run time 148.2 s
\end{lstlisting}

\subsection{\texttt{s06\_prime\_ratios}}
\label{app:out-s06-prime-ratios}

\begin{lstlisting}[style=out]
grid: 40000 points on [2, 1000]; pi_RVM from the first 17061 ordinates tau_k (smooth cut)

ratios at selected x
       x           pi       pi_RVM           li    li/pi_RVM        pi/li    pi/pi_RVM
    10.0       4.0000       3.9876       6.1657     1.546218     0.648747     1.003105
    20.0       8.0000       7.6636       9.9052     1.292496     0.807657     1.043893
    50.0      15.0000      14.7686      18.4695     1.250592     0.812151     1.015670
   100.0      25.0000      26.0813      30.1265     1.155098     0.829835     0.958541
   200.0      46.0000      43.3656      50.1910     1.157393     0.916499     1.060749
   500.0      95.0000      93.7489     101.7975     1.085854     0.933225     1.013346
  1000.0     168.0000     169.9750     177.6097     1.044916     0.945895     0.988381

square-root scale B(x) = sqrt(x) (log log x)^2 / log x on [8, 1000]
max |pi_RVM - li| / B = 2.7213   (at x = 9.73)
max |pi - li| / B     = 2.4490   (at x = 11.00)
max |pi_RVM - li| / B on [50, 1000] = 1.3455
max |pi_RVM - li| / B on [500, 1000] = 0.7216
total run time 1.2 s
\end{lstlisting}

\subsection{\texttt{s07\_product\_identity}}
\label{app:out-s07-product-identity}

\begin{lstlisting}[style=out]
N = 2001052 ordinates; varpi_RVM = 0.9619415807
         K              P_K    P_K - varpi
        10     0.9625640008      6.224e-04
       100     0.9624787401      5.372e-04
      1000     0.9622891707      3.476e-04
     10000     0.9619513787      9.798e-06
    100000     0.9619301865     -1.139e-05
   1000000     0.9619418189      2.382e-07
   2001052     0.9619415390     -4.166e-08
K >=   10000: min P_K = 0.9618023968, max P_K = 0.9620641927, mean = 0.9619418623
K >=  100000: min P_K = 0.9619210208, max P_K = 0.9619605076, mean = 0.9619417591
K >= 1000000: min P_K = 0.9619389488, max P_K = 0.9619439825, mean = 0.9619416633

flow along the lambda-family
Upsilon(0.00) = 0.0436949578
Upsilon(0.25) = 0.0412445375
Upsilon(0.50) = 0.0387987180
Upsilon(0.75) = 0.0363572218
Upsilon(1.00) = 0.0339197739
sum_k 2 eta_k^2 / gamma_k^2 = 0.0098210688
integral_0^1 Upsilon = sum_k log((1/4+tau_k^2)/(1/4+gamma_k^2)) = 0.0388016004
log L = log(1/varpi_RVM)                               = 0.0388015571
total run time 1.5 s
\end{lstlisting}

\begin{thebibliography}{99}

\bibitem{BatemanDiamond1969}
P.~T. Bateman and H.~G. Diamond, Asymptotic distribution of Beurling's generalized prime numbers, in: \emph{Studies in Number Theory}, MAA Studies in Mathematics~6, Mathematical Association of America, 1969.

\bibitem{BenedettoPowellYilmaz2006}
J.~J. Benedetto, A.~M. Powell and \"O.~Y{\i}lmaz, Sigma-delta ($\Sigma\Delta$) quantization and finite frames, \emph{IEEE Trans. Inform. Theory} \textbf{52} (2006), no.~5, 1990--2005.
\href{https://doi.org/10.1109/TIT.2006.872849}{doi:10.1109/TIT.2006.872849}.

\bibitem{Beurling1937}
A.~Beurling, Analyse de la loi asymptotique de la distribution des nombres premiers g\'en\'eralis\'es.~I, \emph{Acta Math.} \textbf{68} (1937), 255--291.
\href{https://doi.org/10.1007/BF02546666}{doi:10.1007/BF02546666}.

\bibitem{BochnerChandrasekharan1956}
S.~Bochner and K.~Chandrasekharan, On Riemann's functional equation, \emph{Ann. of Math. (2)} \textbf{63} (1956), 336--360.
\href{https://doi.org/10.2307/1969614}{doi:10.2307/1969614}.

\bibitem{BrouckeDebruyneVindas2020}
F.~Broucke, G.~Debruyne and J.~Vindas, Beurling integers with RH and large oscillation, \emph{Adv. Math.} \textbf{370} (2020), 107240.
\href{https://doi.org/10.1016/j.aim.2020.107240}{doi:10.1016/j.aim.2020.107240}.

\bibitem{BrouckeVindas2024}
F.~Broucke and J.~Vindas, A new generalized prime random approximation procedure and some of its applications, \emph{Math. Z.} \textbf{307} (2024), Paper No.~62.
\href{https://doi.org/10.1007/s00209-024-03526-4}{doi:10.1007/s00209-024-03526-4}.

\bibitem{Corless1996}
R.~M. Corless, G.~H. Gonnet, D.~E.~G. Hare, D.~J. Jeffrey and D.~E. Knuth, On the Lambert $W$ function, \emph{Adv. Comput. Math.} \textbf{5} (1996), 329--359.
\href{https://doi.org/10.1007/BF02124750}{doi:10.1007/BF02124750}.

\bibitem{DaubechiesDeVore2003}
I.~Daubechies and R.~DeVore, Approximating a bandlimited function using very coarsely quantized data: a family of stable sigma-delta modulators of arbitrary order, \emph{Ann. of Math.} \textbf{158} (2003), 679--710.
\href{https://doi.org/10.4007/annals.2003.158.679}{doi:10.4007/annals.2003.158.679}.

\bibitem{Davenport1980}
H.~Davenport, \emph{Multiplicative Number Theory}, 2nd ed., revised by H.~L. Montgomery, Graduate Texts in Mathematics~74, Springer, New York, 1980.
\href{https://doi.org/10.1007/978-1-4757-5927-3}{doi:10.1007/978-1-4757-5927-3}.

\bibitem{delaValleePoussin1896}
C.-J. de la Vall\'ee Poussin, Recherches analytiques sur la th\'eorie des nombres premiers, \emph{Ann. Soc. Sci. Bruxelles} \textbf{20} (1896), 183--256.

\bibitem{Diamond1970}
H.~G. Diamond, Asymptotic distribution of Beurling's generalized integers, \emph{Illinois J. Math.} \textbf{14} (1970).
\href{https://doi.org/10.1215/ijm/1256053295}{doi:10.1215/ijm/1256053295}.

\bibitem{DiamondMontgomeryVorhauer2006}
H.~G. Diamond, H.~L. Montgomery and U.~M.~A. Vorhauer, Beurling primes with large oscillation, \emph{Math. Ann.} \textbf{334} (2006), 1--36.
\href{https://doi.org/10.1007/s00208-005-0638-2}{doi:10.1007/s00208-005-0638-2}.

\bibitem{DiamondZhang2016}
H.~G. Diamond and W.-B. Zhang, \emph{Beurling Generalized Numbers}, Mathematical Surveys and Monographs~213, American Mathematical Society, Providence, RI, 2016.
\href{https://doi.org/10.1090/surv/213}{doi:10.1090/surv/213}.

\bibitem{DLMF2024}
NIST Digital Library of Mathematical Functions, F.~W.~J. Olver, A.~B. Olde Daalhuis, D.~W. Lozier et al. (eds.), \S4.13 (Lambert $W$-function) and \S5.11 (Stirling series).
\url{https://dlmf.nist.gov/}.

\bibitem{Edwards1974}
H.~M. Edwards, \emph{Riemann's Zeta Function}, Academic Press, New York, 1974.

\bibitem{FrancaLeClair2014}
G.~Fran\c{c}a and A.~LeClair, Statistical and other properties of Riemann zeros based on an explicit equation for the $n$-th zero on the critical line, arXiv:1307.8395.
\href{https://doi.org/10.48550/arXiv.1307.8395}{doi:10.48550/arXiv.1307.8395}.

\bibitem{GrosswaldSchnitzer1978}
E.~Grosswald and F.~J. Schnitzer, A class of modified $\zeta$ and $L$-functions, \emph{Pacific J. Math.} \textbf{74} (1978), no.~2, 357--364.
\href{https://doi.org/10.2140/pjm.1978.74.357}{doi:10.2140/pjm.1978.74.357}.

\bibitem{GunturkThao2003}
C.~S. G\"unt\"urk and N.~T. Thao, Ergodic dynamics in sigma-delta quantization: tiling invariant sets and spectral analysis of error, arXiv:math/0308223.
\href{https://doi.org/10.48550/arXiv.math/0308223}{doi:10.48550/arXiv.math/0308223}.

\bibitem{Hadamard1896}
J.~Hadamard, Sur la distribution des z\'eros de la fonction $\zeta(s)$ et ses cons\'equences arithm\'etiques, \emph{Bull. Soc. Math. France} \textbf{24} (1896), 199--220.
\href{https://doi.org/10.24033/bsmf.545}{doi:10.24033/bsmf.545}.

\bibitem{Hamburger1921}
H.~Hamburger, \"Uber die Riemannsche Funktionalgleichung der $\zeta$-Funktion, \emph{Math. Z.} \textbf{10} (1921), 240--254.
\href{https://doi.org/10.1007/BF01211612}{doi:10.1007/BF01211612}.

\bibitem{KaczorowskiPerelli1999}
J.~Kaczorowski and A.~Perelli, On the structure of the Selberg class, I: $0\le d\le1$, \emph{Acta Math.} \textbf{182} (1999), 207--241.
\href{https://doi.org/10.1007/BF02392574}{doi:10.1007/BF02392574}.

\bibitem{Knopp1994}
M.~I. Knopp, On Dirichlet series satisfying Riemann's functional equation, \emph{Invent. Math.} \textbf{117} (1994), 361--372.
\href{https://doi.org/10.1007/BF01232248}{doi:10.1007/BF01232248}.

\bibitem{LapidusvanFrankenhuysen2000}
M.~L. Lapidus and M.~van Frankenhuysen, \emph{Fractal Geometry and Number Theory: Complex Dimensions of Fractal Strings and Zeros of Zeta Functions}, Birkh\"auser, Boston, 2000.
\href{https://doi.org/10.1007/978-1-4612-5314-3}{doi:10.1007/978-1-4612-5314-3}.

\bibitem{OdlyzkoTeRiele1985}
A.~M. Odlyzko and H.~J.~J. te~Riele, Disproof of the Mertens conjecture, \emph{J. reine angew. Math.} \textbf{357} (1985), 138--160.
\href{https://doi.org/10.1515/crll.1985.357.138}{doi:10.1515/crll.1985.357.138}.

\bibitem{Nakamura2020}
T.~Nakamura, Dirichlet series with periodic coefficients, Riemann's functional equation and real zeros of Dirichlet $L$-functions, preprint, arXiv:2008.02570 (2020).
\href{https://doi.org/10.48550/arXiv.2008.02570}{doi:10.48550/arXiv.2008.02570}.

\bibitem{Perelli2016}
A.~Perelli, Converse theorems: from the Riemann zeta function to the Selberg class, arXiv:1605.02354.
\href{https://doi.org/10.48550/arXiv.1605.02354}{doi:10.48550/arXiv.1605.02354}.

\bibitem{PlattTrudgian2021}
D.~J. Platt and T.~S. Trudgian, The Riemann hypothesis is true up to $3\cdot10^{12}$, \emph{Bull. London Math. Soc.} \textbf{53} (2021); arXiv:2004.09765.
\href{https://doi.org/10.48550/arXiv.2004.09765}{doi:10.48550/arXiv.2004.09765}.

\bibitem{Riemann1859}
B.~Riemann, \"Uber die Anzahl der Primzahlen unter einer gegebenen Gr\"osse, \emph{Monatsberichte der K\"oniglich Preussischen Akademie der Wissenschaften zu Berlin} (1859), 671--680.

\bibitem{Siegel1922}
C.~L. Siegel, Bemerkung zu einem Satz von Hamburger \"uber die Funktionalgleichung der Riemannschen Zetafunktion, \emph{Math. Ann.} \textbf{86} (1922), 276--279.
\href{https://doi.org/10.1007/BF01457989}{doi:10.1007/BF01457989}.

\bibitem{Titchmarsh1939}
E.~C. Titchmarsh, \emph{The Theory of Functions}, 2nd ed., Oxford University Press, 1939.

\bibitem{Titchmarsh1986}
E.~C. Titchmarsh, \emph{The Theory of the Riemann Zeta-Function}, 2nd ed., revised by D.~R. Heath-Brown, Clarendon Press, Oxford, 1986.

\bibitem{vonMangoldt1895}
H.~von Mangoldt, Zu Riemanns Abhandlung ``\"Uber die Anzahl der Primzahlen unter einer gegebenen Gr\"osse'', \emph{J. reine angew. Math.} \textbf{114} (1895), 255--305.
\href{https://doi.org/10.1515/crll.1895.114.255}{doi:10.1515/crll.1895.114.255}.

\end{thebibliography}
\end{document}
